\documentclass[a4paper,11pt]{ltjsarticle}
\usepackage[margin=25mm,headheight=16pt]{geometry}
\usepackage{amsmath,amssymb,amsthm,mathtools,bm,mathrsfs}
\usepackage{booktabs,longtable,array,enumitem}
\usepackage{xcolor}
\usepackage[unicode,colorlinks=true,linkcolor=blue!40!black,citecolor=blue!40!black,urlcolor=blue!45!black]{hyperref}
\usepackage{fancyhdr}
\usepackage{microtype}
\setlength{\parindent}{1em}
\setlength{\parskip}{0.25em}
\setlist{itemsep=0.2em,topsep=0.4em}
\allowdisplaybreaks[2]
\pagestyle{fancy}
\fancyhf{}
\fancyhead[L]{\small 3次超曲面のPoisson構造と非可換化}
\fancyhead[R]{\small 2026-11-09 / v6}
\fancyfoot[C]{\thepage}
\newcommand{\CC}{\mathbb C}
\newcommand{\ZZ}{\mathbb Z}
\newcommand{\PP}{\mathbb P}
\newcommand{\OO}{\mathcal O}
\newcommand{\cK}{\mathcal K}
\newcommand{\cP}{\mathcal P}
\newcommand{\cD}{\mathcal D}
\newcommand{\cC}{\mathcal C}
\newcommand{\cE}{\mathcal E}
\newcommand{\cF}{\mathcal F}
\newcommand{\cM}{\mathcal M}
\newcommand{\fm}{\mathfrak m}
\newcommand{\Perf}{\operatorname{Perf}}
\newcommand{\Ku}{\operatorname{Ku}}
\newcommand{\MF}{\operatorname{MF}}
\newcommand{\HH}{\operatorname{HH}}
\newcommand{\QGr}{\operatorname{QGr}}
\newcommand{\qgr}{\operatorname{qgr}}
\newcommand{\Dsg}{D_{\mathrm{sg}}}
\newcommand{\RHom}{\operatorname{RHom}}
\newcommand{\REnd}{\operatorname{REnd}}
\newcommand{\Hom}{\operatorname{Hom}}
\newcommand{\Ext}{\operatorname{Ext}}
\newcommand{\Ind}{\operatorname{Ind}}
\newcommand{\Def}{\operatorname{Def}}
\newcommand{\thick}{\operatorname{thick}}
\newcommand{\Spec}{\operatorname{Spec}}
\newcommand{\Spf}{\operatorname{Spf}}
\newcommand{\Id}{\operatorname{Id}}
\newcommand{\sym}{\operatorname{sym}}
\newcommand{\Cone}{\operatorname{Cone}}
\newcommand{\rank}{\operatorname{rank}}
\newcommand{\Pic}{\operatorname{Pic}}
\newcommand{\otL}{\mathbin{\otimes^{\mathbf L}}}
\newcommand{\DocDate}{2026年11月9日}
\newcommand{\DocVersion}{論文分割版 v6}
\theoremstyle{plain}
\newtheorem{theorem}{定理}[section]
\newtheorem{proposition}[theorem]{命題}
\newtheorem{lemma}[theorem]{補題}
\newtheorem{corollary}[theorem]{系}
\theoremstyle{definition}
\newtheorem{definition}[theorem]{定義}
\newtheorem{problem}[theorem]{問題}
\newtheorem{example}[theorem]{例}
\theoremstyle{remark}
\newtheorem{remark}[theorem]{注意}
\renewcommand{\proofname}{証明}
\hypersetup{pdfauthor={},pdfsubject={3次超曲面のPoisson構造・非可換量子化・rotation付き残余圏}}

\hypersetup{pdftitle={非可換三次超曲面のrotation付き可換復元}}
\title{非可換三次超曲面のrotation付き可換復元\\
\large 具体的方程式・一次変形・存在の障害}
\author{}
\date{2026年11月10日・論文C（研究草稿 v21）}
\fancyhead[R]{\small 2026-11-10 / v21}
\begin{document}
\maketitle
\begin{abstract}
滑らかな可換cubicのPoisson量子化とその中心三次商から得る
残余dg圏について、標準次数シフトに対応するrotationを保持した
可換モデルの存在を研究する。
形式的なmarked復元、実在するsmooth proper代数的族でのétale局所実現、
明示的な楕円Sklyanin一ブロックと一方有限中心の多ブロック比較を整理する。
一方、Calabi--YauなZhang twistの例は任意のAS-regular中心三次商の
rotation付き可換復元が偽であることを示す。
一般混合項・二つ以上の無限位数ブロック・相対saturationは未解決である。
偏極付きambient復元（論文B）からrotation付き存在は自動的に従わない。
本研究の究極の目標は、一般次元でPoissonスキームの任意の既約成分の一般点を扱うことである。
各成分が零Poisson構造を含むことと、次元合同類によるrotationの復元可能性を整理する。
étale局所実現には、同値双加群の障害をrotation不変部分で消去する独立の証明を与える。
一般混合型では可換模型の存在は未解決である。
その方程式をrotationの低次数積から抽出する有限手順と、
cubic threefoldの純Poisson方向の一次消滅を示す。
分離型の既知のHessian公式を圏比較に結び付け、
最小混合試験族の三次中心元を全パラメータの閉公式で構成する。
候補の具体化と、残余対角のperfect性を用いる存在証明を並行して進める。
\end{abstract}
\tableofcontents
\section{序論と対象の区別}\label{sec:paperC-introduction}
本稿では、五つ以上の次数一生成元を持つKoszul AS-regular代数
\(A_\hbar\) と三次中心正則元 \(\Omega_\hbar\) の商
\(B_\hbar=A_\hbar/(\Omega_\hbar)\) を扱う。
論文Bでは逆に周囲を \(B\)、中心三次商を \(A\) と書くため、
両論文の記号を混同しない。
本稿の対象は残余dg圏 \(\mathcal K_\hbar\) と
次数シフトから定まる自己同値 \(\rho_\hbar\) の対である。
\[
 (\mathcal K_\hbar,\rho_\hbar)
 \stackrel{?}{\simeq}
 (\operatorname{Ku}(Y_\hbar),\rho_{Y_\hbar}).
\]
Torelli型定理が示す可換模型の一意性と、
模型の存在を論理的に区別する。
一楕円ブロックの良い複素ファイバーでの明示的な存在結果と、
相対族全体の同値の未完成部分を分けて記述する。
背景のPoisson構造・中心量子化は論文A、
一般17次元中心三次対の偏極付き復元は論文Bを参照する。
一般17次元のPoisson中心対から元の可換cubicの変形への忘却は、
分離点では7次元の核を持ち、中心平面の移動6方向とPoissonスカラー1方向に分かれる。
固定した元の可換cubicと平面の下では、スカラーの\(\mathbb G_m\)から
滑らかな楕円平面切断の\(\operatorname{Pic}^0\)へ非定数の一価有理写像は存在しない。
従って、一般のrotation付き可換復元は、元の可換cubicと平面を同時に
固定した素朴なPicardパラメータ化ではなく、可換cubicまたはmarkingの
非自明な変化を含む必要があり得る（論文Bの対応する命題を参照）。
これは形式復元・多価対応・実際の相対復元の不存在を意味しない。
\subsection{方程式の抽出と同値の存在}
可換模型の具体式を作ることを、本稿の存在問題と並行して扱う。
第\ref{sec:explicit-jacobian-recovery}節では、rotationの次数一・二の
導来自然変換の積から二次関係を取り出し、可積分性の線形方程式を解いて
射影同型を除いて唯一の可換候補を求める。低次数データの一致だけでは圏同値は従わない。
第\ref{pf:sec:first-order}節では、cubic threefoldの純Poisson方向が
Kuznetsov成分の一次変形として消えることを、Todd補正を含めて示す。
第\ref{em:sec:comparison}節では、Mori--SmithのHessian公式を
分離型の可換模型へ適用し、中心元の単純な可換化が二次で異なることを示す。
第\ref{mt:sec:family}節では、最小混合試験族の三次中心元と
全Artin切断でのPBW性を証明し、復元先の二次項を求める入力を固定する。
いずれも一般混合型の存在や相対smoothnessを無条件に結論するものではない。

\section{rotationの低次数積からの可換方程式の有限復元}
\label{sec:explicit-jacobian-recovery}

本節では、可換模型の\emph{方程式の抽出}を、その模型との
\emph{dg圏同値の存在}から分離する。
前者には有限次元線形代数による必要十分な判定がある。
これを非可換cubicの残余圏に適用するには、まずその圏について
以下の低次数積を実際に計算しなければならない。
本節は一般の非可換cubicの可換模型の存在を証明するものではない。

先行研究との関係を明確にしておく。
Huybrechts--Rennemoはrotationから得られるHochschild環と
Jacobian環を比較し\cite{HR}、Lin--Zhangはgraded matrix factorizationを用いて
一般のFano超曲面に対するrotation付きTorelliを証明している
\cite[Theorem~3.5, Remark~3.6, Section~5]{LZ}。
勾配空間から多項式を復元する問題はUeda--Yoshinagaの研究
\cite[Lemma~3]{ExplicitUedaYoshinaga}に現れ、
積分の非一意性はHwangのsymmetrizer群
\cite[Theorems~1.5--1.6]{ExplicitHwang}によって説明される。
以下の有限検査器はこれらの既知の構造を本研究に適用するための整理であり、
復元原理そのものの新規性を主張しない。
具体的な検査器および反例の初出は未確認である。

\subsection{計算すべき量は次数一と二の積である}

$\cC$を小さな$\CC$-線形dg圏、$\rho$をdg自己同値とする。
通常の三角圏上の自然変換ではなく、dg双加群による\emph{導来自然変換}を用いる。
対角双加群を$\Delta$、$\rho$の双加群を$P$と書き、
\[
 V:=H^0\RHom_{\cC^e}(\Delta,P),\qquad
 W:=H^0\RHom_{\cC^e}(\Delta,P\otimes^{\mathbf L}_{\cC}P)
 \tag{EJ1}\label{eq:ej-low-degree}
\]
とおく。$\alpha,\beta:\Id\to\rho$の積を
$\rho(\beta)\circ\alpha:\Id\to\rho^2$で定める。
これにより$\mu:V\otimes V\to W$を得る。
一般のpairではこの積の対称性は自動ではない。
従って最初に
\[
 \mu(v,w)-\mu(w,v)=0\quad(v,w\in V)
 \tag{EJ2}\label{eq:ej-commutativity}
\]
を確認し、その後でのみ
$\bar\mu:\operatorname{Sym}^2V\to W$および
$K:=\ker\bar\mu$を定義する。
以下、これらのベクトル空間は有限次元であるとする。

\begin{proposition}[可換cubicにおける低次数の同定]
\label{prop:ej-geometric-low-degree}
$F\in\CC[x_0,\ldots,x_{N-1}]_3$が
滑らかなcubic $X\subset\PP^{N-1}$を定め、$N=n+2\ge5$とする。
$(\cC,\rho)=(\Ku(X),\rho_X)$に対して
\[
 V\simeq\operatorname{Jac}(F)_1=\CC\langle x_0,\ldots,x_{N-1}\rangle,
 \qquad W\simeq\operatorname{Jac}(F)_2
\]
であり、$\mu$はJacobian環の積である。特に積は対称で、
\[
 \dim V=N,\quad \dim W=\binom N2,\quad
 \bar\mu\text{ は全射},\quad
 K=\langle\partial_0F,\ldots,\partial_{N-1}F\rangle
 \subset\operatorname{Sym}^2V.
 \tag{EJ3}\label{eq:ej-jacobian-kernel}
\]
\end{proposition}
\begin{proof}
Orlovのdg比較をrotationも含めて用いると、$\rho_X$はgraded matrix
factorizationの内部次数シフトとなる。
Lin--Zhangの\cite[Theorem~2.2, Theorem~3.5, Remark~3.6, Section~5]{LZ}
にある対角kernelの計算を適用する。
恒等sectorは$\operatorname{Jac}(F)_t$を与える。
三次式の非恒等sectorは全$N$変数を動かすので、
cohomological degree零の$\Hom(\Id,\rho^t)$に寄与できるのは
$N$が偶数で$t=N/2$の場合に限る。
$N$が奇数なら寄与はなく、$N$が偶数でも$N\ge6$なので$t=1,2$では寄与しない。
従って表示した二つの同定を得る。
このsector計算には$\gcd(N,3)=1$を必要としない。
同論文Proposition~3.7の互いに素という条件は、内部シフトを
Serre関手とcohomological shiftで表示してSerre代数へ埋め込むための条件であり、
ここでは$\rho$自体を入力している。

恒等sector上の合成は多項式の積である
\cite[Remark~3.6]{LZ}。
Jacobianイデアルは二次式$\partial_iF$で生成され、次数一部分は零である。
滑らかさにより$\partial_iF$は一次独立である。
実際、定数係数の関係があれば線形座標変換後に$F$は一つの変数に依らず、
その変数軸上の射影点が特異点になる。
従って次数二の核と全射性、および次元公式が従う。
\end{proof}

\begin{remark}[低次数だけを使う理由]
全ての$\Hom(\Id,\rho^t)$をそのままJacobian環と同一視してはならない。
例えばcubic fourfoldでは$t=3$に非恒等sectorが現れる。
命題\ref{prop:ej-geometric-low-degree}は、$n\ge3$に共通する
$t=1,2$だけを取り出すことでこの問題を避けている。
\end{remark}

\subsection{二次関係の積分は有限線形方程式である}

$V$の基底を$x_0,\ldots,x_{N-1}$とし、
$S=\operatorname{Sym}V=\CC[x_0,\ldots,x_{N-1}]$とする。
$N$次元部分空間$K\subset S_2$の基底を$q_0,\ldots,q_{N-1}$と書く。
次の線形部分空間を定義する：
\[
 \mathcal L_K:=\left\{A=(A_{ia})\in\operatorname{Mat}_N(\CC):
 \sum_a A_{ia}\partial_jq_a=\sum_a A_{ja}\partial_iq_a
 \text{ for all }i,j\right\}.
 \tag{EJ4}\label{eq:ej-integrability}
\]
各等式の両辺は一次式なので、係数比較により
高々$N\binom N2$本の斉次線形方程式を$N^2$個の未知数について解くだけでよい。

\begin{theorem}[方程式抽出の必要十分条件]
\label{thm:ej-effective-recovery}
$K\subset S_2$を$N$次元部分空間とする。
次の二条件は同値である。
\begin{enumerate}
 \item 滑らかなcubic $F\in S_3$が存在し、
 $K=\langle\partial_0F,\ldots,\partial_{N-1}F\rangle$である。
 \item $K$の共通零点は$\mathbb A^N$で原点だけであり、
 $\mathcal L_K$は可逆行列$A$を含む。
\end{enumerate}
第二条件が成り立つとき、任意の可逆$A\in\mathcal L_K$について
\[
 \boxed{\quad
 F_A=\frac13\sum_{i,a}x_i A_{ia}q_a,
 \qquad \partial_iF_A=\sum_aA_{ia}q_a
 \quad}
 \tag{EJ5}\label{eq:ej-euler-recovery}
\]
が求める滑らかなcubicを与える。
\end{theorem}
\begin{proof}
第一条件が成り立つなら、$\partial_iF$は$K$の別の基底なので
$\partial_iF=\sum_a A_{ia}q_a$となる可逆行列$A$がある。
混合偏微分の交換性が$A\in\mathcal L_K$を与える。
三次のEuler等式により、偏微分の共通射影零点は
$V(F)$の特異点に一致するから、$K$には射影共通零点がない。

逆に$g_i=\sum_aA_{ia}q_a$と置く。
式\eqref{eq:ej-integrability}と二次のEuler等式から
\[
 \partial_jF_A
 =\frac13\left(g_j+\sum_i x_i\partial_jg_i\right)
 =\frac13\left(g_j+\sum_i x_i\partial_i g_j\right)=g_j.
\]
$A$が可逆なので、勾配は$K$を張る。
その共通射影零点がないことから$V(F_A)$は滑らかである。
\end{proof}

\begin{corollary}[有限の判定手順]
\label{cor:ej-finite-algorithm}
$K$の係数が計算可能な標数零の体にあれば、
代数閉包上で定理\ref{thm:ej-effective-recovery}の条件を有限計算で判定できる。
\end{corollary}
\begin{proof}
まず線形方程式\eqref{eq:ej-integrability}の核の基底
$A_1,\ldots,A_r$を求める。
$\mathcal L_K$に可逆行列がある条件は
\[
 D_K(z_1,\ldots,z_r)=\det\left(\sum_{\nu=1}^r z_\nu A_\nu\right)
 \not\equiv0
 \tag{EJ6}\label{eq:ej-determinant-test}
\]
であり、有限の多項式恒等判定である。
非零なら、例えば$\{0,1,\ldots,N\}^r$の有限格子のいずれかで非零になる。
これは$D_K$の各変数の次数が高々$N$であることと一変数の零点数から帰納的に従う。
その値を使って式\eqref{eq:ej-euler-recovery}を適用すればよい。
共通零点条件は、斉次イデアル$(K)$の商が零次元であることと同値であり、
Gr\"obner基底で判定できる。
この方法は実行可能性を保証する有限手順であって、最適な計算量の主張ではない。
\end{proof}

\subsection{候補は射影同型を除いて一意だが、式は一意とは限らない}

\begin{proposition}[滑らかな候補の射影的一意性]
\label{prop:ej-projective-uniqueness}
同じ$K\subset S_2$を勾配空間に持つ二つの滑らかなcubicは射影同型である。
ただし固定座標での方程式がscalar倍に限られるとは限らない。
\end{proposition}
\begin{proof}
$F,G$をその二つのcubicとする。
それぞれの勾配は$K$の基底なので$\nabla G=A\nabla F$を満たす可逆行列$A$がある。
$F$を、$F(x)=T(x,x,x)$を満たす対称三重線形形式$T$で表す。
$a=A^{\mathsf t}$とすれば、$G$の対称性は
\[
 T(au,v,w)=T(u,av,w)=T(u,v,aw)
 \tag{EJ7}\label{eq:ej-symmetrizer}
\]
を意味する。一般にこの式を満たす線形写像全体を$\mathfrak s_T$とする。

$a,b\in\mathfrak s_T$について、三つの引数の間で写像を移すと
\[
 \begin{aligned}
 T(abu,v,w)&=T(bu,av,w)=T(u,av,bw)\\
 &=T(au,v,bw)=T(au,bv,w)=T(bau,v,w).
 \end{aligned}
\]
$F$が滑らかなので$u\mapsto T(u,-,-)$は単射であり、$ab=ba$が従う。
さらに式\eqref{eq:ej-symmetrizer}から$ab$もsymmetrizerである。
従って$\mathfrak s_T$は単位を持つ可換部分代数である。

この代数に非零nilpotentがあれば、その適当な冪を取って
$0\ne b\in\mathfrak s_T$かつ$b^2=0$とできる。
$u=bv\ne0$を選ぶと
\[
 T(u,u,w)=T(bv,bv,w)=T(v,b^2v,w)=0
 \quad\text{for every }w.
\]
これは$[u]$が$V(F)$の特異点であることを意味し、矛盾する。
Jordan分解のnilpotent部分は元の写像の多項式であるから、
従って$\mathfrak s_T$の全ての元は半単純で、特に$a$は対角化できる。
固有空間分解を$\CC^N=\bigoplus_\lambda U_\lambda$とすると、
式\eqref{eq:ej-symmetrizer}により異なる固有空間を含む混合項は零になる。
よって
\[
 F=\sum_\lambda F_\lambda,
 \qquad G=\sum_\lambda\lambda F_\lambda,
 \qquad F_\lambda\in\operatorname{Sym}^3 U_\lambda^*.
\]
$a$が可逆なので各$\lambda\ne0$である。
各$U_\lambda$の座標を$\lambda^{1/3}$倍すれば$F$を$G$へ移せる。
この線形変換が射影同型を与える。
\end{proof}

\begin{example}[Fermatの場合の積分解]
\label{ex:ej-fermat}
$K=\langle x_0^2,\ldots,x_{N-1}^2\rangle$では
式\eqref{eq:ej-integrability}は$A$が対角行列であることと同値である。
従って$\dim\mathcal L_K=N$であり、
\[
 F_A=\frac13\sum_i a_i x_i^3\qquad(a_i\ne0)
\]
は全て解である。これらは射影同型であるが、scalar倍だけの一意性はない。
特にFermat cubic threefoldでは積分解空間の次元は$5$であり、$1$ではない。
\end{example}

\subsection{正しいHilbert級数を持っても積分できない二次関係}

次の反例は、次数一・二の次元や完全交差性だけでは
可換cubicの候補すら作れないことを示す。
非可換cubicの残余圏からこの例が生じるとは主張しない。

\begin{proposition}[五変数の積分不可能な完全交差]
\label{prop:ej-nonintegrable-ci}
添字を$\ZZ/5\ZZ$で読み、
\[
 q_i=x_i^2+x_{i+1}x_{i+2},\qquad
 K=\langle q_0,\ldots,q_4\rangle
 \subset\CC[x_0,\ldots,x_4]_2
 \tag{EJ8}\label{eq:ej-cycle-example}
\]
とする。このとき以下が成り立つ。
\begin{enumerate}
 \item $K$には共通射影零点がなく、$S/(K)$はArtin完全交差で
 Hilbert級数$(1+t)^5$を持つ。
 \item $\mathcal L_K=0$である。
 特に$S/(K)$はどの滑らかなcubic threefoldのgraded Jacobian環とも同型でない。
\end{enumerate}
\end{proposition}
\begin{proof}
$q_i(x)=0$とする。ある$x_i=0$なら
$q_{i-2}=x_{i-2}^2+x_{i-1}x_i=0$から$x_{i-2}=0$である。
$2$と$5$は互いに素なので全ての座標が零になる。
全ての$x_i$が非零なら五つの式$x_i^2=-x_{i+1}x_{i+2}$の積から
$\prod_i x_i^2=-\prod_i x_i^2$となり矛盾する。
従って共通零点は原点だけである。
五つの$q_i$は各々固有の平方項を持つため一次独立であり、
多項式環のsystem of parametersなので正則列である。
よってHilbert級数は$(1-t^2)^5/(1-t)^5=(1+t)^5$である。

次に$A\in\mathcal L_K$とし、$g_i=\sum_aA_{ia}q_a$とおく。
直接微分すると
\[
 \partial_jg_i=2A_{ij}x_j+A_{i,j-1}x_{j+1}+A_{i,j-2}x_{j-1}.
 \tag{EJ9}\label{eq:ej-cycle-derivative}
\]
$j=i+1$の可積分性で$x_{i+2}$と$x_{i-1}$の係数を比較すると
$A_{ii}=0$および$A_{i+1,i-2}=0$である。
添字を動かせば$A_{i,i+2}=0$も得る。
$j=i+2$で$x_{i+3}$と$x_i$の係数を比較すると
$A_{i,i+1}=0$および$A_{i+2,i}=0$、従って$A_{i,i+3}=0$を得る。
残り得る成分を$a_i=A_{i,i-1}$とすると、
$j=i+1$で$x_i$の係数を比較して$a_i=2a_{i+1}$となる。
五回繰り返すと$a_i=32a_i$なので$a_i=0$である。従って$A=0$。

もしgraded Jacobian環と同型なら、その次数一部分の同型は
多項式環の線形座標変換に持ち上がる。
次数二の核から、この座標変換後の$K$は滑らかな三次式の勾配空間になるはずである。
しかし可積分性は座標変換に不変であり、$\mathcal L_K=0$と矛盾する。
\end{proof}

\subsection{非可換cubicに適用するときの結論と残る課題}

\begin{corollary}[唯一の可換候補を抽出する検査器]
\label{cor:ej-candidate}
非可換cubicのenhanced残余pair $(\cC,\rho)$に対して、
式\eqref{eq:ej-low-degree}の$V,W$および積$\mu$が計算されたとする。
このpairが滑らかな可換cubic $n$-foldと同値なら、$N=n+2$として
次の全てを満たす：
\begin{enumerate}
 \item $\dim V=N$、$\dim W=\binom N2$、積の対称性と全射性。
 \item $K=\ker\bar\mu$は$N$次元で、共通射影零点を持たない。
 \item $\mathcal L_K$には可逆行列が存在する。
\end{enumerate}
いずれかが失敗すれば、このpairの可換cubic模型は存在しない。
全て成功すれば、式\eqref{eq:ej-euler-recovery}は
\emph{そのような模型が存在するときの唯一の候補}を射影同型を除いて与える。
この成功だけから$(\cC,\rho)\simeq(\Ku(V(F_A)),\rho_{F_A})$を結論してはならない。
\end{corollary}
\begin{proof}
必要性は命題\ref{prop:ej-geometric-low-degree}と
定理\ref{thm:ej-effective-recovery}による。
候補の射影的一意性は命題\ref{prop:ej-projective-uniqueness}による。
これらの証明は低次数の導来自然変換の積しか使わず、
pair間の同値双加群やその逆を構成していない。
従って圏同値の存在は独立に証明すべき事項として残る。
\end{proof}

この検査器により、今後の具体的作業は明確になる。
まず第\ref{mt:sec:family}節の最小混合試験族について、
続いて非楕円二成分について残余dg代数とrotation双加群を明示し、
$\RHom_{\Gamma^e}(\Gamma,P)$と
$\RHom_{\Gamma^e}(\Gamma,P\otimes_\Gamma^{\mathbf L}P)$の
次数零コホモロジーおよび合成を求める。
これは線束Homの階数表、Poissonの一次項、ambientの二次関係を
そのまま代入する操作ではない。
得られた$K$を積分して候補方程式を抽出した後、
実際の同値双加群の構成、または必要なsmoothnessを確認した上での
\'etale実現定理との比較を行う。
相対対角perfect性(SM)も、この低次数検査だけでは証明されない。

\paragraph{計算の再現範囲.}
\texttt{computations/verify\_explicit\_jacobian\_recovery.py}は有理数上で、
Fermat三次式、混合項を持つ滑らかな五変数三次式、
および命題\ref{prop:ej-nonintegrable-ci}の完全交差を検査する。
Jacobianの次数二商の積から核を取り直して積分し、
勾配との一致、線形系の階数、共通零点、Hilbert関数を確認する。
反例では非零の最大小行列式も証明書に記録する。
入力は既知の可換多項式または明示的な二次関係であり、
非楕円量子化の残余圏の導来自然変換を計算したものではない。

% Current mathematical statement; 2026-11-10.
\section{cubic threefoldにおけるPoisson方向の一次消滅}
\label{pf:sec:first-order}
本節では、可換復元先の方程式を求める際に必要な最初の制約を示す。
対象は滑らかな\emph{cubic threefold}に限定する。
純粋なPoisson方向はKuznetsov成分の一次変形では消えるため、
可換模型が存在するなら、その方程式の本質的な変化は二次以降に現れる。
ここで証明するのは一次接写像の消滅であって、可換模型の存在、
残余smoothness (SM)、全次数での局所自明性ではない。

用いる外部結果は、Calaque--Rossi--Van den BerghのTodd補正付きHKRと
加群構造の両立性\cite[Theorem 1.4]{PFCRVdB}、および
Jacovskis--Lin--Liu--ZhangのHochschild作用の可換図式と
cubic threefoldに対する単射性\cite[Theorems 3.11 and 4.4(2)]{JLLZER}である。
以下の消滅命題はこれらの既知結果の適用であり、その帰結としての新規性は未確認である。

\subsection{通常のHKRとTodd補正の比較}
$X\subset\PP^4_{\CC}$を滑らかな三次超曲面とし、
\[
 \operatorname{HT}^2(X)
 =H^0(X,\wedge^2T_X)\oplus H^1(X,T_X)\oplus H^2(X,\OO_X)
\]
と書く。通常のHKR同型を
$I_{\mathrm{HKR}}:\operatorname{HT}^2(X)\xrightarrow{\sim}\HH^2(X)$、
加群構造と両立するTodd補正付き同型を$I^K$とする。
純粋なPoisson方向とは、一次変形類が
$I_{\mathrm{HKR}}(\pi)$、$\pi\in H^0(X,\wedge^2T_X)$、
$[\pi,\pi]=0$であることをいう。
同時に動かす可換複素構造の成分をこの条件に含めない。

\begin{lemma}[bivectorによる第一Chern類の縮約]
\label{pf:lem:chern}
滑らかな複素代数多様体$Z$と
$\pi\in H^0(Z,\wedge^2T_Z)$に対し、
\[
 \pi^\sharp c_1(K_Z)=0\quad\text{in }H^1(Z,T_Z).
\]
この主張にPoisson条件$[\pi,\pi]=0$は不要である。
\end{lemma}
\begin{proof}
$K_Z$を自明化するアフィン開被覆$\{U_i\}$と、
各$U_i$で消えない最高次微分形式$\mu_i$を取る。
重なりで$\mu_j=g_{ij}\mu_i$、$g_{ij}\in\OO^\times(U_i\cap U_j)$と書く。
$c_1(K_Z)$は、定数倍と符号の規約を除き、
$\{d\log g_{ij}\}$で表される$H^1(\Omega_Z^1)$の類である。

局所ベクトル場$D_i$を
\[
 \iota_{D_i}\mu_i=d(\iota_\pi\mu_i)
\]
で定める。最高次微分形式との縮約は
$T_{U_i}\simeq\Omega^{\dim Z-1}_{U_i}$という同型なので、$D_i$は正則である。
積の微分により
\begin{align*}
 \iota_{D_j-D_i}\mu_i
 &=g_{ij}^{-1}d\bigl(g_{ij}\iota_\pi\mu_i\bigr)
       -d(\iota_\pi\mu_i)\\
 &=d\log g_{ij}\wedge\iota_\pi\mu_i.
\end{align*}
右辺は、縮約と$\pi^\sharp$の符号を固定すれば
$\pm\iota_{\pi^\sharp(d\log g_{ij})}\mu_i$である。
従って$\{\pi^\sharp(d\log g_{ij})\}$は
$\{\pm D_i\}$の\v{C}ech coboundaryであり、そのコホモロジー類は零となる。
符号の選択は消滅の結論に影響しない。
\end{proof}

\begin{lemma}[cubic threefoldの純bivector上ではTodd補正が消える]
\label{pf:lem:todd}
任意の$\pi\in H^0(X,\wedge^2T_X)$について
\[
 I^K(\pi)=I_{\mathrm{HKR}}(\pi).
\]
\end{lemma}
\begin{proof}
\cite[Theorem 1.4]{PFCRVdB}の規約では
$I^K=I_{\mathrm{HKR}}\circ\iota_{\sqrt{\operatorname{td}(X)}}$である。
\[
 \sqrt{\operatorname{td}(X)}
 =1+\tfrac14c_1(T_X)+\theta_2+\cdots,
 \qquad\theta_2\in H^2(X,\Omega_X^2)
\]
を$\pi$に縮約すると、第一補正は
$\tfrac14\pi^\sharp c_1(T_X)\in H^1(T_X)$であり、
$c_1(T_X)=-c_1(K_X)$と補題\ref{pf:lem:chern}から零である。
第二補正は$H^2(X,\OO_X)=0$に値を取る。
三次以上の形式はbivectorには縮約できない。
よって補正は全て消える。逆Todd類を用いる規約でも同じ理由で結論は変わらない。
\end{proof}

\subsection{残余圏への制限写像の核}
$\mathcal K=\Ku(X)$とし、enhancementの完全忠実な包含から誘導される写像を
\[
 \operatorname{res}_X:\HH^2(X)\longrightarrow\HH^2(\mathcal K)
\]
と書く。

\begin{proposition}[Poisson方向の一次消滅]
\label{pf:prop:vanishing}
滑らかなcubic threefold $X$に対し、
\[
 \operatorname{res}_X\circ I_{\mathrm{HKR}}
 \big|_{H^0(X,\wedge^2T_X)}=0.
 \tag{PV}\label{pf:eq:vanishing}
\]
従って、実在する純Poisson一次変形からenhanced残余変形が誘導され、
そのKodaira--Spencer類が上の制限写像で与えられるなら、
残余変形の$\HH^2(\mathcal K)$における一次類は零である。
\end{proposition}
\begin{proof}
\cite[Theorem 3.11]{JLLZER}によれば、
$\alpha'=\operatorname{res}_X\circ I^K$と
\[
 \gamma:\HH^2(\mathcal K)\longrightarrow
 \operatorname{Hom}_{\CC}
 \bigl(H^1(X,\Omega_X^2),H^2(X,\Omega_X^1)\bigr)
\]
の合成はpolyvectorの縮約作用である。
ここで同論文のHochschild homologyの符号では、
$H\Omega_{-1}(X)=H^1(X,\Omega_X^2)$、
$H\Omega_1(X)=H^2(X,\Omega_X^1)$となる。
これら以外の該当するHodge成分は、cubic threefoldのLefschetz定理と
$H^i(X,\OO_X)=0$（$i>0$）により零である。

$\pi\in H^0(X,\wedge^2T_X)$の縮約は
\[
 H^1(X,\Omega_X^2)\xrightarrow{\ \iota_\pi\ }
 H^1(X,\OO_X)=0
\]
なので$\gamma\alpha'(\pi)=0$である。
cubic threefoldについて$\gamma$は単射である
\cite[Theorem 4.4(2)]{JLLZER}。従って$\alpha'(\pi)=0$を得る。
補題\ref{pf:lem:todd}で$I^K$を通常のHKRに戻せば
\eqref{pf:eq:vanishing}となる。
最後の主張は、仮定されたKodaira--Spencer類の制限との両立性による。
\end{proof}

\subsection{明示的復元式に対する制約}
\begin{corollary}[可換復元先の方程式は一次で動かない]
\label{pf:cor:second-order}
$X=V(F)\subset\PP^4$、$R=\CC[[h]]$とし、
特殊ファイバーが$\Ku(X)$であるenhanced残余族$\mathcal K_R$を考える。
その$R/(h^2)$への還元が、命題\ref{pf:prop:vanishing}の意味で
純Poisson一次変形から誘導されていると仮定する。
さらに滑らかな可換cubic族$Y_R=V(F_h)$とenhanced同値
\[
 \mathcal K_R\simeq\Ku(Y_R/R),\qquad Y_0=X
\]
が\emph{既に存在する}と仮定する。
$F_h=F+hG+O(h^2)$と書けば
\[
 [G]=0\quad\text{in }
 J(F)_3,
 \qquad J(F)=\CC[z_0,\ldots,z_4]/(\partial_0F,\ldots,\partial_4F).
 \tag{T0}\label{pf:eq:tangent-zero}
\]
従って$h=0$で恒等となる線形座標変換と方程式の単元倍により
$F_h=F+O(h^2)$とできる。
同じ結論は一次底$\CC[h]/(h^2)$上だけでも成立する。
\end{corollary}
\begin{proof}
命題\ref{pf:prop:vanishing}より、$\mathcal K_R$の一次類は零である。
enhanced同値はこの類を$\Ku(Y_R/R)$の一次類に送るので、
幾何的Kodaira--Spencer類$\kappa_Y\in H^1(X,T_X)$について
\[
 \eta(\kappa_Y)=0,\qquad
 \eta:H^1(X,T_X)\longrightarrow\HH^2(\Ku(X))
\]
となる。$H^1(T_X)$上のTodd補正も$H^2(\OO_X)=0$で消えるので、
ここでも通常のHKRと補正付きHKRは一致する。$\eta$は同型\cite[Theorem 4.4(2)]{JLLZER}であるから
$\kappa_Y=0$である。ここでは特殊ファイバーの同値が恒等である必要はなく、
Hochschild類の零は任意のenhanced同値で零に移る。

超曲面の法束完全列と射影空間のEuler列により、
$H^1(X,T_X)=J(F)_3$で、$\kappa_Y$は$[G]$に対応する。
従って一次形式$\ell_i$を用いて
$G=\sum_{i=0}^4\ell_i\partial_iF$と書ける。
座標変換$z_i\mapsto z_i-h\ell_i(z)$を代入すると一次項が消える。
方程式のスカラー変化もEuler恒等式
$\sum_i z_i\partial_iF=3F$によって同じJacobian部分に含まれる。
\end{proof}

\begin{remark}[仮定と計算方針]
\label{pf:rem:scope}
量子代数の交換子の半古典極限が$\pi$であることだけでは、
射影化した圏の一次変形が純粋に$I_{\mathrm{HKR}}(\pi)$であるとは断定できない。
可換方程式やgluingを一次で同時に動かすと、$H^1(T_X)$成分が加わり得る。
具体的な量子化族に系\ref{pf:cor:second-order}を適用する際は、
この純粋性とenhanced残余変形への制限との両立性を確認する。

確認後は、可換模型の最初の本質的な係数を
$F_h=F+h^2G_2+h^3G_3+\cdots$の$[G_2]\in J(F)_3$から探せばよい。
本節は$G_2$を計算しておらず、その非零性も、全ての奇数次項の消滅も主張しない。
また$\HH^2$での消滅はunderlying残余圏の一次類についての主張であり、
指定rotationを含む全coherentデータの自明化を結論しない。
存在・SM・高次同値の証明は、それぞれ独立に残る。
\end{remark}

\section{可換復元の定式化と主結果}\label{sec:intro}

滑らかなcubicのKuznetsov成分を
\[
 D^b(X)=\langle\Ku(X),\OO_X,\OO_X(1),\ldots,\OO_X(n-2)\rangle
\]
の規約で取る。rotationは
\[
 \rho_X=\mathbf L_{\OO_X}\circ(-\otimes\OO_X(1))
\]
である。すべての圏はdg enhancementと冪等完備化を保って扱う。
可換cubicについては、\((\Ku(X),\rho_X)\) が \(X\) の射影同型類を決定する
\cite[Theorem 1.7]{LZ}。このTorelli定理は可換モデルの一意性を与えるが、
任意の非可換な圏の組がその像に属することを述べない。
幾何側のenhanced同値は、包含・射影と合成してFourier--Mukai核で表されるため、
同定理のFourier--Mukai条件を満たす。
本稿は、この存在問題の形式近傍、飽和な代数的族のétale近傍、
および明示的な楕円ブロックの族を扱う。

\subsection{形式的な存在と一意性}

\begin{theorem}[次数付き変形の形式的可換モデル]\label{thm:main}
 \(X=V(F)\subset\PP^{n+1}_{\CC}\) を滑らかなcubic、\(n\ge3\) とする。
 \(R\) を剰余体が \(\CC\) の有限次元可換局所Artin代数とし、
 \(B_R=\bigoplus_{d\ge0}(B_R)_d\) を中心的な \(R\)-代数で、
 \((B_R)_0=R\)、各次数部分が有限自由、かつ指定された次数付き同型
 \[
 B_R\otimes_R\CC\simeq\CC[x_1,\ldots,x_{n+2}]/(F)
 \]
 を持つものとする。次を定義する。
 \[
 \cP_R=D(\QGr B_R)^c,\qquad
 \cK_R=\langle\OO_R,\ldots,\OO_R(n-2)\rangle^\perp\subset\cP_R,\qquad
 \rho_R=\mathbf L_{\OO_R}\circ(1).
 \]
 このとき \((\cK_R,\rho_R)\) はsmooth properな \(R\)-linear dg圏と自己同値の組である。
 あるmarkedな滑らかなcubic変形
 \(X_R\subset\PP^{n+1}_R\) と、markingを保つ同値
 \[
 (\cK_R,\rho_R)\simeq(\Ku(X_R),\rho_{X_R})
 \]
 が存在する。\(X_R\) のmarked射影同型類は一意である。
 これらはArtin底の変更と両立する。
\end{theorem}

ここで \(\QGr\) は全次数付き加群をtorsion加群で割ったGrothendieck圏であり、
上付き \(c\) はcompact対象を表す。Artin底上では
\(\cP_R\) を \(D^b(\qgr B_R)\) 全体に置き換えない。
例えば可換な定数族でも、底 \(R\) が非正則なら両者は異なる。
後述する体上の有限パラメータでは、残余圏を通常のgraded matrix factorization圏、
すなわちOrlov同値で \(D^b(\qgr B)\) に埋め込まれる成分として扱う。
特異なファイバーで、この成分と \(D(\QGr B)^c\) の直交成分を混同しない。

\begin{corollary}[形式族]\label{cor:formal}
 \(B_\hbar\) が次数ごとに \(\CC[[\hbar]]\)-flatな、定理の斉次座標環の形式的変形なら、
 \[
 F_\hbar\in\CC[[\hbar]][x_1,\ldots,x_{n+2}]_3,\qquad F_\hbar\bmod\hbar=F
 \]
 が存在し、各 \(\CC[\hbar]/(\hbar^{r+1})\) 上で残余圏の組を
 \(V(F_\hbar)\) の組と同一視できる。これらの同値は互換的に選べる。
 形式的cubicのmarked同型類は一意である。
\end{corollary}

これは形式的な主張である。係数の収束や、指定された代数的パラメータ空間への
降下、任意の非零パラメータでの可換モデルの存在を含まない。
また \(F_\hbar=F\) を意味しない。

\subsection{代数的族における局所的存在}

定理\ref{thm:etale}は、有限型複素底 \(S\) 上のsmooth properな代数的dg圏の族
\(\cK_S\) と、可逆perfect bimoduleで表された自己同値 \(\rho_S\) を入力とする。
ある点 \(s\) で \((\cK_s,\rho_s)\simeq(\Ku(X),\rho_X)\) なら、
\(s\) のétale近傍 \(T\) 上に滑らかなcubic \(Y/T\) と
\[
 (\cK_T,\rho_T)\simeq(\Ku(Y/T),\rho_{Y/T})
\]
が存在する。従って、可換cubicで実現される複素点はZariski開部分をなす。
これはgraded flatな代数の族から、相対smooth properなdg模型を作る課題を
省略するものではない。

また \(n\ne4\) では \(\HH^2(\Ku(X))=J(F)_3\)、\(\HH^1=0\) であり、
rotationを忘れた形式的復元も成立する（定理\ref{thm:unpolarized-formal}）。
同じ条件下では、代数的dg圏の族についてもrotationを最初から指定せずに
étale局所モデルを作れる（系\ref{cor:etale-unpolarized}）。
これらは局所的な実現の結果であり、rotationなしの大域的な一意性とは分けて扱う。

\subsection{有限パラメータを含む明示的族}

\[
 S_h=\CC\langle x,y,z\rangle/
 ([x,y]-hz^2,\ [y,z]-hx^2,\ [z,x]-hy^2)
\]
と置き、\(T=\sym(xyz)\) を六つの順列の平均、
\(G_h=x^3+y^3+z^3-2h^2T\) とする。
中心変数 \(u_1,\ldots,u_m\) と斉次三次式 \(P(u)\) に対して
\[
 B_{h,P}=S_h[u_1,\ldots,u_m]/(G_h+P(u))
\]
を考える。\(P=\sum_i u_i^3\) の場合を \(B_{h,m}\) と略記する。

\begin{theorem}[一ブロックと任意の中心三次式の可換モデル]\label{thm:global}
 \(m\ge2\)、\(n=m+1\)、\(h\in\CC\)、\(h^6\ne-27\) とする。
 \(P\in\CC[u_1,\ldots,u_m]_3\) のアフィン臨界点が原点だけであると仮定する。
 \(B_{h,P}\) の残余dg圏とrotationは、滑らかなcubic
 \[
 Y_{h,P}=V\!\left(x^3+y^3+z^3+h^2xyz+P(u)\right)
 \subset\PP^{m+2}
 \]
 の \((\Ku(Y_{h,P}),\rho_{Y_{h,P}})\) と同値である。
 この \(Y_{h,P}\) は、同じrotation付き残余圏を持つ滑らかな可換cubicとして
 射影同型を除いて一意である。
\end{theorem}

この定理はsixfoldでは \(m=5\)、sevenfoldでは \(m=6\) を含む。
中心変数どうしの任意の三次混合項を許すが、一つのSklyaninブロックと中心部分が
分離されているという仮定を保持する。\(P=\sum u_i^3\) はその特殊例である。
複数の非可換ブロックについても、対応する可換cubicの残余圏からの完全忠実埋込みを得る
（命題\ref{prop:multiblock-embedding}）。次の部分的な有限位数条件の下では、その像が全体を生成する。

\begin{theorem}[一つの任意位数ブロックを許す多ブロック可換モデル]\label{thm:multiblock-model}
 \(r\ge2\)、\(m\ge0\)、\(n=3r+m-2\) とする。
 各 \(j\) について \(h_j^6\ne-27\) とし、楕円曲線
 \[
 E_{h_j}=V(x_j^3+y_j^3+z_j^3+h_j^2x_jy_jz_j),\qquad
 o_j=[1:-1:0],\quad p_j=[1:-1:-h_j]
 \]
 上の点 \(p_j\) は、高々一つを除いて有限位数であると仮定する。
 すなわち因子を並べ替えた後、\(j\ge2\) にのみ有限位数を仮定し、\(p_1\) の位数には条件を課さない。
 \(m>0\) の場合、\(P\in\CC[u_1,\ldots,u_m]_3\) は原点だけに臨界点を持つとし、
 \(m=0\) の場合は中心部分を省く。
 異なるSklyaninブロックの元を可換とした代数
 \[
 B_{\mathbf h,P}=
 (S_{h_1}\otimes\cdots\otimes S_{h_r})[u_1,\ldots,u_m]
 \big/\left(\sum_jG_{h_j}+P\right)
 \]
 のrotation付き残余dg圏は、滑らかなcubic
 \[
 Y_{\mathbf h,P}=V\!\left(
 \sum_{j=1}^r(x_j^3+y_j^3+z_j^3+h_j^2x_jy_jz_j)+P(u)
 \right)\subset\PP^{3r+m-1}
 \]
 の \((\Ku(Y_{\mathbf h,P}),\rho_{Y_{\mathbf h,P}})\) と同値である。
 この可換モデルは射影同型を除いて一意である。
\end{theorem}

ここで「有限パラメータ」は形式変数でなく実際の複素数であることを意味し、
楕円曲線上の点が「有限位数」であることとは異なる。
定理は一つの無限位数ブロックを許す。
具体的に \(h_1=2\) は無限位数（補題\ref{lem:order-infinite}）、
\(h_j^6=5\ (j\ge2)\) は位数五（補題\ref{lem:order-five}）となる。
従って二ブロックのsixfold、三ブロックのsevenfoldにも、真に無限位数の因子を含む例が得られる。
二つ以上の無限位数ブロックを任意に許す主張は、本定理には含まれない。
証明は第\ref{sec:finite-center}節に置く。

\begin{proposition}[例外パラメータのhomological singularity]\label{prop:bad-intro}
 \(m\ge0\) とする。Fermat中心部分の族 \(B_{h,m}\) で \(h^6=-27\) とすると、\(S_h[u]\) はAS-regularで
 \(G_h+\sum u_i^3\) は中心正則元であるが、そのgraded matrix factorization圏は
 properではなく、さらに
 \[
 \operatorname{hdim}(\qgr B_{h,m})=\infty
 \]
 である。ここでhomological dimensionは、非零のExtが存在する次数の上限である。
 従って、滑らかな射影cubicのKuznetsov成分とは同値にならない。
\end{proposition}

これはAS-regularな周囲、三次の中心正則元、期待されるHilbert series、
および \(\rho^3\simeq[2]\) だけでは、大域的な可換モデルの存在を保証しないことを示す。
本稿の可換点からのPoisson量子化族に限定した、良い有限パラメータ全体での存在問題は別に残る。
一般のAS-regular中心三次商については、smooth properかつCalabi--Yauな例にも
定理\ref{thm:zhang-cy-fourfold-counterexample}の反例がある。

\subsection{先行研究と本稿の位置}

Orlov \cite{Orlov} とCassidy--Conner--Kirkman--Moore \cite{CCKM} は、
中心正則元による商、特異圏、matrix factorization、残余成分の比較を与える。
体上のwindowのdg模型にはCăldăraru--Tu \cite[Theorem 5.8]{CT} を用いる。
本稿の形式的な構成では、これらの体上の同値をArtin底上へ無条件に延長せず、
De Deken--Lowen \cite{DDL} のabelian deformationとcompact generationを用いる。
smoothnessとpropernessの持ち上げはnilpotent descent
\cite{Mathew}による。

残余圏のHochschild cohomologyはBallard--Favero--Katzarkov \cite{BFK}、
幾何学的Kodaira--Spencer写像はHuybrechts--Rennemo \cite{HR} と、
三次元の場合にはJacovskis--Lin--Liu--Zhang \cite{JLLZ} を用いる。
四次元の余分な二方向に対するrotationの作用は
Huybrechts \cite[Proposition 3.12]{Huy} から確認する。
actual dg変形をHochschild形式変形と同一視する際のcurvature問題
\cite{KL}は、Lurie \cite{Lurie} とBlanc--Katzarkov--Pandit
\cite{BKP} の比較写像の単射性で処理する。

明示的族の楕円ブロックはEtingof--Ginzburg \cite{EG} と
Sklyaninの点曲線の公式 \cite{Walton} に依拠する。
本稿では次数シフトの共役、中心変数を加える際の生成性、
有限群による通常次数への降下を証明する。
可換な係数環についてのThom--Sebastiani定理を、
非可換な係数環に対して仮定を確認せず引用することはしない。

有限中心代数の外部積生成については、Elagin--Lunts--Schn\"urer
\cite[Theorem 4.18]{ELS} が、perfect field上のnoetherian J-2 schemesに載る
coherentな代数に対し、ungradedな有界導来圏の生成定理を証明している。
本稿では、\(\ZZ^r\) の任意の粗化次数を保ち、さらに一方の係数代数には
有限中心性を仮定しない形を証明する（定理\ref{thm:one-sided-boxes}）。
\cite[Theorem 4.18]{ELS} は双方の因子をcoherentな構造代数とするため、
任意のnoetherian係数代数を許すここでの主張は、その定理の直接の適用ではない。
新たな入力はBeilinsonの対角分解 \cite{Beilinson} を非可換係数付きで用いる
多項式拡大の補題である。中心のNoether正規化と分離冪等元によるretractを組み合わせ、
斉次中心局所化の核を次元について帰納する。
対象の細分次数へのliftや、係数代数の有限global dimensionは仮定しない。
有限位数THCRの有限中心性 \cite[Corollary 2.3]{SmithTate} をこの片側のみに適用し、
一つの任意位数ブロックを含む多ブロックのMF比較を得る。
He--Ma--Ye \cite{HMY} のtensor productに対する一般化Kn\"orrer周期性は、
中心元の次数が二である非可換quadricとClifford変形を扱う。
本稿の三次中心元への直接の適用ではない。
周囲のnoetherian性については、Chan \cite{ChanMulti} の多重Serre同値と
面・帯による議論を任意個の楕円因子へ展開し、有限位数の仮定なしで証明する。
Keeler \cite{KeelerMulti} のTheorem 2.7は任意個の可換な可逆bimoduleの
NC-amplenessを判定するが、同論文のTheorem 3.3とCorollary 3.5のnoetherian性は
二因子についての主張である。本稿では任意個の積に必要な帰納法を記す。
基礎となる対角分解と有限中心のungraded生成定理は既知である。
一方有限中心の次数付き定式化および本適用の先行研究上の新規性は確認継続中である。

étale局所復元では、Toën--Vaquié \cite{TV} のperfect objectのmoduliと
smoothnessの無限小判定 \cite{Stacks} を用いる。
新たに示すのは、同値bimoduleとrotation intertwinerを同時にパラメータ化したスタックが、
形式的復元定理によって底上smoothになることである。
Pirozhkov \cite{Pirozhkov} およびLin--Zhang \cite{LZ} のTorelliは
二つの可換モデルの一意性に関する結果であり、この局所的な存在の証明とは役割が異なる。
Huybrechts \cite[Section 9]{HuyUpdate} の四次元での一般的な非可換変形の問題提起と、
本稿のrotation付きの変形問題との接写像の詳細な比較は研究ノートに残す。

\paragraph{主張の状態.}
本稿の定理には以下で証明を与える。これらを一つの一般次元の復元定理として
整理したことの先行研究上の新規性は未確認であり、
既知の基礎定理自体を新しい成果として扱わない。
本稿で構成する可換点からのPoisson量子化族について、
良い有限パラメータ全体での存在は未解決として残す。

\section{次数付き変形から相対残余圏を作る}\label{sec:relative}

\subsection{nilpotentな底変更によるcompactnessの検出}

\begin{lemma}\label{lem:descent}
 \(R\) を上の局所Artin代数、\(k=R/\fm\) とする。
 \(A\) を \(R\)-linear dg代数、\(M\) を \(A\)-加群とし、
 \(A_0=A\otL_R k\)、\(M_0=M\otL_R k\) と書く。
 \(M_0\) が \(A_0\)-compactなら \(M\) は \(A\)-compactである。
 また導来底変更 \(-\otL_R k\) は保守的である。
\end{lemma}
\begin{proof}
 \(\fm\) の冪による有限濾過から
 \(R\in\thick_{D(R)}(k)\) である。
 \(A\)-加群の任意の族 \((N_\alpha)\) に対し、
 \[
 \bigoplus_\alpha\RHom_A(M,P\otL_RN_\alpha)
 \longrightarrow
 \RHom_A\!\left(M,P\otL_R\bigoplus_\alpha N_\alpha\right)
 \]
 が同値となる \(P\in D(R)\) の全体はthickである。
 \(P=k\) では導来随伴と \(M_0\) のcompactnessによって同値となるので、
 \(P=R\) でも同値となる。保守性も、
 \(P\otL_RM=0\) となる \(P\) のthick部分圏を使って同様に従う。
 これは \cite[Propositions 3.28, 3.35; Remark 3.29]{Mathew}
 の特別な場合である。
\end{proof}

\begin{remark}
 この補題は \(k\) が \(R\)-perfectであるとは仮定していない。
 また普通の商 \(A/\fm A\) だけを調べる命題ではなく、
 必ず導来ファイバー \(A\otL_Rk\) を使う。
\end{remark}

\subsection{abelian deformationと明示的な生成子}

\begin{proposition}\label{prop:qgr}
 定理\ref{thm:main}の \(B_R\) について、
 \(\QGr B_R\) は \(\operatorname{Qcoh}(X)\) のflat abelian deformationである。
 \[
 E_R=\bigoplus_{i=0}^{n+1}\OO_R(i)
 \]
 は \(D(\QGr B_R)\) のcompact generatorであり、
 \[
 \Lambda_R=\REnd(E_R),\qquad
 \Lambda_R\otL_Rk\simeq\REnd_X(E_0)
 \tag{\(\ast\)}\label{eq:derived-basechange}
 \]
 がdg代数として成り立つ。
\end{proposition}
\begin{proof}
 まず \(B_R\) は左右noetherianである。
 \(\fm^aB_R/\fm^{a+1}B_R\) は、
 flatnessにより \(B_0\) の有限個の直和である。
 有限濾過と \(B_0\) のnoetherian性を使えばよい。
 次数付き生成元を持ち上げる同じ議論から、
 \(B_R\) は有限生成で、付随する \(\ZZ\)-代数のtailsも通常のtorsion商で記述できる。

 \(X\) の射影正規性により \(B_0\) の \(\ZZ\)-代数は
 \((\OO_X(i))_{i\in\ZZ}\) のdirected Hom-categoryである。
 この列は \(\ZZ\)-generatingであり、\(i\le j\) に対して
 \[
 \Ext_X^{1,2}(\OO_X(i),\OO_X(j))
 =H^{1,2}(X,\OO_X(j-i))=0.
 \]
 hypersurfaceの完全列と射影空間のcohomologyからこの消滅が従う。
 よって \cite[Theorem 4.5]{DDL} を適用できる。
 \(\OO_R(i)\) は明示的なflat liftである。
 実際 \(B_R(i)\) との有限 \(R\)-加群のテンソル積は完全で、
 \(\operatorname{Gr}B_R\to\QGr B_R\) は完全である。
 これらの対象はabelian圏も生成するので、必要なflat resolutionを取れる。

 閉埋め込み \(j:X\hookrightarrow\PP^{n+1}\) を考える。
 \(E_0=j^*(\bigoplus_{i=0}^{n+1}\OO_{\PP^{n+1}}(i))\) はcompactである。
 \(\RHom_X(E_0,M)=0\) なら、随伴とBeilinsonの生成定理
 \cite{Beilinson}から \(Rj_*M=0\)、従って \(M=0\) となる。
 したがって \(E_0\) はcompact generatorである。

 \(E_R\) は有界flat liftなので、
 \cite[Propositions 5.8, 5.9]{DDL} によりcompact generatorとなる。
 同論文Proposition 5.5と導来随伴は
 \[
 k\otL_R\RHom(E_R,E_R)
 \simeq\RHom_X(k\otL_RE_R,k\otL_RE_R)
 \]
 を与える。同論文Theorem 5.13の証明でも同じ底変更が使われている。
 dg enhancement上の導来還元関手からこの比較を作るため、
 比較は合成と単位を保つ。必要ならhomotopically \(R\)-flatなdg模型を選ぶ。
 これで \eqref{eq:derived-basechange} は複体だけでなくdg代数として成立する。
\end{proof}

\begin{corollary}\label{cor:smooth-proper}
 \(\cP_R=D(\QGr B_R)^c\) はsmooth properであり、
 \(\cP_R\otL_Rk\simeq\Perf(X)\) である。
\end{corollary}
\begin{proof}
 導来Morita理論 \cite{Keller} によって
 \(\cP_R\simeq\Perf(\Lambda_R)\) である。
 \(X\) のpropernessから \(\Lambda_0\) は \(k\)-perfectなので、
 補題\ref{lem:descent}で \(\Lambda_R\) は \(R\)-perfectとなる。
 次に
 \[
 A=\Lambda_R^{\mathrm{op}}\otL_R\Lambda_R,\qquad M=\Lambda_R
 \]
 とする。この導来ファイバーは \(\Lambda_0\) のenveloping algebraと対角加群である。
 \(X\) のsmoothnessにより対角加群はその上でperfectであり、
 同じ補題から \(\Lambda_R\) の対角もperfectとなる。
\end{proof}

\subsection{例外列、mutation、rotation}

\begin{proposition}\label{prop:relative-pair}
 \(\OO_R,\ldots,\OO_R(n-2)\) は \(\cP_R\) 内のrelative exceptional collectionである。
 その右直交 \(\cK_R\) はsmooth properで、導来特殊ファイバーは \(\Ku(X)\) である。
 \(\rho_R=\mathbf L_{\OO_R}\circ(1)\) は \(\cK_R\) の自己同値となる。
 これらの構成はArtin底変更と両立する。
\end{proposition}
\begin{proof}
 Hom複体の導来底変更と補題\ref{lem:descent}から
 \[
 \RHom(\OO_R(i),\OO_R(i))\simeq R,\qquad
 \RHom(\OO_R(j),\OO_R(i))=0\quad(j>i,\ 0\le i,j\le n-2)
 \]
 である。自己Homについては単位写像のconeを還元する。
 propernessにより評価・余評価はperfect対象を保つため、各例外成分はadmissibleである。
 評価のconeによるmutationを順に取って
 \[
 p_R=\mathbf L_{\OO_R}\cdots\mathbf L_{\OO_R(n-2)}
 \]
 と置くと、\(G_R=p_R(E_R)\) は \(\cK_R\) を生成する。
 mutationの導来底変更は \cite[Lemma 5.14]{DDL} による。
 \(G_R\) の導来特殊ファイバーは \(\Ku(X)\) の生成子なので、
 その自己準同型dg代数へ補題\ref{lem:descent}を適用すればsmooth proper性も従う。

 \(T=(1)\) は次数付き射影圏の自己同値で、\(T\OO_R(i)=\OO_R(i+1)\) である。
 \(M\in\cK_R\)、\(1\le i\le n-2\) に対して
 \[
 \RHom(\OO_R(i),TM)=\RHom(\OO_R(i-1),M)=0.
 \]
 \(\OO_R\) によるmutationはこの消滅を保ち、\(i=0\) の直交性も与える。
 よって \(\rho_R\) は \(\cK_R\) を保つ。
 その特殊ファイバーは通常の \(\rho_X\) である。
 Homの比較写像は還元後に同値なので、nilpotentな保守性から
 \(\rho_R\) は完全忠実である。

 \(\rho_R(G_R)\) の特殊ファイバーは生成子である。
 これに直交するInd対象を導来還元すると零になり、
 同じ保守性で元の対象も零になる。従って \(\rho_R(G_R)\) もcompact generatorである。
 冪等完備な圏からの完全忠実完全関手の像はthickなので、\(\rho_R\) は本質的全射でもある。

 底変更については、次数付き導来テンソル関手による生成子の自己Homの自然な比較を使う。
 剰余体上でこの比較は恒等写像となるため、変更後のArtin底上で同値となる。
 評価写像とそのconeも底変更と両立するので、残余圏とrotationまで両立する。
\end{proof}

\section{rotation付きの形式的復元}\label{sec:formal}

\subsection{正確な変形問題}

\begin{definition}
 \((\Ku(X),\rho_X)\) のmarked pair deformationとは、
 smooth properな \(R\)-linear dg圏 \(\cK_R\)、invertible bimoduleで表される
 自己同値 \(\rho_R\)、および特殊ファイバーでのenhanced pairの指定された同定からなる。
 同値にはこのmarkingとの両立を要求する。
 \(\rho_R^3\simeq[2]\) の指定されたcoherenceは定義に含めない。
\end{definition}

\begin{theorem}[pairの形式的復元]\label{thm:pair}
 \(X\) を滑らかなcubic \(n\)-fold、\(n\ge3\) とする。
 marked cubic deformationからmarked pair deformationへの写像は、
 各古典的局所Artin \(\CC\)-代数上で同型類の全単射を与える。
 互換的な形式族についても同型類の一意性と互換的な実現が成り立つ。
\end{theorem}

この主張はfull higher deformation groupoidの同値ではない。
pairには自然変換のscalar automorphismがあり、実現する圏同値も一意とは限らない。

\subsection{actual変形とHochschild形式変形}

\(\cC=\Ind(\Ku(X))\) とする。
actual compactly generated deformationの空間と、その形式moduli problemの間には
\[
 \theta_{\cC}:\operatorname{CatDef}_{\cC}^{c}
 \longrightarrow\operatorname{CatDef}_{\cC}^{\wedge}
\]
がある。右辺の接複体は \(CC^\bullet(\Ku(X))[2]\) である。
actual Morita変形が任意のHochschild classを実現するとは一般にはいえない
\cite{KL}。本稿ではその全射性を使わない。

\begin{lemma}\label{lem:theta}
 上の \(\theta_{\cC}\) は \((-1)\)-truncatedである。
 従って、右辺で同値となった二つのactual変形の同値はactual側へ戻せる。
 この性質はrotationによる \(\ZZ\)-homotopy fixed pointを取った後にも成り立つ。
\end{lemma}
\begin{proof}
 \(\Ku(X)\) はproperでcompact generatorを持つためtameである。
 \cite[Proposition 5.3.21, Theorem 5.1.9]{Lurie} を適用する。
 \cite[Definition 3.3, Section 3.3]{BKP} にも、
 \(\pi_0\) で単射、\(\pi_i\), \(i\ge1\), で同型という形でこの帰結が明記されている。
 比較写像は圏同値に自然であり、markingを移すrotationの作用と両立する。
 空間のmonomorphismはhomotopy limitで保たれる。
\end{proof}

一つのenhanced autoequivalenceを持ち上げる問題は、
markingに作用する \(\ZZ\) のhomotopy fixed pointである。
\(\ZZ\) は一つの可逆射で自由生成されるので、追加の有限群関係を課していない。
\[
 P=\left(\operatorname{CatDef}_{\cC}^{\wedge}\right)^{h\ZZ}
\]
と置くと、これはformal moduli problemであり、
\[
 T_P=R\Gamma(\ZZ,CC^\bullet(\Ku(X))[2]).
\]
\(\ZZ\) の自明加群の長さ一の自由分解から、次の完全列を得る。
\[
 0\longrightarrow
 \operatorname{coker}(1-\rho:\HH^1\to\HH^1)
 \longrightarrow H^0(T_P)
 \longrightarrow
 \ker(1-\rho:\HH^2\to\HH^2)
 \longrightarrow0.                                      \tag{1}\label{eq:tangent}
\]

\subsection{Hochschild計算とKodaira--Spencer写像}

\(N=n+2\) および
\[
 J(F)=\CC[x_1,\ldots,x_N]/(\partial_1F,\ldots,\partial_NF)
\]
を用いる。

\begin{lemma}\label{lem:HH}
 \(n\ge3\) のとき
 \[
 \HH^{<0}(\Ku(X))=0,\quad \HH^0(\Ku(X))=\CC,\quad
 \HH^1(\Ku(X))=0,
 \]
 \[
 \HH^2(\Ku(X))=
 \begin{cases}
 J(F)_3& n\ne4,\\
 J(F)_3\oplus\CC^2& n=4.
 \end{cases}
 \]
 幾何学的Kodaira--Spencer写像は同型
 \[
 H^1(X,T_X)\xrightarrow{\sim}
 \HH^2(\Ku(X))^{\rho_X}=J(F)_3
 \]
 を与える。その次元は \(\binom{n+2}{3}\) である。
\end{lemma}
\begin{proof}
 \cite[Theorem 5.39]{BFK} のKoszul cohomologyによる公式を
 三次のequivariant matrix factorizationへ適用する。
 恒等sectorは偶数次数 \(2r\) に \(J(F)_{3r}\) を与える。
 非自明な \(g\in\mu_3\) では固定点集合は原点だけであり、
 conormalのrankは \(N\)、determinant characterは \(-N\) である。
 cohomological degree \(q\) で残るweightは
 \[
 N+3(q-N)/2=(3q-N)/2.
 \]
 従って二つの非自明sectorは、\(3\mid N\) のときに限り、
 \(q=N/3\) に各一次元を与える。これで列記した群が得られる。

 BFKの恒等sectorのKoszul模型では、\(G\) の乗法が \([G]\) を表す。
 \(F+\varepsilon G\) の実際の変形類との一致は次の比較による。
 \(n\ge4\) では、periodicityから作るkernel morphismと
 Atiyah--Kodaira--Spencer morphismの一致を
 \cite[Lemma 2.19]{HR} とその前後の構成で確認できる。
 ここで三次に必要な次数条件 \(3\le(n+2)/2\) は満たされる。
 \(n=3\) では \cite[Theorem 4.4(2)]{JLLZ} が必要な接写像の同型を与える。
 normal sequenceとEuler sequenceは \(H^1(T_X)=J(F)_3\) を与える。

 幾何学的変形はrotationを持つので、この部分空間は不変である。
 \(n\ne4\) ではそれが \(\HH^2\) 全体である。
 \(n=4\) では \cite[Proposition 3.12]{Huy} により、
 rotationはMukai格子の \(A_2^\perp\) を固定し、\(A_2\) の三つの根を巡回置換する。
 invariantな2-Calabi--Yau volumeによる
 \(\HH^2\simeq\HH_0\) の下で、余分な二方向の固有値は
 \(\zeta,\zeta^2\) となり、不変部分はない。
 最後に偏微分は二次の正則列であるから
 \(H_{J(F)}(t)=(1+t)^N\) であり、次数三の次元は \(\binom N3\) である。
\end{proof}

\subsection{持ち上げの差と自己同型}

\begin{lemma}[Artin帰納法]\label{lem:induction}
 \(D\) をsmoothかつinfinitesimally homogeneousなmarked変形関手で、
 古典的Artin底上ではset-valuedとし、
 \(P\) をformal moduli problemとする。
 自然な写像 \(f:D\to P\) が接空間で同型を誘導すると仮定する。
 さらに、各small extension \(R'\to R\) と \(x'\in D(R')\) に対し、
 \(f(x'\vert_R)\) の任意の自己同型類が \(f(x')\) 上へ持ち上がると仮定する。
 このとき \(D(R)\to\pi_0P(R)\) は各古典的Artin底上で全単射である。
\end{lemma}
\begin{proof}
 Artin長について帰納法を行う。核 \(I\) のsmall extension上に
 \(y'\in P(R')\) があり、その還元を \(f(x)\) と同一視できたとする。
 \(D\) のsmoothnessにより \(x\) を \(x'\) に持ち上げる。
 固定した還元の同定を持つ \(f(x')\) と \(y'\) は、
 \(P(R')\to P(R)\) の同じhomotopy fibreに属する。
 このfibreが空でなければ、その連結成分は
 \(H^0(T_P)\otimes I\) のtorsorである。
 接写像の全射性によって \(x'\) を修正すれば差を消せる。

 単射性では、\(f(x')\simeq f(z')\) の還元を考える。
 帰納法により幾何学的還元は一致するが、
 圏同値の還元はその像の自己同型となり得る。
 仮定によってこの自己同型を上へ持ち上げ、その逆を合成する。
 すると還元を恒等に固定したfibreで二つのliftを比較できる。
 差の接ベクトルの像は零であり、接写像の単射性で元の差も零となる。
\end{proof}

\begin{lemma}\label{lem:automorphisms}
 cubicから得たpairのmarked deformationにおいて、
 small extensionに沿う上の自己同型持ち上げ条件が成り立つ。
\end{lemma}
\begin{proof}
 smooth properな圏ではHochschild複体はperfectで、その導来底変更が成り立つ。
 補題\ref{lem:HH}と局所Artin環上のminimal perfect complexを使うと
 \[
 \HH^0(\cK_R/R)=R,\qquad \HH^1(\cK_R/R)=0
 \]
 である。実際特殊ファイバーの負次数と次数一が零なので、
 minimal complexにはそれらの次数の項がなく、次数零は単位で生成される。

 恒等関手に還元されるinvertible bimoduleは相対対角と同型である。
 perfect bimoduleのliftの差はsmall extensionの核 \(I\) に対して
 \(\Ext^1(\Delta,\Delta)\otimes I=\HH^1(\Ku(X))\otimes I=0\)
 で測られ、対角自身がliftを与えるためである。
 この同型は各段階で指定した特殊ファイバーのmarkingを保つように選ぶ。
 従ってpairの自己同型の基礎となる圏同値は恒等と同型であり、
 残るのは \(\rho_R\Rightarrow\rho_R\) のscalar unitである。
 markingを固定すればこのscalarは \(1+\fm_R\) に入り、必ず \(R'\) へ持ち上がる。
 scalarな二次自己同型も持ち上がる。
 formal moduli problem側の自己同型をactual bimoduleで扱えることは、
 補題\ref{lem:theta}の \(\pi_1\) の同型による。
\end{proof}

\begin{proof}[定理\ref{thm:pair}の証明]
 cubicのmarked変形関手はsmoothである。方程式の係数を持ち上げればよく、
 \(H^1(\OO_X)=H^2(\OO_X)=0\) によって \(\OO_X(1)\) は一意に延長する。
 infinitesimalな射影自己同型もない。実際、二次の偏微分の正則列には
 一次係数の非自明なsyzygyがなく、Euler scalarを除けば \(H^0(T_X)=0\) である。
 これはmarkedな幾何学的変形をset-valuedとして扱うことを正当化する。

 式\eqref{eq:tangent}と補題\ref{lem:HH}により接写像は同型である。
 補題\ref{lem:automorphisms}と補題\ref{lem:induction}を適用し、
 最後に補題\ref{lem:theta}でactual pairの同値へ戻す。
 帰納法では各段階で指定した還元の同値を保てるので、
 \(\CC[\hbar]/(\hbar^{r+1})\) に沿った実現も互換的に選べる。
\end{proof}

\begin{proof}[定理\ref{thm:main}と系\ref{cor:formal}の証明]
 命題\ref{prop:relative-pair}の組へ定理\ref{thm:pair}を適用する。
 互換的な幾何学的liftの方程式の係数の逆極限が \(F_\hbar\) を与える。
 この最後の操作は形式族を作るものであり、非零複素数で評価する操作ではない。
\end{proof}

\begin{corollary}\label{cor:period}
 上で構成したArtin pairについて、指定を伴わない同型
 \(\rho_R^3\simeq[2]\) が存在する。
\end{corollary}
\begin{proof}
 \(\rho_R^3[-2]\) は恒等に還元されるinvertible bimoduleである。
 補題\ref{lem:automorphisms}で使った \(\HH^1\)-rigidityを適用する。
 または可換モデルからこの同型を移してもよい。
 特定のmatrix factorization表示が与えるperiodicity同型とのcoherenceの一致は、
 この系の主張に含めない。
\end{proof}

\begin{remark}
 sevenfoldでは \(\HH^3(\Ku(X))\simeq\CC^2\) があるが、
 上の帰納法は \(\HH^3=0\) を要求しない。
 比較している二つのliftは既に存在し、その差を接空間で修正する。
 higher operationsを無視したformalityの仮定も使っていない。
\end{remark}

\subsection{四次元以外ではrotationを忘れた形式復元}\label{sec:unpolarized-formal}

四次元以外では、形式近傍に限れば、復元の入力からrotationを省くことができる。
ここでいうmarked dg変形は、特殊ファイバーを \(\Ku(X)\) と同定した
smooth properな \(R\)-linear dg圏であり、同値にもこのmarkingとの両立を要求する。

\begin{theorem}[rotationを忘れた形式復元]\label{thm:unpolarized-formal}
 \(X\subset\PP^{n+1}_{\CC}\) を滑らかなcubicとし、\(n\ge3\)、\(n\ne4\) とする。
 任意の古典的局所Artin \(\CC\)-代数 \(R\) について、
 marked cubic変形から \(\Ku(X)\) のmarked smooth proper dg変形への写像は、
 同型類の全単射を与える。従って、任意のそのような \(\cK_R\) は
 一意なmarked cubic変形 \(X_R\) の \(\Ku(X_R)\) と同値である。

 また、固定した \(\cK_R\) 上で、\(\rho_X\) を持ち上げるinvertible bimodule
 \(\rho_R\) が存在し、そのmarked同型類は一意である。
 従ってrotationを忘れる写像
 \[
  \{\text{marked smooth proper pair変形}\}/\simeq
  \ \longrightarrow\
  \{\text{marked smooth proper dg変形}\}/\simeq
 \]
 も全単射である。これらの主張はArtin底変更と互換的な形式族についても成り立つ。
\end{theorem}
\begin{proof}
 \(Q=\operatorname{CatDef}_{\Ind\Ku(X)}^{\wedge}\) と置く。
 補題\ref{lem:HH}の幾何学的Kodaira--Spencer写像は、\(n\ne4\) では
 \[
  H^1(X,T_X)\xrightarrow{\sim}
  \HH^2(\Ku(X))=H^0(T_Q)
 \]
 である。cubicのmarked変形関手はsmoothかつset-valuedである。
 補題\ref{lem:induction}を適用するため、幾何学的な変形の像の自己同型を調べる。

 \(\Ku(X_R)\) の恒等に還元される自己同値を表すinvertible bimoduleは、
 \(\HH^1(\Ku(X))=0\) により相対対角bimoduleと同型である。
 これは補題\ref{lem:automorphisms}と同じsmall extensionの帰納法で従う。
 特殊ファイバー上の同型を指定しても、そこでのscalarは \(R^\times\) へ持ち上がるので、
 同型をmarkingと両立するように選べる。
 従ってmarkedな自己同値の同型類は恒等だけであり、
 必要な自己同型持ち上げ条件は満たされる。
 \(Q\) 側の自己同型についても、補題\ref{lem:theta}の比較写像が
 \(\pi_1\) 上の同型を誘導するため同じことが成り立つ。

 補題\ref{lem:induction}から、cubic変形と \(\pi_0Q(R)\) の全単射が得られる。
 幾何学的変形の像はactual smooth proper変形であり、
 補題\ref{lem:theta}によりactual変形間の同値を検出できる。
 よって最初の主張が従う。この議論はHochschild形式変形の一般的なactual実現性を
 仮定せず、この場合には幾何学的なliftで全てを実現している。

 得られた同値 \(\Ku(X_R)\simeq\cK_R\) を通じて
 \(\rho_{X_R}\) を移せば、\(\rho_X\) のliftが存在する。
 二つのmarkedなlift \(\rho_R,\rho'_R\) があれば、
 \(\rho'_R\rho_R^{-1}\) は恒等に還元されるinvertible bimoduleである。
 再び \(\HH^1=0\) による対角のrigidityを用いると、
 \(\rho'_R\rho_R^{-1}\simeq\mathrm{id}\) となり、markingを保つ同型も選べる。
 これがliftの同型類の一意性とforgetful mapの単射性を与える。
 各small extensionで還元上の同値を保って構成できるため、
 形式族での互換性も従う。
\end{proof}

\begin{remark}\label{rem:unpolarized-scope}
 この定理の一意性は同型類についてのものである。
 rotationとのintertwinerにはscalar自由度があり、
 rotationを忘れる写像がfull higher deformation groupoidの同値だとは主張しない。
 実際、\(\HH^0=\CC\)、\(\HH^1=0\) とrotationのscalar上の自明作用から、
 pairの接複体には \(H^{-1}(T_P)\simeq\CC\) がある一方、
 \(H^{-1}(T_Q)=0\) である。

 \(n=4\) では幾何学的な20方向に加えて \(\HH^2\) に二方向があり、
 rotationの固定部分を取って初めて接写像が同型となる。
 この次元でrotationを省く結論を上の証明から導くことはできない。
 ここでは、Hochschildの全22方向が各Artin底上でactual変形として
 実現されるとは主張していない。

 \(n=7\) では \(\HH^3\simeq\CC^2\) が残るが、
 幾何学的なliftと既に存在するliftを比較する帰納法により結論は成り立つ。
 障害を収め得るcohomology群が非零であることと、
 この変形問題で実際に障害が現れることは同じではない。
 一般の有限パラメータの圏からrotationを省けるという大域的主張も含めない。
\end{remark}

\section{飽和な代数的族に対するétale局所復元}\label{sec:algebraization}

形式的な復元を有限パラメータへ移す際には、圏の族そのものの代数性が必要である。
本節ではそれを明示的に仮定し、可換モデルと圏同値の両方をétale近傍上で構成する。
各次数がflatであることだけから、この仮定が自動的に従うとは主張しない。

\begin{definition}\label{def:algebraic-family}
 \(S\) を有限型複素スキームとする。
 本節で\emph{smooth properな代数的dg圏の族}とは、
 étale局所的に \(\Perf(A)\) と書け、\(A\) が底環上のdg代数で、
 \(A\) 自身が底上perfect、対角 \(A\) が \(A^{\mathrm{op}}\otL A\) 上perfectである族をいう。
 圏は小さく、pretriangulatedかつ冪等完備とし、すべての底変更を導来的に取る。
 自己同値 \(\rho\) の代数的族は、可逆なperfect bimoduleによって与える。
\end{definition}

\begin{theorem}[étale局所的な可換復元]\label{thm:etale}
 \(S\) を有限型複素スキーム、\(s\in S(\CC)\) とし、
 \((\cK_S,\rho_S)\) を定義\ref{def:algebraic-family}の組とする。
 滑らかなcubic \(X^n\subset\PP^{n+1}\)、\(n\ge3\)、およびenhanced pair同値
 \[
 \phi_s:(\cK_s,\rho_s)\xrightarrow{\sim}(\Ku(X),\rho_X)
 \]
 が与えられたとする。このとき、あるpointed étale近傍
 \((T,t)\to(S,s)\)、滑らかな相対cubic \(Y\subset\PP^{n+1}_T\)、および
 \[
 \Phi_T:(\cK_T,\rho_T)\xrightarrow{\sim}(\Ku(Y/T),\rho_{Y/T})
 \]
 が存在し、\(t\) 上で \(Y_t=X\)、\(\Phi_t\simeq\phi_s\) とできる。
 この同値はperfect bimoduleとrotationを結ぶ同型によって与えられる。

 従って、可換cubicのpairで実現できる複素点は \(S\) のZariski開部分をなす。
 その各複素点で、実現する滑らかなcubicの射影同型類は一意である。
\end{theorem}

この定理は、底 \(S\) 全体上のcubicの一意な降下や、同値の一意性を述べない。
また、\(S\) が連結であることだけでは、上の開部分が \(S\) 全体であるとはいえない。
\(n\ne4\) ではrotationを最初から与えなくても、同じ方法でétale局所的な
可換モデルが存在する（系\ref{cor:etale-unpolarized}）。

\subsection{指定された還元同値を保つ持ち上げ}

\begin{lemma}\label{lem:relative-realization}
 \(R'\twoheadrightarrow R\) を剰余体 \(\CC\) の局所Artin代数のsmall extensionとする。
 \((\cK_{R'},\rho_{R'})\) をsmooth proper pairとし、滑らかなcubic \(Y_R\) と同値
 \[
 \phi_R:(\cK_R,\rho_R)\simeq(\Ku(Y_R),\rho_{Y_R})
 \]
 を固定する。このとき \(Y_R\) を延長する \(Y_{R'}\) と、
 還元後に \(\phi_R\) と同型になるpair同値 \(\phi_{R'}\) が存在する。
 \(Y_R\) の射影空間への埋込みも指定してよい。
\end{lemma}
\begin{proof}
 特殊ファイバーでの \(\phi_R\) をmarkingとして固定する。
 まずcubicの係数を持ち上げ、幾何側に一つのlift \(Y^0_{R'}\) を取る。
 \(P\) を第\ref{sec:formal}節のpairのformal moduli problemとし、
 \(I=\ker(R'\to R)\) と置く。
 与えられた圏とこの幾何学的liftを、\(\phi_R\) によって
 \(P(R')\to P(R)\) の同じhomotopy fibre内で比較する。
 存在するliftの差は \(H^0(T_P)\otimes I\) のtorsorで測られる。
 幾何学的Kodaira--Spencer写像
 \(H^1(T_{Y_0})\simeq H^0(T_P)\) により \(Y^0_{R'}\) を修正し、その差を消す。
 補題\ref{lem:theta}の完全忠実性で、指定した還元同値を保つactual同値へ戻る。
 途中で還元上の自己同型が現れる表示を用いた場合には、
 補題\ref{lem:automorphisms}でそれを持ち上げて補正できる。

 抽象的な偏極cubicのliftに対しても、\(\OO(1)\) の延長と
 \(H^0(\OO(1))\) の自由な基底の持ち上げにより、指定された埋込みを延長できる。
 これは \(\operatorname{PGL}_{n+2}\) のsmoothnessを用いることと同じである。
\end{proof}

\subsection{同値とrotationの同型を表す代数的スタック}

\(U\subset\PP(\CC[x_1,\ldots,x_{n+2}]_3)\) を滑らかなcubicの方程式空間とし、
 \(\mathcal X\to U\) を普遍cubicとする。
 相対的な例外列とmutationにより
 \((\Ku(\mathcal X/U),\rho_U)\) はsmooth properな代数的pairの族となる。
 相対生成子として \(\bigoplus_{i=0}^{n+1}\OO_{\mathcal X}(i)\) の残余射影を使える。
 この生成性は、射影空間の相対Beilinson生成子の引戻しと、
 閉埋込みの押出しが保守的であることから従う。

\begin{proposition}\label{prop:equivalence-stack}
 \(S\)-スキーム \(T\) に対し、
 \[
 (f:T\to U,\ P,\ \alpha:P\circ\rho_T\xrightarrow{\sim}\rho_f\circ P)
 \]
 を対象とするスタック \(\mathcal E\) を考える。
 ここで \(P:\cK_T\simeq\Ku(\mathcal X_f/T)\) はperfect bimoduleで表される同値、
 射は \(\alpha\) と両立するbimoduleの同型である。
 このスタックは \(S\times U\) 上、従って \(S\) 上で局所有限表示のArtin 1-stackである。
\end{proposition}
\begin{proof}
 すべての構成はétale降下と両立するため、底をaffineにして生成子を選んでよい。
 \(\cD=\cK_S^{\mathrm{op}}\otL\Ku(\mathcal X/U)\) は \(S\times U\) 上smooth properである。
 Toën--VaquiéのTheorem 3.6とCorollary 3.15 \cite{TV}により、
 そのperfect bimoduleのmoduliは局所幾何的で局所有限表示である。
 同論文は一般の可換底環上で定式化されているので、定数族だけに限定されない。

 \(P\) が同値を与える条件は開である。
 実際、右随伴を表すperfect bimoduleを \(P^R\) とすると、
 単位と余単位
 \[
 \Delta_{\cK}\longrightarrow P^R P,
 \qquad PP^R\longrightarrow\Delta_{\Ku}
 \]
 のconeがともに零となることが必要十分である。
 生成子でdg代数を選んだ模型で、これらのconeを評価する。
 smooth proper性によりその底上の複体はperfectであり、消滅はそこで検出される。
 従って消滅座は開である。
 これは \cite[Lemma 3.25とCorollary 3.24の証明]{TV} の二つの圏に対する同じ議論である。

 同値 \(P\) の幾何ファイバーで
 \[
 \Ext^i(P,P)=\HH^i(\Ku(X_f)),\qquad i\in\ZZ.
 \]
 補題\ref{lem:HH}から負次数は消える。
 perfectな自己Hom複体のTor振幅の開性により、これは任意の底変更後にも負次数が
 消える開部分として扱える。\cite[Corollary 3.21]{TV}により、
 同値を表す部分は通常のArtin 1-stackとなる。

 最後に \(P\rho_T\) と \(\rho_fP\) の間の同型を付加する。
 射のmoduliの局所有限表示性は \cite[Corollary 3.19]{TV} による。
 その中でconeが零となる部分を取れば同型のmoduliを得る。
 同型が存在する点では両側は可逆bimoduleであり、
 負次数のHomは再び \(\HH^{<0}\) の消滅によって零である。
 従ってこの同型の選択は代数的空間で表され、
 \(\mathcal E\) はArtin 1-stackのままである。
 この構成には \(\rho^3\simeq[2]\) の選択や、そのcoherenceは課していない。
\end{proof}

\subsection{smoothnessとétale局所section}

\begin{proof}[定理\ref{thm:etale}の証明]
 \(\phi_s\) と \(X\) の方程式を選んで \(e\in\mathcal E(\CC)\) を得る。
 \(V\to\mathcal E\) をsmooth atlasとし、\(e\) の上に複素点 \(v\) を選ぶ。
 \(v\) における \(V\to S\) のsmall extensionの持ち上げ問題を考える。
 補題\ref{lem:relative-realization}により、対応する \(\mathcal E\) の対象を
 指定された還元を保って持ち上げられる。
 さらに \(V\to\mathcal E\) のsmoothnessにより、それを \(V\) へ持ち上げる。
 底が局所Artin環で剰余体が \(\CC\) なので、atlasへのétale局所持ち上げはsectionを持ち、
 この持ち上げを同じ底上で取れる。

 \(V\to S\) は局所有限型であるから、Noetherian底上の無限小判定
 \cite[Tag 02HX]{Stacks}により \(v\) でsmoothである。
 \(v\) の近傍を縮め、smooth morphismのétale局所section
 \cite[Tag 055Uとその証明中のLemma 37.38.5]{Stacks}を使うと、
 \(\kappa(v)=\kappa(s)=\CC\) であることから、指定した \(v\) を通る
 \((T,t)\to(S,s)\) と \(T\to V\to\mathcal E\) が得られる。
 これが求める \(Y\)、\(P\)、\(\alpha\) を同時に与える。

 \(T\to S\) の像はZariski開であり、その各複素点の上には複素点を取れる。
 従って可換モデルが存在する複素点の集合は開部分の複素点集合である。
 各点でのモデルの一意性は \cite[Theorem 1.7]{LZ} による。
\end{proof}

\begin{corollary}\label{cor:generic-realization}
 定理\ref{thm:etale}の仮定で \(S\) が既約なら、\(s\) を含む空でないZariski開部分
 \(S^\circ\) の全複素ファイバーに可換cubicモデルが存在する。
 幾何学的generic fibreにも、有限型étale底変更を通じてモデルが存在する。
\end{corollary}
\begin{proof}
 上のétale近傍の像を \(S^\circ\) とする。既約性からこれはdenseである。
 \(S\) を必要なら被約化してgeneric pointを考えると、
 \(T\) の支配的な成分のgeneric residue fieldは \(\CC(S)\) の有限分離拡大である。
 その上のcubicと同値を代数閉包へ底変更すればよい。
\end{proof}

\begin{remark}\label{rem:algebraization-boundary}
 本節では形式的なpower seriesの収束を証明していない。
 指定された代数的圏の族に対し、代数的な同値のmoduliのsmoothnessから
 別途、実際の代数的近傍と同値を作っている。
 その完備化は形式的復元定理のモデルとmarked同型類で一致する。
 しかし、任意に選ばれた形式的座標表示や形式的bimoduleそのものが
 元の底上へ一意に降下する、という結論ではない。
\end{remark}


\begin{corollary}[rotationを指定しない局所的存在]\label{cor:etale-unpolarized}
 定理\ref{thm:etale}で \(n\ne4\) とする。
 \(\rho_S\) は仮定せず、\(\cK_s\simeq\Ku(X)\) というenhanced同値だけを与えても、
 pointed étale近傍 \(T\) 上にcubic \(Y/T\) と、指定された特殊ファイバーの同値を
 延長する \(\cK_T\simeq\Ku(Y/T)\) が存在する。
\end{corollary}
\begin{proof}
 命題\ref{prop:equivalence-stack}から \(\alpha\) の指定を除いた同値のスタックを用いる。
 定理\ref{thm:unpolarized-formal}の証明を、還元同値を指定したhomotopy fibre内で行えば、
 補題\ref{lem:relative-realization}のrotationなしの版が得られる。
 後は定理\ref{thm:etale}と同じsmooth atlasとétale局所sectionの議論を適用する。
 得られたモデルからrotationを移せるが、その \(T\) 全体上での一意性は主張しない。
\end{proof}

\section{Sklyaninブロックと次数シフト}\label{sec:elliptic}

\subsection{中心元と点曲線}

\begin{lemma}\label{lem:center}
 \(G_h=x^3+y^3+z^3-2h^2\sym(xyz)\) は、すべての \(h\) で
 \(S_h\) の非零な中心正則元である。
\end{lemma}
\begin{proof}
 \(D=x^3+y^3+z^3\) とする。関係式から
 \[
 [D,x]=h(zy^2z-yz^2y)=h[zy,yz]=h^2[zy,x^2].
 \]
 また六つの順列を同じ順序へ移すと
 \[
 \sym(xyz)=xyz+\frac h2(-x^3+y^3-z^3)
 \]
 となる。この式を \(x\) と交換させ、再び二次関係式を使えば
 \(2[\sym(xyz),x]=[zy,x^2]\) を得る。
 従って \([G_h,x]=0\)。巡回置換で \(y,z\) についても従う。

 非零性は特殊化で失われない形で確認できる。
 自由代数の次数三部分の線形汎関数 \(\ell\) を
 \[
 \ell(x^3)=\ell(y^3)=\ell(z^3)=1,\quad
 \ell(xyz)=\ell(yzx)=\ell(zxy)=h/2,
 \]
 \[
 \ell(xzy)=\ell(zyx)=\ell(yxz)=-h/2
 \]
 とし、ちょうど二文字が等しい語では零とする。
 三つの二次関係式の左右に一文字を掛けた18個の式を \(\ell\) はすべて殺す。
 例えば \(\ell(z(xy-yx-hz^2))=h/2+h/2-h=0\) である。
 従って \(\ell\) は \((S_h)_3\) に降り、\(\ell(G_h)=3\) となる。
 \((a,b,c)=(1,-1,-h)\) はSklyanin代数の12個の非正則パラメータに属さず、
 \(S_h\) はAS-regular domainである \cite[Introduction]{CGWZ}。
 \(h=0\) は多項式環である。よって \(G_h\) は左右正則である。
\end{proof}

標準Sklyanin記法では \((a,b,c)=(1,-1,-h)\) である。
標準中心元の式 \cite[Equation (2.16)]{Walton}
\[
 g=c(a^3-c^3)x^3+a(b^3-c^3)xyz+b(c^3-a^3)yxz+c(c^3-b^3)y^3
\]
に代入し、\(yxz=xyz-hz^3\) を用いると \(g=-hG_h\) となる。
\(h\ne0\)、\(h^6\ne-27\) では点曲線の公式
\cite[Theorem 2.14]{Walton}から
\[
 E_h=V(x^3+y^3+z^3+h^2xyz)\subset\PP^2
\]
を得る。これは滑らかな楕円曲線であり、
\[
 S_h/(G_h)\simeq B(E_h,L,\sigma_h),\qquad L=\OO_{E_h}(1),
\]
と書ける。ここで \(\sigma_h=t_p\) は平行移動である。
原点をinflection point \(o=[1:-1:0]\) に取ると、
多重線形化した関係式を第一成分 \(o\) で評価した核は
\(p=[1:-1:-h]\) である。従ってこの規約で \(\sigma_h=t_p\) となる。
この中心商と楕円曲線の比較は \cite[Sections 3.5--3.6]{EG} も参照。
\(h=0\) は通常のFermat平面三次曲線の斉次座標環である。
この場合は \(p=o\)、\(\sigma_0=\Id\) とする。

\subsection{rotationに含まれる平行移動を消す}

\begin{lemma}\label{lem:elliptic-conjugacy}
 \(E\) を滑らかな楕円曲線、\(L\) を次数三の直線束、
 \(\sigma=t_p\) とする。
 \(C=B(E,L,\sigma)\) のgraded singularity categoryを \(D^b(E)\) と同一視したとき、
 その次数シフトは可換平面三次曲線のrotation
 \(\rho_0=T_{\OO_E}\circ(-\otimes L)\) と共役である。
\end{lemma}
\begin{proof}
 右加群の規約を
 \[
 \Gamma_*(F)_d=H^0(E,F\otimes L_d),\qquad
 L_d=L\otimes\sigma^*L\otimes\cdots\otimes(\sigma^{d-1})^*L
 \]
 と取る。tailsの次数シフトは \(H(F)=\sigma_*(F\otimes L)\) である。
 Gorenstein parameter零のOrlov windowを比較すると
 \[
 \rho_{\mathrm{nc}}=T_{\OO_E}\circ H=t_{-p}^*\rho_0.
 \tag{2}\label{eq:rho-nc}
 \]
 このwindow公式を確認する。特異圏への商を \(q\) とし
 \(\Psi_i=qR\Gamma_{\ge i}\) と置くと、
 \((1)\Psi_1=\Psi_0H\)、\(q(k)=\Psi_0(\OO_E)[1]\)、
 \[
 \Psi_1F\longrightarrow\Psi_0F
 \longrightarrow\RHom(\OO_E,F)\otimes q(k)
 \]
 がある。楕円曲線のSerre双対から
 \(\Psi_1=\Psi_0T_{\OO_E}^{-1}\) となり、\eqref{eq:rho-nc}を得る。

 原点を \(o\)、\(P_q=\OO_E(q-o)\)、\(A_q=(-\otimes P_q)\) とする。
 spherical twistの核は \(I_\Delta[1]\) なので、
 \[
 K_{T_{P_q}}=(P_q^{-1}\boxtimes P_q)\otimes I_\Delta[1].
 \]
 see-sawを両方向のファイバーへ適用すると
 \[
 (P_q^{-1}\boxtimes P_q)\otimes I_\Delta
 \simeq I_{\{y=x-q\}}.
 \]
 よって \(T_{P_q}=t_q^*T_{\OO_E}\) であり、
 \[
 A_q\rho_0A_q^{-1}=t_q^*\rho_0,\qquad
 \rho_{\mathrm{nc}}=A_{-p}\rho_0A_{-p}^{-1}.
 \tag{3}\label{eq:conjugacy}
 \]
 これらは核とwindowによるdg同値として実現される。
 三次周期性の二つの同型の差は \(\HH^0(D^b(E))^\times=\CC^\times\) のscalarであり、
 共役の自然同型をその三乗根で調整すれば周期性も両立させられる。
 \(\HH^{<0}(D^b(E))=0\) なので、この自然同型の等式に追加の高次ホモトピー障害はない。
 後で用いる有限次数差の作用は、この周期性から構成する。
\end{proof}

\begin{corollary}\label{cor:block}
 \(h^6\ne-27\) のとき、三次周期性と次数シフトを保つdg Morita同値
 \[
 \MF^{\ZZ}(S_h,G_h)
 \simeq
 \MF^{\ZZ}\bigl(\CC[x,y,z],x^3+y^3+z^3+h^2xyz\bigr)
 \]
 が存在する。
\end{corollary}
\begin{proof}
 \cite[Theorem 3.6.2]{EG}、Orlov window、
 \cite[Theorem 5.8]{CCKM}によって両辺を \(D^b(E_h)\) と比較し、
 補題\ref{lem:elliptic-conjugacy}を使う。
 MFと特異圏の比較は、factorizationを商環上の両側無限の完全周期分解へ送る
 dg関手で実現する。中心正則元の商はGorensteinであり、complete resolutionは
 MCM加群の安定圏のdg模型を与える。
 CCKMの完全忠実性をすべてのshiftに適用すると、Hom複体の比較がquasi-isomorphismとなる。
 周囲がAS-regularであることによる高次syzygyのfactorization表示が本質的全射性を与える。
 特異圏からwindowへのdg比較は \cite[Section 5, Theorem 5.8]{CT} による。
 \(h=0\) では恒等の可換比較を用いる。
\end{proof}

\section{中心三次式を加える比較と通常次数への降下}\label{sec:TS}

\subsection{二つの次数ブロックと外部積}

中心変数側の三次式をFermat型に制限する必要はない。
この節では \(T=\CC[u_1,\ldots,u_m]\)、\(m\ge1\)、とし、
\(P\in T_3\) は
\[
 J(P)=\left(\frac{\partial P}{\partial u_1},\ldots,
                 \frac{\partial P}{\partial u_m}\right),\qquad
 \dim_\CC T/J(P)<\infty                                      \tag{4a}\label{eq:central-isolated}
\]
を満たすものとする。これは斉次三次式 \(P\) のアフィン臨界点が
原点だけであるという条件である。
\[
 \Gamma=(\ZZ e_0\oplus\ZZ e_1)/\langle3e_0-3e_1\rangle,
 \qquad w=3e_0=3e_1                                       \tag{4b}\label{eq:central-gamma}
\]
と置く。\(S_d\) の次数を \(de_0\)、すべての \(u_i\) の次数を \(e_1\) とする。
したがって中心変数どうしの混合項を含む \(P\) も次数 \(w\) である。

\begin{lemma}[任意の孤立中心三次式に対する外部積]\label{lem:TS}
 \(S\) を連結非負次数付きnoetherian AS-regular \(\CC\)-代数、
 \(g\in S_3\) を中心正則元とする。
 \eqref{eq:central-isolated}の下で、外部積はpretriangulated hullと
 冪等完備化の後にdg Morita同値
 \[
 \MF^\ZZ(S,g)\otL_\CC\MF^\ZZ(T,P)
 \simeq
 \MF^\Gamma(S\otimes_\CC T,g+P)                             \tag{4c}\label{eq:central-TS}
 \]
 を与える。この同値は二つの因子の次数シフトと、
 両者の三次内部シフトを \([2]\) と同一視する周期性を保つ。
\end{lemma}
\begin{proof}
 \(R=S\otimes_\CC T=S[u_1,\ldots,u_m]\)、\(W=g+P\)、
 \(A=R/(W)\)、\(C=S/(g)\) と置く。
 まず \(W\) は中心正則元である。
 実際、\(u\) に関する単項式順序を固定すると、\(P\ne0\) の最高項の係数は
 非零のscalarであり、任意の非零多項式との積の最高項は消えない。
 以下では中心正則元の商に対するMFとgraded singularity categoryの比較
 \cite[Theorem 5.8]{CCKM}を用いる。
 \(\Gamma\to\ZZ\)、\(e_0,e_1\mapsto1\)、は有限の核を持ち、
 正の整数次数を与えるので、有限次数付き自由分解による同じ比較は
 \(\Gamma\)-次数についても適用できる。

 \paragraph{外部積と完全忠実性.}
 factorizationの規約を
 \[
 Q=(Q_1\xrightarrow{\varphi}Q_0
                  \xrightarrow{\psi}Q_1(w)),\qquad
 \psi\varphi=g\Id,\quad\varphi\psi=g\Id
 \]
 とする。これを \(Q^{i+2}=Q^i(w)\) を満たす曲がった複体に展開すると、
 \(Q^0=Q_0\)、\(Q^1=Q_1(w)\) である。
 二つのfactorizationの外部積は、二つの奇微分をKoszul符号付きで足すことにより定める。
 奇微分の外部和の平方は \(g+P\) である。
 二つの \([2]\) を同じ内部シフトに対応させる関係が
 \(3e_0=3e_1\) であり、得られるfactorizationは有限自由である。

 有限自由対象 \(Q,Q'\) と \(K,K'\) に対して、
 係数を \(S\)-次数と \(T\)-次数に分解すると
 \[
 \Hom_{\MF^\Gamma}(Q\boxtimes K,Q'\boxtimes K')
 \simeq
 \operatorname{Tot}\!\left(
   \Hom_{\MF^\ZZ(S,g)}(Q,Q')\otimes_\CC
   \Hom_{\MF^\ZZ(T,P)}(K,K')\right)                         \tag{4d}\label{eq:central-hom}
 \]
 を得る。具体的には、左辺の次数 \(r\) の斉次写像は、
 \(a+b=r\) を満たす右辺の次数 \((a,b)\) の写像の和に一意に分解する。
 二つの内部次数の差を移す操作は \((3,-3)\) の整数倍だけであり、
 これは二つのcohomological degreeを \((2,-2)\) だけ変える操作に対応する。
 微分と合成の符号も外部和のKoszul符号に一致する。
 各因子のHom複体は、有限個の自由加群の次数シフトと係数代数の非負次数性から
 下に有界である。従って各全cohomological degreeに寄与する項は有限個であり、
 \eqref{eq:central-hom}に完備化は必要ない。
 これが外部積のdg完全忠実性を与える。

 \paragraph{中心変数を零とした商への生成性の帰着.}
 \(\MF^\Gamma(R,W)\) の有限自由対象の奇微分を \(D\) とする。
 \(u_i\) は中心変数なので、係数行列の \(u_i\)-微分は \(R\) 上の導分であり、
 \[
 D\,\frac{\partial D}{\partial u_i}
 +\frac{\partial D}{\partial u_i}\,D
 =\frac{\partial P}{\partial u_i}\Id
 \]
 が成り立つ。内部シフトを付けると、これは
 \(\partial_iP:N\to N(2e_1)\) の零ホモトピーである。
 斉次イデアル \(J(P)\) は \((u_1,\ldots,u_m)\)-primaryなので、
 ある整数 \(d\ge2\) に対してすべての \(u_i^d\) が \(J(P)\) に属する。
 斉次係数でこれらを偏微分の線形結合に書けば、
 \(u_i^d:N\to N(de_1)\) も特異圏で零となる。

 \(N\) を有限生成 \(A\)-加群の有界複体で表し、
 \(u_1^d,\ldots,u_m^d\) に沿うKoszul複体を \(K(u^d;N)\) と書く。
 零写像のconeは元の対象を直和因子に持つため、\(N\) は
 \(K(u^d;N)\) の直和因子である。
 より正確には、一変数ごとのKoszul cone関手は完全関手でretractを保つので、
 各 \(u_i^d\) の零作用から得られるretractを順に合成すればよい。
 このKoszul複体のcohomologyは \((u_1^d,\ldots,u_m^d)\) に殺され、
 \[
 (u_1,\ldots,u_m)^{m(d-1)+1}\subset(u_1^d,\ldots,u_m^d)
 \]
 である。従って各cohomology加群は、有限個の \(C=A/(u_1,\ldots,u_m)\)-加群を
 商とする有限filtrationを持つ。標準切断の三角形も使うと、
 右辺の特異圏は \(C\)-加群の制限とその次数シフトで厚く生成される。

 ここで制限が特異圏へ降りることを確認する。
 \(u_1,\ldots,u_m\) は \(A\) の中心正則列である。
 一つの変数 \(u_i\) について、\(R\) で \(u_i a=Wb\) とする。
 \(u_i=0\) とした商では、\(\overline W=g+P|_{u_i=0}\) は正則である。
 特殊化した \(P\) が非零なら最高項のscalar係数を使い、
 零なら \(g\) の正則性を使えばよい。
 従って \(\overline b=0\)、すなわち \(b=u_i b_1\) であり、
 多項式環 \(R\) で \(u_i\) を消去して \(a=Wb_1\) を得る。
 先行する変数を順に零とした同じ議論が正則列を与える。
 よって \(C\) は有限自由Koszul分解を持つ \(A\)-perfect加群であり、
 制限は \(\Dsg(C)\to\Dsg(A)\) を誘導する。

 \paragraph{制限加群の安定化の明示.}
 Euler公式から
 \[
 q_i=\frac13\frac{\partial P}{\partial u_i},\qquad
 P=\sum_{i=1}^{m}u_iq_i
 \]
 と書ける。\(T\) 上の通常のKoszul分解を
 \[
 K_j=T(-je_1)\otimes_\CC\Lambda^j\CC^m,
 \qquad d_K=\sum_i u_i\iota_i
 \]
 とする。\(\iota_i\) は外積代数の縮約、\(\varepsilon_i\) は対応する基底による左外積である。
 \[
 h_K=\sum_iq_i\varepsilon_i:K_j\longrightarrow K_{j+1}(w)
 \]
 と置くと
 \[
 d_Kh_K+h_Kd_K=P\Id,\qquad h_K^2=0.                         \tag{4e}\label{eq:central-koszul}
 \]
 ここで \(q_i\) の次数は \(2e_1\) であり、
 \(K_{j+1}(w)\) のシフトによりすべての表示写像は内部次数零である。

 有限Koszul複体を偶奇で折り畳み、
 \[
 (K_P)_0=\bigoplus_{j\ge0}K_{2j}(jw),\qquad
 (K_P)_1=\bigoplus_{j\ge0}K_{2j+1}(jw)
 \]
 と置く。\(d_K+h_K\) は
 \[
 (K_P)_1\longrightarrow(K_P)_0
             \longrightarrow(K_P)_1(w)
 \]
 を与え、\eqref{eq:central-koszul}によってこれは \(P\) の有限自由factorizationである。

 \(M\) を \(C\)-加群とし、その特異圏での像を表すfactorizationを
 \(Q=(Q_1\xrightarrow\varphi Q_0\xrightarrow\psi Q_1(w))\) とする。
 有限回のsyzygyとcohomological shiftを取れば
 \(M=\operatorname{coker}\varphi\) としてよい。
 実際、元の加群との差に現れる \(C\)-perfect対象は制限後も \(A\)-perfectである。
 \(\psi\varphi=g\Id\) と \(g\) の正則性から \(\varphi\) は単射であり、
 \[
 0\longrightarrow Q_1\xrightarrow\varphi Q_0\longrightarrow M\longrightarrow0
 \]
 は \(S\)-自由分解である。
 これをhomological degree \(1,0\) に置き、
 \(d_Q=\varphi\)、\(s_Q:Q_0\to Q_1(w)\) を \(\psi\) とする。
 \(s_Q\) は \(Q_1\) 上では零とする。

 \(u_iM=0\) として \(M\) を \(R\)-加群に制限したときの有限自由分解は
 \[
 F_a=\bigoplus_{b+c=a}Q_b\otimes_\CC K_c,
 \qquad d=d_Q\otimes1+(-1)^b\otimes d_K
 \]
 である。\(b\in\{0,1\}\) は \(Q\) 側のhomological degreeである。
 \[
 s=s_Q\otimes1+(-1)^b\otimes h_K:F_a\longrightarrow F_{a+1}(w)
 \]
 と置くと、異なる因子の奇作用の反可換性から
 \[
 ds+sd=W\Id,\qquad s^2=0.                                  \tag{4f}\label{eq:central-shamash}
 \]
 が成り立つ。

 このデータの周期化が実際に \(A\)-自由分解を与えることも確認しておく。
 \(F_a=0\) for \(a<0\) とし、
 \[
 \mathcal F_n=
 \bigoplus_{j\ge0}(F_{n-2j}\otimes_R A)(-jw),\qquad n\ge0
 \]
 と置く。微分は同じ \(j\) 内の \(d\) と、\(j\ge1\) について
 \[
 s:(F_{n-2j}\otimes_R A)(-jw)
 \longrightarrow(F_{n-2j+1}\otimes_R A)(-(j-1)w)
 \]
 の和である。\eqref{eq:central-shamash}により平方は零になる。
 各 \(n\) で和は有限である。

 \(j\) による増大濾過を使うと、第一ページは
 \(H_*(F\otimes_R A)\) の内部・homological shiftの和となる。
 \(W\) が正則で \(WM=0\) なので、
 \[
 H_i(F\otimes_R A)=
 \begin{cases}
 M&i=0,\\
 M(-w)&i=1,\\
 0&i\ne0,1.
 \end{cases}
 \]
 次の微分は、\(j\ge1\) の \(H_0\) を \(j-1\) の \(H_1\) へ送る同型である。
 これを \(ds+sd=W\Id\) から直接見られる。
 \(a\in F_0\) が \(M\) の元を表すなら、\(s(a)\) の \(A\) 上の像はcycleであり、
 \(d s(a)=Wa\) である。
 短完全列
 \[
 0\longrightarrow F(-w)\xrightarrow{W}F
      \longrightarrow F\otimes_R A\longrightarrow0
 \]
 のconnecting mapはそのcycleを元の \(a\) に送る。
 従ってこの微分は、上の \(H_1\simeq M(-w)\) の標準同定の下で恒等である。
 隣接する二項はすべて相殺し、\(j=0\) の \(H_0=M\) だけが残る。
 濾過は各全homological degreeで有限なので収束し、
 \(\mathcal F\) は \(M\) の \(A\)-自由分解となる。

 \(F\) は有限複体なので、この分解の十分先は二周期的である。
 そのfactorizationは
 \[
 \widehat F_0=\bigoplus_{j\ge0}F_{2j}(jw),\qquad
 \widehat F_1=\bigoplus_{j\ge0}F_{2j+1}(jw),
 \qquad d+s
 \]
 による折り畳みであり、Koszul符号と内部シフトを含めて
 \(Q\boxtimes K_P\) に一致する。
 実際、\(\widehat F_0\) は
 \(Q_0\otimes(K_P)_0\oplus Q_1\otimes(K_P)_1(w)\)、
 \(\widehat F_1\) は
 \(Q_1\otimes(K_P)_0\oplus Q_0\otimes(K_P)_1\) である。
 従って制限加群の安定化は
 \[
 \operatorname{stab}_A(M)\simeq Q\boxtimes K_P.               \tag{4g}\label{eq:restriction-mf}
 \]
 ここで全体のshiftにも曖昧さはない。
 十分大きい \(\ell\) について
 \(\mathcal F_{2\ell}=(\widehat F_0\otimes_R A)(-\ell w)\)、
 \(\mathcal F_{2\ell+1}=(\widehat F_1\otimes_R A)(-\ell w)\) であり、
 そのcokernelは \(\Omega_A^{2\ell}M\) である。
 \(\Omega_A^{2\ell}M\simeq M[-2\ell]\) と \((w)\simeq[2]\) によって
 この内部シフトはちょうど打ち消される。

 最後に、\(\Gamma/\ZZ e_0\simeq\ZZ/3\) であるため、
 任意の \(\Gamma\)-次数付き \(C\)-加群は、
 三つの次数類 \(0,e_1,2e_1\) に属する部分に分解する。
 各部分は通常の \(\ZZ\)-次数付き \(C\)-加群の内部シフトであり、
 \eqref{eq:restriction-mf}の右辺の \(T\) 側を \((j)\)、\(j=0,1,2\)、
 だけシフトすればすべて得られる。
 前段のKoszul filtrationと合わせて、外部積の像は右辺全体を厚く生成する。
 これで本質的全射性が従う。
 周期性および二つの次数シフトとの両立は、ここで用いた外部和の構成そのものから従う。
\end{proof}

\begin{remark}[一変数の場合の明示行列]\label{rem:central-one-variable}
 \(m=1\)、\(P=u^3\) では \(K_P=(u,u^2)\) である。
 上の分解とホモトピーは
 \[
 Q_1(-e_1)\xrightarrow{(-u,\varphi)^t}
 Q_1\oplus Q_0(-e_1)\xrightarrow{(\varphi,u)}Q_0,
 \]
 \[
 s_0=(\psi,u^2)^t,\qquad s_1=(-u^2,\psi)
 \]
 となる。折り畳んだfactorizationは
 \[
 \alpha=
 \begin{pmatrix}\varphi&u\\-u^2&\psi\end{pmatrix}:
 Q_1\oplus Q_0(-e_1)\longrightarrow Q_0\oplus Q_1(w-e_1),
 \]
 \[
 \beta=
 \begin{pmatrix}\psi&-u\\u^2&\varphi\end{pmatrix}:
 Q_0\oplus Q_1(w-e_1)\longrightarrow
 Q_1(w)\oplus Q_0(w-e_1).
 \]
 中心性により \(\alpha\beta=\beta\alpha=(g+u^3)\Id\) である。
 これらの式は、一般の \(K_P\) による安定化の符号と次数シフトの具体的な検算になる。
\end{remark}

\subsection{有限群による整数次数の回復}

\(\Gamma\to\ZZ\)、\(e_0,e_1\mapsto1\)、の核は
\(T_\Gamma\simeq\ZZ/3\) であり、\(e_1-e_0\) で生成される。
分裂 \(1\mapsto e_0\) の下で、その双対群
\(H_\Gamma=\Hom(T_\Gamma,\CC^\times)\simeq\mu_3\) は
\(S\) 上恒等、すべての \(u_i\) を同時に三乗根倍する作用である。
\(P\) は斉次三次式なのでこの作用は \(W=g+P\) を保ち、
\[
 \MF^\Gamma(R,W)
 \simeq\bigl(\MF^\ZZ(R,W)\bigr)^{H_\Gamma}.                \tag{5}\label{eq:eq}
\]
有限群の位数は可逆である。
同変な有限次数付き自由加群は、連結な整数次数と有限群表現の半単純性により
同変な斉次基底を持つので、この比較は有限自由factorizationのdg模型で実現する。
さらに双対群 \(T_\Gamma\) で同変化すると、元の通常次数の圏が回復される
\cite[Theorem 4.2, Section 8, Lemma 8.6, Theorem 8.9]{Elagin}。
冪等完備化はこの比較の一部として取る。

外部積 \eqref{eq:central-TS} の左辺で双対作用の生成元は
\[
 \rho_T\rho_S^{-1}
\]
である。ここで \(\rho_S,\rho_T\) はそれぞれの因子の内部次数一のシフトである。
その三乗を恒等と同一視する自然同型は、両因子の
\((3)\simeq[2]\) から得られる。
二つのperiodicityを含めて \(\rho_S\) を保つdg同値はこの作用と両立し、
\eqref{eq:eq}と双対同変化により通常次数へ降りる。
この際 \(e_0\) のシフトも保存されるので、通常次数のdegree shift \((1)\) が保存される。
有限群作用のcoherenceはfactorizationの内部シフトのcoherenceから引き継ぎ、
単に自己同値の三乗が同型になるという存在命題だけから群作用を作るのではない。

\subsection{Sklyaninブロックへの適用}

\begin{proof}[定理\ref{thm:global}の証明]
 補題\ref{lem:TS}を \((S_h,G_h)\) と
 \((\CC[x,y,z],x^3+y^3+z^3+h^2xyz)\) のそれぞれに適用する。
 系\ref{cor:block}は第一因子の次数シフトと三次周期性を保ち、
 \(\MF^\ZZ(T,P)\) 側では恒等を用いる。
 前節の降下によって通常次数のdegree shiftを保つdg Morita同値
 \[
 \MF^\ZZ(S_h[u],G_h+P)
 \simeq
 \MF^\ZZ\!\left(\CC[x,y,z,u],x^3+y^3+z^3+h^2xyz+P\right)
 \]
 を得る。体上のOrlov同値と \cite{CCKM} により、これは残余圏の比較である。

 可換方程式の \(u_i\)-偏微分がすべて零なら、
 \eqref{eq:central-isolated}によって全 \(u_i=0\) である。
 \(h^6\ne-27\) の下では残る平面三次式のアフィン臨界点も原点だけである。
 よって射影超曲面 \(Y_{h,P}\) は滑らかである。
 次元は \(m+1\ge3\) なので、一意性は \cite[Theorem 1.7]{LZ} による。
\end{proof}



\begin{corollary}[平行移動を伴う一般の楕円ブロック]\label{cor:elliptic-central-cubic}
 \(S\) をHilbert series \((1-t)^{-3}\) を持つnoetherianな
 二次三次元AS-regular \(\CC\)-代数、\(g\in S_3\) を中心正則元とする。
 滑らかな平面三次曲線 \(E\)、\(L=\OO_E(1)\)、および
 ある原点に関する平行移動 \(\sigma=t_p\) に対して、次数付き同型
 \[
 S/(g)\simeq B(E,L,\sigma)
 \]
 が存在すると仮定する。
 \(f_E\in\CC[x,y,z]_3\) を \(E\subset\PP^2\) の方程式とし、
 \(P\in\CC[u_1,\ldots,u_m]_3\)、\(m\ge2\)、は
 \eqref{eq:central-isolated}を満たすとする。
 このとき \(S[u]/(g+P)\) のrotation付き残余圏の滑らかな可換モデルは
 \[
 V(f_E(x,y,z)+P(u))\subset\PP^{m+2}
 \]
 であり、射影同型を除いて一意である。
\end{corollary}
\begin{proof}
 Gorenstein parameter零のOrlov同値により
 \(\MF^\ZZ(S,g)\) は \(D^b(E)\) と同値であり、
 次数シフトは補題\ref{lem:elliptic-conjugacy}の
 \(T_{\OO_E}\sigma_*(-\otimes L)\) に対応する。
 同補題の直線束による共役とperiodicityのscalar調整により、
 \(\MF^\ZZ(\CC[x,y,z],f_E)\) との、次数シフトと三次周期性を保つdg同値を得る。
 あとは補題\ref{lem:TS}と有限次数差による降下を適用すればよい。
 滑らかさと一意性は定理\ref{thm:global}と同じ証明による。
 方程式 \(f_E\) の非零scalar倍による違いは、
 分離された \(x,y,z\) の同時三乗根倍によって射影同型に吸収される。
\end{proof}

\begin{remark}[この拡張の範囲]
 定理\ref{thm:global}では中心変数間の任意の三次混合項を許すが、
 非可換ブロックの生成元と中心変数の混合項は含めない。
 別の非可換ブロックについては、その生成元による通常の偏微分が代数の導分とは限らず、
 本節の \(J(P)\)-primaryな支持への帰着は適用できない。
 この点は独立した生成性の問題として残る。
 本節の拡張の先行研究上の新規性は未確認である。
\end{remark}

\section{複数ブロックの比較と生成性の帰着}\label{sec:multiblock}

\subsection{生成性を仮定しない外部積}

\(A_1,\ldots,A_s\) を、\((A_i)_0=\CC\) を満たす局所有限な
非負次数付き代数、\(f_i\in Z(A_i)_3\) を中心的な三次斉次元とする。
この小節の \(\MF\) は、有限自由次数付きmatrix factorizationのdg圏の
pretriangulated hullを取り、さらに冪等完備化したものを表す。
この定義にはtensor product algebraのnoetherian性を仮定しない。
\[
 \Gamma_s=\ZZ^s/\langle3e_i-3e_1:2\le i\le s\rangle,
 \qquad w=3e_1,
 \qquad A=\bigotimes_{i=1}^s A_i,
 \qquad W=\sum_{i=1}^s f_i
\]
とする。\(A_i\) の次数 \(d\) の斉次元には \(de_i\) を割り当てる。

\begin{lemma}[外部積の完全忠実性]\label{lem:multiblock-ff}
 外部積はdg完全忠実関手
 \[
 \Phi:\Perf\!\left(\bigotimes_{i=1}^s\MF^{\ZZ}(A_i,f_i)\right)
 \longrightarrow\MF^{\Gamma_s}(A,W)
 \]
 を定める。\(i\) 番目の因子のdegree shiftは右辺の \((e_i)\) に対応し、
 \((3e_i)=[2]\) の標準同型とも両立する。
\end{lemma}
\begin{proof}
 まず有限自由factorization \(P_i,Q_i\) を取る。
 tensor productの奇微分はKoszul符号付きの外部和であり、
 その二乗は \(W\) となる。
 各因子のHom複体を \(C_i^\bullet\) と書くと、
 自由加群のHomを行列係数ごとに分解することにより
 \[
 \Hom_{\MF^{\Gamma_s}(A,W)}^\bullet
       (\boxtimes_iP_i,\boxtimes_iQ_i)
 \cong
 \operatorname{Tot}\!\left(\bigotimes_i C_i^\bullet\right)
 \tag{\(\dagger\)}\label{eq:multiblock-hom}
 \]
 を得る。ここでは次数群の関係 \(3e_i=3e_1\) が、
 各因子の二つのcohomological shiftを他の因子へ移す操作に正確に対応する。
 従って右辺の次数 \(q\) の和は
 \(q_1+\cdots+q_s=q\) にわたる通常のtensor product複体の和であり、
 追加の完備化や直積を取らない。

 この和の有限性を確認する。各 \(P_i,Q_i\) は有限個の内部次数shiftからなり、
 \(A_i\) は非負次数なので、\(C_i^q=0\) は十分小さい \(q\) で成り立つ。
 したがって固定した全次数 \(q\) に寄与する \((q_1,\ldots,q_s)\) は有限個である。
 同様に、行列係数を元の \(\ZZ^s\)-次数へ分解したときも、
 固定した \(\Gamma_s\)-次数と非負次数の条件の下では有限個の成分しか現れない。
 これにより \eqref{eq:multiblock-hom} は実際のHom複体の同型である。
 differentialとcompositionの一致はtensor productのKoszul符号から従う。
 完全忠実性は有限個のcone、shift、直和因子を取っても保たれるので、主張を得る。
 内部次数shiftとperiodicityとの両立はこの構成自体に含まれる。
\end{proof}

\subsection{任意個の楕円THCRの積と周囲の正則性}

Chanの多重Serre同値 \cite[Theorem 3.4]{ChanMulti} は、代数のnoetherian性を
仮定せず、right ampleな可逆bimoduleの組に適用できる。
同論文のCorollaries 5.8--5.9は二因子のTHCRと三次元AS-regular代数の積を扱う。
以下では、その証明で用いる面と帯の議論をfactor数について帰納し、任意個の楕円因子を扱う。
Keeler \cite[Theorem 2.7, Corollary 2.8]{KeelerMulti} によっても、
この積のbimoduleの組の左右NC-amplenessを確認できる。
各translationはnumerical Picard groupに自明に作用し、全ての座標が正の外部積はampleだからである。
ただし多重amplenessだけではnoetherian性は自動的に従わず、以下で帯の有限性を別に扱う。

\begin{lemma}[楕円THCRの有限テンソル積]
\label{lem:elliptic-thcr-tensor-noeth}
 \(E_1,\ldots,E_r\) を滑らかな複素射影楕円曲線、\(L_i\) を正次数の可逆層、
 \(\sigma_i\) をtranslationとする。このとき
 \[
 C=\bigotimes_{i=1}^r C_i,\qquad C_i=B(E_i,L_i,\sigma_i)
 \]
 は左右noetherianである。
\end{lemma}
\begin{proof}
 \(X=\prod_iE_i\) とし、各factorから引き戻した可逆bimoduleを
 \(\mathcal L_i\) と書く。異なるfactorの作用は可換であり、
 \(\mathcal L^{\mathbf d}=\mathcal L_1^{d_1}\cdots\mathcal L_r^{d_r}\) と置けば
 \[
 C_{\mathbf d}=H^0(X,\mathcal L^{\mathbf d})\qquad(\mathbf d\in\mathbb N^r)
 \]
 は \(C\) のテンソル積多重次数と乗法を与える。
 可逆bimoduleの逆を用いて \(\mathcal L^{\mathbf d}\) は
 \(\mathbf d\in\mathbb Z^r\) に対しても定義する。

 \emph{Ampleness.}
 外部積のvery ample可逆層 \(H=\boxtimes_iH_i\) を固定する。
 任意のcoherent層 \(\mathcal F\) を、有限個の \(H\) の負冪の直和からの全射で
 \(\dim X\) 段まで近似する。
 各 \(\mathcal L_i^{d_i}\) の台となる可逆層の次数は \(d_i\deg L_i\) なので、
 その各近似項を \(\mathcal L^{\mathbf d}\) でtwistした高次cohomologyは、
 \(d_i\) が全て十分大きいと消える。
 ここでは楕円曲線上の正次数の可逆層の \(H^1\) の消滅とK\"unneth公式を用いる。
 最後の核にはGrothendieck消滅を適用し、dimension shiftingによって
 \[
 H^q(X,\mathcal F\mathcal L^{\mathbf d})=0
 \quad(q>0,\ \mathbf d\gg0)
 \]
 を得る。従ってこの組はright ampleである。
 \(\mathcal F\otimes H^{-q}\) にもこの消滅を適用すると、
 Castelnuovo--Mumford regularityにより
 \(\mathcal F\mathcal L^{\mathbf d}\) の大域生成も \(\mathbf d\gg0\) で従う。
 自動同型によるpushforwardはcohomologyと大域生成を変えないので、
 この議論は可逆bimoduleの左右の規約に依存しない。

 \emph{面と有限幅の帯.}
 代数 \(C\) の右noetherian性と、全てのcoherent \(\mathcal F\)、\(q\ge0\) に対する
 \[
 Q^q(\mathcal F)=\bigoplus_{\mathbf d\in\mathbb N^r}
 H^q(X,\mathcal F\mathcal L^{\mathbf d})
 \]
 の有限生成性を、\(r\) について同時に示す。
 \(r=0\) の場合は有限次元ベクトル空間の主張である。
 \(r-1\) 個の因子まで成立すると仮定し、
 \(C_{\widehat i}=\bigotimes_{j\ne i}C_j\)、\(X_{\widehat i}=\prod_{j\ne i}E_j\) と置く。
 座標 \(d_i=a\) を固定した \(Q^q(\mathcal F)\) の面は、\(C_{\widehat i}\)-加群である。
 射影 \(\pi_i:X\to X_{\widehat i}\) と
 \(\mathcal F_i=\mathcal F\mathcal L_i^a\) に対するLerayスペクトル系列の
 全多重次数の和は
 \[
 E_2^{p,b}=Q^p_{X_{\widehat i}}
       (R^b\pi_{i*}\mathcal F_i)
 \ \Longrightarrow\
 \bigoplus_{\mathbf d_{\widehat i}\ge0}
 H^{p+b}(X,\mathcal F_i\mathcal L_{\widehat i}^{\mathbf d_{\widehat i}})
 \]
 となる。射影公式と残るtranslationに関する自然性により、これは
 \(C_{\widehat i}\)-加群のスペクトル系列である。
 帰納法で各 \(E_2\) 項はnoetherianな \(C_{\widehat i}\) 上有限生成である。
 従って、有限個の核・商・拡張で得られる各面もnoetherianである。

 次の帯に関する観察を使う。
 \(\mathbb N^r\)-次数付き \(C\)-加群 \(M\) の各面 \(M_{d_i=a}\) が
 \(C_{\widehat i}\) 上noetherianであり、\(M_{\ge\mathbf b}=0\) とする。
 \[
 M^{(i)}=\bigoplus_{d_1\ge b_1,\ldots,d_i\ge b_i}M_{\mathbf d}
 \quad(0\le i\le r)
 \]
 は \(M=M^{(0)}\supset\cdots\supset M^{(r)}=0\) という \(C\)-部分加群のfiltrationである。
 各 \(M^{(i-1)}/M^{(i)}\) をさらに \(d_i\ge a\) でfiltrationすると、
 有限個の商は固定面の \(C_{\widehat i}\)-部分加群になる。
 その商では \((C_i)_{>0}\) の作用は零なので、これらは \(C\)-加群としてもnoetherianである。
 よって \(M\) はnoetherianである。この観察は \(C\) 自身のnoetherian性を使わない。
 各面がnoetherianな加群の部分商についても、同じ観察を適用できる。

 \emph{切断加群.}
 \(q>0\) のとき、amplenessにより \(Q^q(\mathcal F)_{\ge\mathbf b}=0\) となる
 \(\mathbf b\) がある。従って前段階から
 \(Q^q(\mathcal F)\) は\emph{noetherian}である。
 特に、後で用いる \(Q^1\) の任意の部分加群も有限生成である。

 \(q=0\) について、\(\mathcal F\mathcal L^{\mathbf a}\) を大域生成する
 \(\mathbf a\ge0\) を選び、逆bimoduleをtensorして
 \[
 0\longrightarrow\mathcal K\longrightarrow
       (\mathcal L^{-\mathbf a})^{\oplus t}
       \longrightarrow\mathcal F\longrightarrow0
 \]
 を得る。\(\mathcal K\) はcoherentである。
 この積の場合、\(\mathbf v\) の一つでも成分が負なら
 \(H^0(X,\mathcal L^{\mathbf v})=0\) である。
 実際、外部積のそのfactorの次数が負であり、K\"unneth公式を適用できる。
 従って切断を全非負多重次数で取った列は
 \[
 C(-\mathbf a)^{\oplus t}\longrightarrow Q^0(\mathcal F)
       \longrightarrow Q^1(\mathcal K)
 \]
 と正確である。最初の写像の像は有限生成であり、余核は既にnoetherianと示した
 \(Q^1(\mathcal K)\) の部分加群なので、\(Q^0(\mathcal F)\) は有限生成である。
 また任意の \(\mathbf a\ge0\) について
 \[
 Q^0(\mathcal F)_{\ge\mathbf a}(\mathbf a)
 \cong Q^0(\mathcal F\mathcal L^{\mathbf a})
 \]
 であるから、第一象限によるtruncationも有限生成である。

 \emph{イデアルの有限生成.}
 多重斉次右イデアル \(I\subset C\) を取る。
 \cite[Theorem 3.4]{ChanMulti} により、sheafificationはexactであり、
 \(\widetilde I\subset\widetilde C=\OO_X\) となる。
 \(X\) はnoetherianなので、このquasi-coherent部分層はcoherentである。
 随伴のunit \(\eta_I:I\to Q^0(\widetilde I)\) の自然性と
 \(\eta_C:C\xrightarrow{\sim}Q^0(\OO_X)=C\) により、\(\eta_I\) は単射である。
 その余核は多重tailsに関するtorsionであり、\(Q^0(\widetilde I)\) の商だから有限生成である。
 有限生成torsion加群は一つの第一象限上で零となるため、ある \(\mathbf a\ge0\) に対して
 \[
 I_{\ge\mathbf a}=Q^0(\widetilde I)_{\ge\mathbf a}
 \]
 となる。右辺の有限生成性は前段階で証明した。
 一方、\(I/I_{\ge\mathbf a}\) は \(C/C_{\ge\mathbf a}\) の部分加群である。
 \(C\) の固定面は有限次元空間 \((C_i)_a\) と \(C_{\widehat i}\) のtensor productなので
 \(C_{\widehat i}\)-noetherianである。
 従って帯の観察により \(C/C_{\ge\mathbf a}\) はnoetherianであり、
 \(I/I_{\ge\mathbf a}\) も有限生成となる。これで \(I\) の有限生成性を得た。

 加法と両立する \(\mathbb N^r\) の整列順序、例えば全次数付き辞書式順序で
 任意の右イデアルの先頭多重斉次項を取る。
 その斉次イデアルの有限生成元を元のイデアルへ持ち上げ、
 先頭次数を下げる帰納法を行えば、任意の右イデアルも有限生成である。
 これで通常の右noetherian性と帰納法が完了する。
 opposite algebraには逆translationを伴う正次数の可逆bimoduleが対応するので、
 同じ証明により左noetherian性も得る。
\end{proof}

\begin{lemma}[正次数中心正則元からのnoetherian性の持ち上げ]
 \label{lem:central-noetherian-lift}
 \(D\) を連結局所有限な非負次数付き代数、
 \(g\in Z(D)\) を正次数の斉次正則元とする。
 \(D/(g)\) がnoetherianなら \(D\) もnoetherianである。
\end{lemma}
\begin{proof}
 \(g\)-adic associated gradedは \((D/(g))[t]\) でありnoetherianである。
 任意の斉次右イデアル \(J\) に誘導された \(g\)-adic filtrationの
 associated gradedはこの環の右イデアルなので有限生成である。
 その二重斉次生成元を \(J\) に持ち上げる。
 内部次数 \(d\) の元のinitial termをこれらの持ち上げで消すと、
 \(g\)-adic orderが真に増える。
 しかし \(g^nD\) の内部次数は少なくとも \(n\deg(g)\) なので、
 固定した \(d\) ではこの操作は有限回で終わる。
 よって全ての斉次右イデアルが有限生成である。
 任意の右イデアルについては通常の内部次数の最高次項を取り、
 生成元を持ち上げて次数に関する帰納法を行えば有限生成性を得る。
 左側も同じである。
\end{proof}

\begin{proposition}[周囲のテンソル積]
\label{prop:tensor-sklyanin-ambient}
 \(h_i^6\ne-27\)、\(S_i=S_{h_i}\)、\(g_i=G_{h_i}\) とし、
 \[
 A=(\bigotimes_{i=1}^rS_i)[u_1,\ldots,u_m],\qquad
 N=3r+m,\qquad W=\sum_i g_i+P(u),\qquad R=A/(W)
 \]
 と置く。ここで \(P\) は任意の中心斉次三次式であり、\(m=0\) なら \(P=0\) とする。
 \(A\) は左右noetherianなKoszul AS-regular代数で、
 \[
 \operatorname{gldim}A=N,\qquad H_A(t)=(1-t)^{-N},\qquad
 \operatorname{Ext}^j_A(\CC,A)=
 \begin{cases}\CC(N)&j=N,\\0&j\ne N\end{cases}
 \]
 を満たす。\(W\) は中心正則元であり、\(R\) はinjective dimension \(N-1\)、
 Gorenstein parameter \(N-3\) のAS-Gorenstein代数である。
 この主張に \(\sigma_i\) の有限位数や \(P\) の孤立性は必要ない。
\end{proposition}
\begin{proof}
 \(S_i/(g_i)\) の斉次基底の持ち上げ \(V_i\subset S_i\) を選ぶと、
 \(g_i\) の正則性と次数に関する帰納法により
 \(S_i\cong V_i\otimes\CC[g_i]\) は中心 \(\CC[g_i]\)-加群として自由である。
 \(g_1,\ldots,g_r\) は \(\bigotimes_iS_i\) 上の中心正則列をなし、その商は
 補題\ref{lem:elliptic-thcr-tensor-noeth}の \(C\) である。
 補題\ref{lem:central-noetherian-lift}により、正次数中心正則元について
 商のnoetherian性を元の連結次数付き代数へ持ち上げられる。
 これを反復し、中心変数についてのHilbert基底定理を用いると \(A\) は左右noetherianとなる。

 各 \(S_i\) の長さ三のKoszul双加群分解と、多項式環のKoszul双加群分解をtensorすると、
 \(A\) の対角双加群の長さ \(N\) の有限射影分解を得る。
 よって \(\operatorname{gldim}A\le N\) である。
 各factorの \(\CC\) のKoszul分解とそのdualのK\"unneth公式は、
 表示したExt公式と左側の同じ式を与える。
 従ってglobal dimensionはちょうど \(N\) であり、AS-Gorenstein条件も成立する。
 Koszul性とHilbert seriesはtensor分解から従い、GK dimensionは \(N\) である。

 \(A\) は \(\CC[g_1,\ldots,g_r,u_1,\ldots,u_m]\) 上自由である。
 この中心多項式環で \(W\) は \(g_1\) に関してmonic一次式なので、\(A\) 上でも左右正則である。
 中心正則元によるchange of ringsにより、商のinjective dimensionは一つ減り、
 Gorenstein parameterは三減る。
\end{proof}

\subsection{多ブロックHesse模型からの埋込み}

\(r\ge1\)、\(m\ge0\) とし、独立した変数 \(x_j,y_j,z_j\) に対して
\[
 S_{h_j}=\CC\langle x_j,y_j,z_j\rangle/
 ([x_j,y_j]-h_jz_j^2,[y_j,z_j]-h_jx_j^2,[z_j,x_j]-h_jy_j^2),
\]
\[
 G_j=x_j^3+y_j^3+z_j^3-2h_j^2\sym(x_jy_jz_j),\qquad
 F_j=x_j^3+y_j^3+z_j^3+h_j^2x_jy_jz_j
\]
と置く。異なるブロックの元は可換とし、
\[
 A_{\mathbf h,m}=(S_{h_1}\otimes\cdots\otimes S_{h_r})[u_1,\ldots,u_m],
 \qquad W_{\mathbf h,m}=\sum_{j=1}^rG_j+\sum_{a=1}^m u_a^3,
\]
\[
 F_{\mathbf h,m}=\sum_{j=1}^rF_j+\sum_{a=1}^m u_a^3,
 \qquad Y_{\mathbf h,m}=V(F_{\mathbf h,m})\subset\PP^{3r+m-1}
\]
とする。

\begin{proposition}[多ブロックの完全忠実比較]\label{prop:multiblock-embedding}
 全ての \(j\) について \(h_j^6\ne-27\) と仮定する。
 degree shiftと三次periodicityを保つdg完全忠実関手
 \[
 \MF^{\ZZ}(\CC[x_j,y_j,z_j,u_a],F_{\mathbf h,m})
 \longrightarrow
 \MF^{\ZZ}(A_{\mathbf h,m},W_{\mathbf h,m})
 \]
 が存在する。\(n=3r+m-2\ge3\) なら、左辺は
 \((\Ku(Y_{\mathbf h,m}),\rho_{Y_{\mathbf h,m}})\) である。
 この命題は、右辺の全対象が像に属するとは主張しない。
\end{proposition}
\begin{proof}
 \(s=r+m\) とし、各非可換ブロックと各一変数ブロックに別の \(e_i\) を割り当てる。
 系\ref{cor:block}は各 \(j\) について
 \[
 \MF^{\ZZ}(\CC[x_j,y_j,z_j],F_j)
 \simeq\MF^{\ZZ}(S_{h_j},G_j)
 \]
 をdegree shiftとperiodicityを保って与える。
 これらをtensorし、\(\MF^{\ZZ}(\CC[u_a],u_a^3)\) には恒等を用いる。
 非可換側へは補題\ref{lem:multiblock-ff}の完全忠実関手を用いる。

 可換側の外部積は本質的全射でもある。
 これは可換なsmooth affine ambientに対するequivariant Thom--Sebastiani定理
 \cite[Theorem 5.15, Corollary 5.18]{BFK}を適用したものである。
 二因子の場合は \(G=H=\mathbb G_m\)、potentialのcharacterを双方 \(t\mapsto t^3\)
 と取る。異なる二因子上でのcharacterの差 \((-3,3)\) は非torsionであり、
 \(G\times_{\mathbb G_m}H\) のcharacter groupは
 \(\ZZ^2/\langle3e_2-3e_1\rangle\) となる。
 各potentialは斉次なので、Euler公式により和のcritical locusは
 個々の零点の積に含まれる。従って同定理の支持条件を除くことができる。
 反復して \(\Gamma_s\) を得る。
 ここで引用している定理のambientは全て可換であり、非可換ambientへ拡張して引用してはいない。

 これで \(\Gamma_s\)-次数付きの完全忠実関手を得た。
 \(\Gamma_s\to\ZZ\)、\(e_i\mapsto1\) の核は
 \(T_s\cong(\ZZ/3)^{s-1}\) である。
 細分次数の圏から通常次数の圏へ戻す双対同変化では、
 \(T_s\) の生成元は因子のdegree shiftの差
 \(\rho_i\rho_1^{-1}\) として作用する。
 全ての比較はこれらと三次periodicityを保つので、
 有限群関係とそのcoherenceを保つ同変関手となる。
 同変対象間のHomは元のHomの有限群不変部分であり、標数零では不変部分を取る操作は完全である。
 従って完全忠実性は同変化で保たれる。
 \cite[Theorem 4.2 and Section 8]{Elagin}の双対同変化を用いて通常次数へ降ろせば、
 求める関手を得る。\((e_1)\) との両立は降下後の \((1)\) との両立になる。

 最後に \(Y_{\mathbf h,m}\) のsmoothnessを確認する。
 特異点では全ての \(u_a=0\) であり、各三変数ブロックの偏微分が同時に零となる。
 \(F_j\) の射影零点がsmoothなので、そのaffine critical locusは原点だけである。
 従って全座標が零となり、射影特異点は存在しない。
 Orlov同値により最後の主張を得る。
\end{proof}

\begin{remark}
 命題\ref{prop:tensor-sklyanin-ambient}により、この独立ブロック族の周囲は
 平行移動の位数によらずnoetherian AS-regularである。
 中心正則元の商についてのMF--特異圏比較とOrlov windowにより、
 右辺は \(A_{\mathbf h,m}/(W_{\mathbf h,m})\) の残余成分と同一視できる。
 この周囲の正則性は、残余圏そのもののpropernessや比較の本質的全射性とは別である。
\end{remark}

\subsection{Euler導分と中心支持への帰着}

\begin{lemma}[ブロックpotentialの零ホモトピー]\label{lem:block-euler}
 任意の有限自由factorization \(P\) of \(W_{\mathbf h,m}\) 上で、
 各 \(G_j:P\to P(w)\) と各 \(u_a^2:P\to P(2e_{r+a})\) は零ホモトピックである。
 この主張は \(\Gamma_s\)-次数付きでも通常の整数次数付きでも成り立つ。
\end{lemma}
\begin{proof}
 \(j\) 番目のブロックの次数を数え、他のブロックと \(u_a\) では零となるEuler導分を
 \(\delta_j\) と書く。Sklyanin関係式はそのブロックについて斉次なので、
 \(\delta_j\) は自由代数から商代数へ降りる。
 \(\delta_j(W_{\mathbf h,m})=3G_j\) である。
 factorizationの有限自由加群に斉次基底を選び、その微分行列 \(d\) の係数に
 \(\delta_j\) を作用させる。積のLeibniz則から
 \[
 d\,\delta_j(d)+\delta_j(d)\,d
 =\delta_j(d^2)
 =3G_j\,\Id_P.
 \]
 従って \(\frac13\delta_j(d)\) が所要のホモトピーである。
 \(\delta_j\) は代数の元の \(\Gamma_s\)-次数を変えない。
 \(3e_j=w\) という次数群の関係があっても、
 元のtensor factorsから定義した導分の存在は失われない。
 同じ計算を中心変数の導分 \(\partial_{u_a}\) に対して行えば
 \[
 d\,\partial_{u_a}(d)+\partial_{u_a}(d)\,d=3u_a^2\Id_P
 \]
 を得る。基底に依存するのはホモトピーの選択であり、零ホモトピックという結論ではない。
\end{proof}

\begin{proposition}[生成性の中心商への帰着]\label{prop:block-support}
 全ての \(h_j^6\ne-27\) とし、\(A=A_{\mathbf h,m}\) と置く。
 命題\ref{prop:tensor-sklyanin-ambient}により、この \(A\) はnoetherian AS-regularである。さらに
 \[
 R=A/(W_{\mathbf h,m}),\qquad
 C_j=S_{h_j}/(G_j),\qquad C=\bigotimes_{j=1}^r C_j
 \]
 と置く。\(C\) は \(R\) の商とみなす。
 次の等式が、冪等完備化した \(\Gamma_s\)-次数付き特異圏で成り立つ。
 \[
 \Dsg^{\Gamma_s}(R)
 =\thick\!\left(
 \operatorname{Res}\bigl(D^b(\operatorname{gr}^{\Gamma_s}C)\bigr)
 \right).
 \tag{\(\ddagger\)}\label{eq:block-support}
 \]
 また、\(C\) は \(R\)-perfectなので、restrictionは特異圏の間の関手にも降りる。
\end{proposition}
\begin{proof}
 各 \(G_j\) は中心正則元である。\(S_{h_j}/G_jS_{h_j}\) の斉次基底を持ち上げた
 ベクトル部分空間を \(V_j\subset S_{h_j}\) とすると、次数に関する帰納法と正則性から
 \[
 V_j\otimes_\CC\CC[t_j]\xrightarrow{\sim}S_{h_j},\qquad
 v\otimes t_j^b\longmapsto vG_j^b
 \]
 は中心 \(\CC[t_j]\)-加群の同型である。
 従って \(A\) は中心多項式環
 \(\CC[t_1,\ldots,t_r,u_1,\ldots,u_m]\) 上自由である。
 \(W\) はこの環の \(\sum_jt_j+\sum_au_a^3\) に対応するので正則であり、
 \(R\) は \(\CC[t_1,\ldots,t_{r-1},u_1,\ldots,u_m]\) 上自由となる。
 よって
 \[
 u_1,\ldots,u_m,G_1,\ldots,G_{r-1}
 \]
 は \(R\) の中心正則列で、その商は \(C\) である。
 Koszul分解は \(C\) の有限自由 \(R\)-分解を与える。
 これでrestrictionがperfect対象をperfect対象へ送ることも従う。

 まず冪等完備化前の \(\Dsg^{\Gamma_s}(R)\) の対象 \(N\) を取る。
 中心正則元によるAS-regular代数の商に対するMF比較
 \cite[Theorem 5.8]{CCKM}を使えば、\(N\) はfinite-free factorizationのcokernelのshiftで表せる。
 この比較は通常の整数次数で記述されているが、
 \(\Gamma_s\simeq\ZZ\oplus\Gamma_{s,\mathrm{tors}}\) と標数零を用い、
 各有限自由分解とホモトピーを斉次成分ごとに取る同じ構成で
 \(\Gamma_s\)-次数を保った比較を得る。
 補題\ref{lem:block-euler}より、\(N\) は
 \[
 K(G_1,\ldots,G_{r-1},u_1^2,\ldots,u_m^2;N)
 \]
 のshiftと内部twistの直和因子となる。
 ここで \(K\) は通常の中心Koszul複体である。
 この複体の各cohomologyは \(G_1,\ldots,G_{r-1},u_1^2,\ldots,u_m^2\) に殺される。
 さらに \((u_1,\ldots,u_m)^{m+1}\subset(u_1^2,\ldots,u_m^2)\) なので、
 各cohomologyの有限filtrationの商は全 \(u_a\) と全 \(G_j\) に殺される \(C\)-加群である。
 \(G_r\) の消滅は \(W=0\) から従う。
 noetherian性によりこれらは有限生成であり、有限個のtruncation triangleとfiltrationから
 未完備の特異圏の各対象が右辺に入る。
 最後に直和因子閉包を取れば、\eqref{eq:block-support} を得る。
\end{proof}

\subsection{必要な生成性の正確な内容}

命題\ref{prop:block-support}は、全ての \(C\)-加群のrestrictionが
\(C_1,\ldots,C_r\) の外部積から生成されることまでは述べていない。
十分条件の一つは
\[
 D^b(\operatorname{gr}^{\Gamma_s}C)
 =\thick\left\{
 \operatorname{coarsen}(M_1\boxtimes\cdots\boxtimes M_r)(\gamma)
 :M_j\in D^b(\operatorname{gr}^{\ZZ}C_j),\ \gamma\in\Gamma_s
 \right\}
 \tag{G}\label{eq:block-generation}
\]
である。\(\gamma\) の全shiftは、中心変数が零となった後にも残る次数類を含めるために必要である。
以下で用いるsyzygyへの置換を正当化するため、
\(A_{\ne i}=(\bigotimes_{j\ne i}S_{h_j})[u_1,\ldots,u_m]\) とし、
他因子の外部積に全ての \(u_a\) が零で作用する \(A_{\ne i}\)-加群を \(M_{\ne i}\) とする。
\(S_{h_i}\) の \(\CC[G_i]\) 上の自由性と \(W=0\) による \(G_i\) の消去から、
\(R\) は \(A_{\ne i}\) 上自由であり、
\[
 \operatorname{Res}(C_i\boxtimes M_{\ne i})
 \simeq M_{\ne i}\otL_{A_{\ne i}}R.
\]
各 \(S_{h_j}\) 上の有限自由分解と \(u_a\) のKoszul分解を外部tensorすると、
\(M_{\ne i}\) は \(A_{\ne i}\)-perfectであり、上のrestrictionは \(R\)-perfectである。
従って一因子が \(C_i\)-perfectである項はrestriction後に零となり、
各因子を有限syzygyへ置き換えてよい。有界複体の場合もtruncation triangleから従う。
各因子の有限syzygyを取りMFで表した後、外部tensor productの有限ambient分解を
\(W\) の零ホモトピーで周期化すると、そのrestrictionの安定化は外部積factorizationになる
（全体のcohomological shiftは生成性に影響しない）。
したがって (G) が成立すれば、命題\ref{prop:block-support}と
補題\ref{lem:multiblock-ff}により非可換側の外部積もMorita同値となり、
命題\ref{prop:multiblock-embedding}を圏同値へ強められる。
より弱く、(G)の等式をrestriction後の \(\Dsg^{\Gamma_s}(R)\) だけで示しても十分である。

しかし (G) を、multigraded tailsの可換な積の幾何だけから推論することはできない。
\(C\) は元来 \(\ZZ^r\)-次数を持つが、ここで必要なのは
\(\Gamma_s\)-次数付き加群である。\(\ZZ^r\to\Gamma_s\) の核は
\(3(e_i-e_1)\) からなる無限格子を含み、有限群降下ではない。
\(\Gamma_s\)-次数付き加群は一般には \(\ZZ^r\)-次数へ持ち上がらない。
最も単純には、\(\CC[a,b]\) に \(\deg a=e_1\)、\(\deg b=e_2\)、
\(3e_1=3e_2\) を入れると、\(\CC[a,b]/(a^3-b^3)\) は \(\Gamma_2\)-次数付きだが
\(\ZZ^2\)-次数へ持ち上がらない。そのannihilatorが \(\ZZ^2\)-斉次でないためである。
この例は (G) 自体を反証するものではなく、個々の対象を細分次数へ持ち上げる証明法を反証する。

さらにmultigraded tailsでは各座標軸上のtorsionが消える。
生成性を完成させるには、これらのtorsionと上記の非liftable加群を、
restriction後も含めて制御する必要がある。
第\ref{sec:finite-center}節では、高々一つを除く因子の有限中心性でこの両方を扱い、
一つの任意位数ブロックを許して本質的全射性を証明する。
二つ以上の無限位数因子の一般の場合は未解決として保持する。

\begin{remark}[別の十分条件]
 \(\mathcal A\to\cD\) をsmooth proper圏からのdg完全忠実関手とする。
 \(\cD\) がproperでSerre関手を持ち、埋込みがSerre関手と両立すれば、
 \(\mathcal A\) のsaturated性から左右随伴が存在し、その直交成分は左右一致する。
 したがって \(\cD\) に非自明な直交直和分解がなければ埋込みは同値である。
 これは上の生成性とは別の完成条件となる。
 ただし複数ブロックの全MFについて、必要なpropernessと直交分解不能性を
 ここで証明したわけではない。
\end{remark}

\section{一方有限中心の外部積生成性と多ブロックの可換モデル}\label{sec:finite-center}

この節では生成条件 \eqref{eq:block-generation} を、高々一つを除くブロックが
有限位数である場合に証明する。任意位数の一因子は係数代数として残し、
有限中心の因子を一つずつ追加する。
非可換係数付きBeilinson分解、中心のNoether正規化、分離的な稠密開集合と
その補集合の帰納法を用いる。
\(\ZZ^r\)-次数へ持ち上がらない対象と軸上のtorsionも初めから含む。
有限中心の場合の基礎的な局所化・生成法については \cite{ELS} と比較されたい。

\subsection{斉次中心局所化と稠密開集合}

\(\operatorname{gr}^{\Lambda}D\) は有限生成の \(\Lambda\)-次数付き右 \(D\)-加群の圏とする。
ここで \(\Lambda\) は有限生成アーベル群であり、torsionを持っていてよい。
全てのthick閉包は、有限直和、shift、cone、直和因子に関する閉包を表す。

\begin{lemma}[斉次中心局所化]\label{lem:graded-localization}
 \(D\) をnoetherianな \(\Lambda\)-次数付き代数、
 \(f\in Z(D)\) を斉次元とする。
 有界導来圏の局所化は
 \[
 D^b(\operatorname{gr}^{\Lambda}D)\big/
 D^b_{f\text{-tor}}(\operatorname{gr}^{\Lambda}D)
 \simeq D^b(\operatorname{gr}^{\Lambda}D_f)
 \]
 を与える。左辺の核は、各cohomologyが \(f\) の冪に殺される対象からなる。
 この核は \(D/(f)\) 上の有限生成 \(\Lambda\)-次数付き加群のrestrictionで生成される。
\end{lemma}
\begin{proof}
 局所化は中心元についての平坦な局所化であるからcohomologyと可換である。
 有限生成加群が局所化後に零なら、その有限個の生成元を同時に殺す \(f\) の冪がある。
 \(f\)-adic filtrationと有界truncation triangleにより核についての主張を得る。

 局所化後の有限生成次数付き加群は、有限個の斉次生成元の
 \(D\)-spanを取ることによって、局所化前の有限生成次数付き加群から得られる。
 射については有限生成加群の有限生成自由分解の必要な有限部分を使う。
 中心局所化と \(\operatorname{Ext}^q\) は可換であり、
 次数零の射は \(f\) の冪を分母としてclearできる。
 同じことは有界複体にも有限個のtruncationによって成立する。
 次数を保つ分母の射は \(f^N:M(-N\deg f)\to M\) と書く。
 そのconeは \(f\)-torsionである。
 従って上記Verdier商は局所化後の圏へ完全忠実に写り、
 加群とその有限個のconeを全て含むので本質的全射である。
 この議論で次数を細分する必要はない。
\end{proof}

\begin{lemma}[有限中心代数の良い稠密開集合]\label{lem:generic-separable}
 \(k\) を標数零の体とする。
 \(D\) を連結非負次数付き \(k\)-代数で、有限生成の可換次数付き中心部分代数
 \(Z\subset Z(D)\) 上有限なものとする。
 \(Z\) がreducedであると仮定する。
 ある斉次元 \(f\in Z\) と斉次両側イデアル \(I\subset D\) が存在し、次を満たす。
 \begin{enumerate}
 \item \(f\) は \(Z\) の全てのminimal primeを避ける。
       特に \(Z/(f)\ne0\) なら \(\dim Z/(f)<\dim Z\) である。
 \item \(Z_f\) は \(k\) 上smoothな有限型可換環である。
 \item \(I_f\) は冪零であり、\(Q_f=(D/I)_f\) は
       \(Z_f\) 上有限射影かつ分離的な代数である。
 \end{enumerate}
 \(\dim Z=0\) の場合には \(Z=k\)、\(f=1\) と取れる。
\end{lemma}
\begin{proof}
 \(K\) を \(Z\) の全商環とする。
 \(Z\) はreduced noetherianなので、\(K\) はminimal primeの分数体の有限積である。
 \(D_K\) は \(K\) 上有限次元であり、そのJacobson radical \(J\) は冪零である。
 \(I\) を \(D\to D_K\) による \(J\) の逆像とする。
 次数を与える \(\mathbb G_m\)-作用はminimal primeの集合と \(J\) を保つので、
 \(I\) は斉次である。
 \(I^N\) は十分大きい \(N\) で \(K\) 上零となり、
 有限生成性から、ある全商環の分母を反転すれば \(I^N=0\) となる。

 \(Q_K=D_K/J\) は半単純である。
 標数零なので、各simple factorの中心は対応する成分体の有限分離拡大であり、
 \(Q_K\) は \(K\) 上分離的である。
 generic freenessを用いて、ある稠密開集合上で \(Q\) を有限射影とする。
 分離性は、積の写像 \(Q\otimes_ZQ^{\mathrm{op}}\to Q\) の
 \(Q\)-双加群としてのsection、すなわち分離冪等元の存在で表せる。
 generic fiberにある分離冪等元と、それが満たす有限個の等式の分母を払えば、
 さらに小さい稠密開集合上でこの分離冪等元が存在する。
 また \(Z\) のsmooth locusは全てのgeneric pointを含む。

 以上の良い開集合は \(\mathbb G_m\)-作用で不変に取れる。
 実際、\(I^N\) の消滅、局所自由性、分離性、smoothnessは全て
 次数作用によって保たれる開条件である。
 補集合のイデアルは斉次で、どのminimal primeにも含まれない。
 \(Z_0=k\) と \(k\) の無限性により、そのイデアル内から
 全てのminimal primeを避ける一つの斉次元 \(f\) を選べる。
 必要なら候補を冪にして次数を揃え、一般の線型結合を取ればよい。
 \(\dim Z>0\) のとき零次元の既約成分はない
 （全ての既約成分は連結な錐の頂点を含む）ので、この操作に例外はない。
 \(Z=k\) の場合は有限次元代数のradicalと半単純商を用いればよい。
\end{proof}

\begin{lemma}[分離代数と射影分解]\label{lem:separable-perfect}
 \(Z\) をsmooth affineな有限型 \(k\)-代数、
 \(Q\) を \(Z\) 上有限射影かつ分離的な代数とする。
 このとき \(Q\) は有限global dimensionを持つ。
 特に \(Q\) に任意のアーベル群次数を与えたとき、
 \(D^b(\operatorname{gr}^{\Lambda}Q)\) は \(Q(\lambda)\)
 \((\lambda\in\Lambda)\) で生成される。
\end{lemma}
\begin{proof}
 任意の右 \(Q\)-加群 \(M\) は、分離冪等元により
 \(M\otimes_ZQ\to M\) の \(Q\)-線型sectionを持つ。
 \(Z\) は有限global dimensionを持ち、\(Q\) は \(Z\)-flatなので、
 \(M\) の \(Z\)-射影分解を \(-\otimes_ZQ\) で送ると
 \(M\otimes_ZQ\) の有限 \(Q\)-射影分解を得る。
 従ってその直和因子 \(M\) も有限射影次元を持つ。

 次数付き有限生成加群に有限生成斉次自由分解を作ると、十分高いsyzygyは
 ungradedには射影的である。
 自由加群からこのsyzygyへの次数零の全射のungraded sectionを取り、
 sectionの次数零成分を抜き出すと次数付きsectionになる。
 よってsyzygyは次数付き有限自由加群の直和因子であり、
 元の加群は有限個の \(Q(\lambda)\) から生成される。
\end{proof}

\subsection{非可換係数の多項式拡大と一方有限中心の生成定理}

ここでは係数代数の有限中心性を仮定しない。
これが一つの無限位数ブロックを許すための要点である。
以下、\(k\) は標数零の体、\(\Lambda\) は有限生成アーベル群とする。

\begin{lemma}[非可換係数を持つ多項式拡大]\label{lem:polynomial-coefficient}
 \(A\) を左右noetherianな \(\Lambda\)-次数付き \(k\)-代数とする。
 \(P=k[t_1,\ldots,t_d]\) の各変数に任意の \(\Lambda\)-次数を与え、
 \(A\otimes_kP\) には和の次数を与える。このとき
 \[
 D^b(\operatorname{gr}^{\Lambda}(A\otimes_kP))
 =\thick\{M\otimes_kP:\ M\in\operatorname{gr}^{\Lambda}A\}.
 \]
 \(P\) の斉次元 \(f\) を反転した \(P_f\) についても同じ結論が成立する。
 \(A\) が有限global dimensionを持つという仮定はない。
\end{lemma}
\begin{proof}
 中心多項式拡大のHilbert基底定理により \(A\otimes P\) は左右noetherianである。
 \(d=0\) は自明なので \(d>0\) とする。
 有限生成 \(\Lambda\)-次数付き \(A\otimes P\)-加群 \(T\) の有限斉次表示を取る。
 射影用の独立な整数次数を導入し、\(t_0,t_1,\ldots,t_d\) の射影次数を全て1、
 \(A\) の射影次数を0とする。\(\Lambda\)-次数は
 \(\deg_\Lambda(t_0)=0\)、\(\deg_\Lambda(t_i)=\deg_\Lambda(t_i\in P)\) とする。
 表示の各関係を \(t_0\) で斉次化すると、\(\PP^d_k\) 上の
 \(\mathcal A=A\otimes_k\OO_{\PP^d}\)-加群 \(\cF\) が得られる。
 これは \(\mathcal A\) 上有限表示で、\(t_0\ne0\) への制限が \(T\) である。
 各関係の斉次化では \(\Lambda\)-次数を変えていない。
 以後の \(\Lambda\)-次数は対角化可能群 \(D(\Lambda)\) による線型化としても解釈できる。

 必要なcohomologyの有限性を確認する。
 \(\mathcal A\) 上のcoherent加群には有限個の
 \(\mathcal A(-a)(\lambda)\) からの全射があり、核もcoherentである。
 これは中心変数による有限斉次表示とnoetherian性から従う。
 標準affine被覆の\v Cech複体は長さ \(d\) なので、cohomologyは次数 \(0,\ldots,d\)
 に限られる。一方、\(H^q(\PP^d,\mathcal A(a))
 =A\otimes_kH^q(\PP^d,\OO(a))\) は有限生成 \(A\)-加群である。
 上記全射とその核の長完全列を用い、\(q=d,d-1,\ldots,0\) の順に帰納すれば、
 全てのcoherent \(\mathcal A\)-加群の \(H^q\) が有限生成 \(A\)-加群と分かる。
 この議論は線型化と両立するので、\(\cF\) に有限階数線束をテンソルしたものの
 導来大域切断は \(D^b(\operatorname{gr}^{\Lambda}A)\) に属する。
 ここでcoherentは \(\mathcal A\) 上の意味であり、\(\OO\) 上有限とは仮定しない。

 Beilinsonの対角分解 \cite{Beilinson} を、因子の順序を適切に選んで
 \[
 0\longrightarrow\Omega^d(d)\boxtimes\OO(-d)
 \longrightarrow\cdots\longrightarrow
 \Omega^1(1)\boxtimes\OO(-1)
 \longrightarrow\OO\boxtimes\OO
 \longrightarrow\OO_\Delta\longrightarrow0
 \]
 と書く。これは \(D(\Lambda)\)-線型化を保つ分解である。
 \(\mathcal A\) の係数を保ったまま \(\cF\) に作用させると、\(\cF\) は
 \[
 R\Gamma(\PP^d,\cF\otimes_{\OO}\Omega^i(i))\otimes_k\OO(-i),
 \qquad 0\le i\le d,
 \]
 の有限個のconeから作られる。
 対角のKoszul分解は各射影に関してflatであるから、\(\cF\) の
 \(\OO\)-flatnessは必要ない。非可換な \(A\)-作用も全ての操作と可換である。
 係数複体は前段落により \(D^b(\operatorname{gr}^{\Lambda}A)\) に属する。
 \(t_0\ne0\) に制限すると、\(t_0\) の \(\Lambda\)-次数が0であるため
 \(\OO(-i)\) は次数を保って自明化される。
 よって \(T\) は \(M\otimes P\) から生成される。
 係数複体を有限個のcohomology加群に分けるには通常のtruncationを用いればよい。
 有界複体全体についても同じtruncationから結論する。

 \(P_f\) の場合は補題\ref{lem:graded-localization}による局所化を適用する。
 有限生成次数付き加群の有限表示は局所化前に持ち上がるので、
 局所化後の生成対象はちょうど \(M\otimes P_f\) である。
\end{proof}

\begin{lemma}[多項式中心上の分離的な稠密開集合]\label{lem:polynomial-separable}
 \(D\) を連結局所有限な非負次数付き \(k\)-代数で、有限型可換斉次中心部分代数
 \(Z\) 上有限なものとする。\(d=\dim Z\) と置く。
 正次数の斉次変数を持つ多項式部分代数 \(P=k[t_1,\ldots,t_d]\subset Z\)、
 非零斉次元 \(f\in P\)、斉次両側イデアル \(J\subset D\) を選んで、
 \(D\) が \(P\) 上有限、\(J_f\) が冪零、\(Q=(D/J)_f\) が
 \(P_f\) 上有限射影かつ分離的となるようにできる。
 分離冪等元は元の整数次数で次数0に取れる。
\end{lemma}
\begin{proof}
 斉次Noether正規化によって \(Z\) は正次数変数の多項式部分代数 \(P\) 上有限となる。
 \(Z\) のreducednessは必要ない。\(K=\operatorname{Frac}P\) とすると、
 \(D_K\) は有限次元 \(K\)-代数である。
 そのradicalの \(D\) への逆像を \(J\) とする。
 次数のスカラー作用は \(P\)、\(K\) とradicalを保つので \(J\) は斉次である。
 十分大きい \(a\) で \((J^a)_K=0\) だから、有限生成性によって
 ある非零元を反転すると \(J^a=0\) となる。
 半単純商 \(D_K/\operatorname{rad}(D_K)\) は標数零の体上分離的である。
 generic freenessと分離冪等元の有限個の等式の分母払いにより、
 ある稠密開集合上で \(D/J\) は有限射影かつ分離的となる。

 これらの条件の成立する開集合は次数作用に関して不変である。
 その補集合のイデアルは斉次であり、\(P\) は整域なので、そこから
 非零斉次元 \(f\) を一つ選べる。
 \(d=0\) の場合は \(P=k\)、\(f=1\) とし、有限次元代数のradicalを使う。
 最後に、分離冪等元を元の整数次数で分解し、次数0の成分を取る。
 \(\mu(e)=1\) と \((q\otimes1)e=e(1\otimes q)\) は斉次 \(q\) ごとに次数に分解されるから、
 次数0成分も同じ条件を満たす。これが必要な次数付き分離冪等元である。
\end{proof}

\begin{theorem}[一方だけが有限中心である外部積生成]\label{thm:one-sided-boxes}
 \(A\) を左右noetherianな \(\Lambda\)-次数付き \(k\)-代数とする。
 \(D\) を連結局所有限な非負次数付き \(k\)-代数で、
 有限型可換斉次中心部分代数 \(Z\) 上有限なものとする。
 任意の準同型 \(\beta:\ZZ\to\Lambda\) により \(D\) の次数を粗化し、
 \(A\otimes_kD\) には和の次数を与える。このとき
 \begin{equation}\label{eq:one-sided-boxes}
 \begin{split}
 D^b(\operatorname{gr}^{\Lambda}(A\otimes_kD))
 =\thick\{\,M\boxtimes\operatorname{coarsen}_\beta N:\qquad\qquad\qquad\quad\\
 M\in\operatorname{gr}^{\Lambda}A,\quad
 N\in\operatorname{gr}^{\ZZ}D\,\}.
 \end{split}
 \end{equation}
 特に \(A\) の有限中心性、有限global dimension、あるいは
 左辺の対象の細分次数へのliftは仮定しない。
\end{theorem}
\begin{proof}
 補題\ref{lem:polynomial-separable}の \(P\subset Z\) に対して
 \(A\otimes D\) は \(A\otimes P\) 上左右有限である。
 後者はnoetherianなので前者もnoetherianである。
 右辺のthick部分圏を \(\mathcal B\) とし、\(d=\dim Z\) に関して帰納する。
 \(d=0\) では \(D\) は有限次元であり、連結性から \(D_{>0}\) は冪零で商は \(k\) である。
 \(D_{>0}\)-adic filtrationの各層は有限生成 \(\Lambda\)-次数付き \(A\)-加群であるから、
 \(M\boxtimes k\) からの生成を得る。

 \(d>0\) とし、同補題の \(f,J,Q\) を取る。
 \(f=1\) と取れる場合も以下の開集合の議論はそのまま適用でき、閉集合側は空になる。
 \(J_f\)-adic filtrationにより、\(A\otimes D_f\) 上の全加群は
 \(A\otimes Q\) 上の有限生成加群から生成される。
 そのような加群 \(T\) は \(A\otimes P_f\) 上も有限生成である。
 次数0の分離冪等元 \(e=\sum_\nu x_\nu\otimes y_\nu\) により、乗法写像
 \[
 \operatorname{Res}T\otimes_{P_f}Q\longrightarrow T
 \]
 は \(A\otimes Q\)-線型かつ \(\Lambda\)-次数0のsection
 \(t\mapsto\sum_\nu tx_\nu\otimes y_\nu\) を持つ。
 \(A\) の作用は \(Q\) と可換なので、このsectionは \(A\)-線型でもある。

 補題\ref{lem:polynomial-coefficient}により \(\operatorname{Res}T\) は
 \(M\otimes P_f\) から生成される。
 \(Q\) は \(P_f\) 上射影なので \(-\otimes_{P_f}Q\) はexactであり、
 上のsectionと合わせて \(T\) は \(M\otimes Q\) から生成される。
 \(M\otimes Q\) は \(M\boxtimes(D/J)\) の局所化であるから、
 開集合上の全圏は \(\mathcal B\) の像で生成される。
 ここで \(A\otimes Q\) が有限global dimensionを持つとは主張していない。

 局所化の核は補題\ref{lem:graded-localization}により、
 \(A\otimes(D/(f))\) 上の有限生成次数付き加群から生成される。
 \(D/(f)\) は \(P/(f)\) の像という有限型可換斉次中心上有限である。
 \(f\) が非単元ならこの中心の次元は \(d-1\) 以下であるから、帰納法を適用できる。
 得られる外部積の第二因子は \(D/(f)\)-加群であり、元の \(D\)-加群としても
 元の整数次数を保っている。従って局所化の核は \(\mathcal B\) に含まれる。
 開集合と核の生成をVerdier局所化によって合わせれば \eqref{eq:one-sided-boxes} を得る。
 全てのconeと直和因子は粗化後の圏で取っており、対象の次数のliftは用いない。
\end{proof}

\subsection{任意の粗化次数と有限中心因子の反復}

\begin{theorem}[有限中心代数の外部積生成]\label{thm:finite-center-boxes}
 \(D_1,\ldots,D_r\) を連結局所有限な非負次数付き \(k\)-代数で、
 各 \(D_i\) は有限型可換斉次中心部分代数上有限とする。
 任意の準同型 \(\alpha:\ZZ^r\to\Lambda\) に対し、
 \(D=\bigotimes_iD_i\) に粗化次数を与えると
 \[
 D^b(\operatorname{gr}^{\Lambda}D)
 =\thick\{\operatorname{coarsen}_{\alpha}
 (M_1\boxtimes\cdots\boxtimes M_r)(\lambda):
 M_i\in\operatorname{gr}^{\ZZ}D_i,\ \lambda\in\Lambda\}.
 \tag{\(G_{\mathrm{fin}}\)}\label{eq:finite-center-generation}
 \]
 \(\alpha\) は単射・全射である必要はなく、無限の核を持っていてもよい。
\end{theorem}
\begin{proof}
 \(A=k\) から出発し、\(\beta_i(1)=\alpha(e_i)\) として
 定理\ref{thm:one-sided-boxes}を順に適用する。
 各段階の係数代数は有限型可換中心上有限なので左右noetherianである。
 \(k\) の有限生成 \(\Lambda\)-次数付き加群は有限個の \(k(\lambda)\) の和であり、
 反復によって表示の外部積を得る。
 複体の外部積は体 \(k\) 上exactなので、前段階のthick生成を次段階へ送れる。
\end{proof}

\begin{remark}[nonliftable対象と軸torsion]
 外部積のconeには粗化後の全ての次数零の射を用いる。
 例えば \(k[a,b]/(a^3-b^3)\) は \(3e_1=3e_2\) の次数では
 \(a^3-b^3\) による二つのfree shiftの間のconeで生成される。
 同じ加群が \(\ZZ^2\)-次数へ持ち上がる必要はない。
 また \(f\)-adic filtrationは正次数の中心元に殺される軸上の加群も保持する。
 multigraded tailsへの移行によってこれらを捨てる操作は用いない。
\end{remark}

\subsection{高々一つの無限位数ブロックへの帰結}

次の有限中心性は \cite[Corollary 2.3]{SmithTate} に含まれる既知の結果である。
必要な斉次中心を明示するため、証明を記す。

\begin{lemma}[有限位数THCRの有限中心性]\label{lem:finite-order-thcr}
 \(E\) をsmooth elliptic curve、\(L\) を正次数の線束とし、
 \(\sigma\in\operatorname{Aut}(E)\) が有限位数 \(q\) を持つとする。
 このとき \(B=B(E,L,\sigma)\) は有限型可換斉次中心部分代数上有限である。
\end{lemma}
\begin{proof}
 \(L=\mathcal O_E(D)\) となるdivisor \(D\) を選び、
 \(D_d=D+\sigma^*D+\cdots+(\sigma^{d-1})^*D\) と置く。
 通常のskew Laurent表示により、\(B\) を
 \[
 B_d=H^0(E,\mathcal O_E(D_d))t^d
 \subset k(E)[t,t^{-1};\sigma^*]
 \]
 と書ける。\(\sigma^q=1\) なので \(D_q\) は \(\sigma\)-不変であり、
 \(q\)-Veroneseは可換なsection ring
 \[
 R=B^{(q)}=\bigoplus_{n\ge0}H^0(E,\mathcal O_E(nD_q))t^{qn}
 \]
 になる。\(R\) 上では \(\sigma\) を係数に作用させ、\(t^q\) を固定する。
 \(Z=R^{\langle\sigma\rangle}\) は \(B\) の中心に含まれる。
 実際、\(a\in H^0(E,\mathcal O_E(nD_q))\) が \(\sigma\)-不変なら、
 任意の \(bt^d\in B_d\) に対して
 \[
 (at^{qn})(bt^d)=ab\,t^{qn+d}
 =b\sigma^{d*}(a)t^{qn+d}=(bt^d)(at^{qn})
 \]
 である。

 \(D_q\) は正次数なのでampleである。従って \(R\) は有限型可換代数であり、
 各 \(0\le d<q\) について
 \[
 \bigoplus_{n\ge0}H^0(E,\mathcal O_E(nD_q+D_d))t^{qn+d}
 \]
 は \(R\) 上有限なsection moduleである。
 これらの有限直和が \(B\) なので \(B\) は \(R\) 上有限である。
 有限群の不変式環について \(R\) は有限型環 \(Z\) 上有限であり、
 従って \(B\) も \(Z\) 上有限である。
\end{proof}




\begin{corollary}[一つの任意位数因子を含む生成条件]\label{cor:PI-generation}
 全ての \(j\) で \(h_j^6\ne-27\) とし、並べ替えた後の
 \(\sigma_j=t_{p_j}\ (j\ge2)\) のみ有限位数と仮定する。
 このとき \(C_j=S_{h_j}/(G_j)\)、\(C=\bigotimes_jC_j\) に対して、
 条件 \eqref{eq:block-generation} は成立する。
 中心変数を零とした後に残る \(\Gamma_s\) の全shiftも含む。
\end{corollary}
\begin{proof}
 \(C_1\) はnoetherianであり、その \(\Gamma_s\)-次数は
 単射 \(\ZZ e_1\hookrightarrow\Gamma_s\) によって与えられる。
 任意の有限生成 \(\Gamma_s\)-次数付き \(C_1\)-加群は
 \(\Gamma_s/\ZZ e_1\) の剰余類ごとの部分加群に分解する。
 有限個の生成元の次数が関わる剰余類は有限個だけであるから、
 各剰余類の代表によるshiftを戻せば、元の整数次数を持つ \(C_1\)-加群の有限和となる。
 ここで細分しているのは一因子上の加群だけであり、積の一般の加群ではない。

 \(j\ge2\) では補題\ref{lem:finite-order-thcr}によって \(C_j\) は有限中心上有限である。
 定理\ref{thm:one-sided-boxes}を
 \(A=C_1\)、\(D=C_2\)、\(\beta(1)=e_2\) に適用し、次に係数代数を
 \(C_1\otimes C_2\) として反復する。
 各段階で左右noetherian性は同定理の有限中心多項式拡大によって保証される。
 外部積は体上exactなので、前段階で得たthick生成をそのまま次段階へ送れる。
 これで全ての \(C_j\) について条件 \eqref{eq:block-generation} が得られる。
\end{proof}

以下では定理\ref{thm:multiblock-model}の仮定の下で
\[
 \begin{gathered}
 N=3r+m,\qquad A=\left(\bigotimes_{j=1}^rS_{h_j}\right)[u_1,\ldots,u_m],\\
 W=\sum_jG_{h_j}+P(u),\qquad R=A/(W)=B_{\mathbf h,P}.
 \end{gathered}
\]
と置く。\(m=0\) の場合は \(P=0\) とし、中心変数の因子を省く。

\begin{proof}[定理\ref{thm:multiblock-model}の証明]

 \emph{周囲の代数と中心元の正則性.}
 命題\ref{prop:tensor-sklyanin-ambient}により、\(A\) は左右noetherianな
 Koszul AS-regular代数で、global dimensionとGorenstein parameterは \(N\) である。
 \(W\) は中心正則元である。以下の支持帰着のために、中心上の自由性も具体的に用いる。

 \(C_j\) の斉次基底の持ち上げを \(V_j\subset S_{h_j}\) とすると、
 正則性と次数に関する帰納法により
 \[
 V_j\otimes\CC[t_j]\xrightarrow{\sim}S_{h_j},\qquad
 v\otimes t_j^a\longmapsto vG_j^a
 \]
 は中心 \(\CC[t_j]\)-加群としての同型である。
 従って \(A\) は中心多項式環
 \(Z_0=\CC[t_1,\ldots,t_r,u_1,\ldots,u_m]\) 上自由である。
 \(W\) は \(\sum_jt_j+P(u)\in Z_0\) の作用なので、左右正則である。
 その商 \(R\) は
 \(\CC[t_1,\ldots,t_{r-1},u_1,\ldots,u_m]\) 上自由であり、
 \[
 u_1,\ldots,u_m,G_1,\ldots,G_{r-1}
 \]
 は \(R\) の中心正則列で、商は \(C\) である。

 \emph{全生成性.}
 \(m>0\) では中心多項式環全体 \((T,P)=(\CC[u],P)\) を一つの因子とし、
 \[
 \Gamma=\ZZ^{r+1}/\langle3e_j-3e_0:1\le j\le r\rangle
 \]
 を用いる。\(\deg(u_a)=e_0\)、Sklyanin因子の次数は \(e_j\) である。
 \(m=0\) では \(e_0\) を省き、\(\Gamma=\Gamma_r\) とする。
 finite-free factorizationの微分を \(d\)、第 \(j\) ブロックのEuler導分を
 \(\delta_j\) とすれば、補題\ref{lem:block-euler}と同じ計算により
 \[
 d\delta_j(d)+\delta_j(d)d=3G_j\operatorname{id}.
 \]
 中心変数についても行列係数を微分することで
 \[
 d\,\partial_{u_a}(d)+\partial_{u_a}(d)\,d
 =\partial_{u_a}P\operatorname{id}
 \]
 を得る。これらは斉次のnull-homotopyである。

 \(J_P=(\partial_{u_1}P,\ldots,\partial_{u_m}P)\) とする。
 孤立性より、ある \(b\) について \((u_1,\ldots,u_m)^b\subset J_P\) である。
 命題\ref{prop:block-support}の証明で \(u_a^2\) を \(\partial_{u_a}P\) に置き換えると、
 各対象は \(G_1,\ldots,G_{r-1},\partial_{u_1}P,\ldots,\partial_{u_m}P\)
 に関するKoszul複体のshift・内部twistの直和因子になる。
 各cohomologyの \((u)\)-adic有限filtrationの商は全ての \(u_a\) に殺される。
 それらは \(G_1,\ldots,G_{r-1}\) にも殺され、\(W=0\) と \(P(0)=0\) から
 \(G_r\) にも殺される。従って
 \[
 \Dsg^\Gamma(R)
 =\thick\operatorname{Res}D^b(\operatorname{gr}^\Gamma C).
 \]
 \(m=0\) では中心変数についてのKoszul複体とfiltrationを省けばよい。
 全ての \(\Gamma\)-次数付き \(C\)-加群が
 系\ref{cor:PI-generation}により外部積で生成される。

 この外部積のrestrictionのstabilizationを説明する。
 各 \(C_j=S_{h_j}/(G_j)\) の加群を十分高いsyzygyに置き換えて二項ambient自由分解で表し、
 \(m>0\) では \(T/(P)\) 上の残余体 \(k\) についても同じ操作をする。
 一つのhypersurface因子においてperfectな項は、restriction後に \(R\)-perfectとなる。
 この点を右加群の規約で明示する。選んだambientを \(A_i\)、potentialを \(f_i\)、
 その商を \(D_i\)、他のambientの積を \(A'\)、potentialの和を \(g'\) と書く。
 \(A_i\cong V_i\otimes\CC[f_i]\) と \(f_i=-g'\) の消去により、
 \(R=(A_i\otimes A')/(f_i+g')\) は \(A'\) 上左右とも自由である。
 他因子の外部積 \(M'\) は \(M'g'=0\) を満たし、各ambient上の有限自由分解の
 tensor productによって \(A'\)-perfectである。従って
 \[
 \begin{aligned}
 M'\otL_{A'}R&\simeq M'\otimes_{A'}R
       \xrightarrow{\sim}\operatorname{Res}(D_i\boxtimes M'),\\
 m'\otimes[a\otimes b]&\longmapsto\bar a\otimes m'b
 \end{aligned}
 \]
 は次数を保つ同型であり、右辺は \(R\)-perfectである。
 有限cone、shift、直和因子を取ることで、\(D_i\) を任意の \(D_i\)-perfect複体に
 置き換えても成立する。
 中心因子についても \(m>0\) では孤立性から \(P\ne0\) であり、
 \(T\cong V_0\otimes\CC[P]\) となるので同じ議論を適用できる。
 従ってsyzygyへの置換で失う項は特異圏で零である。

 二項ambient分解の微分を \(d_i\)、potential倍のnull-homotopyを \(s_i\) とする。
 外部tensor積上のKoszul符号付きの和 \(d=\sum d_i\)、\(s=\sum s_i\) は
 \(ds+sd=W\)、\(s^2=0\) を満たす。
 そのShamash周期化のtailはfactorizationの外部積に一致する。
 よって各boxのrestrictionのstabilizationは外部積関手のthick imageに入る。
 全 \(\Gamma\)-shiftも各因子のdegree shiftで実現されるので、
 補題\ref{lem:multiblock-ff}の完全忠実関手はMorita同値である。

 \emph{可換側との比較と整数次数への降下.}
 各Sklyanin因子に系\ref{cor:block}を用い、中心因子には恒等を用いる。
 可換側の外部積には \cite[Theorem 5.15, Corollary 5.18]{BFK} を適用する。
 命題\ref{prop:multiblock-embedding}と同様に、\(\Gamma\to\ZZ\) の有限核の双対による
 同変化と双対同変化を用いる。因子ごとのdegree shiftと \((3)=[2]\) のcoherenceは
 各比較で保存されるので、通常の整数次数へ降下すると
 \[
 \MF^{\ZZ}(A,W)\simeq
 \MF^{\ZZ}\!\left(\CC[x_j,y_j,z_j,u_a],
 \sum_j(x_j^3+y_j^3+z_j^3+h_j^2x_jy_jz_j)+P(u)\right)
 \tag{\(*\)}\label{eq:PI-model}
 \]
 を得る。

 \emph{射影模型と一意性.}
 中心正則元の商 \(R\) はnoetherian AS-Gorensteinで、injective dimensionは \(N-1\)、
 Gorenstein parameterは \(N-3\) である。
 \cite[Theorem 5.8]{CCKM} のMF--特異圏比較と
 \cite[Theorem 2.5]{Orlov} の残余windowを適用できる。
 \(N>3\) ならそのwindowは
 \[
 D^b(\qgr R)=\langle\mathcal K_R,\OO,\OO(1),\ldots,\OO(N-4)\rangle
 \]
 を与え、\(N=3\) なら例外対象はなく \(\mathcal K_R=D^b(\qgr R)\) である。
 degree shiftはwindowを通してrotationに対応する。
 可換potentialの勾配は各変数群に分離しており、各 \(F_{h_j}\) と \(P\) の
 アフィン臨界点は原点だけであるから、その和も原点以外に臨界点を持たない。
 従って \(Y_{\mathbf h,P}\) は滑らかで、同じOrlov比較は可換側を
 \((\Ku(Y_{\mathbf h,P}),\rho_{Y_{\mathbf h,P}})\) と同定する。
 \(N-2\ge3\) における射影同型類の一意性は
 \cite[Theorem 1.7]{LZ} による。
\end{proof}

\begin{lemma}[位数五の明示的パラメータ]\label{lem:order-five}
 \(h^6=5\) とする。原点 \(o=[1:-1:0]\) を持つ
 \(E_h=V(x^3+y^3+z^3+h^2xyz)\) 上で、\(p=[1:-1:-h]\) はちょうど位数五である。
 従って対応するSklyaninの平行移動 \(\sigma=t_p\) も位数五である。
\end{lemma}
\begin{proof}
 \(o\) における接線は \(3x+3y-h^2z=0\) であり、
 この直線上で \(F_h=(1+h^6/27)z^3\) となるので、\(o\) はinflection pointである。
 逆元は \([x:y:z]\mapsto[y:x:z]\) である。実際、この二点と \(o\) は一直線上にある。
 \(a=h^3\) とし、\(q=[a+1:a-1:-2h]\) と置く。
 \(F_h(q)=0\) であり、\(p\) における接線
 \[
 (3+a)x+(3-a)y+2h^2z=0
 \]
 は \(-q=[a-1:a+1:-2h]\) を通る。
 \(h\ne0\) では \(-q\ne p\) であるから、接線の加法則により \(q=2p\) である。
 \(q\) における接線の係数は
 \[
 [a^2+8a+3:a^2-8a+3:h^2(a^2+11)]
 \]
 であり、これを \(p\) で評価すると \(a(5-a^2)=h^3(5-h^6)\) となる。
 \(p\ne q\) は \(q=2p\) と \(p\ne o\) から従う。
 \(h^6=5\) では \(p\) がこの接線上にあるので \(2q+p=0\)、従って \(5p=0\) である。
 \(p\ne o\) かつ五は素数なので、位数はちょうど五である。
\end{proof}

\begin{remark}[二・三ブロックのsixfold・sevenfold]
 全ての \(h_j^6=5\) とする。
 \(r=2,m=2,P=u_1^3+u_2^3\) は二つの非可換ブロックを持つcubic sixfoldの
 rotation付き残余圏の完全な可換模型を与える。
 \(r=3,m=0\) は三つの非可換ブロックを持つcubic sevenfoldの完全な可換模型を与える。
 両者の方程式は定理\ref{thm:multiblock-model}の \(Y_{\mathbf h,P}\) であり、
 元の対象の一部分ではなく全残余圏を実現する。
\end{remark}

\begin{lemma}[無限位数の明示的パラメータ]\label{lem:order-infinite}
 \(h=2\) とする。楕円曲線
 \[
 E_2=V(x^3+y^3+z^3+4xyz),\qquad o=[1:-1:0]
 \]
 上の点 \(p_2=[1:-1:-2]\) は無限位数である。
\end{lemma}
\begin{proof}
 曲線と点は \(\mathbb Q\) 上定義されている。
 \(2^6+27=91\) は5,11のどちらでも割れないので、Hesseの勾配による判定から
 5と11で良い還元を持つ。
 \(z=1\)、\(z=0,y=1\)、\([1:0:0]\) の三つのchartに分けて列挙すると
 \[
 \#E_2(\mathbb F_5)=9,\qquad \#E_2(\mathbb F_{11})=12.
 \]
 実際、\(z=1\) の点数はそれぞれ8と11、\(z=0\) にはどちらも一点だけ存在する。
 この列挙は整数の合同計算で再現でき、\texttt{verify\_algebra.py} に全座標を記録する。
 良い還元では剰余標数と素なtorsionの還元写像は単射である
 \cite[Reduction of torsion points]{SnowdenDVR}。
 従って \(p_2\) の位数 \(N\) が有限なら、\(N\) の5と素な部分は9を割り、
 11と素な部分は12を割る。後者により5-primary部分はなく、前者により11-primary部分もなく、
 残りの素数について両方を比較すれば \(N\mid3\) となる。

 \(p_2\ne o\) である。一方、接線による倍点公式から
 \[
 2p_2=[9:7:-4],\qquad -p_2=[1:-1:2]
 \]
 であり、これらは異なる射影点である。
 従って \(3p_2\ne0\) であり、\(N\mid3\) は不可能である。
 \(p_2\) は無限位数である。
\end{proof}

\begin{remark}[真に無限位数を含むsixfoldとsevenfold]
 \(h_1=2\)、\(h_j=5^{1/6}\ (j\ge2)\) と取る。
 第一ブロックの平行移動は無限位数、他のブロックは位数5である。
 \((r,m)=(2,2)\)、\(P=u_1^3+u_2^3\) では可換sixfoldは
 \[
 V\bigl(x_1^3+y_1^3+z_1^3+4x_1y_1z_1
       +x_2^3+y_2^3+z_2^3+5^{1/3}x_2y_2z_2+u_1^3+u_2^3\bigr)
 \subset\PP^7
 \]
 となる。\((r,m)=(3,0)\) では第二・第三ブロックの係数を \(5^{1/3}\) とした
 cubic sevenfoldを得る。いずれも全残余圏とrotationの同値である。
 さらに他の有限位数パラメータを固定すれば、\(h_1^6\ne-27\) を満たす
 全ての \(h_1\in\CC\) について同じ結論が成り立つ。
 これは各ファイバーでの主張であり、同値の族をこのパラメータ直線全体に
 代数的に降下させたという主張ではない。
\end{remark}

\subsection{証明範囲と残る課題}

有限位数条件は \(j\ge2\) のみに用いる。
従ってより一般に、\(C_1\) がnoetherian、\(j\ge2\) の各 \(C_j\) が
有限型可換斉次中心上有限で、周囲が命題\ref{prop:block-support}の仮定を満たせば、
同じ生成と安定化の議論を適用できる。
定理\ref{thm:one-sided-boxes}では任意のnoetherianな係数代数を保てる一方、
新たに加える因子の有限中心性は必要である。
従って二つ以上の独立な無限位数THCRを持つ場合は、この証明だけでは解けない。
これはその場合の非存在を意味せず、条件 (G) が未解決であるという意味である。

次の焦点は、二つの無限位数因子の条件 (G) を証明することである。
それが得られれば、残りの有限中心因子は定理\ref{thm:one-sided-boxes}で追加できる。
必要なのは、THCRに対しても任意のnoetherian係数を許す外部積生成、
または粗化次数を保った対角の有限構成である。
可換ambientのThom--Sebastiani定理を非可換ambientへ無条件に適用したわけではない。
先行研究上の新規性は未確認である。

\section{適用範囲と残る問題}\label{sec:scope}

\subsection{例外パラメータでpropernessが失われる}

\begin{proof}[命題\ref{prop:bad-intro}の証明]
 \(h_0^6=-27\) とする。\(E_{h_0}\) は三本の直線からなる三角形である。
 まず
 \[
 S_{h_0}/(G_{h_0})\simeq B(E_{h_0},L,\sigma_{h_0})                 \tag{6}\label{eq:triangle}
 \]
 を確認する。楕円曲線の場合だけの公式を、この特殊値へ無条件に適用しない。
 \(h_0\) の近傍で、関係式を多重線形化した行列
 \[
 M_h(x,y,z)=
 \begin{pmatrix}-y&x&-hz\\-hx&-z&y\\z&-hy&-x\end{pmatrix}
 \]
 は \(E_h\) 上でrank二である。実際その行列式は
 \(-h(x^3+y^3+z^3+h^2xyz)\) であり、rank一なら小行列式から
 \(yz+hx^2=yz-hx^2=0\) と巡回的な式を得て、\(x=y=z=0\) となる。
 核は相対的な射 \(\sigma_h\) を与える。グラフを転置すると \(h\) を \(-h\) に
 変える同じ関係式になるので、この射は相対自己同型である。

 \(\mathcal L_d=\bigotimes_{j=0}^{d-1}(\sigma^j)^*\OO_E(1)\) は、
 滑らかなファイバーでは次数 \(3d\)、三角形では各成分次数 \(d\) を持つ。
 \(d>0\) なら \(H^1(E_h,\mathcal L_d)=0\)、\(h^0=3d\) であり、
 相対切断加群はlocally freeで底変更と両立する。
 自然な次数付き写像 \(S_h\to\bigoplus_dH^0(E_h,\mathcal L_d)\) の下で、
 \(G_h\) は一般の滑らかなファイバーで零なので、局所自由性から恒等的に零となる。
 三角形のtwisted homogeneous coordinate ringが次数一で生成されることを確認する。
 \(L\) のbasepoint-free pencilに対するpencil trickを使うと、
 多次数 \((d,d,d)\) の \(M\) について
 \(H^0(L)\otimes H^0(M)\to H^0(LM)\) は
 \(H^1(ML^{-1})=0\) なら全射である。
 \(d\ge2\) では各成分上の正値性からこの消滅が従う。
 \(d=1\) で \(ML^{-1}\) が非自明なら、多次数零の直線束の切断は
 三角形の三成分の接合条件により零で、Euler数も零なので同じ消滅が成り立つ。
 \(M\simeq L\) なら、平面三次曲線の斉次イデアルに二次式がないことから全射である。
 従って \(S_{h_0}/(G_{h_0})\to B(E_{h_0},L,\sigma_{h_0})\) は全射である。
 補題\ref{lem:center}の正則性から左辺のHilbert係数は \(1,3d\) であり、
 右辺の係数と一致する。よってこの全射は同型で、\eqref{eq:triangle}が従う。
 \(\mathcal L_d\) の各成分上の次数が \(d\) とともに増大するので
 \(\sigma\)-amplenessも成り立ち、tailsは \(\operatorname{coh}E_{h_0}\) と同値になる。

 Gorenstein parameter零のOrlov同値と \cite[Theorem 5.8]{CCKM}から
 \[
 \MF^\ZZ(S_{h_0},G_{h_0})\simeq D^b(\operatorname{coh}E_{h_0}).
 \]
 三角形の節点 \(p\) のskyscraper sheafは、任意に高い次数で非零の自己Extを持つ。
 節点の完備局所環 \(\CC[[s,t]]/(st)\) の剰余体の最小自由分解は無限であり、
 局所大域Extスペクトル系列では高次Extの支持が \(p\) だけなので、
 この非零性は大域自己Extにも残る。従ってこのdg圏はproperでない。

 \(m=0\) の場合の非propernessはこれで示された。以下では \(m\ge1\) とする。
 \(P=\sum_i u_i^3\) とし、第\ref{sec:TS}節の二ブロック次数 \(\Gamma\) と
 Koszul factorization \(K_P\) を使う。
 \(K_P\) の全微分行列は \((u_1,\ldots,u_m)\) に属する。
 もし恒等写像が零ホモトピーなら、係数を原点で評価して非零自由加群上の恒等写像が
 零となるので矛盾する。従って \(H^0\REnd(K_P)\ne0\) である。
 節点の対象を \(N\) とすると、外部積の完全忠実性と体上のKünneth公式により
 \[
 H^r\REnd(N)\otimes H^0\REnd(K_P)
 \subset H^r\REnd(N\boxtimes K_P)
 \]
 が直和因子として残る。よって \(\Gamma\)-次数の全体圏はnonproperである。
 もし通常の整数次数付き圏がproperなら、その有限群 \(H_\Gamma\simeq\mu_3\) による
 同変圏も、Hom複体の不変部分が直和因子となることからproperとなる。
 これは式\eqref{eq:eq}に反する。

 さらに、この議論は単なるHom総次元の非有限性より強い。
 \(\Gamma\)-次数の一つの対象に、任意に高い正次数の自己Homが残っている。
 有限群同変化から忘却すると、同変Homは通常のHomの直和因子なので、
 通常次数でも同じ正次数非消滅を持つ一つの対象が得られる。
 \(m>0\) ではOrlovの完全忠実windowによりこれを \(D^b(\qgr B_{h_0,m})\) へ送れ、
 \(m=0\) では上の三角形との同値を用いる。
 もし \(\qgr B_{h_0,m}\) のhomological dimensionが有限なら、
 その有界導来圏の任意の固定された二対象の間の正次数Homは十分高次で消える。
 これは有界複体のcohomologyとExtからなる有限スペクトル系列で分かる。
 今得た対象はこれに反するから、homological dimensionは無限である。

 一方、周囲の \(S_{h_0}[u]\) はAS-regularであり、三次の中心元の正則性は
 \(m>0\) なら最後の中心変数についてmonicであること、\(m=0\) なら
 補題\ref{lem:center}から従う。従ってHilbert seriesは依然として
 \((1-t^3)/(1-t)^{m+3}\) であり、degree shiftの三乗も \([2]\) である。
 滑らかな射影多様体のadmissible成分はproperなので、求める非存在が従う。
\end{proof}

この例は、量子化を構成する開集合をさらに縮めて例外値を除くことを妨げない。
従って、零を含む十分小さな開集合での存在問題に対する反例ではない。

\subsection{量子化原稿への適用}

三次の中心元によるAS-regularな周囲の商
\[
 B_\hbar=A_\hbar/(\Omega_\hbar),\qquad B_0=\CC[x]/(F)
\]
を、次数ごとにflatな族として構成できていれば、
その形式完備化には定理\ref{thm:main}が適用できる。
この適用にはPoisson構造の分類を再度用いず、各Artin次数でのgraded flatnessだけを使う。
従って、量子化の構成自体の証明と、本稿の相対残余圏・復元の証明は独立に検証できる。
論文Aの代数的量子化定理への適用では、そこで構成する次数ごとのflatnessを用いる。
代数の量子化と圏の復元は異なる段階であり、それぞれの仮定を区別する。

形式的な可換モデルは、半古典極限の固定された \(X\) と同じとは限らない。
定理\ref{thm:global}の代表式では \(h^2xyz\) が現れる。
このモデルを初期の \(F\) と取り違えない。

\subsection{一般の有限パラメータと次の核心}

定理\ref{thm:etale}により、smooth properな代数的pairの族が与えられれば、
可換cubic点のétale近傍で存在問題は解ける。
一般の量子化族に残る最初の課題は、その相対saturationの構成である。
各Artin切断のsmooth proper性だけから、有限表示の代数的模型の存在を結論しない。

独立な複数ブロックでは、命題\ref{prop:multiblock-embedding}で得た完全忠実な像が
全MF圏を生成するかが、二つ以上の無限位数因子の場合の核心である。
命題\ref{prop:block-support}と条件\eqref{eq:block-generation}は、その具体的な帰着を与える。
高々一つの無限位数因子を持つ場合は定理\ref{thm:multiblock-model}で解決した。
中心三次式内の混合項は定理\ref{thm:global}と定理\ref{thm:multiblock-model}に含まれるが、
非可換ブロックと中心変数の混合項は含まれない。
また、étale局所モデルの元の底への降下と、実現開部分の外のファイバーの分類は
別の大域問題として残る。


\section{可換cubic模型の存在を妨げるrotationの不変量}
\label{sec:rotation-obstruction}

本節では、次数付き代数の条件と、可換cubicのrotationを持つという条件を区別する。
三次中心正則元を持つAS-regular代数の商というだけでは、可換模型の存在は従わない。
さらに、商の射影圏が滑らかな可換cubicのcoherent圏と同値であり、周囲の代数が
Calabi--Yauであっても、与えられたrotationを保つ可換模型が存在しない例がある。
以下の二つの反例は有限位数のZhang twistで構成する。
これらが本稿で構成する特定のPoisson量子化族のファイバーであることは主張しない。

\subsection{Serre関手とrotationの必要条件}

\begin{proposition}[Serre--rotation障害]
\label{prop:serre-rotation-obstruction}
 \(\cC\) をSerre関手 \(S_{\cC}\) を持つproperなdg圏、
 \(\rho\) をその自己同値とする。\(N=n+2\ge5\) とし、
 滑らかなcubic \(Y^n\subset\PP^{N-1}\) とのenhanced pair同値
 \[
 (\cC,\rho)\simeq(\Ku(Y),\rho_Y)
 \]
 が存在するなら、
 \begin{equation}\label{eq:serre-rotation-necessary}
 S_{\cC}\simeq\rho^{\,3-N}[N-2],\qquad
 \Theta_N(\cC,\rho):=S_{\cC}\rho^{N-3}[2-N]\simeq\Id
 \end{equation}
 である。特に \(\rho^3\simeq[2]\) だけでは十分でない。
\end{proposition}
\begin{proof}
 可換斉次座標環 \(R_Y=\CC[z_1,\ldots,z_N]/(F_Y)\) のKrull次元は
 \(N-1\)、標準加群は \(R_Y(3-N)\) である。
 滑らかさによって原点以外の錐は正則であり、次数付きAuslander--Reiten双対性は
 \(\Dsg^{\ZZ}(R_Y)\) のSerre関手を
 \((3-N)[N-2]\) と与える
 \cite[Theorem 7, p.~161]{UeyamaIsolated}。
 Orlovのwindowで内部次数シフト \((1)\) は \(\rho_Y\) に対応するので、
 \(S_{\Ku(Y)}\simeq\rho_Y^{3-N}[N-2]\) を得る。
 Serre関手は同値によって保存されるから、pair同値を通してこの式を移せばよい。
\end{proof}

\subsection{有限位数Zhang twistと次数シフト}

以後 \(\zeta\) は \(\zeta^3=1\)、\(\zeta\ne1\) を満たす。
\(S=\CC[x_0,\ldots,x_{N-1}]\)、\(F=\sum_i x_i^3\)、
\(X=V(F)\subset\PP^{N-1}\) とする。
\(q_i\in\{1,\zeta,\zeta^2\}\) を選び、
\(\tau(x_i)=q_i x_i\) と置く。\(\tau^3=\Id\)、\(\tau(F)=F\) である。
\(f:X\to X\) を \(f^*=\tau\) に対応する射影自己同型とする。
右Zhang twist \(A=S^{\tau}\) の積を、斉次元に対して
\begin{equation}\label{eq:right-zhang-product}
 a*b=a\,\tau^{\deg a}(b)
\end{equation}
で定義する。右辺の積は元の可換多項式環の積である。

\begin{lemma}\label{lem:zhang-cubic-structure}
 \(A\) は左右noetherianなKoszul AS-regular domainで、
 global dimensionは \(N\)、Hilbert級数は \((1-t)^{-N}\) である。
 同じ次数付きベクトル空間内の元 \(\Omega=F\in A_3\) は中心正則元であり、
 \[
 B=A/(\Omega)\simeq (S/(F))^{\tau},\qquad
 H_B(t)=\frac{1-t^3}{(1-t)^N}.
 \]
 さらに
 \begin{equation}\label{eq:zhang-qgr-pair}
 \qgr B\simeq\operatorname{coh}X,\qquad
 (\mathcal K_B,\rho_B)\simeq
 (\Ku(X),f_*\rho_X)
 \end{equation}
 である。従って \(\mathcal K_B\) はsmooth properで、
 \(\rho_B^3\simeq[2]\) を満たす。
\end{lemma}
\begin{proof}
 積\eqref{eq:right-zhang-product}は結合的であり、生成元の関係式は
 \[
 x_i*x_j=(q_j/q_i)x_j*x_i\qquad(i<j)
 \]
 となる。従って \(A\) はscalarによる可逆なOre拡大の反復である。
 順序付き単項式 \(x_0^{a_0}*\cdots*x_{N-1}^{a_{N-1}}\) は
 元の可換単項式の非零scalar倍なので基底をなし、domain性とHilbert級数が従う。
 Ore拡大のHilbert基底定理から左右noetherian性を得る。
 多項式環のKoszul分解をtwistすると、\(k\) の長さ \(N\) の線形自由分解になる。
 その双対もtwistされたKoszul分解であるから
 \[
 \operatorname{Ext}^j_A(k,A)=
 \begin{cases}k(N)&j=N,\\0&j\ne N\end{cases}
 \]
 となり、左側にも同じ式が成り立つ。
 同じtwistを双加群Koszul分解に施した有限分解、またはOre拡大のglobal dimensionの公式により、
 global dimensionは \(N\) である。これはZhang twistがAS-regular性を保つことの
 この場合の具体的な確認である \cite{ZhangTwist,RRZ}。

 任意の斉次 \(a\in A_d\) に対して
 \[
 a*F=a\tau^d(F)=aF,
 \qquad F*a=F\tau^3(a)=Fa
 \]
 であるから、\(F\) は中心的である。
 domain性から正則であり、\((F)\) の次数付きベクトル空間は元の多項式環の
 \((F)\) と一致する。このため商も右twistであり、表示したHilbert級数を得る。
 なお \(x_i*x_i*x_i=q_i^3x_i^3=x_i^3\) なので、
 \(\Omega\) はtwist後の生成元の三乗和としても同じ元である。

 次数付き右 \(R=S/(F)\)-加群 \(M\) に、
 \(m*a=m\tau^{\deg m}(a)\) という作用を入れた加群を \(\mathsf T(M)\) と書く。
 \(\mathsf T\) は逆twistを逆関手に持つ完全同値で、有限生成性とtorsionを保つ。
 従って \(\operatorname{gr}R\simeq\operatorname{gr}B\) および
 \(\operatorname{coh}X\simeq\qgr B\) を与える。

 次数シフトの規約を \(M(1)_d=M_{d+1}\) とする。
 \(\mathsf T(M)(1)\) をuntwistすると、次数 \(d\) の元 \(m\in M_{d+1}\) に対して
 \[
 m\cdot a=m*\tau^{-d}(a)=m\tau(a)
 \]
 となる。従って、untwist後の \((1)\) は \(M(1)_{\tau}\) であり、
 coherent圏では \(f_*(-\otimes\OO_X(1))\) に対応する。
 特に \(\OO_B(i)\) は \(\OO_X(i)\) に対応する。
 \(f_*\OO_X\simeq\OO_X\) なので、\(f_*\) は \(\OO_X\) によるleft mutationと可換である。
 よって残余圏とrotationについて\eqref{eq:zhang-qgr-pair}が従う。
 \(f_*\) は \(\rho_X\) と可換で、\(f_*^3\simeq\Id\) であるため、
 \((f_*\rho_X)^3\simeq\rho_X^3\simeq[2]\) となる。
\end{proof}

\begin{remark}[右twistにおけるNakayamaの規約]
\label{rem:zhang-nakayama}
 上の右twistのNakayama自己同型は
 \begin{equation}\label{eq:right-twist-nakayama}
 \mu_A=\tau^{-N}\xi_{\det\tau},\qquad
 \mu_A(x_i)=\frac{\prod_jq_j}{q_i^N}x_i,
 \end{equation}
 と取れる。ここで \(\xi_c\) は次数 \(d\) 上で \(c^d\) 倍する自己同型である。
 \cite[Theorem 5.4]{RRZ}は左twist
 \(a\circ b=\sigma^{\deg b}(a)b\) に対して
 \(\mu=\sigma^N\xi_{\det\sigma}^{-1}\) を与える。
 \(a\mapsto\tau^{-\deg a}(a)\) は右twistから \(\sigma=\tau^{-1}\) の左twistへの同型であり、
 この公式を移すと\eqref{eq:right-twist-nakayama}になる。
 このように左右の規約を固定してから用いる。
\end{remark}

\subsection{三次元の滑らかな反例}

\begin{theorem}[Serre--rotation条件を破るsmoothな非可換cubic]
\label{thm:zhang-threefold-counterexample}
 \(N=5\)、\(\tau=\operatorname{diag}(\zeta,1,1,1,1)\) とし、
 補題\ref{lem:zhang-cubic-structure}の \(B\) を取る。
 この \(B\) はnoetherianなKoszul AS-regular五次元代数の三次中心正則商であり、
 \(\qgr B\simeq\operatorname{coh}X\) は滑らかなFermat cubic threefoldのcoherent圏と同値である。
 しかし、いかなる滑らかな可換cubic threefold \(Y\) に対しても
 \[
 (\mathcal K_B,\rho_B)\not\simeq(\Ku(Y),\rho_Y).
 \]
\end{theorem}
\begin{proof}
 代数の諸性質とsmooth proper性は補題\ref{lem:zhang-cubic-structure}で示した。
 \(T=f_*|_{\Ku(X)}\) と書くと、同補題により
 \(\rho_B=T\rho_X\)、\(T^3\simeq\Id\)、\(T\rho_X\simeq\rho_XT\) である。
 underlying圏は \(\Ku(X)\) なので、そのSerre関手は
 \(S_X\simeq\rho_X^{-2}[3]\) である。
 もし可換pair模型が存在すれば、命題\ref{prop:serre-rotation-obstruction}により
 \[
 \rho_X^{-2}[3]\simeq (T\rho_X)^{-2}[3]
       \simeq T^{-2}\rho_X^{-2}[3]
 \]
 となり、\(T^2\simeq\Id\) を得る。\(T^3\simeq\Id\) と合わせると \(T\simeq\Id\) となる。

 これはHochschild homologyへの作用に反する。
 例外対象のHochschild homologyは次数零だけにあるため、加法性とHKR同型から
 \[
 \HH_1(\Ku(X))\simeq\HH_1(X)\simeq H^{2,1}(X)
 \]
 である。Griffithsのresidue同型\cite{Griffiths}は
 \[
 J(F)_1\xrightarrow{\sim}H^{2,1}(X),\qquad
 a\longmapsto\operatorname{Res}_X\frac{a\Omega_{\mathrm{proj}}}{F^2},
 \quad
 \Omega_{\mathrm{proj}}=
 \sum_{i=0}^4(-1)^i x_i\,dx_0\wedge\cdots\wedge\widehat{dx_i}\wedge\cdots\wedge dx_4
 \]
 と書ける。\(\tau(F)=F\)、\(\tau^*\Omega_{\mathrm{proj}}=(\det\tau)\Omega_{\mathrm{proj}}
 =\zeta\Omega_{\mathrm{proj}}\) である。
 \(J(F)=\CC[x_0,\ldots,x_4]/(x_0^2,\ldots,x_4^2)\) だから、
 \(f^*\) の \(H^{2,1}(X)\) 上の固有値は
 \[
 \zeta^2\ \text{が1個},\qquad \zeta\ \text{が4個}
 \]
 となる。\(f_*\) の作用はその逆であり、いずれにしても恒等ではない。
 従って \(T\not\simeq\Id\) であり、矛盾する。
\end{proof}

\begin{remark}
 この例では \(\rho_B^3\simeq[2]\) とsmooth proper性の両方が成立する。
 一方、\eqref{eq:right-twist-nakayama}のNakayama自己同型は非自明であり、
 周囲の \(A\) はCalabi--Yauではない。
 次の例は、この最後の条件を加えても一般的な存在命題が成立しないことを示す。
\end{remark}

\subsection{Calabi--Yauな周囲を持つ四次元の反例}

\begin{theorem}[Calabi--Yauな非可換cubicのrotation障害]
\label{thm:zhang-cy-fourfold-counterexample}
 \(N=6\)、\(\tau=\operatorname{diag}(\zeta,\zeta^2,1,1,1,1)\) とし、
 補題\ref{lem:zhang-cubic-structure}の \(A\)、\(B\) を取る。
 このとき \(A\) は六次元Calabi--Yauなnoetherian Koszul AS-regular代数であり、
 \(\Omega\) は三次中心正則元である。
 \(\qgr B\) は滑らかなFermat cubic fourfold \(X\) のcoherent圏と同値であり、
 \(\mathcal K_B\) はsmooth properな2-Calabi--Yau圏である。
 \(\rho_B^3\simeq[2]\) および\eqref{eq:serre-rotation-necessary}も満たす。
 しかし、いかなる滑らかな可換cubic fourfold \(Y\) に対しても
 \[
 (\mathcal K_B,\rho_B)\not\simeq(\Ku(Y),\rho_Y).
 \]
 より具体的に、Hochschild \emph{homology}の次数零部分について
 \begin{equation}\label{eq:cy-fourfold-hh0-obstruction}
 \dim\HH_0(\mathcal K_B)^{\rho_B}=8,
 \qquad
 \dim\HH_0(\Ku(Y))^{\rho_Y}=20
 \end{equation}
 である。
\end{theorem}
\begin{proof}
 \(\det\tau=1\)、\(\tau^6=\Id\) だから、
 \eqref{eq:right-twist-nakayama}によって \(\mu_A=\Id\) である。
 AS-regular性とNakayamaの自明性はCalabi--Yau性を与える
 \cite[Lemma 1.2, Theorem 5.4]{RRZ}。
 残りの代数的性質と \(\mathcal K_B\simeq\Ku(X)\) は
 補題\ref{lem:zhang-cubic-structure}による。
 cubic fourfoldのKuznetsov成分は2-Calabi--Yauなので
 \(S_{\mathcal K_B}\simeq[2]\) であり、
 \(\rho_B^{-3}[4]\simeq[2]\) からSerre--rotation条件も成立する。

 Hochschild homologyの作用を計算する。
 Mukai--HKR同定と残余成分の射影によって
 \begin{equation}\label{eq:cubic-fourfold-hh0-split}
 \HH_0(\Ku(X))\simeq
 H^{2,2}(X)_{\mathrm{prim}}\oplus (A_2\otimes_{\ZZ}\CC)
 \end{equation}
 と書ける。ここで前者は20次元、後者は2次元である。
 \(A_2\) はcubic fourfoldのKuznetsov成分に標準的に含まれる二つのMukai classの張る格子である。
 \cite[Proposition 3.12]{Huy}により、\(\rho_X\) はprimitive部分上で恒等、
 \(A_2\) 上で三つの根を巡回置換するので、後者の固有値は \(\zeta,\zeta^2\) である。
 この記述は一般の滑らかなcubic fourfoldにも成り立つため、
 \(\dim\HH_0(\Ku(Y))^{\rho_Y}=20\) を得る。

 射影自己同型 \(f\) の作用は分解\eqref{eq:cubic-fourfold-hh0-split}を保ち、
 \(A_2\) 上で恒等である。
 実際、その二つの標準classは直線 \(\ell\subset X\) 上の
 \(\OO_\ell(1)\)、\(\OO_\ell(2)\) の残余射影から作られる。
 \(f\) は直線を直線に送り、全ての直線のcohomology classは同じである。
 残余射影のcohomology上の作用は \(\OO_X,\OO_X(1),\OO_X(2)\) から定まり、
 これらを \(f\) が固定するので、二つのMukai classも固定される。
 Mukai--HKR同定に入るTodd classは射影自己同型で不変であるから、
 これは \(\HH_0\) 上の作用についても成立する。

 primitive部分ではresidue表示
 \[
 J(F)_3\xrightarrow{\sim}H^{2,2}(X)_{\mathrm{prim}},\qquad
 a\longmapsto\operatorname{Res}_X\frac{a\Omega_{\mathrm{proj}}}{F^3}
 \]
 を用いる。ここでresidue表示はpole filtrationの商を取った
 \(J(F)_3\simeq\operatorname{Gr}_F^2H^4(X)_{\mathrm{prim}}\) と、
 Hodge分解による \(\operatorname{Gr}_F^2\simeq H^{2,2}_{\mathrm{prim}}\) の合成を意味する。
 射影自己同型は両方の構造を保つ。
 \(\det\tau=1\)、\(\tau(F)=F\) なので、ここで \(f^*\) の作用は
 \(J(F)_3\) 上の通常の置換 \(a\mapsto\tau(a)\) に一致する。
 FermatのJacobian環は \(J(F)=\CC[x_0,\ldots,x_5]/(x_0^2,\ldots,x_5^2)\) であり、
 次数三の基底は平方自由な三変数単項式20個である。
 その固有値を分類すると
 \[
 \begin{array}{c|c|c}
 \text{単項式に含まれる変数}&\text{個数}&\tau\text{の固有値}\\\hline
 x_0,x_1\text{の両方}&4&1\\
 x_0,x_1\text{のどちらも含まない}&4&1\\
 x_0\text{だけ}&6&\zeta\\
 x_1\text{だけ}&6&\zeta^2
 \end{array}
 \]
 となる。従って \(f_*\) のprimitive部分の固定空間も8次元である。

 補題\ref{lem:zhang-cubic-structure}から \(\rho_B=f_*\rho_X\) である。
 primitive部分で \(\rho_X\) は恒等なので、\(\rho_B\) の固定空間は8次元である。
 \(A_2\) 上では \(f_*\) が恒等であり、\(\rho_X\) に固有値1はない。
 よって \eqref{eq:cy-fourfold-hh0-obstruction}が成立する。
 enhanced pair同値は \(\rho\) の \(\HH_0\) 上の作用を共役するので、
 固定空間の次元の不一致が同値の不存在を示す。
\end{proof}

\subsection{smooth properな連結族で保存される不変量}

前の反例と可換cubicのpairとの違いは、連結な圏の族を通じても検出できる。
Hochschild homologyの各次数の次元が個別に一定であることを仮定せず、
perfect複体のEuler標数を用いる。

\begin{proposition}[rotation不変Euler標数の局所一定性]
\label{prop:rotation-euler-locally-constant}
 \(S\) を複素スキームとし、\(\mathcal C_S\) をsmooth properな代数的dg圏の族とする。
 \(\rho_S\) は可逆perfect bimoduleで与えられる自己同値の族とし、
 整数 \(d>0\)、\(a\in\ZZ\) と自然同型 \(\rho_S^d\simeq[2a]\) を仮定する。
 幾何学的点 \(s\) に対して
 \[
 \chi_{\rho}(\mathcal C_s)
   =\sum_{i\in\ZZ}(-1)^i
        \dim_{\kappa(s)}\HH_i(\mathcal C_s)^{\rho_s}
 \]
 と定義する。この和は有限であり、\(s\mapsto\chi_{\rho}(\mathcal C_s)\) は局所一定である。
 特に \(S\) が連結なら一定である。
\end{proposition}
\begin{proof}
 相対Hochschild複体 \(H=\HH(\mathcal C_S/S)\) は \(S\)-perfectであり、
 導来底変更と両立する \cite[Lemmas 3.4, 3.6]{PerryIHC}。
 実際、étale局所的に \(\mathcal C_S=\Perf(E)\) と書けば、
 \(H=E\otimes^{\mathbf L}_{E^{\mathrm{op}}\otimes^{\mathbf L}_S E}E\) である。
 smoothnessによって対角は有限cell双加群の直和因子として表され、propernessによって
 各cellをtraceした複体は \(S\)-perfectとなる。perfect性は直和因子で保たれる。
 これらの構成から底変更も従う。

 \(\rho_S\) が誘導する \(H\) の自己同型を \(r\) と書く。
 Hochschild homologyの加法性により \([1]\) の誘導作用は \(-1\)、従って
 \([2a]\) の作用は恒等である \cite[Section 4, Example 5.6]{CisinskiTabuada}。
 よって \(r^d=1\) であり、
 \[
 e=\frac1d\sum_{j=0}^{d-1}r^j\in\operatorname{End}_{D(S)}(H)
 \]
 は冪等元である。\(\Perf(S)\) は冪等完備なので、その像 \(H^{\mathrm{inv}}\) はperfectである。
 幾何ファイバーで \(e\) の作用は有限巡回群の不変部分への射影であるから、
 \[
 H^{-i}(H^{\mathrm{inv}}\otimes_S^{\mathbf L}\kappa(s))
     =\HH_i(\mathcal C_s)^{\rho_s}.
 \]
 perfect複体のファイバーEuler標数は、局所的な有限階数ベクトル束複体の
 階数の交代和であるため局所一定である。これが主張である。
\end{proof}

\begin{corollary}[Calabi--Yau反例と可換pairの連結成分の分離]
\label{cor:zhang-cy-component-obstruction}
 定理\ref{thm:zhang-cy-fourfold-counterexample}のpairは、
 \(\rho^3\simeq[2]\) を相対的に備えた連結smooth proper代数的dg圏の族によって、
 滑らかな可換cubic fourfoldの標準pairに結ぶことができない。
 実際、両者の不変Euler標数はそれぞれ
 \[
 \chi_{\rho_B}(\mathcal K_B)=10,
 \qquad
 \chi_{\rho_Y}(\Ku(Y))=22
 \]
 である。
\end{corollary}
\begin{proof}
 cubic fourfoldのKuznetsov成分のHochschild homologyは次数 \(-2,0,2\) にのみあり、
 次数 \(-2,2\) は各一次元である。
 標準rotationはsymplecticなので、この二つの部分に恒等として作用する
 \cite[Proposition 3.12]{Huy}。
 定理\ref{thm:zhang-cy-fourfold-counterexample}の \(f\) も、
 \(\det\tau=1\)、\(\tau(F)=F\) により \(H^{3,1}(X)\) のresidueを固定し、
 その複素共役 \(H^{1,3}(X)\) も固定する。
 従って \(f_*\rho_X\) についても \(\HH_{\pm2}\) の不変部分は各一次元である。
 次数零の不変部分は\eqref{eq:cy-fourfold-hh0-obstruction}で計算したので、
 Euler標数はそれぞれ \(1+8+1=10\)、\(1+20+1=22\) となる。
 命題\ref{prop:rotation-euler-locally-constant}を適用する。
\end{proof}

\begin{remark}[量子化族との論理的関係]
\label{rem:counterexample-quantization-scope}
 定理\ref{thm:zhang-cy-fourfold-counterexample}は、
 「多項式環と同じHilbert級数を持つCalabi--Yau AS-regular代数の三次中心正則商であり、
 射影圏がsmooth properである」という条件による一般的な可換pair復元を反証する。
 またSerre--rotation条件を追加するだけでも十分でない。
 underlyingなKuznetsov成分には既に可換模型 \(X\) があるが、指定されたrotationが異なるのである。

 一方、この代数が本稿の特定のPoisson構造から作る、零パラメータで可換な
 中心三次量子化族のファイバーであることは示していない。
 \(\tau\) の有限位数性は中心元の構成に本質的であり、
 \(\zeta\) を任意の複素パラメータに置き換えても三次中心元が残るわけではない。
 従って本例から、その特定の量子化族に対する存在問題の否定は導かない。
 研究目標では、扱う量子化族と残余pairの相対構成を指定し、
 その上での存在問題を述べる必要がある。
\end{remark}

\section{完全忠実比較を同値へ強める条件}
\label{sec:serre-generation}

多ブロック族では、既に構成した完全忠実関手の像が全体を生成するかが問題になる。
外部積による直接的な生成証明とは別に、properness、Serre関手および連結性から
本質的全射性を判定できる。ここでの \(\HH^0\) はcohomologyであり、
前節の反例に用いたhomologyの \(\HH_0\) とは異なる。

\begin{proposition}[Serreと連結性による全生成判定]
\label{prop:serre-connected-generation}
 \(\mathcal A\) を非零のsmooth properな、冪等完備pretriangulated dg圏、
 \(\mathcal D\) をproperな、冪等完備pretriangulated dg圏とする。
 \(\mathcal D\) はSerre関手 \(S_{\mathcal D}\) を持つとし、
 dg完全忠実関手 \(\Phi:\mathcal A\to\mathcal D\) が
 \[
 \Phi S_{\mathcal A}\simeq S_{\mathcal D}\Phi
 \]
 を満たすと仮定する。
 さらに環 \(\HH^0(\mathcal D)\) が \(0,1\) 以外の冪等元を持たないなら、
 \(\Phi\) はdg Morita同値である。
 特に \(\HH^0(\mathcal D)=\CC\) なら、この連結性条件は満たされる。
 \(\mathcal D\) のsmoothnessは仮定しなくてよい。
\end{proposition}
\begin{proof}
 まず \(\Phi\) の像がadmissibleであることを示す。
 \(\mathcal A\) はsmooth properなので、compact generator \(G\) と
 smooth proper dg代数 \(E=\REnd_{\mathcal A}(G)\) を用いて
 \(\mathcal A\simeq\Perf(E)\) と書ける。
 \(D\in\mathcal D\) に対し、\(\RHom_{\mathcal D}(\Phi G,D)\) は
 \(\mathcal D\) のpropernessにより \(\CC\)-perfectである。
 smoothなdg代数上では、このようなpseudo-perfect加群はperfectである。
 実際、対角 \(E\) を有限cell双加群の直和因子として表し、この加群をtensorすると、
 \(\CC\)-perfectな係数を持つ有限cell \(E\)-加群の直和因子として表される。
 したがって \(\mathcal A\) の対象で表現され、\(\Phi\) の右随伴を与える。
 opposite圏で同じ議論を行えば左随伴も得る。
 この構成はdg加群の表現可能性として行われるので、随伴もenhancedである。

 \(\mathcal B=\Phi(\mathcal A)^\perp\) とする。
 \(D\in\mathcal B\)、\(A\in\mathcal A\) に対し、Serre双対性から
 \[
 \Hom_{\mathcal D}(D,\Phi A[i])
 \simeq
 \Hom_{\mathcal D}(S_{\mathcal D}^{-1}\Phi A[i],D)^*
 \simeq
 \Hom_{\mathcal D}(\Phi S_{\mathcal A}^{-1}A[i],D)^*=0
 \]
 となる。従って右直交と左直交は一致する。
 admissibilityによる各対象の分解三角形は、二方向のHom消滅によって分裂する。
 よって \(\mathcal D\) は \(\Phi(\mathcal A)\) と \(\mathcal B\) の直交直和である。
 全てのcohomological shiftのHomが消えるため、これはdg Moritaの意味でも直和分解となる。

 この直和分解は対角双加群の分解を与え、第一成分への射影が
 \(\HH^0(\mathcal D)=\operatorname{End}(\Delta_{\mathcal D})\) の冪等元になる。
 \(\mathcal A\ne0\) なのでこの冪等元は零でない。
 もし \(\mathcal B\ne0\) なら恒等でもないため、仮定に反する。
 従って \(\mathcal B=0\) であり、\(\Phi\) は本質的全射である。
\end{proof}

\begin{corollary}[任意位数の独立ブロックに対する十分条件]
\label{cor:arbitrary-order-conditional-model}
 \(N=3r+m\ge5\) とし、全ての \(h_j^6\ne-27\) とする。
 \(P(u)\) は原点のみを臨界点に持つ中心三次式とし、\(m=0\) では省く。
 独立ブロックの代数
 \[
 B=\left(\bigotimes_{j=1}^rS_{h_j}\right)[u]
     \Big/\left(\sum_jG_{h_j}+P(u)\right)
 \]
 の通常次数付きMF圏を \(\mathcal K_B\)、次数シフトを \(\rho_B\) とする。
 次の三条件を仮定する。
 \begin{enumerate}
 \item \(\mathcal K_B\) はproperである。
 \item Serre関手が存在し、\(S_{\mathcal K_B}\simeq\rho_B^{3-N}[N-2]\) である。
 \item \(\HH^0(\mathcal K_B)\) は非自明な冪等元を持たない。
 \end{enumerate}
 このとき、平行移動の有限位数性を仮定せずに、
 \[
 (\mathcal K_B,\rho_B)\simeq
 \left(\Ku\!\left(V\!\left(\sum_j
 (x_j^3+y_j^3+z_j^3+h_j^2x_jy_jz_j)+P(u)\right)\right),\rho_Y\right)
 \]
 が成り立つ。
\end{corollary}
\begin{proof}
 各楕円ブロックのdegree shiftを保つ同値、外部積の完全忠実性、
 および有限群による通常次数への降下から、
 表示した滑らかな可換cubic \(Y\) について
 \(\Phi:\Ku(Y)\to\mathcal K_B\) が得られ、\(\Phi\rho_Y\simeq\rho_B\Phi\) である。
 中心三次式 \(P\) は一つの可換因子として扱えばよく、
 可換側のequivariant Thom--Sebastianiと非可換側の外部積の完全忠実性は
 Fermat型への分解を必要としない。
 \(Y\) の滑らかさは各変数群の勾配の孤立性から従う。
 \(\Ku(Y)\) は非零のsmooth proper圏であり、
 \eqref{eq:serre-rotation-necessary}と第2条件はSerre関手との互換性を与える。
 命題\ref{prop:serre-connected-generation}を適用すればよい。
\end{proof}

\begin{remark}[この十分条件の検証状態]
 この系は、二つ以上の無限位数ブロックの全生成性を無条件に証明したものではない。
 その場合に残る課題を、三つの明示した圏論的条件へ帰着する結果である。
 与えられた量子化族の全ファイバーについて、properness、Serre公式、
 および \(\HH^0\) の連結性を確認することは依然必要である。
 前節の反例が示すように、一般のAS-regular中心商に対しては、
 比較する完全忠実関手の存在も自動的には得られない。
\end{remark}

\section{相対saturationへの二つの還元}
\label{sec:relative-saturation-reduction}
本節では滑らかな可換cubicを含むPoisson量子化族の相対saturationを、
全超曲面圏の対角の有限証明書へ還元する。
相対Beilinson仮定(RB1)--(RB2)、相対例外列、相対Orlov比較は
元論文の量子化定理からは自動的に従わない。
gluingに関する以下の命題は既知の一般原理の相対版であり、
新規な分類定理とは主張しない\cite[Theorem 3.25]{OrlovGlue}。

\subsection{例外列を除く前後の相対smoothness}
\begin{proposition}[相対smoothnessの同値]
\label{prop:saturation-reduction}
$R$ を可換Noether環、$\mathcal C$ を小さいidempotent completeな
$R$-linear dg圏とし、
$\RHom_{\mathcal C}(X,Y)\in\Perf(R)$ が任意の対象$X,Y$に対し成り立つとする。
有限個の相対例外対象 $E_1,\ldots,E_r$ と
dg半直交分解
\[
\mathcal C=\langle\mathcal K,E_1,\ldots,E_r\rangle
\]
が与えられ、さらに$\mathcal K$がcompact generatorを持つならば、
\[
\mathcal C\text{ は }R\text{-smooth}
\quad\Longleftrightarrow\quad
\mathcal K\text{ は }R\text{-smooth}.
\]
ゆえに相対propernessが既に分かっていれば、
両者の相対saturationは同値な条件である。
\end{proposition}
\begin{proof}
例外対象を一つずつ追加する帰納法により一個の場合に帰着する。
$\mathcal K$ のcompact generatorを$G$、例外対象を$E$とする。
この半直交分解のdg enhancementは$\mathcal K$と$\Perf(R)$を
\[
M=\RHom_{\mathcal C}(G,E)
\]
により三角dg gluingしたものとして表される（右左の規約によるoppositeを適宜取る）。
$\Gamma=\REnd_{\mathcal K}(G)$ と置けば、propernessより
$M$は$R$-perfectである。
$\mathcal K$がsmoothならば
$\Gamma\in\Perf(\Gamma^{\mathrm{op}}\otL_R\Gamma)$。
恒等式 $M\simeq M\otL_\Gamma\Gamma$ に対角の有限cell表示を適用すると、
その像は $M\otL_R\Gamma$ の有限回のshift、cone、直和因子から構成される。
$M$が$R$-perfectだから$M$は$\Gamma$-perfect、従って
gluing bimoduleがperfectとなる。三角dg圏の対角の標準triangleにより
$\mathcal C$はsmoothとなる。
逆に$\mathcal C$がsmoothなら、同じgluingの対角triangleを
両成分へ制限して$\mathcal K$の対角がperfectである。
以上は有限cell・導来テンソル・直和因子のみを使うので
$R$上でも成立する。体上の定理は\cite[Theorem 3.25]{OrlovGlue}。
\end{proof}

\begin{corollary}[必要な対角証明書は一つでよい]
\label{cor:only-ambient-diagonal}
量子化 $B_R=A_R/(\Omega_R)$ が研究ノートの相対Beilinson仮定
(RB1)--(RB2)と相対例外列を満たし、
$E=\bigoplus_{i=0}^{N-1}\mathcal O_{B_R}(i)$ がcompact generatorであるとする。
$\Lambda_R=\REnd(E)$は$R$-perfectであり、
その残余圏$\mathcal K_R$は、次の条件と同値に$R$-smoothである：
\[
\Lambda_R\in
\Perf(\Lambda_R^{\mathrm{op}}\otL_R\Lambda_R).
\tag{S}
\]
残余圏のEnd代数$\Gamma_R$と$\Lambda_R$の二つの対角を
別々に持ち上げる必要はない。
\end{corollary}
\begin{proof}
研究ノートの相対proper模型によって$\Perf(\Lambda_R)$はproperで、
$\mathcal K_R$と有限相対例外列による半直交分解を持つ。
命題\ref{prop:saturation-reduction}を適用する。
\end{proof}

\begin{remark}[閉ファイバーのsmoothnessだけでは足りない]
$R=\CC[\hbar]_{(\hbar)}$とし、
$\mathcal C_0\simeq\Perf(X)$が滑らかな可換cubicの圏でも、
その対角の有限cell分解を$R$へ持ち上げることは別の問題である。
$P\in\Perf(\Lambda_R^{\mathrm{op}}\otL_R\Lambda_R)$と
$P\to\Lambda_R$を構成し、$\hbar=0$で擬同型であることを証明できれば、
coneの$R$-perfect性とsupportの閉性により
零を含むZariski開近傍で(S)が成り立つ。
この有限証明書はまだ一般量子化族について構成されていない。
\end{remark}

\subsection{graded matrix factorizationの三次周期性}
\begin{lemma}[中心三次元の形式的周期性]
\label{lem:mf-three-period}
$A$を可換環$R$上の整数次数付き結合代数、
$\Omega\in Z(A)_3$を中心元とする。
有限射影次数付き加群のgraded matrix factorizationの
内部次数シフトを$(1)$とすると、自然なdg自己同値の同型
\[
(1)^3\simeq[2]
\]
があり、任意の基底変更と両立する。
\end{lemma}
\begin{proof}
因子分解を
\[
P_0\xrightarrow{d_0}P_1\xrightarrow{d_1}P_0(3),
\quad
d_1d_0=\Omega,\quad d_0(3)d_1=\Omega
\]
と書く。懸垂$[1]$は二項を入れ替え、片側を$(3)$移し、
二つの微分に負号を付ける。二回の懸垂では負号が消え、
両項および微分が一斉に$(3)$移る。
従って$[2]\simeq(3)=(1)^3$である。
次数移動と符号のみを用いるので基底変更とも可換である。
\end{proof}

\begin{corollary}[相対Orlov比較の下でのrotation]
\label{cor:relative-rotation-period}
$\mathcal K_R$と$\MF^{\mathrm{gr}}(A_R,\Omega_R)$の間に、
基底変更と両立する相対Orlov同値があり、
$(1)$が標準rotation$\rho_R$に対応すると仮定する。
すると$\rho_R$は自己同値で、相対的に
\[
\rho_R^3\simeq[2]
\]
が成立する。この段階では相対smoothnessを別途仮定しない。
\end{corollary}
\begin{proof}
補題\ref{lem:mf-three-period}を相対同値に沿って移す。
\end{proof}

\subsection{可換復元に向けた残る問題}
上の議論により、少なくとも次の四点を区別できる：
(RB1)--(RB2)を指定量子化族について証明すること、
全超曲面圏$\Perf(\Lambda_R)$の対角の有限証明書を可換点から持ち上げること、
相対Orlov比較でrotationを同定すること、
そして得られたsmooth proper pairが可換cubicのpairの代数的変形像に入ること。
前の三点が解けても最後の可換モデルの存在は別問題であり、
Zhang twistの障害を無視することはできない。

\section{問題と主結果}\label{er:sec:intro}

全体を通じて基礎体は$\CC$である。
三次中心正則元$g$を持つ次数付きAS-regular代数$S$から
$S[u_1,\ldots,u_m]/(g+P(u))$を作るとき、その残余圏を
可換三次超曲面から復元できるかを考える。
中心変数$u_i$は$S$の全ての元と可換とし、$P$は次数三、
\begin{equation}\label{er:eq:isolated}
 V(\partial_{u_1}P,\ldots,\partial_{u_m}P)=\{0\}\subset\mathbb A^m,
 \qquad m\ge2
\end{equation}
を仮定する。これは$P$のアフィン臨界点が原点のみであるという条件である。

使用する量子化は、$t^3\ne-1$に対する
\[
 f_t(x,y,z)=\frac{x^3+y^3+z^3}{3}+txyz
\]
のJacobian Poisson括弧の量子化である。
$h$を次数零の変形パラメータとし、$S_{h,t}$を
\begin{equation}\label{er:eq:relations}
\begin{aligned}
 [y,z]&=h\bigl(x^2+\tfrac t2(yz+zy)\bigr),\\
 [z,x]&=h\bigl(y^2+\tfrac t2(zx+xz)\bigr),\\
 [x,y]&=h\bigl(z^2+\tfrac t2(xy+yx)\bigr)
\end{aligned}
\end{equation}
で定める。$x,y,z,u_i$の内部次数は一である。
次の多項式と有理関数を用いる：
\begin{equation}\label{er:eq:coefficients}
\begin{aligned}
 b(h,t)&=1-\frac{h^2t^2}{4},&
 a(h,t)&=t+\frac{h^2}{3}\left(1+\frac{t^3}{4}\right),\\
 \vartheta(h,t)&=\frac{a(h,t)}{b(h,t)},&
 \Delta(h,t)&=a(h,t)^3+b(h,t)^3,\\
 U(h,t)&=1+\frac{h^2t^2}{2}+\frac{h^4t}{6}-\frac{h^4t^4}{48},\\
 W(h,t)&=\frac{(4-h^2t^2)\bigl(12t+h^2(t^3-8)\bigr)}{48}.
\end{aligned}
\end{equation}
$\mathscr U_t=\{h:U(h,t)b(h,t)\Delta(h,t)\ne0\}$は$0$を含む開集合である。
$\sym(xyz)$は六つの順序付き積の平均を表し、
\begin{equation}\label{er:eq:center}
 \Omega_{h,t}=\frac{U(h,t)}3(x^3+y^3+z^3)+W(h,t)\sym(xyz)
\end{equation}
と置く。中心性・正則性は第\ref{er:sec:algebra}節で確認する。

\begin{theorem}[一ブロックの明示的可換復元]\label{er:thm:main}
$t^3\ne-1$、$h\in\mathscr U_t$および\eqref{er:eq:isolated}を仮定する。
\[
 B_{h,t,P}=S_{h,t}[u_1,\ldots,u_m]/(\Omega_{h,t}+P(u)),
 \qquad
 Y_{h,t,P}=V\bigl(f_{\vartheta(h,t)}(x,y,z)+P(u)\bigr)
 \subset\PP^{m+2}
\]
とおく。$Y_{h,t,P}$は滑らかな三次超曲面であり、dg enhancementを保つ同値
\begin{equation}\label{er:eq:main-pair}
 (\mathcal K_{B_{h,t,P}},\rho_B)
 \simeq (\Ku(Y_{h,t,P}),\rho_Y)
\end{equation}
が存在する。ここで$\mathcal K_B$はOrlovの残余成分、
$\rho_B=\mathbf L_{\OO_B}\circ(1)$である。
従って$\mathcal K_B$はsmooth properである。
この可換模型は射影同型を除いて一意である。
楕円点曲線上の平行移動の有限位数性は仮定しない。
\end{theorem}

\begin{remark}[主張の範囲]\label{er:rem:scope}
定理は各複素パラメータにおける同値である。
全ての同値を底上の一つの代数的dg同値に選んだとは主張していない。
また$u_i$と$x,y,z$の本質的な混合項や、二つ以上の独立な無限位数楕円ブロックは対象外である。
smooth proper性はこの特定の残余成分について同値から得る結論であり、
任意のAS-regular中心三次商からの一般的帰結ではない。
\end{remark}

\subsection{先行研究と論理的な依存関係}
三次元Sklyanin代数の中心三次元と楕円曲線による商の表示は古典的である
\cite{SmithTate,Walton,DeLaetLeBruyn}。
Orlovの定理はGorenstein parameter零の特異圏を射影圏と同一視し、
正のparameterでは特異圏を残余成分として実現する\cite[Theorem 2.5]{Orlov}。
非可換行列因子分解との比較には\cite[Theorem 5.8]{CCKM}を用いる。
Etingof--Ginzburgの\cite[Theorem 3.6.2]{EG}にも楕円曲線とgraded factorizationの比較が現れる。
本稿では同結果のgenericという仮定を全パラメータへ読み替えず、
各$h\in\mathscr U_t$で中心商とOrlovの適用条件を確認する。

本稿の比較の要点は、楕円曲線の次数シフトに現れる平行移動を直線束のテンソルで共役し、
その三次周期性を保存したまま中心変数を追加することである。
有限群の双対同変化には\cite[Theorem 4.2, Section 8]{Elagin}を用いる。
可換模型の一意性は\cite[Theorem 1.7]{LZ}によるが、これは存在を与える定理ではない。
本稿の結果は既存の一楕円ブロック比較をHesse量子化へ適用し、
その変形方向と一定の可換模型を持つ代数的曲線を明示したものである。
これらの組合せについて先行研究上の初出を確定したとは主張しない。

\subsection{圏と次数の規約}
右加群を用いる。$\MF^\Gamma(A,w)$は有限射影$\Gamma$次数付き行列因子分解のdg圏の
pretriangulated hullを取り冪等完備化したものとする。
ここで行列因子分解自体は$D^2=w$を満たすcurved objectであり、通常の複体ではない。
一方、二対象間のHomには$d(\varphi)=D_Q\varphi-(-1)^{|\varphi|}\varphi D_P$が入り、
$w$の中心性から$d^2=0$となる。
$\Gamma=\ZZ$、$\deg w=3$なら、内部シフトは自然な周期性
\begin{equation}\label{er:eq:period}
 (1)^3=(3)\simeq[2]
\end{equation}
を備える。以下の同値では、この周期性も記録する。

\section{中心三次元と楕円商}\label{er:sec:algebra}
\begin{proposition}\label{er:prop:center}
$h\in\mathscr U_t$では$S_{h,t}$はnoetherianな三次元Koszul AS-regular domainで、
Hilbert級数は$(1-z)^{-3}$である。
\eqref{er:eq:center}は中心正則元であり、$\Omega_{0,t}=f_t$。
さらに滑らかな三次曲線
\[
 E_{h,t}=V\bigl(x^3+y^3+z^3+3\vartheta(h,t)xyz\bigr)
\]
と$L=\OO_{E_{h,t}}(1)$について
\begin{equation}\label{er:eq:thcr}
 S_{h,t}/(\Omega_{h,t})\simeq B(E_{h,t},L,t_{p_{h,t}}),
 \qquad
 p_{h,t}=[1-ht/2:-1-ht/2:-h]
\end{equation}
が成り立つ。原点は$o=[1:-1:0]$である。$h=0$では$t_{p_{0,t}}=\Id$と解釈する。
\end{proposition}
\begin{proof}
$q=ht/2$、$\alpha=1-q$、$\beta=-1-q$、$\gamma=-h$とすると、関係式は
$\alpha yz+\beta zy+\gamma x^2=0$とその巡回置換である。
$h\ne0$では$\alpha\beta\gamma=h b(h,t)\ne0$かつ
\[
 (\alpha^3+\beta^3+\gamma^3)^3-(3\alpha\beta\gamma)^3
 =-27h^3\Delta(h,t)\ne0.
\]
従って楕円型Sklyanin代数の既知の定理を適用できる\cite{SmithTate,Walton}。
その点曲線は
$\alpha\beta\gamma(x^3+y^3+z^3)-
(\alpha^3+\beta^3+\gamma^3)xyz=0$である。
$\alpha^3+\beta^3+\gamma^3=-3h a(h,t)$を代入すると$E_{h,t}$を得る。
$1+\vartheta^3=\Delta/b^3\ne0$なので滑らかである。

中心元の規格化を確認するため、$S_3=x^3+y^3+z^3$、
$C=xyz+yzx+zxy$、$D=yxz+xzy+zyx$と書く。
\cite[Section 2, p.~3]{DeLaetLeBruyn}の中心三次元の表示は
\[
 G=\alpha(\beta^3-\gamma^3)C+
   \beta(\gamma^3-\alpha^3)D+
   \gamma(\alpha^3-\beta^3)S_3.
\]
関係式を$x,y,z$倍して足すと$C-D=hS_3+q(C+D)$である。
この式を用いて直接整理すると
\begin{equation}\label{er:eq:central-identity}
 G=-9h\left(\frac U3 S_3+\frac W6(C+D)\right)=-9h\Omega_{h,t}.
\end{equation}
従って$h\ne0$では$\Omega_{h,t}$は標準中心元の非零スカラー倍であり、
中心正則で、その商が\eqref{er:eq:thcr}となる。
$h=0$では代数は多項式環、$U(0,t)=1$、$W(0,t)=t$であるから
$\Omega_{0,t}=f_t$である。この場合は滑らかな平面三次曲線の通常の斉次座標環であり、
同じ商の記述が成立する。
\end{proof}

\begin{remark}
$h=0$における$S_{0,t}$自体のpoint schemeは$\PP^2$であって、$E_{0,t}$ではない。
$E_{0,t}$は中心三次元で割った商の射影曲線である。
この二つを混同すると零パラメータを含む族の記述を誤る。
また中心元の公式の証明では\cite{DeLaetLeBruyn}のSection 2の表示を使う。
そのarXiv v1の冒頭のTheorem 1に印刷された表示とSection 2の表示には係数の差があるため、
\eqref{er:eq:central-identity}を別途多項式として検算する。
\end{remark}

\section{楕円曲線上のrotationの共役}\label{er:sec:rotation}
この節では$E$を滑らかな平面三次曲線、$L=\OO_E(1)$、$\sigma=t_p$とし、
$S$をnoetherianな三次元二次AS-regular代数、$g\in Z(S)_3$を正則元で
$S/(g)=B(E,L,\sigma)$となるものとする。

\begin{lemma}\label{er:lem:block}
$\MF^\ZZ(S,g)$は$\Perf(E)$とdg Morita同値である。
この同値の下で内部シフトは
\[
 \rho_p=T_{\OO_E}\circ\sigma_*\circ(-\otimes L)
       =t_{-p}^*\rho_0,
 \qquad \rho_0=T_{\OO_E}\circ(-\otimes L)
\]
に対応する。$T$はspherical twistである。
\end{lemma}
\begin{proof}
三次元AS-regular代数のparameterは三であり、三次中心正則元で割ると
Gorenstein parameterは零になる。
\cite[Theorem 5.8]{CCKM}のcokernel関手および
\cite[Theorem 2.5(iii)]{Orlov}のwindowを合成すると
\[
 \MF^\ZZ(S,g)\simeq \Dsg^{\mathrm{gr}}(S/(g))
 \simeq D^b(\qgr(S/(g)))\simeq\Perf(E)
\]
を得る。ここではcokernelとwindowのdg模型を使う。
同じ有限射影分解・dg quotientによってHom複体の比較が得られるので、
三角圏の抽象的同値だけを用いているわけではない。
THCRの右加群の規約では、tailsの次数シフトは
$\sigma_*(-\otimes L)$である。
windowを一つ移動させる際の差は$\OO_E$の評価射のconeであり、
従って残余側の次数シフトは$T_{\OO_E}\sigma_*(-\otimes L)$になる。
最後に、平行移動は$\OO_E$を保つのでそのpullbackは$T_{\OO_E}$と可換である。
\end{proof}

\begin{lemma}[次数零直線束による共役]\label{er:lem:conjugacy}
$P_q=\OO_E(q-o)$、$A_q=(-\otimes P_q)$と置くと
\begin{equation}\label{er:eq:conjugacy}
 A_q\rho_0A_q^{-1}\simeq t_q^*\rho_0.
\end{equation}
特に$\rho_p=A_{-p}\rho_0A_{-p}^{-1}$であり、$A_p$は$(\Perf(E),\rho_p)$から
$(\Perf(E),\rho_0)$への同値を与える。$p$の位数に条件はない。
\end{lemma}
\begin{proof}
テンソル関手との共役により$A_qT_{\OO_E}A_q^{-1}=T_{P_q}$である。
sourceを第一因子、targetを第二因子とするFourier--Mukai規約で、
$T_{P_q}$のkernelは
\[
 (P_q^{-1}\boxtimes P_q)\otimes I_\Delta[1]
\]
である。$E\times E$の二つの方向のfiberへの制限を計算すると、see-sawの原理により
\[
 (P_q^{-1}\boxtimes P_q)\otimes\OO(-\Delta)
 \simeq\OO(-\Gamma_q),\qquad
 \Gamma_q=\{(x,y):y=x-q\}
\]
となる。実際、第一座標$x$を固定すると両辺の次数$-1$の直線束の類は
$-x+q$であり、第二座標$y$を固定すると類は$-y-q$である。
よって$T_{P_q}=t_q^*T_{\OO_E}$。
$A_q$は$(-\otimes L)$と可換なので\eqref{er:eq:conjugacy}を得る。
\end{proof}

\begin{lemma}[周期性を含むブロック比較]\label{er:lem:periodic-block}
$E=V(f_E)$ならば、内部シフトと\eqref{er:eq:period}を保つdg Morita同値
\[
 \MF^\ZZ(S,g)\simeq\MF^\ZZ(\CC[x,y,z],f_E)
\]
が存在する。
\end{lemma}
\begin{proof}
両側を補題\ref{er:lem:block}により$\Perf(E)$と同一視し、
補題\ref{er:lem:conjugacy}の$A_p$を用いる。
これにより内部シフトをintertwineする自然同型を得る。
二つの三次周期性をこの同型に沿って比較した差は、同値関手の可逆な自然自己変換である。
$\HH^0(E)=H^0(E,\OO_E)=\CC$だから、この差は非零スカラーである。
その三乗根を用いてintertwinerを再規格化すれば、三次周期性も一致する。
負次数のHochschild cohomologyは零であり、既にfactorization側で与えられた
内部シフトのcoherenceをこの同値で移す。
従って単に自己同値の同型類を合わせるだけでなく、次節の有限群作用に必要な
周期性のデータも保存できる。
\end{proof}

\section{中心変数の追加と有限次数差の降下}\label{er:sec:suspension}
この節は、ブロックの比較を全非可換cubicへ移すための生成性を示す。
$T=\CC[u_1,\ldots,u_m]$、$A=S\otimes T$、$w=g+P$とする。
次数群を
\[
 \Gamma=(\ZZ e_0\oplus\ZZ e_1)/(3e_0-3e_1)
\]
とし、$S_d$の次数を$de_0$、$T_d$の次数を$de_1$とする。

\begin{lemma}[中心変数に対する外部積生成]\label{er:lem:TS}
$S$が非負次数付きnoetherian AS-regular代数、$g\in Z(S)_3$が正則で、
$P$が\eqref{er:eq:isolated}を満たせば、外部積はdg Morita同値
\begin{equation}\label{er:eq:TS}
 \Perf\bigl(\MF^\ZZ(S,g)\otimes_\CC\MF^\ZZ(T,P)\bigr)
 \simeq\MF^\Gamma(A,w)
\end{equation}
を与える。第一・第二因子の内部シフトはそれぞれ$(e_0)$、$(e_1)$に対応する。
\end{lemma}
\begin{proof}
外部積の微分は$D_Q\otimes1+(-1)^{\deg}1\otimes D_K$であり、その二乗は$g+P$である。
有限自由factorization間のHomを行列成分ごとに計算すると、右辺のHom複体は左辺の
二つのHom複体のtensor totalizationに一致する。
非負の内部次数と固定された有限個の生成元のshiftにより、各Hom複体は下に有界であり、
各total次数では和は有限である。これが完全忠実性を与える。

本質的全射性には$u$方向の微分を使う。
任意のfactorization$F$の微分を$D$とすると、$S$と$u_i$が独立なので
$\partial_{u_i}$は$S[u]$の導分であり、
\[
 D(\partial_{u_i}D)+(\partial_{u_i}D)D
   =(\partial_{u_i}P)\Id_F.
\]
従ってJacobianイデアル$J(P)$は$F$上で零ホモトピックに作用する。
\eqref{er:eq:isolated}より、十分大きい$d$に対し全ての$u_i^d$は$J(P)$に属する。
$u_i^d$による乗法のconeを順に取ったKoszul対象は$F$のshiftを直和因子に持つ。
特異圏でその対象を有界複体として表すと、そのcohomologyは
$u_1^d,\ldots,u_m^d$で殺される。
$(u_1,\ldots,u_m)$の冪による有限filtrationの各商は
$C=S/(g)=A/(w,u)$上の有限生成次数付き加群である。
従って右辺は、$C$-加群のrestrictionをstabilizeした対象で生成される。

ここでrestrictionが特異圏に降りることを確認する。
$u_1,\ldots,u_m$は$A/(w)$上の中心正則列である。
実際、途中で一部の$u$を零にしても$g+$残りの三次式は
$S$次数によるfiltrationで正則な最高記号$g$を持つ。
一つの$u_i$についても、先に$u_i=0$とした商でこの元が正則であることから
$u_i$が非零因子となる。
従って$C$は$A/(w)$上有限射影次元を持ち、restrictionはperfectをperfectへ送る。

$C$-加群$M$の十分高いsyzygyを取れば、その$S$上の射影次元は高々一となる。
実際、自由$C$-加群の$S$上の射影次元は一なので、syzygyを一つ取るごとに
次数二以上の$\Ext_S$は一つdimension shiftし、$S$の有限大域次元により最後には消える。
その二項$S$-射影分解と$g$の零ホモトピーは
$Q\in\MF^\ZZ(S,g)$を与える。
$M$を$A$上で分解するには、この二項分解と$u_1,\ldots,u_m$のKoszul分解をtensorすればよい。
$w$の零ホモトピーには、前者の$g$の零ホモトピーと
\[
 P=\sum_i u_i\,\frac{\partial_{u_i}P}{3}
\]
から得るKoszulホモトピーを足す。
この有界分解を周期化すると、restrictionされた$M$のstabilizationは
\[
 Q\boxtimes K_P,\qquad
 K_P=\left(\bigwedge^\bullet\CC^m\otimes T,
       \sum_i u_i\iota_i+\frac{\partial_{u_i}P}{3}\,\varepsilon_i\right)
\]
に一致する。$\varepsilon_i$は外積、$\iota_i$は縮約で、$K_P$の微分の二乗は$P$である。
二項分解の完全性、Koszul分解の完全性、およびこの式は、周期化のcohomologyを
直接同定し、余分な有限射影尾部は特異圏で零となる。

最後に、$C$は$\ZZ e_0$に次数を持つ一方、
$\Gamma/\ZZ e_0\simeq\ZZ/3$である。
任意の$\Gamma$次数付き$C$-加群をこの三剰余類に分ければ、
上の$\ZZ$次数付き議論と有限個のshiftですべての対象が得られる。
従って粗化次数でのみ現れる加群を落とさず、外部積が全生成する。
\end{proof}

\begin{lemma}[通常次数への降下]\label{er:lem:descent}
補題\ref{er:lem:periodic-block}のように第一因子の内部シフトと三次周期性を保つdg同値は、
\eqref{er:eq:TS}を通じ、通常の整数次数について
\[
 \MF^\ZZ(S[u],g+P)\simeq
 \MF^\ZZ(\CC[x,y,z,u],f_E+P)
\]
を誘導する。この同値は通常次数の$(1)$を保つ。
\end{lemma}
\begin{proof}
$\Gamma\to\ZZ$を$e_0,e_1\mapsto1$で定めると、核は$\ZZ/3$である。
その双対群$H\simeq\mu_3$は、整数次数の代数において$u$だけを三乗根倍する。
$P$の次数が三なのでcurvatureは不変である。
$\Gamma$次数付きfactorizationは、整数次数付きfactorizationの$H$-同変対象と同一視される。

その同変圏には双対群$H^\vee\simeq\ZZ/3$が文字をtensorすることで作用する。
外部積の表示では生成元は$\rho_T\rho_S^{-1}$であり、その三乗の自明化は
二つの$(3)\simeq[2]$の比較である。
第一因子の同値がこのデータを保つので、$H^\vee$-同変化にも同値を誘導する。
有限群の位数が可逆であることと冪等完備性を用い、
双対同変化で元の圏を回復する定理\cite[Theorem 4.2]{Elagin}を適用する。
同論文Section 8のdg模型、あるいは有限群のaveraging idempotentで、この比較はdg同値として行える。
$e_0$のシフトは全過程で保存されるため、通常次数の$(1)$も保存される。
\end{proof}

\subsection{主定理の証明}
\begin{proof}[定理\ref{er:thm:main}の証明]
命題\ref{er:prop:center}により$(S,g)=(S_{h,t},\Omega_{h,t})$は
補題\ref{er:lem:block}の全ての仮定を満たす。
$E=V(f_{\vartheta})$なので、周期性を保つブロック比較と中心変数の追加・降下から
\[
 \MF^\ZZ(S_{h,t}[u],\Omega_{h,t}+P)
 \simeq\MF^\ZZ(\CC[x,y,z,u],f_{\vartheta}+P)
\]
を得る。$S_{h,t}[u]$は次元$m+3$のnoetherian AS-regular代数であり、
$\Omega_{h,t}+P$は中心正則である。
その商のGorenstein parameterは$m$であり、
Orlov windowはこのMF圏を
$\langle\OO_B,\ldots,\OO_B(m-1)\rangle^\perp$へ実現する。
degree shiftはその残余成分の$\mathbf L_{\OO_B}(1)$に移る。
可換側も同じwindowを用いるので\eqref{er:eq:main-pair}を得る。

可換方程式の全偏微分が零なら、まず\eqref{er:eq:isolated}により$u=0$である。
残る$f_{\vartheta}$は$1+\vartheta^3\ne0$だからアフィン臨界点は原点のみである。
従って$Y$は射影的に滑らかである。
$\Ku(Y)$は滑らかな射影多様体のadmissible成分なのでsmooth properである。
最後に、別の滑らかな可換cubicのpairと同値ならば、二つの可換pairのenhanced同値を得る。
これはadmissible成分のFourier--Mukai同値として表され、
\cite[Theorem 1.7]{LZ}から二つのcubicは同型になる。
次元$m+1\ge3$ではその同型は超平面束を保ち、射影同型である。
\end{proof}

\begin{corollary}[任意の滑らかな三変数三次式]\label{er:cor:ternary}
滑らかな三変数三次式$g$を
$g(x)=\lambda f_t(Mx)$、$\lambda\ne0$、$M\in\mathrm{GL}_3(\CC)$とHesse形にする。
$g$のJacobian括弧を量子化し、$h'=\lambda\det(M)h$、
$\Omega_h^g(x)=\lambda\Omega_{h',t}(Mx)$で中心元を正規化する。
このとき$h'\in\mathscr U_t$では、
$S_h^g[u]/(\Omega_h^g+P)$の可換pair模型は
\[
 V\bigl(\lambda f_{\vartheta(h',t)}(Mx)+P(u)\bigr)
\]
である。さらに$Mx$を中心変数の一次形式だけ平行移動した場合も、同じ座標置換で結論を移せる。
\end{corollary}
\begin{proof}
滑らかな平面三次曲線のHesse標準形は\cite[Lemma 1]{ArtebaniDolgachev}による。
$y=Mx$の下で定数三重ベクトルは$\det(M)$倍となり、対称化は線形変換と可換なので
$S_h^g\simeq S_{h',t}$である。
商の方程式を$\lambda$で割り、定理\ref{er:thm:main}を$P/\lambda$に適用する。
最後に方程式を$\lambda$倍して元の座標に戻す。
中心線形シフトは次数付き代数の同型であり、内部シフトも保つ。
\end{proof}

\section{復元されたcubicの二次から始まる変形}\label{er:sec:second-order}
次の主張は、定理\ref{er:thm:main}で明示した可換cubicの族自体に関するものである。
非可換族の全体との相対同値を仮定しない。

\begin{theorem}[markedな二次変形の非自明性]\label{er:thm:second-order}
$t$と$P$を固定する。$Y_{h,t,P}$の方程式は
\begin{equation}\label{er:eq:second}
 f_t+P+\frac{h^2(1+t^3)}{3(1-h^2t^2/4)}xyz
\end{equation}
である。従って$h$に関して一次のmarked変形は自明だが、
$\CC[h]/(h^3)$上のmarked変形は自明でない。
$q=h^2$をパラメータとする族のKodaira--Spencer類は
\[
 \frac{1+t^3}{3}[xyz]\in
 J(f_t+P)_3,
 \qquad J(F)=\CC[x,y,z,u]/(\partial F),
\]
であり、これは非零である。
\end{theorem}
\begin{proof}
\eqref{er:eq:coefficients}から直接
\[
 \vartheta(h,t)-t=
 \frac{h^2(1+t^3)}{3(1-h^2t^2/4)}
\]
を得る。滑らかな射影超曲面のembeddedな一次変形を射影座標変換と方程式のスカラー倍で割った空間は
Jacobian環の三次部分である。Euler等式により$F$自身もJacobianイデアルに含まれる。
次元三以上では、この記述は$H^1(T_Y)$による抽象変形の接空間とも一致する。

$f_t$の三つの偏微分は$x^2+tyz,y^2+tzx,z^2+txy$であり、正則列をなす。
従ってそのJacobian環はHilbert級数$(1+z)^3$を持ち、三次部分の次元は一である。
同環で
\[
 x^2y=-t y^2z=t^2xz^2=-t^3x^2y
\]
なので、$1+t^3\ne0$から$x^2y=0$となり、同様に全ての指数型$(2,1,0)$の単項式が零となる。
また$x^3=y^3=z^3=-txyz$である。
従って三次部分は$xyz$で生成され、次元一であるから$[xyz]\ne0$。
$t=0$でも同じ議論が成立し、$xyz$は三次の唯一の平方自由単項式である。
変数が分離しているので$J(f_t+P)=J(f_t)\otimes J(P)$であり、
$[xyz]\otimes1$は非零のままである。

一次の変形項が零なので、$h^2$での座標変換が消せるのはJacobianイデアルの三次部分だけである。
$h$で始まる座標変換が一次を保つ場合、その一次項は$F$の無限小射影自己同型である。
滑らかな三次超曲面では、偏微分の一次係数syzygyが存在しないこととEuler等式により
それはスカラーEuler方向に限られ、射影的には零である。
よって非零の上記Jacobian類は$\CC[h]/(h^3)$上で消せない。
\end{proof}

\begin{remark}
この結論はmarked変形、または局所moduli stackについてのものとして述べている。
特別なcubicの有限自己同型による商を取ったcoarse moduliの局所座標で、
同じ二次の次数がそのまま最初の非零次数になるとは主張しない。
また、ファイバーごとの\eqref{er:eq:main-pair}のみから、非可換族の相対Kodaira--Spencer類を
上の式と同一視してはいない。
\end{remark}

\section{同じ可換模型を持つ非可換な代数的曲線}\label{er:sec:constant-model}
固定した$t$では復元されたcubicが動く一方、$t$も動かすと正反対の現象が現れる。

\begin{theorem}[ファイバーごとに同じpairを持つ量子化族]\label{er:thm:isocategorical}
$\tau^3\ne-1$と\eqref{er:eq:isolated}を満たす$P$を固定し、
$Y_0=V(f_\tau+P)$とする。
$(h,t)=(0,\tau)$を通る平面曲線
\begin{equation}\label{er:eq:curve}
 \mathscr C_\tau:
 h^2t^3+3\tau h^2t^2+12t+4h^2-12\tau=0
\end{equation}
のその点を含む適切な滑らかな開部分上に、次数別平坦な代数的族
$B_{h,t,P}$がある。次を満たす。
\begin{enumerate}
 \item $h$への射は$(0,\tau)$でétaleであり、全ての複素ファイバーで
 $(\mathcal K_{B_{h,t,P}},\rho_B)\simeq(\Ku(Y_0),\rho_{Y_0})$。
 \item $B_{h,t,P}$の$(0,\tau)$における次数付き結合代数の一次変形は非自明である。
 特にこの族は可換代数の定数族ではない。
 \item この曲線の関数体上で、Sklyaninの平行移動点$p_{h,t}$は無限位数である。
\end{enumerate}
ここで主張しているのはファイバーごとの同値であり、
一つの相対同値による自明化ではない。
\end{theorem}
\begin{proof}
$\vartheta(h,t)=\tau$の分母を払うと\eqref{er:eq:curve}になる。
その$t$による偏微分は$(0,\tau)$で12である。
従ってこの点を含む唯一の既約成分を取り、適切な開近傍に制限すると、
曲線は滑らかで$h$への射はétaleとなる。
さらに$t^3+1$、$U$、$b$、$\Delta$が非零となるよう縮める。
いずれも$(0,\tau)$で非零である。
ここでは$\vartheta=\tau$なので、定理\ref{er:thm:main}の可換模型は常に$Y_0$である。

代数の表示はこの曲線上の正則関数を係数とする有限表示である。
各幾何ファイバーの$S$のHilbert級数は$(1-z)^{-3}$で一定である。
滑らかな曲線上の有限生成加群のtorsionはその台でfiber次元を増加させるので、
各次数成分は局所自由になる。
中心元による各次数の乗法は全fiberで単射であり、その余核も局所自由である。
従って商$B$は次数別平坦である。

局所環の完備化では$t=\tau+O(h^2)$である。
従って$\CC[h]/(h^2)$上で
\[
 [x,y]=h(z^2+\tau xy).
\]
右辺の係数は$B_0=\CC[x,y,z,u]/(f_\tau+P)$で零でない。
実際、唯一の定義方程式は三次であり、非零の二次式を消さない。
可換な定数変形へ同型な代数は可換なので、この一次変形は自明ではない。

最後に、点曲線は一定の$E_\tau$であり、$p_{h,t}$はその上の点である。
原点$o$の局所パラメータ$z/x$で
\[
 (z/x)(p_{h,t})=-\frac h{1-ht/2}=-h+O(h^2)
\]
となる。標数零の楕円曲線の形式群では、任意の正整数$r$について
$[r](v)=rv+O(v^2)$である。
従って$[r]p$の局所パラメータは$-rh+O(h^2)$であり、恒等的に零ではない。
よって$p$は関数体上で無限位数である。
\end{proof}

\begin{corollary}[固定したHesse係数でのgenericな位数]\label{er:cor:nontorsion}
固定した$t^3\ne-1$に対しても、$p_{h,t}$は$\CC(h)$上で無限位数である。
各固定整数$r\ge1$に対し、$[r]p_{h,t}=o$となる$h\in\mathscr U_t$は有限個である。
従って非常に一般的な$h$で$p_{h,t}$は無限位数である。
\end{corollary}
\begin{proof}
滑らかな楕円曲線族の零切断に沿う形式群にも$[r](v)=rv+O(v^2)$が成り立つので、
前の証明と同じ一次係数が非torsion性を与える。
二つの切断$[r]p$と$o$の一致集合は底の閉集合であり、恒等ではない。
一次元の有限型既約底の真の閉集合は有限である。
全$r$の例外集合の和は可算なので、最後の主張が従う。
\end{proof}

\section{相対中心商と、相対復元に残る一段階}\label{er:sec:relative}
この節では$t$を固定し、
$R=\CC[h,(U(h,t)b(h,t)\Delta(h,t))^{-1}]$とする。
$\mathscr E\to\Spec R$は$f_{\vartheta}=0$で定まる滑らかな楕円曲線族、
$\mathcal L=\OO_{\mathscr E}(1)$、$p=p_{h,t}$、$\sigma=t_p$とする。
$S_R$と$\Omega_R$は式\eqref{er:eq:relations}、\eqref{er:eq:center}で定義する。

\begin{proposition}[中心商の相対THCR表示]\label{er:prop:relative-THCR}
$\mathcal L_d=\mathcal L\otimes\sigma^*\mathcal L\otimes\cdots\otimes(\sigma^{d-1})^*\mathcal L$
と置くと、自然な次数付き同型
\begin{equation}\label{er:eq:relative-THCR}
 S_R/(\Omega_R)\simeq
 R\oplus\bigoplus_{d\ge1}H^0(\mathscr E,\mathcal L_d)
\end{equation}
が存在する。右辺の積は$\sigma$で捩った切断の積である。
この同型と両辺の各次数成分の形成は、任意の可換$R$-代数への底変更と両立する。
\end{proposition}
\begin{proof}
各$S_{R,d}$は有限生成$R$-加群であり、全閉fiberにおける次元は一定である。
$R$は$\CC[h]$の局所化であるからPIDであり、torsion成分があればその台でfiber次元が増える。
従って各次数は有限自由である。
同様に$\Omega_R$の乗法の余核も全fiberでHilbert級数
$(1-z^3)/(1-z)^3$を持ち、次数別有限自由である。

$d\ge1$では各fiberの$\mathcal L_d$の次数は$3d>0$である。
楕円曲線の$H^1$の消滅とcohomology and base changeにより、
$H^0(\mathscr E,\mathcal L_d)$は階数$3d$の有限局所自由$R$-加群で、その形成は任意の底変更と可換である。
$d=0$についても$\pi_*\OO_{\mathscr E}=\OO_{\Spec R}$で同じ主張が成立する。

三つの座標切断から$S_R$の右辺への写像を作る。
二次関係式が満たされることは各複素fiberで点曲線のgraphによる表示から従う。
右辺の各次数が$R$-torsion-freeなので、この恒等式は$R$上でも成立する。
同じ理由で$\Omega_R$の像は零である。
得られた商からの写像は$h\ne0$では古典的THCR表示、$h=0$では平面三次曲線の
projective normalityによる通常のsection ring表示であり、全fiberで同型である。
各次数の有限自由加群間の写像にNakayamaの補題を適用すると\eqref{er:eq:relative-THCR}を得る。
各次数が自由であり、切断の形成も底変更と可換なので、最後の主張も従う。
\end{proof}

\begin{remark}[証明できた相対命題と残る問題]\label{er:rem:relative-gap}
命題\ref{er:prop:relative-THCR}は、非可換三次元ブロックの中心商を実在する代数的底上で記述している。
しかし、それだけで
$\MF_R(S_R,\Omega_R)$と$\Perf(\mathscr E/R)$を
任意の底変更と両立して同一視したわけではない。
さらに$P(u)$を加えた残余圏について、対角のperfect性まで得たわけでもない。

正則な底についてのOrlov比較の候補はBurke--Stevenson
\cite[Definition 5.1, Theorem 6.4]{BurkeStevenson}である。
同定理には、全体の代数の左右noetherian性、有限injective dimension、
$\RHom_A(A_0,A)$のGorenstein型、$\chi$条件が必要である。
ファイバーごとのAS-regular性だけを、これらの相対的な条件の証明として扱ってはならない。
本稿はこの経路で相対同値とその底変更compatibilityを全て確認したとは主張しない。

明示的な滑らかな族$Y\to\Spec R$の相対Kuznetsov成分はsmooth properな可換模型として存在する。
これと実際の非可換量子化の相対残余圏を同定することが次の課題である。
一般の中心混合項および複数の無限位数ブロックの全生成は、さらに別の問題として残る。
\end{remark}
\section{具体的な可換模型と中心式の可換化の相違}
\label{em:sec:comparison}

可換復元には、候補となる方程式を取り出すことと、その方程式の
Kuznetsov成分へのenhanced同値を作ることの二段階がある。
一楕円ブロックの分離型では、前者は二次関係式の多重線形化から計算でき、
後者には定理\ref{er:thm:main}のブロック比較を使える。
本節ではこの具体性を明示し、三次中心元の表示から積の順序を消す操作が、
その可換模型を与えないことを証明する。
一般混合型への圏同値は本節の結論に含めない。
Hessianによる点曲線公式はMori--Smithの既知の結果であり、
本節では量子化の規約に合わせた証明を示し、既存の圏比較への接続を明示する。
以下の二次比較と混合項の障害も初等的に証明できる。
それらを先行研究に対して新規なものとする確認は未完了である。

\subsection{三変数対称Jacobian量子化のHessian公式}

三変数の任意の斉次三次式$f\in\CC[x,y,z]_3$に対し、
$\operatorname{Hess}(f)=\det(\partial_i\partial_j f)$と書く。
これはHessian\emph{行列}ではなく、その行列式である。
以下の多項式恒等式は$f$の滑らかさを仮定しない。
点曲線の記述はMori--Smith
\cite[Lemma 7.1, Theorem 7.2, pp.~23--25]{PMoriSmith}にある。
同論文の$A_{f,\lambda}$の規約では、行列式は
$\lambda^3\operatorname{Hess}(f)/8+3\lambda f$である。
ここでは関係式の符号を合わせて、その式を$-3h$で正規化する。

\begin{proposition}[Mori--SmithのHessian公式の量子化表示]
\label{em:prop:hessian-formula}
対称Jacobian関係式
\[
 [y,z]=h\sym(\partial_xf),\qquad
 [z,x]=h\sym(\partial_yf),\qquad
 [x,y]=h\sym(\partial_zf)
\]
を多重線形化すると、$v=(x,y,z)^{\mathsf T}$について係数行列は
\[
 M_f(v)=K(v)-\frac h2 H_f(v),\qquad
 K(v)=\begin{pmatrix}0&-z&y\\z&0&-x\\-y&x&0\end{pmatrix},
 \quad H_f=(\partial_i\partial_j f)
\]
となる。その行列式は
\begin{equation}\label{em:eq:hessian-identity}
 \det M_f=-3h\left(f+\frac{h^2}{24}\operatorname{Hess}(f)\right).
\end{equation}
従って$h\ne0$では、多重線形化の退化三次式は
\begin{equation}\label{em:eq:hessian-candidate}
 f_h^{\mathrm{pt}}=f+\frac{h^2}{24}\operatorname{Hess}(f)
\end{equation}
で与えられる。この式は、退化三次式が恒等的に零になる場合も含む恒等式である。
実際の点曲線との同定には、対応するAS-regular代数と点加群についての追加仮定が必要である。
\end{proposition}
\begin{proof}
二次式$q$の対称化の多重線形化は
$\frac12\sum_j(\partial_jq)(v)v'_j$である。
$q=\partial_i f$と置けば行列表示を得る。
任意の対称$3\times3$行列$H$と$\varepsilon$について、
$K^{\mathsf T}=-K$より
\[
 \det(K-\varepsilon H)=-\det(K+\varepsilon H).
\]
左辺は$\varepsilon$の三次以下の奇多項式であり、
$\operatorname{adj}K=vv^{\mathsf T}$なので一次係数は
$-\operatorname{tr}(\operatorname{adj}K\,H)=-v^{\mathsf T}Hv$である。
三次係数は$-\det H$であり、従って
\[
 \det(K-\varepsilon H)=-\varepsilon v^{\mathsf T}Hv-
 \varepsilon^3\det H.
\]
三次のEuler等式を二回用いると$v^{\mathsf T}H_fv=6f$である。
$\varepsilon=h/2$を代入すれば\eqref{em:eq:hessian-identity}を得る。
\end{proof}

この公式は座標選択とも整合する。実際、
$g(x)=\lambda f(Lx)$、$L\in\mathrm{GL}_3(\CC)$、$\lambda\ne0$なら
\begin{equation}\label{em:eq:hessian-covariance}
 \operatorname{Hess}(g)(x)
 =\lambda^3\det(L)^2\operatorname{Hess}(f)(Lx),\qquad
 g_h^{\mathrm{pt}}(x)=\lambda f_{h'}^{\mathrm{pt}}(Lx),
 \quad h'=\lambda\det(L)h.
\end{equation}
第一式は$H_g=\lambda L^{\mathsf T}H_f(Lx)L$の行列式を取れば従い、
第二式は代入で得られる。$h'$はJacobian Poisson括弧の座標変換則と一致する。
固定した体積形式の下での式であり、Hessianの数値係数を座標変換と独立に扱わない。

\begin{corollary}[一般の滑らかな一ブロックの候補式]
\label{em:cor:hessian-reconstruction}
$g$は滑らかな平面三次曲線を定め、$\Omega_h^g$は
系\ref{er:cor:ternary}の正規化された中心元とする。
同系のパラメータ開集合と$P(u)$の仮定の下では、
定理\ref{er:thm:main}の圏比較を入力とすることにより、
\begin{equation}\label{em:eq:general-hessian-output}
 V\left(g(x)+\frac{h^2}{24}\operatorname{Hess}(g)(x)+P(u)\right)
\end{equation}
が$S_h^g[u]/(\Omega_h^g+P)$のrotation付き可換模型となる。
これは同系の模型の射影同型な表示であり、一般混合型の復元定理ではない。
\end{corollary}
\begin{proof}
$g=\lambda f_t(Lx)$とHesse形に移し、$h'=\lambda\det(L)h$とする。
直接計算で
\[
 \operatorname{Hess}(f_t)=
 2(t^3+4)xyz-2t^2(x^3+y^3+z^3),\qquad
 f_t+\frac{h'^2}{24}\operatorname{Hess}(f_t)
 =b(h',t) f_{\vartheta(h',t)}
\]
を得る。従って\eqref{em:eq:hessian-covariance}により
\eqref{em:eq:general-hessian-output}の方程式は
$\lambda b(h',t) f_{\vartheta(h',t)}(Lx)+P(u)$となる。
$b(h',t)\ne0$なので、三変数を一斉に$b(h',t)^{-1/3}$倍すれば、
系\ref{er:cor:ternary}の方程式
$\lambda f_{\vartheta(h',t)}(Lx)+P(u)$になる。
この座標変換と同系の圏同値により結論を得る。
ここで用いた三乗根は複素ファイバーごとに選べる。
族上の同型を選ぶ際には、必要な有限\'etale被覆を別途指定する。
\end{proof}

\subsection{二次関係式から取り出せる分離型の方程式}

$S_{h,t}$、$\Omega_{h,t}$、$P(u)$および$\mathscr U_t$は
第\ref{er:sec:intro}節のものとする。
$t^3\ne-1$、$h\in\mathscr U_t$、$m\ge2$とし、$P$のアフィン臨界点は原点のみとする。
三次中心商は
\[
 B_{h,t,P}=S_{h,t}[u_1,\ldots,u_m]/(\Omega_{h,t}+P(u))
\]
である。既存のSklyanin中心商と楕円曲線の比較
\cite[pp.~20--21]{SmithTate}\cite[Sections 3.5--3.6]{EG}だけでなく、
内部シフトとの整合性および中心変数を加えた後の生成性を使う点が重要である。

\begin{proposition}[分離型の具体的復元手順]
\label{em:prop:determinant}
$\alpha=1-ht/2$、$\beta=-1-ht/2$、$\gamma=-h$と置く。
関係式を第二の三変数について線形にした行列を
\[
 M(x,y,z)=
 \begin{pmatrix}
 \gamma x&\beta z&\alpha y\\
 \alpha z&\gamma y&\beta x\\
 \beta y&\alpha x&\gamma z
 \end{pmatrix}
\]
とすると
\begin{equation}\label{em:eq:determinant}
 -\det M=\alpha\beta\gamma(x^3+y^3+z^3)
 -(\alpha^3+\beta^3+\gamma^3)xyz.
\end{equation}
$h\ne0$では、$b=1-h^2t^2/4$を用いて
\begin{equation}\label{em:eq:explicit-output}
 F_{\mathrm{rec}}(h)
 =-\frac{\det M}{3hb}+P(u)
 =\frac{x^3+y^3+z^3}{3}+\vartheta(h,t)xyz+P(u)
\end{equation}
を得る。右辺は$h=0$まで正則に延長する。
定理\ref{er:thm:main}の圏比較を入力とすれば、
$Y_h=V(F_{\mathrm{rec}}(h))$は$B_{h,t,P}$のrotation付き可換模型である。
\end{proposition}
\begin{proof}
多重線形化された関係式は
$M(x,y,z)(x',y',z')^{\mathsf T}=0$である。
行列式を六項に展開すると\eqref{em:eq:determinant}を得る。
さらに
\[
 \alpha\beta\gamma=hb,\qquad
 \alpha^3+\beta^3+\gamma^3
 =-3h\left(t+\frac{h^2}{3}(1+t^3/4)\right)
\]
を代入すれば\eqref{em:eq:explicit-output}となる。
零パラメータでは行列式全体が零になるため、共通因子$h$を除いてから延長する。
これは$S_{0,t}$のpoint schemeを楕円曲線と同一視する操作ではない。
最後の圏同値は定理\ref{er:thm:main}そのものであり、
行列式の計算だけから従うものではない。
\end{proof}

一般の三次元Sklyanin関係式
$\alpha yz+\beta zy+\gamma x^2=0$とその巡回置換で与えられた場合も、
滑らかな点曲線および標準中心三次商であることを確認したうえで、
\eqref{em:eq:determinant}が点曲線の方程式を直接与える。
分離型では方程式の非零スカラー倍はブロック変数の一斉拡大で吸収できる。
その拡大に必要な三乗根は$\CC$上で存在するが、一般の基礎体上では底拡大の指定が必要である。
これは元の古典極限の方程式を再利用する手順ではなく、
与えられた量子関係式から変形した楕円曲線を取り出す手順である。

\subsection{中心式の積を可換にすると二次で違う模型になる}

$S_3=x^3+y^3+z^3$と置き、中心元の指定された対称表示
$\Omega_{h,t}=U(h,t)S_3/3+W(h,t)\sym(xyz)$の積の順序を消した多項式を
\[
 F_{\mathrm{sym}}(h)=U(h,t)S_3/3+W(h,t)xyz+P(u)
\]
と定義する。これは選んだ表示に依存する操作であって、
代数の可換化とは異なる。可換化との相違は次の小節で扱う。

\begin{proposition}[二次Jacobian類の不一致]
\label{em:prop:second-order-failure}
$F_0=f_t+P$、$J(F_0)=\CC[x,y,z,u]/(\partial F_0)$とする。
$F_{\mathrm{rec}}(h)$と$F_{\mathrm{sym}}(h)$はいずれも$h$の一次項を持たず、
二次係数のJacobian類はそれぞれ
\begin{equation}\label{em:eq:two-classes}
 \frac{1+t^3}{3}[xyz],\qquad
 -\frac{2(1+t^3)}{3}[xyz]
 \quad\text{in }J(F_0)_3
\end{equation}
である。差$(1+t^3)[xyz]$は非零である。
従って二つの偏極付きcubic族は、$\CC[h]/(h^3)$上で
特殊ファイバーの恒等写像を保つ射影座標変換によって同型にならない。
特に、中心式の積の順序を消す操作は、分離型でもmarkedな可換復元を与えない。
\end{proposition}
\begin{proof}
式\eqref{er:eq:coefficients}を展開すると
\begin{align*}
 \vartheta(h,t)&=t+\frac{1+t^3}{3}h^2+O(h^4),\\
 U(h,t)&=1+\frac{t^2}{2}h^2+O(h^4),\\
 W(h,t)&=t-\frac{t^3+4}{6}h^2+O(h^4).
\end{align*}
$J(f_t)$では
$x^2+tyz=y^2+tzx=z^2+txy=0$なので
$S_3=-3txyz$である。従って$F_{\mathrm{sym}}$の二次係数の類は
\[
 \left[\frac{t^2}{6}S_3-\frac{t^3+4}{6}xyz\right]
 =-\frac{2(1+t^3)}{3}[xyz]
\]
となる。

$[xyz]\ne0$も有限な線形代数で直接確認できる。
次数三の単項式を
\[
 (x^3,x^2y,x^2z,xy^2,xyz,xz^2,y^3,y^2z,yz^2,z^3)
\]
の順に並べる。列を
$(x\partial_xf_t,y\partial_xf_t,z\partial_xf_t,
x\partial_yf_t,y\partial_yf_t,z\partial_yf_t,
x\partial_zf_t,y\partial_zf_t,z\partial_zf_t,xyz)$
とする係数行列の行列式は
\[
 -(1+t^3)^2.
\]
従って$t^3\ne-1$のとき、最初の九列は独立であり、$xyz$はその張る空間に含まれない。
変数の分離から$J(F_0)=J(f_t)\otimes J(P)$であり、
$J(P)_0=\CC$なので$[xyz]\otimes1$も非零である。

最後に、$F_0$は滑らかなcubicなので、その偏微分は二次の正則列をなす。
これらに一次係数のsyzygyはなく、Euler等式を除いて無限小射影自己同型はない。
従って一次項が零である二つのmarked族の間の一次座標変換は射影的に零である。
二次座標変換と方程式の単元倍で変えられる二次項は$(\partial F_0)_3$に限られる。
差のJacobian類が非零なので、求める同型は存在しない。
\end{proof}

\begin{example}[Fermatブロック]
\label{em:ex:fermat}
$t=0$、$P(u,v)=u^3+v^3$では
\begin{align*}
 F_{\mathrm{rec}}(h)&=\frac{x^3+y^3+z^3}{3}+u^3+v^3
                       +\frac{h^2}{3}xyz,\\
 F_{\mathrm{sym}}(h)&=\frac{x^3+y^3+z^3}{3}+u^3+v^3
                       -\frac{2h^2}{3}xyz.
\end{align*}
両式の差は正確に$h^2xyz$である。
ここで$J(F_0)=\CC[x,y,z,u,v]/(x^2,y^2,z^2,u^2,v^2)$なので、
差の非零性は平方自由単項式$xyz$だけで確認できる。
この判定はmarked変形の判定であり、
有限自己同型群で割った粗モジュライの局所座標について、同じ次数の主張を加えていない。
\end{example}

\subsection{代数そのものの可換化は次元さえ保たない}

\begin{proposition}[文字通りの可換化の失敗]
\label{em:prop:abelianization}
$t=0$、$h\ne0$とする。$B_{h,0,P}$を、全ての交換子が生成する両側イデアルで割った
可換代数を$B_{h,0,P}^{\mathrm{ab}}$と書く。このとき
\[
 B_{h,0,P}^{\mathrm{ab}}
 \simeq
 \CC[x,y,z,u_1,\ldots,u_m]
 /(x^2,y^2,z^2,P(u)-2h^2xyz/3).
\]
その被約環は$\CC[u]/(P)$であり、$\operatorname{Proj}$の次元は$m-2$である。
従って、次元$m+1$の復元先cubicをこの可換化から得ることはできない。
\end{proposition}
\begin{proof}
関係式$[y,z]=hx^2$とその巡回置換に可換性を課すと、$h\ne0$から
$x^2=y^2=z^2=0$が従う。
中心方程式の三乗項は消え、表示した最後の関係式が残る。
$x,y,z$は冪零なので、被約化後に残る関係は$P=0$である。
$P$は原点以外に臨界点を持たない非零三次式であり、特に被約である。
よって被約環と次元の主張を得る。
\end{proof}

\subsection{非楕円混合型で次に計算すべきもの}

節点型と尖点型で用いる三変数三次式は
\[
 f_{\mathrm N}=y^3+z^3+xyz,\qquad
 f_{\mathrm C}=xz^2+y^3
\]
であり、全体の可換特殊ファイバーには$ux^2$という混合項がある。
この項は、一ブロック定理で許される中心変数方向への線形平行移動に吸収できない。

\begin{lemma}[混合項$ux^2$は中心シフトで除去できない]
\label{em:lem:essential-mixing}
$\epsilon=\mathrm N,\mathrm C$のいずれについても
$x^2\notin\langle\partial_xf_\epsilon,\partial_yf_\epsilon,
\partial_zf_\epsilon\rangle\subset\CC[x,y,z]_2$である。
従って$f_\epsilon+ux^2$から$u$の一次混合項を
$(x,y,z)\mapsto(x+a u,y+b u,z+c u)$によって除去できない。
\end{lemma}
\begin{proof}
点$(x,y,z)=(1,0,0)$では両方の$f_\epsilon$の三つの偏微分が全て零であり、
$x^2=1$である。よって表示した線形空間に$x^2$は属さない。
中心シフトで$f_\epsilon$の$u$の一次係数に加わるものは
$a\partial_x f_\epsilon+b\partial_y f_\epsilon+c\partial_z f_\epsilon$だけなので、結論を得る。
\end{proof}

従って非楕円二例には、単に三変数のpoint curveを取り出して
定理\ref{er:thm:main}を適用することはできない。
候補式を得るためには、全七変数の量子関係式と中心三次元から残余圏の
rotation付き乗法データを計算するか、混合項を保つ具体的な行列因子分解同値を作る必要がある。
候補cubicが滑らかであることや、有限個のHochschild次元が一致することだけでは同値は証明されない。
また中心変数を通常の偏微分で扱う生成性の議論は、非楕円混合型では
その偏微分が代数導分であることから確かめ直す必要がある。

\subsection{証明入力と厳密検算の範囲}

系\ref{em:cor:hessian-reconstruction}と命題\ref{em:prop:determinant}の圏同値部分は
定理\ref{er:thm:main}と、その証明で用いるdg window、
periodicityを保つブロック比較、中心変数に対する外部積生成、
有限群の双対同変化\cite[Theorem 4.2, Section 8]{Elagin}を入力としている。
行列式や中心係数の一致は、これらの圏論的入力の独立監査を代替しない。
命題\ref{em:prop:second-order-failure}、命題\ref{em:prop:abelianization}、
補題\ref{em:lem:essential-mixing}の証明は上記の初等的計算で完結する。
付属の\texttt{verify\_explicit\_model\_comparison.py}は、
一般の十係数三次式に対するHessian恒等式、Hesse形への特殊化、
多重線形化行列の行列式、二次係数、十次元係数行列の行列式、
Fermat例のJacobian剰余、混合項の評価を有理数上で厳密に検算する。
この検算は圏同値、一般混合型の可換復元、残余smoothnessの計算機証明ではない。

\section{最小混合試験族の三次中心元の閉公式}
\label{mt:sec:family}

一楕円ブロックの分離型から一般混合型へ進むため、
五変数のcubic
\begin{equation}\label{mt:eq:classical}
 F_\lambda=\frac{x^3+y^3+z^3}{3}+u^3+v^3+\lambda uxy
\end{equation}
を固定する。$u,v$を中心、$h,\lambda$を次数零のパラメータとし、
\begin{equation}\label{mt:eq:relations}
 [y,z]=h(x^2+\lambda uy),\qquad
 [z,x]=h(y^2+\lambda ux),\qquad
 [x,y]=hz^2
\end{equation}
で周囲の代数$A_{h,\lambda}$を定める。
本節ではこの族の三次中心元を多項式の閉公式で構成し、
自由代数内の三つの恒等式によって中心性を証明する。
これは可換復元先の二次係数$G_2$の計算ではなく、
それを計算するための非可換代数側の入力である。
中心式の特殊形は既存の混合中心式の有限体検算に現れるが、
ここでは特殊化を用いず、全パラメータに対する証明書を与える。
閉公式の先行研究上の初出は未確認である。

\begin{lemma}[古典特殊ファイバーの滑らかさ]
\label{mt:lem:classical-smooth}
$V(F_\lambda)\subset\PP^4$は$\lambda^3\ne-3$のとき、かつそのときに限り滑らかである。
\end{lemma}
\begin{proof}
全偏微分が零なら$z=v=0$であり、残る方程式は
\[
 x^2+\lambda uy=y^2+\lambda ux=3u^2+\lambda xy=0.
\]
$u=0$なら$x=y=0$となるので射影点は存在しない。
$u=1$とすると、$\lambda=0$では解がなく、$\lambda\ne0$なら$x,y$は非零である。
最初の二式を掛けると$xy=\lambda^2$を得て、第三式は$3+\lambda^3=0$となる。
逆にこの等式が成立すれば、$x^3=3$を満たす$x$と
$y=-x^2/\lambda$を選ぶことで解を得る。
三次のEuler等式により、これらの解は$F_\lambda=0$も満たす。
\end{proof}

\subsection{中心三次式と自由代数の証明書}

$a=\lambda u$と略記し、$T=xyz+xzy+yxz+yzx+zxy+zyx$と置く。

\begin{proposition}[混合試験族の中心元]
\label{mt:prop:center}
\begin{equation}\label{mt:eq:center}
 \boxed{\begin{aligned}
 \Omega_{h,\lambda}={}&
 \frac{x^3+y^3+z^3}{3}-\frac{h^2}{9}T
 +\frac a2(xy+yx)\\
 &+\frac{h^2a^2}{6}z+\frac{h^4a}{18}z^2+u^3+v^3.
 \end{aligned}}
\end{equation}
は$A_{h,\lambda}$の三次中心元である。
これは$\CC[h,\lambda]$上の恒等的な中心性であり、
$h\ne0$への局所化も、有限$h$-jetへの制限も必要としない。
$h=0$で$\Omega_{0,\lambda}=F_\lambda$である。
\end{proposition}
\begin{proof}
まず$a$を独立な中心係数とし、自由代数
$\mathbb Q[h,a]\langle x,y,z\rangle$で計算する。
関係式を
\[
 r_1=xy-yx-hz^2,\quad
 r_2=xz-zx+hy^2+hax,\quad
 r_3=yz-zy-hx^2-hay
\]
と書く。二つの演算を
\[
 \mathcal S_p(r)=rp^2+prp+p^2r,\qquad
 \mathcal D_{p,q}(r)=rpq-prq-qrp+qpr
\]
とし、$\{p,q\}_+=pq+qp$とする。
$\omega=\Omega_{h,\lambda}-u^3-v^3$について、以下の三式が
\emph{関係式で割る前の自由代数}において成立する：
\begin{align}
 [\omega,x]={}&
 -\frac13\mathcal S_y(r_1)-\frac13\mathcal S_z(r_2)
 +\frac{h^2}{18}
   \bigl(\mathcal D_{z,x}(r_1)+\mathcal D_{y,x}(r_2)\bigr)\notag\\
 &+\frac{ah^3}{18}[x,r_1]-\frac a2\{x,r_1\}_+
  +\frac{ah}{6}[z,r_2]+
  \frac h6[yz+zy,r_3],\label{mt:eq:certificate-x}\\
 [\omega,y]={}&
 \frac13\mathcal S_x(r_1)-\frac13\mathcal S_z(r_3)
 -\frac{h^2}{9}\mathcal D_{z,y}(r_1)
 +\frac{ah^3}{18}[y,r_1]+\frac a2\{y,r_1\}_+\notag\\
 &+\frac{ah^2}{18}[r_2,x]+\frac h9[r_2,2xz+zx]
 +\frac{h^2}{9}[r_2,y]y+\frac{ah}{6}[r_3,z],
 \label{mt:eq:certificate-y}\\
 [\omega,z]={}&
 \frac{a(2h^3+9)}{18}[r_1,z]
 +\frac h9[2xy+yx,r_1]+\frac{h^2}{9}[r_1,z]z\notag\\
 &+ayr_2+\frac13\mathcal S_x(r_2)
 -\frac{h^2}{9}\mathcal D_{y,z}(r_2)
 +ar_3x+\frac13\mathcal S_y(r_3).
 \label{mt:eq:certificate-z}
\end{align}
各式は語と$h,a$の冪ごとに展開して確認でき、全係数は有理数である。
右辺の各項は$r_i$を含むので、商代数で全交換子が零となる。
$a=\lambda u$を代入し、中心スカラー$u^3+v^3$を戻せば中心性を得る。
零パラメータでの表示は直接代入による。
\end{proof}

\subsection{Artin切断でのsorted PBW表示}

閉公式の中心性と、周囲の代数のflatnessを別々に確認する。
$u,v$の中心性を先に使い、active語を$z<y<x$の順に並べる。
還元規則は
\begin{equation}\label{mt:eq:rewrite}
 xy\longmapsto yx+hz^2,\qquad
 xz\longmapsto zx-hy^2-hax,\qquad
 yz\longmapsto zy+hx^2+hay
\end{equation}
である。これは一般の非零$h$で通常の語順序に関して停止する規則とは主張しない。

\begin{lemma}[全Artin切断でのPBW性]
\label{mt:lem:artin-pbw}
任意の$N\ge0$について、$h^{N+1}=0$を課した代数は
$\CC[\lambda,u,v,h]/(h^{N+1})$上で
$\{z^iy^jx^k:i,j,k\ge0\}$を基底に持つ。
従って各次数成分は$\CC[\lambda,h]/(h^{N+1})$上自由であり、
周囲のHilbert級数は$(1-t)^{-5}$である。
\end{lemma}
\begin{proof}
$h^r w$の還元に、辞書式に並べた二つの非負整数
$(N-r,\operatorname{inv}(w))$を用いる。
$h$を含まない交換項では反転数が減り、補正項では$h$の冪が増える。
従って$h^{N+1}=0$のもとで全還元が停止する。
active語の重なりは$xyz$の一つだけである。
最初に$xy$を還元した場合と、最初に$yz$を還元した場合は、いずれも有限回の置換で
\begin{equation}\label{mt:eq:overlap}
 zyx+h(x^3-y^3+z^3)
\end{equation}
になる。$a$を含む項は両経路でそれぞれ相殺する。
包含の曖昧さはない。係数に逆元を導入しないmonicな合流的還元なので、
既約語$z^iy^jx^k$が基底となる。
中心係数$a$を$\lambda u$へ置き換え、独立な中心変数$v$を加えても同じ証明が成り立つ。
\end{proof}

この補題は全ての$h$-Artin切断およびその逆極限についてのformal PBW性を示す。
非零$h$での全ファイバーのAS-regular性や代数的族全体のflatnessを、
Artin切断の計算だけから結論していない。
また、中心商は$v$に関してmonicな三次式であるため、
各Artin切断で$1,v,v^2$を、$v$を含まない周囲の代数上の基底に持つ。
従って中心三次商も次数別に平坦で、Hilbert級数は
$(1-t^3)/(1-t)^5$である。

\subsection{二次jetと、復元先に残る計算}

sorted PBW表示に直すと\eqref{mt:eq:center}は
\begin{align}\label{mt:eq:second-jet}
 \Omega_{h,\lambda}\equiv{}&
 \frac{x^3+y^3+z^3}{3}+a\,yx+u^3+v^3
 +\frac{ha}{2}z^2\notag\\
 &+h^2\left(-\frac23zyx+\frac{a^2}{6}z\right)
 \pmod{h^3}.
\end{align}
対称表示では偶数冪だけに見える式が、sorted表示では一次項を持つ。
これは積と基底表示の違いであり、可換復元先の一次変形が非零だという意味ではない。
$h^2a^2z/6$の補正は必要であり、これを省いた元の$x$との交換子は
\[
 -\frac{h^3a^2}{6}(y^2+ax)\pmod{h^4}
\]
となる。従って古典混合項を対称化して加えるだけでは中心にならない。

$\lambda=0$では既存の分離型に戻る。
一方、古典式の混合方向$uxy$は
\[
 J(F_0)_3=\bigl(\CC[x,y,z,u,v]/(x^2,y^2,z^2,u^2,v^2)\bigr)_3
\]
で非零であり、
active変数の中心シフトが作る$u\langle x^2,y^2,z^2\rangle$にも属さない。
従って中心部分空間を固定したactive変数の線形シフトだけで、
この混合方向を分離型へ戻すことはできない。

\subsection{対角対称性による二次復元候補の一方向への制限}

固定した$\lambda$について、有限対角群
\[
 H=\{(\alpha,\beta,\gamma,\delta,\varepsilon)\in\mu_3^5:
       \alpha\beta\gamma=1\},\qquad
 \chi(\alpha,\beta,\gamma,\delta,\varepsilon)=\alpha\beta\delta
\]
を考える。作用は
\[
 (x,y,u,z,v,h)\longmapsto
 (\alpha x,\beta y,\gamma u,\delta z,\varepsilon v,\chi h)
\]
で与える。$\lambda$は固定する。

\begin{lemma}[量子族の同変性とJacobian固有空間]
\label{mt:lem:character-space}
この作用は\eqref{mt:eq:relations}の関係式イデアルを保ち、
$\Omega_{h,\lambda}$および$F_\lambda$を固定する。
$J_\lambda=\CC[x,y,u,z,v]/(\partial F_\lambda)$とすると、
次数三の固有空間は
\begin{equation}\label{mt:eq:character-spaces}
 (J_\lambda)_3^\chi=\CC\cdot[xyz],\qquad
 (J_\lambda)_3^{\chi^2}=0.
\end{equation}
ここで固有空間の計算は、$\lambda^3=-3$を含む全ての$\lambda$について成り立つ。
\end{lemma}
\begin{proof}
例えば最初の関係式は全体が$\beta\delta$倍になる。
実際、$\chi\alpha^2=\beta\delta$、
$\chi\gamma\beta=\beta\delta$である。
他の二関係式も同様である。
中心式の各項の重みはそれぞれ一である。
特に$h^2T$の重みは$\chi^3=1$、
$h^2a^2z$の重みは$\chi^2\gamma^2\delta=1$、
$h^4az^2$の重みは$\chi^4\gamma\delta^2=1$となる。

固有空間は単項式の列挙で求まる。
$\gamma=(\alpha\beta)^{-1}$を消去し、独立な文字を
$(\alpha,\beta,\delta,\varepsilon)$とすると、
$(x,y,u,z,v)$の重みはそれぞれ
\[
 (1,0,0,0),\ (0,1,0,0),\ (2,2,0,0),\
 (0,0,1,0),\ (0,0,0,1)
 \quad\text{in }(\ZZ/3)^4
\]
である。多項式環の次数三の$\chi=(1,1,1,0)$固有空間は
$\langle xyz,u^2z\rangle$だけであり、
Jacobianイデアルの同じ固有空間は
\[
 \langle z\partial_uF_\lambda\rangle
 =\langle\lambda xyz+3u^2z\rangle
\]
だけである。後者の$u^2z$の係数は非零なので、
商は$[xyz]$で生成される一次元空間となる。
同様に$\chi^2=(2,2,2,0)$固有空間は
多項式環では$\langle uz^2\rangle$だけであり、
$u\partial_zF_\lambda=uz^2$がこれを消す。
有限群の標数零の表現は半単純なので、この計算は商の固有空間を正確に与える。
\end{proof}

\begin{corollary}[marked同変復元を仮定した二次係数の制約]
\label{mt:cor:second-character}
$\lambda^3\ne-3$とする。
非可換族に対するmarkedな可換復元が$\CC[h]/(h^3)$上に存在し、
上記の$H$作用が復元先の偏極付きcubic族に持ち上がって、
超平面束の特殊ファイバー上の所定の線形化も持ち上がり、
底の$h\mapsto\chi h$と両立すると仮定する。
その場合、復元先の一次変形類は零であり、
一次項を除去した後の二次変形類は
\begin{equation}\label{mt:eq:unknown-scalar}
 [G_2]=c(\lambda)[xyz]\in (J_\lambda)_3
\end{equation}
という一つの係数で表される。
本系は復元の存在や、その同変性を証明するものではない。
\end{corollary}
\begin{proof}
有限群について平均を取ることにより、
超平面束の切断加群の基底を、特殊ファイバーの上記の対角表現を保つように持ち上げられる。
方程式の生成元も、$F_\lambda$への留数を保って$H$不変に選べる。
従って方程式を
$F_\lambda+hG_1+h^2G_2$と書いたとき、
$G_i$の重みは$\chi^{-i}$である。
補題\ref{mt:lem:character-space}より$[G_1]=0$である。
実際、同変座標では$G_1$は$uz^2$のスカラー倍であり、
$z\mapsto z-hc u$という同変な座標変換で取り除ける。
この変換の同変性は$\chi\gamma=\delta$による。
その後も$G_2$の重みは$\chi^{-2}=\chi$なので、
同じ補題により\eqref{mt:eq:unknown-scalar}を得る。
\end{proof}

量子代数側の同変性は補題で直接確認したが、
可換復元への同変性の持上げは系の明示的な仮定である。
この区別を保ったうえで、残余圏の計算が求める二次変形類を
一つの係数$c(\lambda)$に絞れる。

次に必要なのは、$B_{h,\lambda}=A_{h,\lambda}/(\Omega_{h,\lambda})$の
残余dg圏からrotation付きの乗法を計算し、可換模型の
$G_2\in J(F_\lambda)_3$を取り出すことである。
本節の中心式の$h^2$係数を$G_2$と同一視しない。
残余圏のsmoothness、実際の圏同値、一般混合型の可換復元は本節では証明していない。

証明書の再現は
\texttt{python3 computations/verify\_mixed\_test\_family.py}
で行う。同コードは
\eqref{mt:eq:certificate-x}--\eqref{mt:eq:certificate-z}を
自由代数内の係数ごとに厳密検算する。
別に$h^3,h^4,h^5,h^9$を法とするsorted還元、唯一の重なり、
二次補正を省いた負例も確認し、
\texttt{mixed\_test\_family\_certificate.json}を生成する。
自由代数恒等式による全パラメータの中心性と、有限jetの独立検算を区別して記録する。

\section{同値双加群の障害消去とétale局所実現}\label{er:sec:etale}
本節では、可換模型の存在に使う局所実現原理を証明する。
入力は実在するsmooth properな代数的dg族であり、proper性だけの族ではない。
圏の全変形をHochschild複体から実現するという主張は用いず、
実在するperfect双加群のmoduliと、その完全な障害理論を用いる。
ここでのpairは自己同値一つを指定したもので、$\rho^3\simeq[2]$という
選択された同型やそのcoherenceを追加データとして含めない。

\subsection{perfect加群の小拡大障害}
小拡大は剰余体$k=\CC$を持つ局所Artin代数の全射
$A'\twoheadrightarrow A$で、核$I$が$\mathfrak m_{A'}I=0$を満たすものとする。
$D'$は$A'$上smooth properなdg代数とし、$D=D'\otL_{A'}A$とする。

\begin{lemma}[加群の完全な障害理論と特徴作用]\label{er:lem:obstruction}
$P\in\Perf(D)$の持上げには自然な障害
\[
 o(P;D')\in\operatorname{Ext}^2_{D_k}(P_k,P_k)\otimes_k I
 \tag{O}\label{er:eq:obstruction}
\]
があり、零であることと$\Perf(D')$への持上げの存在は同値である。
障害はMorita関手と対象の同型に関して自然である。
$D'=C'^{\rm op}\otL_{A'}B'$で$P$が$C$から$B$への同値を表すとき、
$B'$を固定し、$C'$のliftをHochschild類$\xi\in\mathrm{HH}^2(C_k)\otimes I$だけ変えると、
障害の変化は、符号を一貫して選べば
\[
 o_{\rm new}-o_{\rm old}=-\chi_P(\xi),\qquad
 \chi_P:\mathrm{HH}^2(C_k)\xrightarrow{\sim}
       \operatorname{Ext}^2_{C_k^{\rm op}\otimes B_k}(P_k,P_k).
 \tag{V}\label{er:eq:variation}
\]
\end{lemma}
\begin{proof}
Toën--Vaquié\cite[Theorem 3.6, Corollaries 3.15, 3.17]{TV}を
可換底$A'$と$D'$に適用する。smoothnessによりpseudo-perfect加群と
perfect加群は一致し、その実際の対象moduliの相対接複体は$\RHom_D(P,P)[1]$である。
導来対象moduliの平方零拡大に関する障害理論は、その接複体の次数一に完全な障害を持つ。
従って障害は$\operatorname{Ext}^2_D(P,P\otL_A I)$に入り、
導来随伴と$\mathfrak m_A I=0$によって\eqref{er:eq:obstruction}に一致する。
これは単なる接空間の計算を代数化の代用にしているのではない。
実在する対象の導来moduliが無限小凝集的であることを使っており、
障害のnull-homotopyはそのmoduliの$A'$-点、すなわち実際のperfect加群を与える。

変化式も説明する。cofibrantな代数模型と加群のbar模型を選び、
構造を次数一のbar余微分$b$で表す。任意のlift$\widetilde b$の曲率
$\widetilde b^{\,2}$は$I$に値を取り、$[b,\widetilde b^{\,2}]=0$である。
持上げを$I$値の$h$だけ変えると、$I^2=0$により曲率は$[b,h]$だけ変わる。
この曲率類が上のmoduliの障害である。
代数側を$\xi$だけ変えた際の曲率の変化は、$\xi$を加群の構造写像に挿入したcochainであり、
それが特徴作用$\chi_P(\xi)$である。符号はsourceの反対代数とtorsorの規約に合わせて
\eqref{er:eq:variation}と取る。必要なのは同じ規約を両式で使うことである。
この記述はLowen\cite[Theorem 4.8]{ERLowen}の特徴作用による障害の記述を、
bar模型でdg加群と任意の小拡大に書いたものである。
$P_k$の逆Morita双加群をテンソルすれば、その自己Homと対角の自己Homが同定されるので、
$\chi_P$は同型である。曲率の構成、またはmoduliの障害の自然性から、最後の自然性も従う。
\end{proof}

\begin{lemma}[還元後の同値とintertwiner]\label{er:lem:lift}
$C',B'$は$A'$上smooth properとする。
$P'\in\Perf(C'^{\rm op}\otL B')$の還元がMorita同値なら$P'$自身も同値である。
さらに$Q'$もperfectで、$\operatorname{Ext}^1(P_k,Q_k)=0$なら、
任意の射$P'_A\to Q'_A$は$A'$上へ持ち上がる。元の射が同型ならliftも同型となる。
\end{lemma}
\begin{proof}
perfect双加群は両側に随伴を持ち、smooth proper性により随伴双加群もperfectである。
unitとcounitのconeは底上perfectで、還元後に零であるから導来Nakayamaで零となる。
射の持上げには三角$Q'\otL I\to Q'\to Q'_A$を$\RHom(P',-)$で評価する。
射の障害は$\operatorname{Ext}^1(P_k,Q_k)\otimes I$に入り、仮定で零である。
持ち上げた射のconeに同じNakayamaを適用すれば同型性も従う。
\end{proof}

\subsection{rotationを保つ障害消去}
$H$は複素数体上smoothな有限型スキーム、$h\in H(\CC)$とする。
$\mathcal A_H$はsmooth proper dg族、$r_H$は可逆perfect双加群で表された自己同値とする。
$\mathcal A=\mathcal A_h$、$r=r_h$と置く。局所的なdg代数模型を取って議論する。

\begin{theorem}[自己同値付き局所実現の判定]\label{er:thm:criterion}
次を仮定する。
\[
 \mathrm{HH}^{<0}(\mathcal A)=0,\qquad
 \mathrm{HH}^{1}(\mathcal A)=0,
 \qquad
 T_hH\longrightarrow\mathrm{HH}^2(\mathcal A)^r
 \text{ が全射}.
 \tag{KS}\label{er:eq:KS}
\]
最後の写像は$\mathcal A_H$のKodaira--Spencer写像である。
$S$は有限型複素スキーム、$(\mathcal B_S,s_S)$はsmooth proper dg族と
可逆perfect双加群のpairとする。$s\in S(\CC)$に指定した同値
$(\mathcal A,r)\simeq(\mathcal B_s,s_s)$があるなら、
pointed étale近傍$V\to(S,s)$、$f:V\to H$、$f(v)=h$と
指定同値を延長するpairの同値
\[
 (f^*\mathcal A_H,f^*r_H)\simeq(\mathcal B_V,s_V)
\]
が存在する。
\end{theorem}
\begin{proof}
まず小拡大$A'\twoheadrightarrow A$上の持上げを証明する。
$\operatorname{Spec}A'\to S$を固定し、$A$上の点$h_A$、同値双加群$P_A$、
intertwiner $\alpha_A:s_AP_A\xrightarrow{\sim}P_Ar_A$が与えられているとする。
$H$はsmoothなので$h_A$の任意のlift$h'_0$を取れる。
両側のliftされた自己同値から、双加群圏の自己同値
\[
 \mathcal F'(P)=s'\circ P\circ(r'_0)^{-1}
\]
を得る。$\alpha_A$は$\mathcal F_A(P_A)\simeq P_A$を与える。
補題\ref{er:lem:obstruction}の自然性より、$P_A$の持上げ障害$o$は
$\mathcal F_k$の作用で不変である。
Morita同定$\operatorname{Ext}^2(P_k,P_k)=\mathrm{HH}^2(\mathcal A)$の下で、
この作用は$r$による共役作用（規約によりその逆）である。
従って
\[
 o\in\mathrm{HH}^2(\mathcal A)^r\otimes I.
 \tag{I}\label{er:eq:invariant}
\]
この段階では$\alpha_A$を$A'$上へ持ち上げたとは仮定していない。

\eqref{er:eq:KS}と\eqref{er:eq:variation}により、$h'_0$を
$T_hH\otimes I$の元だけ修正して$o$を消せる。
修正後のsourceと固定したtargetの間にperfect双加群$P'$が存在し、
補題\ref{er:lem:lift}によって同値となる。
$\alpha_A$を持ち上げる障害は
\[
 \operatorname{Ext}^1(s_kP_k,P_kr_k)\otimes I
       \simeq \mathrm{HH}^1(\mathcal A)\otimes I=0
\]
なので、intertwinerも同型として持ち上がる。
自己同値一つのintertwiner以外に群関係は課していないので、ここで必要なデータは全て揃う。
任意のnilpotent拡大は小拡大に分解できる。

次に代数化する。$H\times S$上で、perfectな
$\mathcal A_H^{\rm op}\otimes\mathcal B_S$-加群$P$と
同型$s_SP\simeq Pr_H$のmoduliを考え、$P$が同値である部分を$\mathcal E$とする。
Toën--Vaquié\cite[Theorem 3.6, Corollaries 3.15, 3.19]{TV}により、
perfect対象とその射の導来moduliは局所幾何的かつ局所有限表示である。
unit、counitおよびintertwinerのconeは底上perfectだから、その消滅は開条件である。
従って$\mathcal E$も局所幾何的かつ局所有限表示である。
同値点での自己Extの負次数は$\mathrm{HH}^{<0}(\mathcal A)=0$であり、
その近傍でこの消滅を保てる。従って古典的切断はその近傍でArtin 1-stackとなる。
ここでは局所dg代数模型を用いてTVの定理を$H\times S$上の可換係数環に適用している。

上の小拡大持上げは、指定点$e=(h,s,P_k,\alpha_k)$で
$\mathcal E\to S$が形式的にsmoothであることを示す。
局所有限表示のArtin stackの無限小判定により、この射は$e$でsmoothである。
$e$を覆うsmooth atlas $U\to\mathcal E$を取り、点の近傍へ縮めると$U\to S$もsmoothである。
滑らかな射のpointed étale局所section\cite[Tags 02HX, 055U]{StacksER}を使い、
$V\to S$と$V\to U\to\mathcal E$を得る。
この普遍対象が所要の族、双加群およびintertwinerを与える。
\end{proof}

\begin{remark}[使わない仮定と先行研究]\label{er:rem:etale-scope}
この証明は$\mathrm{HH}^3=0$、Hochschild cochainのformality、
抽象的なHochschild変形の全てがactual Morita変形になること、または$\rho^3=[2]$の指定を使わない。
Lin--Zhang\cite[Theorem 2.2, Proposition 2.4]{ERLinZhang}の同値核の障害消去と同じ原理を用いるが、
ここではtargetを非幾何的なsmooth proper dg族として扱い、
自己同値が障害を不変部分へ制限することと、intertwinerの持上げを別々に証明した。
同論文の幾何的projectorの仮定を満たさないtargetへその定理を直接適用していない。
この一般定式化の新規性は未確認である。
\end{remark}

\subsection{滑らかなcubicへの適用}
\begin{lemma}[低次Hochschild群と幾何的接写像]\label{er:lem:cubic}
$X=V(F)\subset\mathbb P^{N-1}$を滑らかなcubic、$N=n+2\ge5$とし、
$J(F)=\CC[z_1,\ldots,z_N]/(\partial_1F,\ldots,\partial_NF)$とする。
\[
 \mathrm{HH}^{<0}(\Ku(X))=0,\quad
 \mathrm{HH}^0(\Ku(X))=\CC,\quad \mathrm{HH}^1(\Ku(X))=0,
\]
\[
 \mathrm{HH}^2(\Ku(X))=
 \begin{cases}J(F)_3&n\ne4,\\J(F)_3\oplus\CC^2&n=4,\end{cases}
 \qquad
 H^1(X,T_X)\xrightarrow{\sim}\mathrm{HH}^2(\Ku(X))^{\rho_X}=J(F)_3.
 \tag{H2}\label{er:eq:HH}
\]
$\dim J(F)_3=\binom N3$である。
\end{lemma}
\begin{proof}
BFK\cite[Theorem 5.39]{BFKER}のequivariant matrix factorizationの公式を用いる。
$G=\mathbb G_m$、各座標の重みは1、potentialのcharacterは3である。
恒等sectorのKoszul cohomologyは、偏微分が正則列であるためJacobian商だけで、
cohomological degree $2r$に$J(F)_{3r}$を与える。
非恒等$g\in\mu_3$の固定空間は原点、conormalの次元は$N$、determinantの重みは$-N$である。
次数$q$の重み零条件は$(3q-N)/2=0$であり、従って$3\mid N$かつ$q=N/3$のときだけ
各sectorから一次元が現れる。これで負次数、次数0、1、2の群が得られる。

幾何的Kodaira--Spencer写像の同定は、単なる次元一致ではない。
$n\ge4$では$3\le(n+2)/2$なので、Huybrechts--Rennemo
\cite[Lemma 2.19, \S2.5, \S3.3]{HRER}のAtiyah--Kodaira--Spencerと
三次周期性の比較を適用でき、$F+\epsilon G$の像は恒等sectorの$[G]\in J(F)_3$である。
$n=3$ではJacovskis--Lin--Liu--Zhang\cite[Theorem 4.4(2)]{JLLZER}の
幾何的写像の同型を用いる。normal sequenceとEuler sequenceから$H^1(T_X)=J(F)_3$である。
幾何的族にはrotationがあるのでこの像は不変である。
$n=4$で余分な二方向が不変でないことは、Huybrechts\cite[Proposition 3.12]{HuyER}による。
rotationはsymplecticで、Mukai格子の$A_2$上で三つの根を巡回する。
不変なCalabi--Yau類による$\mathrm{HH}^2\simeq\mathrm{HH}_0$の同定の下で、
その二方向の固有値は$\zeta,\zeta^2$である。よって不変部分は正確に$J(F)_3$となる。
最後に偏微分が二次正則列であることから$H_{J(F)}(t)=(1+t)^N$を得る。
\end{proof}

\begin{theorem}[cubicのrotation付きétale局所実現]\label{er:thm:ER}
$S$は有限型複素スキーム、$\mathcal K_S$は実在するsmooth properな代数的dg族で、
$\rho_S$は可逆perfect双加群により表される自己同値とする。
$s\in S(\CC)$で、$n\ge3$の滑らかなcubic $X$に対する指定同値
$(\mathcal K_s,\rho_s)\simeq(\Ku(X),\rho_X)$があるなら、
pointed étale近傍$V\to(S,s)$、滑らかなcubic族$Y/V$、および指定同値を延長する同値
\[
 (\mathcal K_V,\rho_V)\simeq(\Ku(Y/V),\rho_{Y/V})
 \tag{ER}\label{er:eq:ER}
\]
が存在する。
\end{theorem}
\begin{proof}
滑らかな三次形式の射影パラメータ空間の$[F]$を含むアフィン開を$H$とする。
普遍cubic $\mathcal X/H$のrelative perfect圏はsmooth properで、
$\OO,\ldots,\OO(n-2)$のrelative exceptional列の射影により
smooth properな相対Kuznetsov成分を得る。
射影とrotationは有限個の評価三角で構成され、底変更に可換である。
普遍cubic上の有限個のtwistから生成子を選べるので、局所的なdg代数模型も存在する。
$T_{[F]}H\to H^1(T_X)=J(F)_3$は全射である。
補題\ref{er:lem:cubic}と定理\ref{er:thm:criterion}を適用する。
\end{proof}

\begin{corollary}[rotationを最初に指定しない存在]\label{er:cor:unpolarized}
$n\ne4$なら、前定理で$\rho_S$を入力として指定せずとも、
$\mathcal K_s\simeq\Ku(X)$からétale局所的な可換cubic模型が得られる。
これはrotationなしの大域的な一意性の主張ではない。
\end{corollary}
\begin{proof}
補題\ref{er:lem:cubic}では幾何的接写像が$\mathrm{HH}^2$全体に全射である。
定理\ref{er:thm:criterion}の証明から不変性とintertwinerの段階を省けばよい。
\end{proof}

\section{Hochschildの奇数次障害とrotation不変部分}\label{ho:sec:parity}
可換復元における対角smoothnessを、形式変形の障害と混同しないために、
滑らかな三次超曲面のKuznetsov成分の低次Hochschild群を全次元で整理する。
以下は既知のequivariant matrix factorizationの計算からの帰結であり、
残余dg圏の相対smoothnessに関する新たな一般定理ではない。

\begin{proposition}[奇数次Hochschild群とrotation固有値]\label{ho:prop:sectors}
$X=V(F)\subset\mathbb P^{N-1}_\CC$を滑らかな三次超曲面、$N=n+2\ge5$とし、
$J(F)=\CC[z_1,\dots,z_N]/(\partial F)$とする。
$\operatorname{Ku}(X)$の標準rotationを$\rho$とする。
$\zeta$を原始三乗根と書けば、次の複素ベクトル空間の分解が成り立つ：
\[
 \operatorname{HH}^{q}(\operatorname{Ku}(X))\cong
 \begin{cases}
 J(F)_{3q/2},&q\text{ 偶数},\\ 0,&q\text{ 奇数}
 \end{cases}
 \quad\oplus\quad
 \begin{cases}
 \CC_{\zeta}\oplus\CC_{\zeta^2},&3\mid N,\ q=N/3,\\
 0,&\text{それ以外}.
 \end{cases}
 \tag{HS}\label{ho:eq:sectors}
\]
ここで$J(F)_d=0$（$d<0$）とし、右側の二つの一次元空間は
$\rho$の共役作用の非自明な固有空間である。
$J(F)$から来る項では$\rho$は恒等的に作用する。
特に
\[
 \operatorname{HH}^3(\operatorname{Ku}(X))
 \cong\begin{cases}\CC^2&n=7,\\0&n\ne7,\end{cases}
 \qquad
 \boxed{\operatorname{HH}^3(\operatorname{Ku}(X))^\rho=0\quad(n\ge3).}
 \tag{V3}\label{ho:eq:HH3}
\]
さらに$\operatorname{HH}^2(\operatorname{Ku}(X))^\rho\cong J(F)_3$で、
$\dim_\CC J(F)_3=\binom N3$である。
\end{proposition}
\begin{proof}
Orlovのgraded matrix factorization表示と
Ballard--Favero--Katzarkov\cite[Theorem 5.39, pp.~82--85]{BFKER}の
固定sector公式を用いる。
$G=\mathbb G_m$は全生成変数の重みが$1$となる作用で、
potentialのcharacterは三次である。$K_\chi=\mu_3$とする。
$g=1$において$V_g=\mathbb A^N$、$W_g=0$である。
偏微分$(\partial F)$は原点にのみ共通零点を持つ長さ$N$の正則列であり、
Koszul cohomologyは零次の$J(F)$に集中する。
重み零の条件とcohomological gradingから、恒等sectorは偶数次数$q=2r$に
$J(F)_{3r}$を与え、奇数次数には寄与しない。

$g\in\mu_3\setminus\{1\}$なら固定空間は原点で、$\dim W_g=N$、
$J(F_g)=\CC$である。
Theorem 5.39のtwistとdeterminant characterを代入すると、
$u=(q-N)/2$に対する重み零条件は$3u+N=0$となる。
従って$q=N/3$でのみ各非自明sectorが一次元を与える。
各sectorは標準的な文字twist $(1)$ による共役作用で
$g$の文字（規約によって逆文字）を受ける。
これは同論文のkernelの固定graph分解と
\cite[Theorem 6.5]{BFKER}のcharacter作用の計算と同じものである。
Orlov同値の下で次数shift $(1)$ は$\rho$に対応するため、
二つのsectorの固有値は順序を除いて$\zeta,\zeta^2$となる。
$J(F)$から来る恒等sectorではこの作用は自明である。
これで\eqref{ho:eq:sectors}を得る。

$q=3$は奇数なので、非零になり得るのは$N/3=3$、すなわち$N=9$
（$n=7$）だけである。しかもその二方向には$\rho$-不変ベクトルがない。
$q=2$の不変部分は常に$J(F)_3$である。
滑らかさから偏微分は二次正則列なので、
$H_{J(F)}(t)=(1+t)^N$となり$\dim J(F)_3=\binom N3$を得る。
\end{proof}

\begin{corollary}[非楕円cubic fivefoldの変形障害の所在]\label{ho:cor:fivefold}
本研究の節点型・尖点型の二つの滑らかな可換特殊ファイバー
$X^5$について、共に
\[
 (\dim\operatorname{HH}^0,\dim\operatorname{HH}^1,
 \dim\operatorname{HH}^2,\dim\operatorname{HH}^3,
 \dim\operatorname{HH}^4)=(1,0,35,0,7).
 \tag{F5}\label{ho:eq:fivefold}
\]
従って滑らかな可換特殊ファイバーのdg圏自身の通常のMorita変形に対する
\emph{一次の}Kuranishi障害空間$\operatorname{HH}^3$は零である。
ただし、量子化から構成した$\Gamma_R$の
$\Gamma_R^e$上の対角perfect性を、これだけからは結論できない。
\end{corollary}
\begin{proof}
$N=7$を命題\ref{ho:prop:sectors}に代入する。
$J(F)_0,J(F)_3,J(F)_6$の次元はそれぞれ$1,35,7$で、
それ以外の表示次数$q=1,3$は零となる。
Morita変形の一次障害が$\operatorname{HH}^3$に入ることは
標準的なHochschild変形理論の適用である。
\end{proof}

\begin{remark}[残余対角との論理的境界]\label{ho:rem:barrier}
$\operatorname{HH}^3(\Ku(X))^\rho=0$は重要な形式的情報であるが、
この不変部分の消滅だけで任意のrotation付きdg圏変形の厳密な
$\mathbb Z/3$-equivariant obstruction theoryを構成したとはしない。
また、Artin切断上の変形が可換的に説明できたとしても、
有限型係数環上の\emph{実在する}残余対角のperfect性、
すなわち$\Gamma_R\in\Perf(\Gamma_R^{\mathrm{op}}\otL_R\Gamma_R)$は別の問題である。
所定の量子化のgeneric復元の最優先課題は依然としてこの残余対角のperfect性である。
本節はBFKの既知の公式の具体的な帰結と、その本研究における障害の切分けを与える。
\end{remark}

\section{一般次元・各Poisson成分の支配的chartと可換復元}
\label{sec:all-poisson-components}

本節の目標は、任意の$n\ge3$、滑らかな複素cubic $X^n$と、
アフィンPoissonスキームの各被約既約成分$Z$について、
その幾何学的一般点でrotation付き可換模型を構成することである。
特定の三様体の17次元族や、一つの階数の層に限定しない。
以下では、証明する代数幾何的な還元と、まだ仮定である圏論的な実現原理を区別する。
新規性の包括的な確認は行っておらず、線形代数・錐のchart・完備化の忠実平坦性自体は標準的である。
本節の寄与は、これらを全Poisson成分の問題へ組み合わせ、
\emph{特異な錐の頂点上への降下を、genericな存在定理の前提から外す}ことである。

\subsection{各成分の零点と滑らかな支配的chart}

\begin{proposition}[Poisson錐の成分と零点]
\label{prop:all-components-cone}
$X$を滑らかな複素射影多様体とし、
\[
 E_X=H^0(X,\textstyle\bigwedge^2T_X),\qquad
 \operatorname{Pois}(X)=\{\pi\in E_X:[\pi,\pi]=0\}
\]
とおく。その被約化の各既約成分はスカラー作用で不変で、零を含む。
滑らかなcubicを動かす相対Poissonスキームでも、各成分は
その成分にある$(X,\pi)$に対応する$(X,0)$を含む。
\end{proposition}
\begin{proof}
Jacobi方程式は$\pi$について二次斉次である。
既約成分$Z$に対し、$\mathbb G_m\times Z$の作用像の閉包は既約で$Z$を含む。
$Z$の極大性からその閉包は$Z$であり、スカラー作用は$Z$を保つ。
$t\mapsto t\pi$を$\mathbb A^1$まで延長すると、閉性から$0\in Z$となる。
相対の場合は$t\mapsto(X,t\pi)$とすればよい。
\end{proof}

\begin{theorem}[滑らかなスカラーchart]
\label{cw:thm:chart}
$Z\ne\{0\}$を有限次元複素ベクトル空間$E$内の被約既約アフィン錐とする。
$\mathbb PZ$の滑らかな稠密アフィン開部分$T$と正則写像
$p:T\to Z\setminus\{0\}$を選べて、
\[
 S=T\times\mathbb A^1_h,\qquad
 c:S\longrightarrow Z,\quad (t,h)\longmapsto h p(t)
 \tag{C}\label{cw:eq:contraction}
\]
は次を満たす。
$S$は滑らかな既約スキーム、$c$は支配的双有理、
$c(t,0)=0$であり、$c|_{T\times\mathbb G_m}$は$Z$の稠密開部分への同型である。
とくに$\mathbb C(S)=\mathbb C(Z)$である。
\end{theorem}
\begin{proof}
$\mathbb PZ$は被約既約で、標数零なので滑らかな稠密開部分を持つ。
その非空開部分を、ある線形座標$z_0\ne0$の射影chartの中に取り、さらにアフィンに縮める。
代表を$z_0(p(t))=1$と規格化すれば、$p$は正則である。
$h\ne0$では逆写像は
\[
 z\longmapsto\bigl([z],z_0(z)\bigr)
\]
であり、像は$z_0\ne0$かつ$[z]\in T$という開集合である。
これが双有理性を含む全主張を証明する。
\end{proof}

\begin{remark}[頂点での降下はgenericな目標には不要]
\label{cw:rem:no-cone-descent}
$Z$が頂点で特異でも$S$は滑らかである。
$S$上の量子化族が$Z$へ降下するとは主張しない。
しかし一般点の模型を求めるだけなら、同じ関数体を持つ$S$上での構成で足りる。
$t$を一つ固定した曲線ではなく、\emph{$T$の全方向を含む族}を保つことが重要である。
射影Poisson空間そのものに零点はなく、式\eqref{cw:eq:contraction}の$h=0$が可換ファイバーを与える。
相対Poissonスキームについても、その射影化内の滑らかな稠密chartで
相対tautological lineを自明化すれば同じ構成ができる。
零切断内にのみ含まれる成分は、もともと可換な族として別に扱う。
\end{remark}

\begin{proposition}[ブロック正規形の有限被覆上の固定]
\label{cw:prop:normal-form-cover}
$K=\mathbb C(T)$とし、genericなPoisson構造の三重ベクトル$P$と三次式$F$が
$\overline K$上で、ある$N=3k+m$について
\[
 P=\sum_{s=1}^k\partial_{x_{s1}}\wedge\partial_{x_{s2}}\wedge\partial_{x_{s3}},
 \qquad F=\sum_{s=1}^k f_s(x_s,u)
 \tag{B}\label{cw:eq:block-form}
\]
という正規形を持つと仮定する。
$T$を縮めると、整な有限étale全射$T'\to T$が存在し、
この順序付きブロック、中心変数、座標変換を全て$T'$上で正則に選べる。
$S'=T'\times\mathbb A^1_h\to Z$は支配的でgenerically finiteであり、
$h=0$は可換なcubicの族を与える。
\end{proposition}
\begin{proof}
座標変換、各ブロックの基底、三次式の係数を未知数とすると、
\eqref{cw:eq:block-form}と座標変換の可逆性は有限個の多項式方程式・不等式である。
$\overline K$上の解の有限個の座標は$K$の有限拡大$L$に属する。
標数零なので$L/K$は分離的である。
原始元の最小多項式を取り、係数の分母と判別式を反転すると、
$T$の開部分上に関数体$L$の有限étale被覆$T'$を得る。
解の座標に現れる有限個の分母は、そのノルムをさらに反転して全て可逆にできる。
これで座標変換と式が$T'$全体へ広がる。
有限étale被覆と式\eqref{cw:eq:contraction}の合成は、$h\ne0$で有限étaleである。
\end{proof}
\begin{remark}
この命題の正規形の存在は論文AのPoisson正規形を入力とする。
その証明を本節で再検証したとは扱わない。
$k$は選んだ成分のgenericなブロック数であり、$k$が同じ成分を同一視しない。
また、各$f_s(x_s,0)$が滑らかな平面三次曲線を定めるという仮定は置いていない。
\end{remark}

\subsection{有限行列による中心三次元の代数化}

\begin{lemma}[genericな形式解の有限jetを保つ代数化]
\label{cw:lem:linear-algebraization}
$R$を有限生成整域、$K=\operatorname{Frac}(R)$、
$M(h)\in\operatorname{Mat}_{a\times b}(R[h])$とする。
$\widehat z\in K[[h]]^b$が$M\widehat z=0$を満たすなら、
任意の整数$r\ge0$について、$R$の非零元$s$と
$q(h)\in R_s[h]$、$q(0)\in R_s^\times$、および
\[
 z\in R_s[h,q^{-1}]^b,
 \qquad Mz=0,\qquad z\equiv\widehat z\pmod {h^{r+1}}
 \tag{L}\label{cw:eq:linear-approx}
\]
が存在する。合同は$K[[h]]^b$内で解釈する。
行列$M$のrankが$h=0$で一定であることは仮定しない。
\end{lemma}
\begin{proof}
$\mathcal O=K[h]_{(h)}$とし、$z_{\le r}\in K[h]^b$を$\widehat z$の$r$次切断とする。
$Mz_{\le r}$は$h^{r+1}$で割れるから
\[
 b_r=-h^{-r-1}Mz_{\le r}\in K[h]^a
\]
と置ける。形式解$w=(\widehat z-z_{\le r})/h^{r+1}$は$Mw=b_r$を満たす。
$\mathcal O\to K[[h]]$は忠実平坦である\cite[Tag 00MC]{Stacks}。
従って$\operatorname{coker}(M:\mathcal O^b\to\mathcal O^a)$における$b_r$の像は零で、
$Mv=b_r$を満たす$v\in\mathcal O^b$が存在する。
$z=z_{\le r}+h^{r+1}v$とすれば所要の合同と恒等式を得る。
$v$の共通分母を$q\in K[h]$、$q(0)\ne0$と取る。
$K$の係数に現れる有限個の分母と$q(0)$を$R$で反転すると、
式\eqref{cw:eq:linear-approx}を得る。
\end{proof}

\begin{corollary}[次数三の中心持上げの一様化]
\label{cw:cor:center-algebraization}
$\mathcal A$を$R[h]$上の次数一生成代数とし、
$\mathcal A_3,\mathcal A_4$が有限自由、$\mathcal A/(h)$が可換とする。
$\mathcal A_1$の生成元$x_1,\ldots,x_N$に対し
\[
 M_h:\mathcal A_3\longrightarrow\mathcal A_4^{\oplus N},
 \qquad z\longmapsto h^{-1}([z,x_1],\ldots,[z,x_N])
 \tag{Z}\label{cw:eq:center-matrix}
\]
は有限行列で表される。
$\mathcal A_K$の次数別$h$-進完備化に中心三次元$\widehat\Omega$があり、
$\widehat\Omega\bmod h=F$なら、前補題の局所化後に
$\Omega\in Z(\mathcal A_{R_s[h,q^{-1}]})_3$、$\Omega\bmod h=F$を取れる。
指定した任意の有限jetも保持できる。
Hilbert級数が$(1-t)^{-N}$のPBW族なら、未知数の数と方程式の数はそれぞれ
\[
 b=\binom{N+2}{3},\qquad a=N\binom{N+3}{4}
\]
であり、$h$の打切り次数によって無限に増やす必要はない。
\end{corollary}
\begin{proof}
$\mathcal A/(h)$が可換で$\mathcal A_4$が自由なので、各交換子を$h$で一意に割れる。
中心性は生成元との交換子が全て零であることと同値である。
前補題をこの行列に適用する。
\end{proof}
\begin{remark}[この代数化が証明しないこと]
\label{cw:rem:center-limit}
形式中心元の存在と、適用する周囲の代数$\mathcal A$の構成はこの系の入力である。
係数方向を縮めることは許すが、$q(0)$は可逆なので縮めた基底の零切断は全て残る。
この議論は中心元を\emph{代数的な有理式}にするものであり、
全ての形式係数を一致させること、全Poisson錐上の普遍中心元、
全次数PBW性、AS-regular性、相対saturationを証明するものではない。
\end{remark}

\begin{lemma}[正則中心商の次数別flatness]
\label{cw:lem:flat-quotient}
$D$を有限型基底、$\mathcal A$を$D$上の次数付き代数で各次数が有限局所自由とする。
全ての幾何学的ファイバーが整域で、$\Omega\in Z(\mathcal A)_3$が
各ファイバーで非零であれば、$\mathcal B=\mathcal A/(\Omega)$は各次数で有限局所自由である。
$H_{\mathcal A_s}(t)=(1-t)^{-N}$なら
$H_{\mathcal B_s}(t)=(1-t^3)/(1-t)^N$である。
\end{lemma}
\begin{proof}
各次数$d$で$\Omega:\mathcal A_{d-3}\to\mathcal A_d$は
全ての幾何学的ファイバーで単射である。
有限局所自由加群の行列の極大小行列式を用いると、この写像は局所的に分裂した単射で、
余核は有限局所自由である。従って主張とHilbert級数の式を得る。
\end{proof}

\subsection{一つの可換ファイバーから成分の一般点へ}

\begin{definition}[条件 (ER)]
\label{cw:def:ER}
本節で条件 (ER) とは次の局所実現原理を指す。
有限型複素スキーム$D$上のsmooth properな代数的dg圏と
可逆perfect bimoduleで表されるrotationの族$(\mathcal K_D,\rho_D)$が、
ある複素点$d$で滑らかなcubicの標準pairとenhanced同値なら、
$d$のpointed étale近傍で相対cubicとrotationを保つ同値が存在する。
これは現行論文Cの\texttt{rotation\_models.tex}にある
\texttt{thm:etale}を\emph{入力}とする条件である。
本節ではその形式変形論・実現スタックの証明全体を再証明しない。
Toën--Vaquié \cite[Theorem 3.6]{TV} はperfect対象のmoduliの代数性を与える入力だが、
それ単独が条件 (ER) の証明ではない。
\end{definition}

\begin{theorem}[支配的chart上の一可換点によるgeneric復元]
\label{thm:all-components-conditional}\label{cw:thm:generic}
$X^n$、$n\ge3$を滑らかなcubic、$Z\ne\{0\}$を
$\operatorname{Pois}(X)_{\mathrm{red}}$の任意の既約成分とする。
定理\ref{cw:thm:chart}と命題\ref{cw:prop:normal-form-cover}の
$S'=T'\times\mathbb A^1_h\to Z$を使う。
ある点$s_0=(t_0,0)$を含む開部分$V\subset S'$上に、
実際の次数別flatな量子化中心三次商と、これに対応する
smooth properな代数的残余pair族$(\mathcal K_V,\rho_V)$があり、
$s_0$で$(\mathcal K_{s_0},\rho_{s_0})\simeq(\Ku(X),\rho_X)$とする。
条件 (ER) が成立すれば、$V$の稠密開部分の各幾何学的ファイバーは
滑らかな可換cubicのrotation付きpairを持つ。
特に$\mathbb C(Z)$のある有限分離拡大上で、その構成したgenericな量子化の可換模型が存在する。
\end{theorem}
\begin{proof}
条件 (ER) のétale近傍の像$W\subset V$は$s_0$を含む開集合である。
$S'$は既約だから$W$は稠密で、$W\cap\{h\ne0\}$は非空である。
$S'\to Z$は$h\ne0$で$Z$の稠密開部分への有限étale射なので、その像も開で稠密である。
étale近傍の支配的な成分のgeneric residue fieldは
$\mathbb C(S')$の有限分離拡大であり、
$\mathbb C(S')/\mathbb C(Z)$も有限分離なので結論を得る。
\end{proof}
\begin{remark}[量化と降下]
\label{rem:all-components-gap}
定理は\emph{構成した}量子化族についての条件付き存在であり、
同じPoisson初項を持つ任意の非同値な量子化を同一視しない。
$Z$上に量子化そのものが既に降下している場合には、
$Z$の稠密開の各幾何学的ファイバーについて結論を読むこともできる。
そうでなければ、有限支配被覆上のgeneric模型の存在として読む。
また、多ブロック生成性はこの方法で独立の前提にする必要はないが、
その代わり全体の残余圏についてsmooth proper性を示す必要がある。
\end{remark}

\begin{proposition}[有限対角証明書によるsaturationの開近傍]
\label{cw:prop:diagonal-certificate}
$R$をNoether環、$\Lambda$を$R$上のdg代数で$\Lambda\in\Perf(R)$とする。
$\Lambda^e=\Lambda^{\mathrm{op}}\otimes_R^{\mathbf L}\Lambda$と置く。
実際に
\[
 P\in\Perf(\Lambda^e),\qquad \alpha:P\to\Lambda
 \tag{D}\label{cw:eq:diagonal-certificate}
\]
が構成され、ある点$s$で$\alpha_s$が擬同型なら、$s$のZariski開近傍で
$\Lambda$はsmooth properである。この結論は導来底変更と両立する。
\end{proposition}
\begin{proof}
$\Lambda^e$は$R$-perfectなので、$P$も$R$-perfectである。
従って$C=\operatorname{Cone}(\alpha)$は$R$-perfectで、$C\otimes_R^{\mathbf L}\kappa(s)=0$。
有限個の項からなる有限局所自由複体で表すと、ファイバーの非零コホモロジーの存在座は閉である。
よって$s$の開近傍で$C=0$であり、$\Lambda\simeq P\in\Perf(\Lambda^e)$を得る。
perfect性と擬同型の導来引戻しの両立性から最後の主張も従う
\cite[Tag 0BDK]{Stacks}。
\end{proof}
\begin{remark}
ここで必要なのは\emph{既に持ち上がった}式\eqref{cw:eq:diagonal-certificate}である。
特殊ファイバーの対角がperfectであるだけから、その持上げを結論してはいない。
また$R$-perfectとは底上の複体の性質であって、
有限自由$R$-加群だけからなる\emph{厳密なdg代数模型}の存在を仮定していない。
\end{remark}

\subsection{次元合同類と残る未解決条件}

\begin{proposition}[Serre関手からのrotationの復元]
\label{prop:all-components-rotation-serre}
smooth properなdg圏と自己同値$\rho$が
\begin{equation}\label{eq:all-components-serre}
 \rho^3\simeq[2],\qquad S\simeq\rho^{1-n}[n]
\end{equation}
を満たすとする。
$n\not\equiv1\pmod3$なら、$3a+(1-n)b=1$を満たす整数$a,b$に対して
\begin{equation}\label{eq:all-components-rotation}
 \rho\simeq S^b[2a-nb].
\end{equation}
特に$n\equiv0\pmod3$なら$\rho\simeq S[-n/3]$、
$n\equiv2\pmod3$なら$\rho\simeq S^{-1}[(n+4)/3]$である。
$n\equiv1\pmod3$ならこれらの式から$S\simeq[(n+2)/3]$を得るが、
rotationの同型類は一般には決定されない。
\end{proposition}
\begin{proof}
$\gcd(3,n-1)=1$のときBézout等式を用いて
$\rho=(\rho^3)^a(\rho^{1-n})^b$とする。
$n-1=3m$の場合は$S\simeq[n-2m]$である。
この後者の場合の追加の曖昧さは、論文Cの四次元Zhang twist反例で実際に現れる。
\end{proof}

\begin{example}[全体が滑らかでも純ブロックは特異になり得る]
\label{cw:ex:singular-block}
$\mathbb P^6$内のcubic fivefold
\[
 F=x^3+y^3+u z^2+a^3+b^3+c^3+u^3=0
\]
と、三重ベクトル
$P=\partial_x\wedge\partial_y\wedge\partial_z+
\partial_a\wedge\partial_b\wedge\partial_c$を考える。
$F$は滑らかで、$\iota_{dF}P$は一般階数4のPoisson構造である。
しかし第一ブロックの$u=0$での純三次式$x^3+y^3$は特異である。
\end{example}
\begin{proof}
勾配の零点では$x=y=a=b=c=0$、$uz=0$、$z^2+3u^2=0$なので、$z=u=0$である。
よってアフィン錐の特異点は原点のみである。
二ブロックは互いに可換で$u$はCasimir、各ブロックはJacobian括弧なのでJacobi恒等式を満たす。
$u\ne0$のchartで両ブロックが非零となる点は存在し、そのとき階数は4である。
$x^3+y^3=0$は$[0:0:1]$で特異である。
\end{proof}
\begin{remark}
この例が、そのPoisson既約成分の\emph{一般点}でも特異ブロックを強制するとは主張しない。
全点に対する議論で、全体の滑らかさから各純ブロックの楕円性を仮定してはならないという注意である。
\end{remark}

\begin{remark}[全非可換成分との区別]\label{rem:all-components-nc-scope}
任意のAS-regular中心三次商の全点での復元は既存のZhang twist反例により偽である。
個別の反例点だけではPoisson既約成分の一般点での復元は否定されない。
本節は$n\ge3$を対象とし、曲線・cubic surfaceへ圏論的入力を無条件に拡張しない。
\end{remark}

本節により、generic復元では、Poisson錐の頂点上への量子化の降下、
全成分の分類、全パラメータでの模型の存在を先に解く必要はなくなった。
残る本質的な仕事は、各成分を支配する上のchartで
\emph{実際の}相対PBW代数と中心正則商を構成し、
残余圏のproper模型、有限対角証明書、rotationと相対Orlov比較を得ることである。
形式中心の存在が分かれば中心元だけの代数化は補題\ref{cw:lem:linear-algebraization}で処理できるが、
これを残余圏のsaturationの証明と混同しない。

\subsection{有限計算と証明書の範囲}

添付コードは、$\operatorname{Gr}(2,5)$の標準chartと錐の縮約、
上の特異ブロック例の勾配・Jacobi恒等式、中心行列の線形代数化の例を有理数係数で確認する。
さらに、三変数Fermat量子化の中心元については、
有限特殊化ではなく\emph{$\mathbb Q[h]$上の自由代数恒等式}を証明書として収録した。
$r_1=xy-yx-hz^2$、$r_2=yz-zy-hx^2$、$r_3=zx-xz-hy^2$、
$D=x^3+y^3+z^3$、$T=xyz+xzy+yxz+yzx+zxy+zyx$とすると、
\begin{align*}
 [3D-h^2T,x]={}&-3r_1y^2+h^2r_1zx-h^2xr_1z-3yr_1y-h^2zr_1x
 +h^2xzr_1-3y^2r_1\\
 &+2h^2r_2x^2-3hr_2yz-h^2xr_2x-h^2x^2r_2+3hyzr_2\\
 &+3r_3z^2+3zr_3z+3z^2r_3.
\end{align*}
語ごとの係数比較でこの式を確認でき、巡回置換で$y,z$との交換子も得る。
従って$(D/3-h^2T/9)$はこの三変数商で中心的である。
この証明書は全混合項の中心公式、PBW性、圏同値の検証ではなく、
標準的なFermat入力について特殊値$h=0$を除外しない普遍的検証である。

\section{一般次元の局所代数的模型：仮定と主定理}
\label{lp:sec:main}
本節では、既に構成された次数別flatな量子化から、射影化のproperなdg模型を作る。
Poisson成分のブロック数、純ブロックの楕円性、有限中心性は仮定しない。
ただし、量子化代数そのものの全次数flatnessは入力である。
基底は完備化やArtin切断ではなく、可換点における実在の正則局所環とする。

\begin{definition}[局所入力]\label{lp:def:input}
$(R,\mathfrak m,k)$を有限生成複素代数の局所化である正則局所環とし、
$d=\dim R$、$N\ge5$とする。
$A=\bigoplus_{j\ge0}A_j$は中央$R$-代数、$A_0=R$で、各$A_j$は有限自由$R$-加群、
$A\otimes_R k\simeq k[z_1,\ldots,z_N]$を次数付き代数として満たすとする。
$\Omega\in Z(A)_3$が与えられ、その剰余$F$は滑らかな
$cubic\ X=V(F)\subset\mathbb P^{N-1}_k$を定めるとする。
$B=A/(\Omega)$、$n=N-2$、$a=N-3$と置く。
内部shiftは$M(j)_q=M_{q+j}$、複体はcohomological conventionを用いる。
\end{definition}

\begin{theorem}[局所proper模型と有界Orlov分解]\label{lp:thm:main}
定義\ref{lp:def:input}の下で以下が成立する。
\begin{enumerate}[label=(\arabic*)]
\item $A,B$は左右noetherianである。$A$の$R$上の対角には長さ$N$の有限自由双加群分解があり、
$\operatorname{gldim}A\le N+d$である。$\Omega$は中心正則で、各$B_j$は有限自由、
その形成は任意の通常の底変更と両立する。
\item 左右で
\[
 \operatorname{R\underline{Hom}}_A(R,A)\simeq R(N)[-N],\qquad
 \operatorname{R\underline{Hom}}_B(R,B)\simeq R(a)[-(N-1)].
 \tag{G}\label{lp:eq:G}
\]
$B$は有限左右injective dimensionを持ち、Burke--Stevensonの意味でGorensteinである。
従って通常の有界三角圏について、$D^b(\operatorname{qgr}B)$は
$D_{\rm sg}^b(\operatorname{gr}B)$と$a$個の$\operatorname{Perf}(R)$成分による
Orlov型半直交分解を持つ。特異圏では自然な三次周期性$(3)\simeq[2]$が成立する。
\item $\mathcal O(j)=\pi B(j)$とすると
\[
 \operatorname{Ext}^p_{\operatorname{qgr}B}(\mathcal O(i),\mathcal O(j))\simeq
 \begin{cases}
 B_{j-i}&p=0,\\
 B_{i-j-a}^{\vee}&p=n,\\
 0&\text{その他}
 \end{cases}\qquad (B_q=0\text{ for }q<0).
 \tag{H}\label{lp:eq:H}
\]
ここで$(-)^\vee=\operatorname{Hom}_R(-,R)$である。
この計算の自然な導来Homの底変更写像は、命題\ref{lp:prop:basechange}の通常の底変更で擬同型となる。
\item $\mathcal C_B=\operatorname{thick}_R\{\mathcal O(j):j\in\mathbb Z\}$を
$D^b(\operatorname{qgr}B)$の標準dg enhancement内の冪等完備な部分と定めると、
\[
 E=\bigoplus_{i=0}^{N-1}\mathcal O(i),\qquad
 \mathcal C_B=\operatorname{thick}_R(E)\simeq\operatorname{Perf}(\Lambda_R),
 \quad\Lambda_R=\operatorname{REnd}(E)\in\operatorname{Perf}(R).
 \tag{P}\label{lp:eq:proper-model}
\]
\item $\mathcal K_B^{\rm lin}=\mathcal C_B\cap
\langle\mathcal O,\ldots,\mathcal O(a-1)\rangle^\perp$は
有限個のmutationにより生成対象$G$を持ち、
$\mathcal K_B^{\rm lin}\simeq\operatorname{Perf}(\Gamma_R)$、
$\Gamma_R=\operatorname{REnd}(G)\in\operatorname{Perf}(R)$である。
特殊ファイバーでは$\mathcal C_B\otimes_R^{\mathbf L}k\simeq\operatorname{Perf}(X)$、
$\mathcal K_B^{\rm lin}\otimes_R^{\mathbf L}k\simeq\operatorname{Ku}(X)$を得る。
\end{enumerate}
\end{theorem}

\begin{remark}[主張しないこと]\label{lp:rem:limits}
$\Gamma_R$の底上proper性と$\Gamma_R^e$上の対角perfect性は異なる。
本定理は後者、一般の可換cubic模型の存在、有限型開集合へのdg構造のspread、
任意の導来基底上のOrlov同値を結論しない。
$\mathcal C_B=D^b(\operatorname{qgr}B)$や
$\mathcal K_B^{\rm lin}=D_{\rm sg}^b(\operatorname{gr}B)$も一般には主張しない。
上付き${\rm lin}$は、線束と有限操作から作る部分圏を全coherent残余圏と区別するための記号である。
本稿の結論は条件付きgeneric復元の入力のうち、proper模型・有界Orlov分解を具体化するものである。
\end{remark}

\subsection{先行研究と本節の寄与}
非可換射影圏のsection functorはArtin--Zhangの枠組み\cite{LP-AZ}に属する。
次数零部が体でないGorenstein環の有界Orlov分解には
Burke--Stevenson\cite[Definition 5.1, Lemmas 6.1--6.2, Theorem 6.4]{LP-BS}を用いる。
dg生成対象によるMorita表示はKeller\cite[Theorem 3.8]{LP-Keller}の一般理論である。
これらの定理自体を新しい結果とは扱わない。
以下では、与えられた量子化に対してnoetherian性・Gorenstein型・Homの底変更・有限生成を導き、
既存原稿で別途仮定していた相対proper模型の条件を正則局所基底で検証する。
この組合せの新規性は未確認であり、外部査読や形式証明は未実施である。

\section{局所基底でのnoetherian性と有限Koszul持上げ}
\label{lp:sec:lift}
\begin{lemma}[次数別有限性を用いたnoetherian性]\label{lp:lem:noetherian}
$(R,\mathfrak m)$をnoetherian局所環、$D=\bigoplus_{j\ge0}D_j$を中央$R$-代数とし、
各$D_j$は有限$R$-加群とする。
$D/\mathfrak mD$が右noetherianなら$D$は右noetherianである。左にも同様である。
\end{lemma}
\begin{proof}
$\mathfrak m$の有限生成元を選ぶと、
$(D/\mathfrak mD)[t_1,\ldots,t_e]\twoheadrightarrow\operatorname{gr}_{\mathfrak m}D$
という中央変数の多項式拡大からの全射がある。
従って右辺は右noetherianである。
斉次右イデアル$I\subset D$に誘導filtration $I\cap\mathfrak m^qD$を入れる。
その随伴次数イデアルの有限個の二重斉次生成元を、$f_1,\ldots,f_r\in I$へ持ち上げ、
$J=\sum f_iD$とする。
$x\in I_j$の最低$\mathfrak m$-次数記号を消す操作を繰り返すと、任意の$q$について
$x\in J_j+\mathfrak m^qD_j$を得る。
$D_j/J_j$は有限$R$-加群だからKrull交叉定理より$x\in J_j$である。
従って全ての斉次右イデアルは有限生成である。
任意の右イデアルには内部次数の上昇filtrationを入れ、最高次数項のイデアルの
有限生成元を持ち上げる。内部次数の有限降下により元のイデアルも有限生成である。
この証明は$R$の完備性を使わず、各内部次数における有限性を使う。
\end{proof}

\begin{lemma}[次数別分裂する自由分解の持上げ]\label{lp:lem:resolution}
$D$の各次数が有限自由$R$-加群で、$M$も各次数有限自由とする。
$\overline M$が$\overline D=D/\mathfrak mD$上、有限階数の次数付き自由加群による
有限分解$\overline P_\bullet$を持つなら、同じ項の階数とshiftを持つ
$D$-自由分解$P_\bullet\to M$を持ち上げられる。
分解は各内部次数で$R$-加群の分裂完全列であり、任意の通常の底変更後も完全である。
\end{lemma}
\begin{proof}
$\overline P_0$の生成元の像を$M$の同じ次数へ持ち上げる。
各次数の余核は有限$R$-加群で剰余が零なのでNakayamaの補題で零である。
$P_0\twoheadrightarrow M$は各次数で有限自由加群への全射だから分裂し、
核は各次数有限自由となる。剰余体への底変更後の核も正しく計算される。
核の生成元を$\overline P_1$から同様に持ち上げ、帰納する。
最後の核は各次数有限自由かつ剰余が零であるから零である。
各段階の内部次数別分裂により、底変更後にも完全性が残る。
\end{proof}

\begin{proposition}[全次数の相対Koszul性とambientの対角]\label{lp:prop:koszul}
定義\ref{lp:def:input}の$A$は次数一で生成され、二次関係で定まる。
$R$の左右$A$-加群としての分解と、$A$の$A^e=A\otimes_RA^{\rm op}$-加群としての分解は
それぞれ
\[
 P_j=A(-j)^{\binom Nj},\qquad
 Q_j=A^e(-j)^{\binom Nj}\quad(0\le j\le N)
 \tag{K}\label{lp:eq:koszul}
\]
という項を持ち、全ての通常の底変更と両立する。
特に$A$は$R$上smoothで、$\operatorname{gldim}A\le N+d$である。
\end{proposition}
\begin{proof}
$A_1$の剰余は$N$変数の空間であり、そのliftによる自由代数からの写像は
各次数のNakayamaにより全射となる。
多項式環の$k$のKoszul分解に補題\ref{lp:lem:resolution}を適用すると$P_\bullet$を得る。
第一次関係が次数二で生成されることもこの構成から従う。
同様に$A$を$A^e$-加群として見て、多項式環の対角Koszul分解を持ち上げると$Q_\bullet$を得る。
$A$は$R$上の自由加群（次数別の直和）なので$Q_\bullet$は右$A$-加群としても分裂完全である。
任意の左$A$-加群$M$を右からテンソルすると、項$A\otimes_R M$の有限直和からなる長さ$N$の分解になる。
$\operatorname{gldim}R=d$より各項の$A$-射影次元は高々$d$。
従って$\operatorname{pd}_A M\le N+d$で、右側も同様である。
\end{proof}

\begin{proposition}[相対Gorenstein型と中心商]\label{lp:prop:gorenstein}
式\eqref{lp:eq:G}が成立する。$B$は左右noetherian、各次数有限自由で、
左右injective dimensionは高々$N+d-1$である。
\end{proposition}
\begin{proof}
補題\ref{lp:lem:noetherian}を$A$に適用する。
$P_\bullet$を$A$で双対化した有界複体は、各内部次数で有限自由$R$-複体である。
剰余上ではKoszul双対性により、内部次数$-N$・cohomological degree $N$以外で完全。
有限自由有界複体の剰余が完全なら、単元行列成分の消去とNakayamaにより複体自身も分裂完全である。
内部次数$-N$では最後の項$R$のみ残る。従って$\operatorname{R\underline{Hom}}_A(R,A)=R(N)[-N]$を得る。
左側も同じである。

各次数での$\Omega:A_{j-3}\to A_j$の剰余は非零多項式$F$の乗法で単射。
行列の列数と同じ大きさの小行列式の一つが単元なので、この写像は局所基底変換で分裂単射となる。
よって$\Omega$は正則、$B_j$は自由で、全底変更と両立する。
中央正則元に対するchange-of-ringsは
\[
 \operatorname{R\underline{Hom}}_A(M,A)
 \simeq\operatorname{R\underline{Hom}}_B(M,B)(3)[-1]
 \quad(M\in D(\operatorname{Gr}B))
 \tag{CR}\label{lp:eq:change-rings}
\]
を与える。これは$0\to A(-3)\xrightarrow{\Omega}A\to B\to0$から
$\operatorname{R\underline{Hom}}_A(B,A)=B(3)[-1]$を計算し、導来随伴を使った式である。
$M=R$を代入して$B$の式\eqref{lp:eq:G}を得る。
一般の$M$に適用すると、$A$の有限大域次元から$B$の有限injective dimensionが従う。
商なのでnoetherian性も成り立つ。
\end{proof}

\section{射影化の局所コホモロジーと底変更}
\label{lp:sec:cohomology}
$\tau$を$B_{\ge1}$-torsion functor、$\pi:\operatorname{Gr}B\to\operatorname{QGr}B$を商関手とする。
Burke--Stevenson\cite[Section 4, Lemmas 6.1--6.2]{LP-BS}により、
上記のGorenstein型と$R$-flatnessは$\chi$と必要な有限cohomological dimensionを与える。
従って$D^b(\operatorname{qgr}B)$を$D(\operatorname{QGr}B)$で計算でき、導来sectionの三角
\[
 R\tau(B)\longrightarrow B\longrightarrow R\omega\pi(B)\longrightarrow
 \tag{LC}\label{lp:eq:torsion}
\]
を使用できる。

\begin{lemma}[局所コホモロジーの各次数]\label{lp:lem:local-cohomology}
$d_B=N-1$とすると
\[
 H^q(R\tau B)_t=
 \begin{cases}
  B_{-t-a}^{\vee}&q=d_B,\\
  0&q\ne d_B
 \end{cases}
 \tag{LC'}\label{lp:eq:localcohom}
\]
である。負添字は零とする。
\end{lemma}
\begin{proof}
$M_\ell=B/B_{\ge\ell}$の次数filtrationの商は
$B_j\otimes_R R(-j)$（$0\le j<\ell$）である。
式\eqref{lp:eq:G}と$B_j$の自由性により、各商の$B$への導来Homは次数$d_B$に集中し、
その内部次数は$-a-j$のみである。
長完全列によって$\operatorname{Ext}_B^q(M_\ell,B)=0$（$q\ne d_B$）。
内部次数$t$では$j=-t-a$のただ一層だけが寄与するので、その層が含まれれば群は$B_j^\vee$となる。
$\ell>j$で推移写像は同型である。
$R\tau B=\varinjlim_\ell\operatorname{R\underline{Hom}}_B(M_\ell,B)$のcohomologyを取ればよい。
\end{proof}

\begin{proof}[式\eqref{lp:eq:H}の証明]
$\operatorname{RHom}_{\operatorname{QGr}B}(\pi B,\pi B(t))$
は導来sectionの内部次数$t$に一致する。
式\eqref{lp:eq:torsion}と補題\ref{lp:lem:local-cohomology}を用いると
$H^0=B_t$、$H^{N-2}=B_{-t-a}^\vee$で、それ以外は零となる。
$a>0$なので同じ$t$に対して二群が同時に非零となることはない。
従ってこれはcohomology次元だけでなく、各Hom複体が有限自由加群の一つのshiftであることを示す。
$i$のtwistを外せば式\eqref{lp:eq:H}である。
\end{proof}

\begin{proposition}[線束Homの底変更]\label{lp:prop:basechange}
$R\to R'$を有限大域次元のnoetherian正則環への通常の準同型で、$A'=A\otimes_RR'$が左右noetherianであるものとする。
$B'=B\otimes_RR'$と置けば自然な写像
\[
 \operatorname{RHom}_{\operatorname{qgr}B}(\mathcal O(i),\mathcal O(j))
 \otimes_R^{\mathbf L}R'
 \longrightarrow
 \operatorname{RHom}_{\operatorname{qgr}B'}(\mathcal O'(i),\mathcal O'(j))
 \tag{BC}\label{lp:eq:bc}
\]
は擬同型である。
特に任意の剰余体$R\to\kappa(\mathfrak p)$、および正則noetherian局所環への局所準同型には適用できる。
\end{proposition}
\begin{proof}
$A,B$の有限相対分解と$B$の正則中心商表示は底変更と両立するので、
$A',B'$も対応するGorenstein型、有限injective dimension、$\chi$を持つ。
剰余体については中央の商・局所化によりnoetherian性を得る。
局所準同型の場合は剰余が多項式環なので補題\ref{lp:lem:noetherian}を再適用できる。

まず$R$の有限大域次元により$-\otimes_R^{\mathbf L}R'$は有限Tor振幅を持ち、
有界coherent graded複体を有界coherent複体へ送り、torsionを保つ。
従って射影商の有界導来圏にも底変更が定義される。
自然な底変更写像そのものを確認するため、$M_\ell$を有限階数graded自由$B$-加群の分解$P_\bullet$で計算する。
$M_\ell$は内部次数別$R$-自由なので、そのsyzygyを各次数で分裂させて、分解の全底変更を保てる。
固定した内部次数$t$で$C^\bullet=\operatorname{\underline{Hom}}_B(P_\bullet,B)_t$は
下に有界な有限自由$R$-加群の複体であり、前補題の証明により、
cohomologyは有限自由加群の一つの次数にしか現れない。
$R$の大域次元は有限なので、cohomologyが全て射影なこの複体は、そのcohomology複体にhomotopy同値である。
実際、射影加群からなる完全複体のcycleについてdimension shiftingを$d+1$回適用すると、
全cycleが射影となり複体が分裂する。
従って$C^\bullet$の導来テンソルは通常のテンソルで計算できる。
有限階数自由分解上ではHomと底変更が項ごとに可換であるため、$M_\ell$のHomの底変更を得る。
各$t$で安定するfiltered colimitと式\eqref{lp:eq:torsion}を通せば\eqref{lp:eq:bc}となる。
\end{proof}

\section{有限生成proper部分圏と残余圏}
\label{lp:sec:proper}
\begin{proof}[有界Orlov分解と有限生成の証明]
命題\ref{lp:prop:gorenstein}の結果はBurke--StevensonのDefinition 5.1の仮定を満たす。
$R$は有限大域次元で、$A,B$は中央$R$-flatなので、同論文Lemma 6.2によって$\chi$を満たす。
Theorem 6.4を$A$へ適用すると、$D_{\rm sg}^b(\operatorname{gr}A)=0$であるため
任意の連続した$N$個のtwistが$D^b(\operatorname{qgr}A)$を生成する。
$B$への導来制限は$0\to A(-3)\to A\to B\to0$により有限Tor振幅を持ち、
torsionをtorsionに送るので射影商へ降りる。
この関手は$\mathcal O_A(j)$を$\mathcal O_B(j)$へ送る。
従って全ての$\mathcal O_B(j)$は$E$から有限回のcone・直和因子で作られる。
これが式\eqref{lp:eq:proper-model}の生成性である。

$B$へのTheorem 6.4は、$a>0$より$a$個のtwistとgraded singularity categoryによる半直交分解を与える。
この有界分解の向きは引用定理の規約に従う。以下では線束生成部分における右直交を別に構成する。
この有界三角圏の同値を、任意の導来基底変更と両立するdg同値へ強めたとは主張しない。
\end{proof}

\begin{proposition}[properな残余dg代数]\label{lp:prop:residual}
$\mathcal E=\langle\mathcal O,\ldots,\mathcal O(a-1)\rangle$とする。
各成分は$\operatorname{Perf}(R)$であり、$\mathcal C_B$で
\[
 \mathcal C_B=\langle\mathcal K_B^{\rm lin},\mathcal O,\ldots,\mathcal O(a-1)\rangle
 \tag{SOD}\label{lp:eq:sod}
\]
を持つ。$p=L_{\mathcal O}\circ\cdots\circ L_{\mathcal O(a-1)}$と置くと
$G=p(E)$は$\mathcal K_B^{\rm lin}$を生成する。
$\Gamma_R=\operatorname{REnd}(G)$の底上の複体はperfectであり、
命題\ref{lp:prop:basechange}の底変更でこの構成はMorita同値まで可換である。
\end{proposition}
\begin{proof}
式\eqref{lp:eq:H}から$\operatorname{RHom}(\mathcal O(i),\mathcal O(i))=R$で、
$0\le i<j\le a-1$なら$\operatorname{RHom}(\mathcal O(j),\mathcal O(i))=0$である。
全$\mathcal C_B$の対象間のHomは、$E$間のperfect Homから有限cone・直和因子で作られるので$R$-perfectである。
従って各例外成分の左右随伴はHomとその$R$-双対で定義でき、mutationは有限操作となる。
右端の例外対象から順に左mutationして$p$を作ると、$p(\mathcal C_B)$は正確に$\mathcal E^\perp$である。
$E$の生成性より$G$が生成し、$\operatorname{REnd}(G)$も$R$-perfectとなる。
標準dg enhancementと生成対象のdg Yonedaを用いればMorita表示を得る\cite[Theorem 3.8]{LP-Keller}。

式\eqref{lp:eq:bc}は生成対象間の全導来Homと積に自然なので、有限coneと冪等完備化へ延長すると
線束生成部分圏の底変更で完全忠実性を得る。像は全線束を生成するのでMorita同値である。
評価写像とそのconeも底変更に可換で、$p,G$についても同じである。
特殊ファイバーでは$X$は滑らかで、全twistは$\operatorname{Perf}(X)=D^b(\operatorname{coh}X)$を生成する。
右直交は通常のKuznetsov成分なので主定理(5)が従う。
\end{proof}

\begin{lemma}[局所proper圏の完全忠実性の持上げ]\label{lp:lem:ff-lift}
$\mathcal C,\mathcal D$を$R$-properな冪等完備dg圏とし、対象間のHomが剰余体への導来底変更と可換であるとする。
$\Phi:\mathcal C\to\mathcal D$を底変更と両立するexactなdg関手とする。
$\Phi_k$が完全忠実なら$\Phi$も完全忠実である。
さらに$\Phi_k$が同値で、$\Phi$が右または左随伴を持つなら$\Phi$は同値である。
\end{lemma}
\begin{proof}
任意の対象$M,N$に対し、Hom上の自然な写像のconeは$R$-perfectで、その剰余は零である。
局所環上の導来Nakayamaによりconeは零となるので完全忠実性を得る。
また$Q\in\mathcal D$について$Q_k=0$なら、その自己Homの剰余も零である。
従って$\operatorname{RHom}(Q,Q)=0$、特に恒等射が零なので$Q=0$となる。
$\Phi$の像が片側admissibleであれば、任意の対象は像とその直交成分によるtriangleを持つ。
直交成分の対象$Q$は剰余上で$\Phi_k(\mathcal C_k)$と直交し、$\Phi_k$の同値性から$Q_k=0$。
したがって$Q=0$であり、像は全体である。
\end{proof}

\begin{proposition}[標準rotation候補の完全忠実性]\label{lp:prop:rotation-candidate}
$\widetilde\rho=L_{\mathcal O}\circ(1)$は$\mathcal K_B^{\rm lin}$を保つ完全忠実なexact dg自己関手であり、
特殊ファイバーでは通常のrotationとなる。
上記の底変更と両立する。
本節だけでは全基底上の本質的全射性と$\widetilde\rho^3\simeq[2]$を主張しない。
\end{proposition}
\begin{proof}
$M\in\mathcal E^\perp$なら
$\operatorname{RHom}(\mathcal O(j),M(1))=0$（$1\le j\le a$）。
$L_{\mathcal O}$は$\mathcal O$からのHomを零にする。
$1\le j<a$について$\operatorname{RHom}(\mathcal O(j),\mathcal O)=0$だから、
このmutationは他の必要な直交条件を壊さない。
twistは線束生成部分圏を保ち、mutationも有限操作なので自己関手を得る。
基底変更は前命題の証明と同じである。
特殊ファイバーでは滑らかなcubicの標準rotationであり自己同値なので、
補題\ref{lp:lem:ff-lift}によって$R$上で完全忠実になる。
\end{proof}
\begin{remark}[自己同値化を残す理由]
補題\ref{lp:lem:ff-lift}により、$\widetilde\rho$に有限対象を保つ片側随伴を構成すれば、
本質的全射性も従う。しかしproper性だけから任意のexact関手が随伴を持つとは限らない。
本稿ではこの随伴の存在を仮定なしに付け加えない。
同様に、$\widetilde\rho^3\to[2]$という自然変換を実際に構成し、特殊ファイバーで同型と示せれば、
対象ごとのconeに導来Nakayamaを適用して全体で同型にできる。
自然変換の持上げそのものを、その特殊ファイバーでの存在で代用してはならない。
\end{remark}

\subsection{中心商の導来自己交叉と特異圏の周期性}
\begin{proposition}[三次周期性の代数的な構成]\label{lp:prop:sing-period}
本節の仮定で、graded singularity categoryには自然な同型
\[
 (3)\simeq[2]
 \qquad\text{on }D_{\rm sg}^b(\operatorname{gr}B)
 \tag{PER}\label{lp:eq:period}
\]
が存在する。ここでは射影残余圏のsmoothnessを使わない。
\end{proposition}
\begin{proof}
$Q=B\otimes_A^{\mathbf L}B$をgraded $B$-双加群とする。
中央正則元の長さ一分解から、$H^0(Q)=B$、
$H^{-1}(Q)=(\Omega)/(\Omega)^2\simeq B(-3)$で、他のcohomologyは零である。
この同定は$B$-双加群として成立する。導来双加群の標準truncationは
\[
 B(-3)[1]\longrightarrow Q\longrightarrow B
 \xrightarrow{\eta}B(-3)[2]
\]
を与える。
任意の$M\in D^b(\operatorname{gr}B)$に$M\otimes_B^{\mathbf L}-$を作用させると
\[
 M(-3)[1]\longrightarrow M\otimes_A^{\mathbf L}B\longrightarrow M
 \xrightarrow{\eta_M}M(-3)[2]
\]
となる。$A$はnoetherianで有限大域次元なので$M$は$A$上perfectであり、
中項は$B$上perfectである。従って$\eta_M$はgraded singularity categoryで同型となる。
$\eta$が双加群射から来ることにより$M$に自然である。
内部shiftを三つ戻すと式\eqref{lp:eq:period}を得る。
\end{proof}
\begin{remark}[二種類のrotationを混同しない]
有界Orlov同値で内部shift $(1)$を移せば、全coherent残余圏には
三乗が$[2]$となる自己同値$\rho_{\rm sg}$を定義できる。
他方、命題\ref{lp:prop:rotation-candidate}は線束生成部分に
$\widetilde\rho=L_{\mathcal O}\circ(1)$を具体的に構成する。
この二つの関手を指定したwindowで同定し、全体での底変更とcoherenceを確認する作業は別である。
本稿は通常の中心商から得る特異圏の周期性そのものを未解決とはしないが、
上の同定を省略して線束生成残余部分の相対rotationが完成したとは主張しない。
また$Q$を$B\oplus B(-3)[1]$と双加群として分裂させてはならない。
その拡張類$\eta$こそが周期性を与えている。
\end{remark}

\section{非楕円成分への適用と有限計算}
\label{lp:sec:application}
前稿で構成した二パラメータ非楕円量子化の係数環は
$R_{\mathrm C}^0=\CC[r,s]$、$R_{\mathrm N}^0=\CC[r,s,(4-r^2)^{-1}]$である。
可換点$(r,s)=(0,0)$で局所化すると、どちらも正則局所環
$R=\CC[r,s]_{(r,s)}$になる。
各次数の自由性（局所化前は局所自由性）、七生成元の多項式環への特殊化、
中心三次元は前稿の全次数PBW構成から得られる。
従って主定理は$N=7,n=5,a=4$で適用できる。
この適用は節点・尖点の純ブロックの滑らかさを仮定しない。
既存の中心元・PBW性の証明と再現コードは本bundleの別論文に保持している。

\begin{corollary}[二つの成分の局所proper残余模型]\label{lp:cor:nonelliptic}
上の二つの成分に対し、$R=\CC[r,s]_{(r,s)}$上の実在するproperな残余dg代数
$\Gamma_R$と完全忠実な標準rotation候補を構成できる。
その特殊ファイバーは前稿の滑らかな$cubic\ fivefold$のKuznetsov成分とrotationである。
$B$は有界Orlov分解を持つ。
ただしこの残余dg代数の相対smoothnessは未証明である。
\end{corollary}

\begin{proposition}[線束生成子のcohomology階数]\label{lp:prop:ranks}
$b_q=\binom{N-1+q}{N-1}-\binom{N-4+q}{N-1}$（負次数の項は零）とする。
$\Lambda_R=\operatorname{REnd}(\bigoplus_{i=0}^{N-1}\mathcal O(i))$について
\begin{align*}
 \rank_R H^0(\Lambda_R)&=\sum_{q=0}^{N-1}(N-q)b_q
 =\binom{2N}{N-1}-\binom{2N-3}{N-4},\\
 \rank_R H^{N-2}(\Lambda_R)&=3+2N+\binom{N+1}{2}.
\end{align*}
それ以外のcohomologyは零である。$N=7$では階数は$2838$と$45$である。
\end{proposition}
\begin{proof}
第0次は差$j-i=q\ge0$を持つ組が$N-q$個あることから得る。
二つのbinomial和にhockey-stick公式を二回使うと閉公式になる。
上の次数は$i-j\ge a=N-3$だけが寄与し、差$a,a+1,a+2$を持つ組はそれぞれ3,2,1個である。
$B_0=R$、$\rank B_1=N$、$\rank B_2=\binom{N+1}{2}$を代入する。
\end{proof}

\begin{remark}[cohomologyの表はdg積の代用品ではない]\label{lp:rem:ainfty}
上の数値から$\Lambda_R$や$\Gamma_R$のdg積・高次積を復元したとは主張しない。
特に$H^0(\Lambda_R)\oplus H^{N-2}(\Lambda_R)[-(N-2)]$を
自動的に同じdg代数のsquare-zero模型とみなしてはならない。
計算コードはbinomial恒等式、Homの三角的消滅、例外成分の数を整数で確認するだけで、
noetherian性・Orlov定理・dg Morita同値を形式化したものではない。
\end{remark}

\section{残余smoothnessに残る正確な問題}
\label{lp:sec:remaining}
構成した$\Gamma_R\in\operatorname{Perf}(R)$について、なお必要なのは
\[
 \Gamma_R\in\operatorname{Perf}(\Gamma_R^{\rm op}\otimes_R^{\mathbf L}\Gamma_R).
 \tag{SM}\label{lp:eq:smoothness}
\]
周囲$A$の対角の有限分解は主定理で得たが、式\eqref{lp:eq:smoothness}ではない。
$\Gamma_R\otimes_Rk$がsmoothであることだけから、代数的な有限対角分解の持上げを結論しない。
実際のperfect複体$P$と$P\to\Gamma_R$を構成して一ファイバーで擬同型とすれば、
coneの底上perfect性からその消滅を開近傍へ広げられる。
この判定は有効だが、$P$と射を構成する仕事を代替しない。

同様に、通常の有界Orlov分解があることから、任意の導来基底でのdg同値、
標準rotationの可逆bimodule表示、三次周期性のcoherenceを無条件に結論してはならない。
本稿ではproperな線束生成部分を先に作り、その自己関手候補までを構成した。
全coherent残余圏の比較、標準rotationの本質的全射性と相対周期性の同定は次の課題である。

\begin{remark}[局所環と有限型開集合を区別する]
$R$は可換点の局所環なので、全次数にわたるNakayamaは同一の$R$上で使える。
しかしこれだけで、元の有限型基底の一つのZariski開集合上に
全ての無限次数条件が広がったとはいえない。
有限の自由分解・有限個の行列の個別のspreadと、proper dg代数の全高次積のspreadも区別する。
実際に代数的なsmooth proper pairを得た後に、既存原稿Cのétale局所実現原理(ER)を適用する。
(ER)の圏論的入力は本稿で再証明していない。
\end{remark}

\begin{remark}[同じHom階数でもsmoothとは限らない]
本稿の線束Homの計算は、非零三次式$F$が特異であっても成り立つ。
例えば$F=z_1^2z_2+z_3^3+\cdots+z_N^3$は$[0:1:0:\cdots:0]$で特異だが、
同じHilbert級数と命題\ref{lp:prop:ranks}の階数を持つ。
可換射影超曲面の$\operatorname{Perf}$はこの場合smoothでない。
実際、smooth properならcohomologically finiteなmoduleはperfectであり、
特異点のskyscraperもperfectとなって、その局所環の正則性に矛盾する。
この標準的な対角判定は\cite[Proposition 3.31]{LP-Orlov}のsmooth/properの枠組みに属する。
従って階数表だけをsaturationの証明書として使うことはできない。
\end{remark}
\section{線束生成子の相対最小$A_\infty$模型と対角smoothness}
\label{sec:minimal-ainfty}
本節は、残余圏の対角perfect性(SM)を実際に検証するため、
線束生成子のdg代数へ問題を移し、その相対最小模型の高次演算を制限する。
相対半直交分解の一般定理は既知であり
\cite[Lemma~4.15(3)]{AinftyPerry2019}を参照する。
最小模型の存在には\cite{PetersenTransfer,BurkeTransfer}を用いる。
以下はその適用と整数的な配置の絞り込みであり、SM自体の証明ではない。

\begin{theorem}[残余対角と線束生成子の対角]
\label{thm:ainfty-smooth-equiv}
定義\ref{lp:def:input}の正則局所底$R$上で、
既存の局所proper模型
$\mathcal C_B=\thick_R(E)\simeq\Perf(\Lambda_R)$、
$\mathcal K_B^{\mathrm{lin}}\simeq\Perf(\Gamma_R)$
を考える。ただし
\[
 E=\bigoplus_{i=0}^{N-1}\mathcal O(i),\quad
 \Lambda_R=\REnd(E),\quad N=n+2,\quad n\ge3.
\]
このとき、相対smoothnessについて
\[
 \boxed{\ \Gamma_R\in\Perf(\Gamma_R^e)
 \quad\Longleftrightarrow\quad
 \Lambda_R\in\Perf(\Lambda_R^e)\ },\qquad
 A^e:=A^{\mathrm{op}}\otimes_R^{\mathbf L}A.
 \tag{SM$'$}\label{eq:ainfty-sm-prime}
\]
従ってSMの有限対角証明書を$\Lambda_R$について求めてもよい。
\end{theorem}
\begin{proof}
既存の有限$R$-例外列により
\[
 \mathcal C_B=
 \langle\mathcal K_B^{\mathrm{lin}},
 \mathcal O,\mathcal O(1),\ldots,\mathcal O(n-2)\rangle
\]
は$R$-線形半直交分解である。各相互Homは$R$-perfectであり、
評価・余評価から定める左右mutationが実在するため、各成分はadmissibleである。
$\mathcal O(i)$の生成する成分は$\Perf(R)$と同値であり、$R$上smoothである。

相対半直交分解に関するPerryの補題
\cite[Lemma~4.15(3)]{AinftyPerry2019}により、
admissibleな各成分がsmoothであることと全体がsmoothであることは同値である。
この場合の具体的な理由も記しておく。
全体の対角を残余成分の左右両変数へ制限すると残余対角となり、
包含の両随伴によりrepresentable双加群はrepresentableへ送られるので、
全体の対角perfect性は残余対角perfect性を含意する。
逆に、残余成分と各例外成分がsmoothなら、各貼合せ双加群は
$R$上perfect（全体がproper）であり、
smoothな両成分の包絡dg圏上でもperfectである。
三角dg代数の対角を与える有限三角を使って帰納的に貼り合わせれば
全体の対角がperfectになる。Morita不変性から式\eqref{eq:ainfty-sm-prime}を得る。
\end{proof}

\begin{proposition}[二段階のsquare-zeroコホモロジー環]
\label{prop:ainfty-squarezero}
$H=H^*(\Lambda_R)$、$U=H^0$、$M=H^n$とおく。
通常の積について$H=U\ltimes M[-n]$かつ$M^2=0$である。
各対象指数において
\[
 U(i,j)=
 \begin{cases}B_{j-i}&j\ge i,\\0&j<i,\end{cases}\qquad
 M(i,j)=
 \begin{cases}B_{i-j-(n-1)}^\vee&i-j\ge n-1,\\
 0&\text{otherwise},\end{cases}
 \tag{SQ}\label{eq:ainfty-sq}
\]
である。$U$の積は$B$の積、$M$への作用は線束の合成と
Serre双対表示から定まる。
\end{proposition}
\begin{proof}
既存の線束Hom公式\eqref{lp:eq:H}により$H$は次数$0,n$だけにある。
$m_2$は次数零なので$H^nH^n\subset H^{2n}=0$。
残る成分の式は\eqref{lp:eq:H}そのものである。
\end{proof}

\begin{proposition}[相対最小模型の高次演算の間隙]
\label{prop:ainfty-arity}
$R$-線形でstrictly unitalな最小$A_\infty$模型を$H$上に取れる。
その積$m_k$と、$f_1=\mathrm{id}_H$である最小模型の比較同型の成分$f_k$は
\[
 m_k=0\quad(k\not\equiv2\!\!\!\pmod n),\qquad
 f_k=0\quad(k\not\equiv1\!\!\!\pmod n).
 \tag{AR}\label{eq:ainfty-arity}
\]
特に$m_3,\ldots,m_{n+1}$は零であり、
最初の可能な高次積は$m_{n+2}$、最初の可能な非線形な
比較同型は$f_{n+1}$である。
\end{proposition}
\begin{proof}
Homの各コホモロジーは有限自由$R$-加群で、
Hom複体は$R$-perfectである。従って射影複体模型を選べば
cycle、boundaryの列を分裂でき、単位を保存する
homotopy transferによって相対最小模型を得る。
$m_k$のcohomological degreeは$2-k$である。
次数$n$の入力を$p$個、出力次数を$q\in\{0,n\}$とすれば
$np+2-k=q$、従って$k\equiv2\pmod n$。
比較同型$f_k$の次数は$1-k$であり、同じ議論から
$k\equiv1\pmod n$。strict unital性によって
単位を含む$k\ge3$の入力では$m_k$が消える。
\end{proof}

\begin{proposition}[最初のHochschild類と次の二次障害]
\label{prop:ainfty-hoch}
$H$を$m_2$による次数付き代数とすると、
$m_{n+2}$は内部次数$-n$のHochschild cocycleを定め、
\[
 [m_{n+2}]\in\HH^{n+2}(H,H)_{-n}
 \tag{HC}\label{eq:ainfty-hh}
\]
は$f_1=\mathrm{id}$のstrictly unitalな最小模型同型に依らない。
次に現れ得る積は$m_{2n+2}$で、suspended cochainの規約では
\[
 \delta m_{2n+2}
   =-\frac12[m_{n+2},m_{n+2}].
 \tag{OB}\label{eq:ainfty-ob}
\]
従って任意の候補$m_{n+2}$について
$\frac12[m_{n+2},m_{n+2}]$の
$\HH^{2n+3}(H,H)_{-2n}$での類の消滅は延長可能性の必要条件である。
実在する$\Lambda_R$の模型ではこの類は消えるが、
$m_{2n+2}=0$とは結論しない。
\end{proposition}
\begin{proof}
$m_3,\ldots,m_{n+1}$が零なので、
長さ$n+3$のStasheff恒等式は$\delta m_{n+2}=0$である。
同型条件の最初の非自明な式では、$m_{n+2}$の差が
$\delta f_{n+1}$となるので類は不変である。
長さ$2n+3$の式は$m_2,m_{2n+2}$の合成と
$m_{n+2}$同士の挿入からなり、表示したGerstenhaber関係を得る。
後者の内部次数は$-2n$である。
\end{proof}

\begin{proposition}[最初の高次積の対象指数制限]
\label{prop:ainfty-indices}
$m_{n+2}$の非零候補は、次数$n$の入力を一個持ち
次数零を出力するもの、または次数$n$の入力を二個持ち
次数$n$を出力するものだけである。
前者が$\operatorname{End}^0(\mathcal O(i))=R$を出力するなら、
指数列は
\[
0\to1\to\cdots\to N-1\to0
\tag{CY}\label{eq:ainfty-cyclic}
\]
の$N$個の巡回移動に限る。
後者では$n$個の正の指数上昇の合計は
$n,n+1,n+2,n+3$のいずれかであり、
各上昇幅は高々$4$である。
\end{proposition}
\begin{proof}
式\eqref{eq:ainfty-sq}より、非単位な次数零の射は指数が
少なくとも一つ上昇し、次数$n$の射は指数が
$n-1,n,n+1$だけ下降する。
$m_{n+2}$の次数条件から高次入力数は一個または二個となる。
出力が同じ指数の場合、$n+1$個の上昇幅の和は
一個の下降幅に等しい。各上昇幅が$\ge1$、
下降幅が$\le n+1$なので全ての上昇幅は一、
下降は$N-1$から$0$に限られる。
二個の下降幅を$d_1,d_2$、出力の下降幅を$d_0$とすると
$n$個の上昇幅の和は$d_1+d_2-d_0$で、
各$d_i\in\{n-1,n,n+1\}$より$n\le u\le n+3$。
一つでも上昇幅が$5$以上なら$u\ge n+4$となり矛盾する。
\end{proof}

\begin{remark}[cubic fivefoldの有限計算の範囲]
$n=5,N=7$なら高次積候補は$m_7,m_{12},m_{17},\ldots$。
$H^0$と$H^5$の全$R$-階数は既存の命題\ref{lp:prop:ranks}で
それぞれ$2838,45$である。
strictly unitalな指数列に限ると、
$m_7$の対象指数配置は出力$H^0$側$2478$通り、
出力$H^5$側$303$通りで、前者の対角指数への出力は$7$通り。
これはHomのベクトル空間基底を数えた値ではない。
計算は\texttt{computations/verify\_minimal\_ainfty\_arity.py}で再現できる。
縮約$(i,p,h)$を実際に選べば、$m_7$は七葉の平面二分木
$C_6=132$個による有限和として計算できる。
しかし高次積の係数・相対bar右逆をまだ得ていない。

$m_7$を二度挿入する許容指数列として
\[
 0\to1\to2\to3\to4\to0\to1\to2\to3
 \to4\to0\to1\to2\to3
\]
がある。二つの$4\to0$が次数$5$で、それ以外は次数零。
最初の七射の$m_7$出力は$0\to2$の次数零に、
その出力を含む外側の$m_7$は$0\to3$の次数零に
入り得る。従って$[m_7,m_7]$は指数条件だけでは消えない。
非零の値やHochschild障害類を計算した主張ではない。
特に、ここから$\Gamma_R\in\Perf(\Gamma_R^e)$を結論してはならない。
\end{remark}

\begin{thebibliography}{99}
\bibitem{PHartshorne}
R.~Hartshorne,
\emph{Algebraic Geometry}, Graduate Texts in Mathematics 52,
Springer, 1977.
\href{https://doi.org/10.1007/978-1-4757-3849-0}{doi:10.1007/978-1-4757-3849-0}.

\bibitem{PMayanskiy}
E.~Mayanskiy,
Poisson cohomology of two Fano threefolds,
\emph{Journal of Algebra} \textbf{424} (2015), 21--45.
\href{https://doi.org/10.1016/j.jalgebra.2014.08.049}{doi:10.1016/j.jalgebra.2014.08.049};
\href{https://arxiv.org/abs/1303.5916}{arXiv:1303.5916}.
本文で用いる Grassmann 記述は Theorem~1。

\bibitem{PEtingofLie}
P.~Etingof,
\emph{Lie Groups and Lie Algebras}, MIT lecture notes, 2024 course version.
\href{https://live.ocw.mit.edu/courses/18-755-lie-groups-and-lie-algebras-ii-spring-2024/mit18_755_s24_lec_full.pdf}{MIT OpenCourseWare 公開原稿}.
本文では Proposition~16.9、\S19.4、Theorem~25.17、\S32.2 を参照する。

\bibitem{PMoriSmith}
I.~Mori and S.~P.~Smith,
The classification of 3-Calabi--Yau algebras with 3 generators and 3 quadratic relations,
\href{https://arxiv.org/abs/1502.07403v2}{arXiv:1502.07403v2} (2016).
Theorems~1.3, 1.4, 4.2 and Corollaries~1.6, 2.7, 2.9.
点曲線のHessian公式はLemma~7.1, Theorem~7.2, pp.~23--25。

\bibitem{PBergerTaillefer}
R.~Berger and R.~Taillefer,
Poincar\'e--Birkhoff--Witt deformations of Calabi--Yau algebras,
\emph{Journal of Noncommutative Geometry} \textbf{1} (2007), 241--270.
\href{https://arxiv.org/abs/math/0610112}{arXiv:math/0610112}.
本文で用いる potential による PBW 変形は Theorem~1.1(i)。

\bibitem{PKontsevich}
M.~Kontsevich,
Deformation quantization of Poisson manifolds, I,
\href{https://arxiv.org/abs/q-alg/9709040}{arXiv:q-alg/9709040} (1997).
本文では \S4.6 および \S6 の affine formality を用いる。

\bibitem{PDolgushev}
V.~Dolgushev,
The Van den Bergh duality and the modular symmetry of a Poisson variety,
\href{https://arxiv.org/abs/math/0612288v3}{arXiv:math/0612288v3} (2007).
Theorem~3 および Proposition~1。

\bibitem{Beilinson}
A. A. Beilinson,
\emph{Coherent sheaves on \(\PP^n\) and problems of linear algebra},
Functional Analysis and Its Applications \textbf{12} (1978), 214--216.
\url{https://doi.org/10.1007/BF01681436}.
原著Section 4の対角Koszul分解を確認。
\url{https://www.mathnet.ru/eng/faa2009}.
\bibitem{DDL}
O. De Deken and W. Lowen,
\emph{Abelian and derived deformations in the presence of \(\ZZ\)-generating geometric helices},
Journal of Noncommutative Geometry \textbf{5} (2011), 477--505.
Theorem 4.5; Propositions 5.5, 5.8, 5.9; Lemma 5.14.
\url{https://doi.org/10.4171/JNCG/83}.
\bibitem{Keller}
B. Keller, \emph{Deriving DG categories},
Annales scientifiques de l'École Normale Supérieure \textbf{27} (1994), 63--102.
\url{https://doi.org/10.24033/asens.1689}.
\bibitem{Mathew}
A. Mathew, \emph{The Galois group of a stable homotopy theory},
Advances in Mathematics \textbf{291} (2016), 403--541.
Proposition 3.28, Remark 3.29, Proposition 3.35.
\url{https://arxiv.org/abs/1404.2156}.
\bibitem{Orlov}
D. Orlov,
\emph{Derived categories of coherent sheaves and triangulated categories of singularities},
arXiv:math/0503632, v3 (2009), Theorem 2.5.
\url{https://arxiv.org/abs/math/0503632}.
\bibitem{CCKM}
T. Cassidy, A. Conner, E. Kirkman and W. F. Moore,
\emph{Periodic free resolutions from twisted matrix factorizations},
Journal of Algebra \textbf{455} (2016), 137--163.
Proposition 5.7, Theorem 5.8.
\url{https://doi.org/10.1016/j.jalgebra.2016.01.037}.
\bibitem{CT}
A. Căldăraru and J. Tu, \emph{Curved \(A_\infty\) algebras and Landau--Ginzburg models},
New York Journal of Mathematics \textbf{19} (2013), 305--342.
Section 5、特にdg Orlov windowのTheorem 5.8.
\url{https://nyjm.albany.edu/j/2013/19-16v.pdf}.
\bibitem{BFK}
M. Ballard, D. Favero and L. Katzarkov,
\emph{A category of kernels for equivariant factorizations and its implications for Hodge theory},
Publications Mathématiques de l'IHÉS \textbf{120} (2014), 1--111.
Theorem 5.39のKoszul cohomologyによる公式、Theorem 5.15とCorollary 5.18のequivariant外部積.
\url{https://doi.org/10.1007/s10240-013-0059-9}.
確認した公刊本文 \url{https://www.numdam.org/article/PMIHES_2014__120__1_0.pdf}.
\bibitem{LZ}
X. Lin and S. Zhang,
\emph{Serre algebra, matrix factorization and categorical Torelli theorem for hypersurfaces},
arXiv:2310.09927, v2 (2024), Theorem 1.7.
\url{https://arxiv.org/abs/2310.09927}.
\bibitem{HR}
D. Huybrechts and J. V. Rennemo,
\emph{Hochschild cohomology versus the Jacobian ring, and the Torelli theorem for cubic fourfolds},
Algebraic Geometry \textbf{6} (2019), 76--99.
Lemma 2.19およびSections 2--3.
\url{https://doi.org/10.14231/AG-2019-005}.
\bibitem{Huy}
D. Huybrechts, \emph{The K3 category of a cubic fourfold},
Compositio Mathematica \textbf{153} (2017), 586--620.
Proposition 3.12およびSection 3.5.
\url{https://doi.org/10.1112/S0010437X16008137}.
\bibitem{JLLZ}
A. Jacovskis, X. Lin, Z. Liu and S. Zhang,
\emph{Infinitesimal categorical Torelli theorems for Fano threefolds},
arXiv:2203.08187, Theorem 4.4(2).
\url{https://arxiv.org/abs/2203.08187}.
\bibitem{Lurie}
J. Lurie, \emph{Derived Algebraic Geometry X: Formal Moduli Problems},
2011. Theorems 5.1.9, 5.3.16およびProposition 5.3.21.
\url{https://www.math.ias.edu/~lurie/papers/DAG-X.pdf}.
\bibitem{BKP}
A. Blanc, L. Katzarkov and P. Pandit,
\emph{Generators in formal deformations of categories},
Compositio Mathematica \textbf{154} (2018).
本文で用いるのはarXiv版のDefinition 3.3およびSection 3.3.
\url{https://arxiv.org/abs/1705.00655}.
\bibitem{KL}
B. Keller and W. Lowen,
\emph{On Hochschild cohomology and Morita deformations},
International Mathematics Research Notices (2009).
\url{https://webusers.imj-prg.fr/~bernhard.keller/publ/KellerLowenHochDefo.pdf}.
\bibitem{EG}
P. Etingof and V. Ginzburg,
\emph{Noncommutative del Pezzo surfaces and Calabi--Yau algebras},
Journal of the European Mathematical Society \textbf{12} (2010), 1371--1416.
Sections 3.5--3.6, Theorem 3.6.2, Proposition 8.1.1(i) およびその証明。
\url{https://doi.org/10.4171/JEMS/235}.
\href{https://arxiv.org/abs/0709.3593v4}{arXiv:0709.3593v4}.
\bibitem{Walton}
C. Walton, \emph{Representation theory of three-dimensional Sklyanin algebras},
Nuclear Physics B \textbf{860} (2012), 167--185.
Theorem 2.14および式(2.15)--(2.16).
\url{https://doi.org/10.1016/j.nuclphysb.2012.02.015}.
\bibitem{CGWZ}
K. Chan, J. Gaddis, R. Won and J. J. Zhang,
\emph{Ozone groups of Artin--Schelter regular algebras satisfying a polynomial identity},
arXiv:2312.13014, v2 (2024). Introductionの非退化Sklyanin代数の記述。
\url{https://arxiv.org/abs/2312.13014}.
\bibitem{Elagin}
A. Elagin, \emph{On equivariant triangulated categories},
arXiv:1403.7027, Theorem 4.2.
\url{https://arxiv.org/abs/1403.7027}.
\bibitem{ELS}
A. Elagin, V. A. Lunts and O. M. Schn\"urer,
\emph{Smoothness of Derived Categories of Algebras},
Moscow Mathematical Journal \textbf{20} (2020), no. 2, 277--309.
確認した一次資料はarXiv:1810.07626v2 (22 March 2019),
Theorem 4.4 (Verdier局所化), Proposition 4.10とTheorem 4.15 (稠密開集合と生成子),
Theorem 4.18 (外部積生成), Theorem 5.1 (有限中心代数の有界導来圏のsmoothness).
\url{https://arxiv.org/abs/1810.07626};
\url{https://www.mathjournals.org/mmj/2020-020-002/2020-020-002-003.html}.

\bibitem{TV}
B. Toën and M. Vaquié,
\emph{Moduli of objects in dg-categories},
Annales scientifiques de l'École Normale Supérieure (4) \textbf{40} (2007), 387--444.
Theorem 3.6; Corollaries 2.13, 3.15, 3.19, 3.21; Lemma 3.25 and proof of Corollary 3.24.
\url{https://arxiv.org/abs/math/0503269};
確認した著者版 \url{https://www.math.univ-toulouse.fr/~btoen/ttt.pdf}.
\bibitem{Stacks}
The Stacks Project Authors, \emph{The Stacks Project}.
Tag 02HX (Lemma 37.12.1, Noetherian底上の無限小smoothness判定),
Tag 055U (Lemma 37.38.6, smooth morphismのétale局所section).
\url{https://stacks.math.columbia.edu/tag/02HX};
\url{https://stacks.math.columbia.edu/tag/055U}.
参照日2026年11月4日。
\bibitem{Pirozhkov}
D. Pirozhkov, \emph{A categorical Torelli theorem for hypersurfaces},
Bulletin of the London Mathematical Society \textbf{56} (2024), 3075--3089.
\url{https://doi.org/10.1112/blms.13117}.
公刊版のabstractで、cubicに対する仮定 \(n>3\) を確認。
本稿の全次元の一意性には \cite[Theorem 1.7]{LZ} を用いる。
\bibitem{HuyUpdate}
D. Huybrechts, \emph{The K3 category of a cubic fourfold: an update},
Beiträge zur Algebra und Geometrie / Contributions to Algebra and Geometry (2025).
Section 9: Unknown deformations.
\url{https://doi.org/10.1007/s13366-025-00788-9}.

\bibitem{ChanMulti}
D. Chan, \emph{Twisted multi-homogeneous coordinate rings},
Journal of Algebra \textbf{223} (2000), 438--456.
Theorem 3.4; proof of Theorem 5.2, Claim 5.3 and Lemmas 5.4--5.6;
Corollaries 5.8--5.9.
\url{https://doi.org/10.1006/jabr.1999.7938}.
確認した著者版 \url{https://web.maths.unsw.edu.au/~danielch/paper/tbcr.ps}, pp.8--9,14--17.
\bibitem{SmithTate}
S. P. Smith and J. Tate,
\emph{The center of the 3-dimensional and 4-dimensional Sklyanin algebras},
K-Theory \textbf{8} (1994), 19--63.
Corollary 2.3 (p.~26); Theorem 3.7(c) (p.~38).
\url{https://doi.org/10.1007/BF00962089};
\url{https://sites.math.washington.edu/~smith/Research/tate.pdf}.
\bibitem{KeelerMulti}
D. S. Keeler, \emph{Noncommutative ampleness for multiple divisors},
Journal of Algebra \textbf{265} (2003), no. 1, 299--311.
Theorem 2.7, Corollary 2.8 (任意有限個のbimoduleのNC-ampleness);
Hypothesis 3.1, Theorem 3.3, Corollary 3.5 (二因子のnoetherian性).
番号はarXiv一次TeXソースで確認し、刊行版との番号一致は未照合。
\url{https://doi.org/10.1016/S0021-8693(03)00126-1};
\url{https://arxiv.org/abs/math/0210417};
\url{https://arxiv.org/e-print/math/0210417}.
\bibitem{WWY3}
C. Walton, X. Wang and M. Yakimov,
\emph{Poisson geometry of PI three-dimensional Sklyanin algebras},
Proceedings of the London Mathematical Society \textbf{118} (2019), 1471--1500.
Theorem 1.3(1),(2a), verified in arXiv:1704.04975v3, p.~3.
\url{https://doi.org/10.1112/plms.12220};
\url{https://arxiv.org/abs/1704.04975}.
\bibitem{ToenFinitude}
B. Toën, \emph{Finitude homotopique des dg-algèbres propres et lisses},
Proceedings of the London Mathematical Society (3) \textbf{98} (2009), 217--240.
Lemmas 3.13--3.14. Lemma 3.14の結論は部分前層のquasi-compactnessであり、一般の開性ではない。
\url{https://doi.org/10.1112/plms/pdn027};
確認した著者版 \url{https://www.math-beb.fr/wp-content/uploads/2026/10/finitesatdgalg.pdf}.

\bibitem{SnowdenDVR}
A. Snowden, \emph{Lecture 8: Elliptic curves over DVRs},
Math 679 lecture notes, University of Michigan, 2013.
Section ``Reduction of torsion points'', good reductionでのprime-to-residue-characteristic torsionの還元の単射性。
\url{https://websites.umich.edu/~asnowden/teaching/2013/679/L08.html}.
参照日2026年11月6日。
\bibitem{HMY}
J.-W. He, X.-C. Ma and Y. Ye, \emph{Generalized Kn\"orrer's Periodicity Theorem},
arXiv:2107.06438v2 (16 August 2021).
対象は中心二次元の商であり、本稿の三次の場合には直接適用しない。
\url{https://arxiv.org/abs/2107.06438}.

\bibitem{OrlovGlue}
D. Orlov,
\emph{Smooth and proper noncommutative schemes and gluing of DG categories},
Advances in Mathematics \textbf{302} (2016), 59--105.
Theorem 3.25 (smoothness of dg gluing; Proposition 3.24はsmoothnessからregularityのみ)。
\url{https://doi.org/10.1016/j.aim.2016.07.014};
\url{https://arxiv.org/abs/1402.7364}.

\bibitem{BurkeStevenson}
J. Burke and G. Stevenson,
\emph{The derived category of a graded Gorenstein ring},
arXiv:1507.00830v1 (2015), Definition 5.1 and Theorem 6.4.
\url{https://arxiv.org/abs/1507.00830}.

% Current verified additions, 2026-11-08.
% Insert inside the manuscript's thebibliography environment.

\bibitem{UeyamaIsolated}
K. Ueyama,
\emph{Noncommutative graded Gorenstein isolated singularities},
in \emph{Proceedings of the 45th Symposium on Ring Theory and Representation Theory}
(Shinshu University, September 7--9, 2012),
K. Koike (ed.), February 2013, 158--163.
Theorem 7, printed p.~161.
\url{https://ring-theory-japan.com/ring/oldmeeting/2012/report2012/proceedings.pdf}.

\bibitem{ZhangTwist}
J.~J. Zhang,
\emph{Twisted graded algebras and equivalences of graded categories},
Proceedings of the London Mathematical Society (3) \textbf{72} (1996), no.~2, 281--311.
Theorems 1.1 and 1.3.
\url{https://doi.org/10.1112/plms/s3-72.2.281}.
\url{https://deepblue.lib.umich.edu/bitstream/handle/2027.42/135651/plms0281.pdf}.

\bibitem{RRZ}
M. Reyes, D. Rogalski and J.~J. Zhang,
\emph{Skew Calabi--Yau algebras and homological identities},
Advances in Mathematics \textbf{264} (2014), 308--354.
Theorems 0.3 and 5.4; Example 5.5.
\url{https://doi.org/10.1016/j.aim.2014.07.010}.
\url{https://arxiv.org/abs/1302.0437}.

\bibitem{Griffiths}
P.~A. Griffiths,
\emph{On the periods of certain rational integrals.~I},
Annals of Mathematics (2) \textbf{90} (1969), no.~3, 460--495.
Section 8, especially Theorems 8.1 and 8.3 and formula (8.6).
\url{https://doi.org/10.2307/1970746}.
\url{https://www.cimat.mx/~luis/seminarios/Teoria-de-Hodge/period_integral_I.pdf}.

\bibitem{KuznetsovFCY}
A. Kuznetsov,
\emph{Calabi--Yau and fractional Calabi--Yau categories},
Journal f\"ur die reine und angewandte Mathematik \textbf{753} (2019), 239--267.
Lemmas 2.6, 3.13 and 3.19; Corollary 4.1.
\url{https://doi.org/10.1515/crelle-2017-0004}.

\bibitem{CHL}
C.-H. Cho, H. Hong and S.-C. Lau,
\emph{Noncommutative homological mirror functor},
Memoirs of the American Mathematical Society \textbf{271} (2021), no.~1326.
Chapter 5, especially Theorems 5.15 and 5.16.
\url{https://doi.org/10.1090/memo/1326}.
\url{https://arxiv.org/abs/1512.07128}.

\bibitem{Lee}
S. Lee,
\emph{Noncommutative homological mirror symmetry of elliptic curves},
Kyoto Journal of Mathematics \textbf{61} (2021), no.~3, 723--743.
Theorem 6.3 and the corollary in the introduction.
\url{https://arxiv.org/abs/2004.09145}.

\bibitem{CisinskiTabuada}
D.-C. Cisinski and G. Tabuada,
\emph{Lefschetz and Hirzebruch--Riemann--Roch formulas via noncommutative motives},
Journal of Noncommutative Geometry \textbf{8} (2014), 1171--1190.
Section 4 and Example 5.6.
\url{https://doi.org/10.4171/JNCG/183}.
\url{https://www.math.univ-toulouse.fr/~dcisinsk/ncLHRR.pdf}.

\bibitem{MathewHH}
A. Mathew,
\emph{Kaledin's degeneration theorem and topological Hochschild homology},
Geometry \& Topology \textbf{24} (2020), no.~6, 2675--2708.
Theorem 1.3.
\url{https://doi.org/10.2140/gt.2020.24.2675}.
\url{https://arxiv.org/abs/1710.09045}.

\bibitem{PerryIHC}
A. Perry,
\emph{The integral Hodge conjecture for two-dimensional Calabi--Yau categories},
Compositio Mathematica \textbf{158} (2022), no.~2, 287--333.
Lemmas 3.4 and 3.6; Theorem 3.5.
\url{https://doi.org/10.1112/S0010437X22007266}.

\bibitem{FuHwangCI}
B. Fu and J.-M. Hwang,
\emph{Isotrivial VMRT-structures of complete intersection type},
Asian Journal of Mathematics \textbf{22} (2018), 333--354.
Lemma 4.1 (Br\"uckmannの消滅定理を引用).
\url{https://doi.org/10.4310/AJM.2018.v22.n2.a10}.
\bibitem{BorelRemmert}
A. Borel and R. Remmert,
\emph{\"Uber kompakte homogene K\"ahlersche Mannigfaltigkeiten},
Mathematische Annalen \textbf{145} (1962), 429--439, Theorem I.
\url{https://wglobke.github.io/publications/BorelRemmert1962.pdf}.

\bibitem{DeLaetLeBruyn}
K. De Laet and L. Le Bruyn,
\emph{The geometry of representations of 3-dimensional Sklyanin algebras},
Algebras and Representation Theory \textbf{18} (2015), 761--776.
arXiv:1405.1158v1, Theorem 1 and Section 2 (canonical central cubic formula);
the result is a reformulation of the Artin--Schelter and Smith--Tate formula.
\url{https://arxiv.org/abs/1405.1158}.

\bibitem{ArtebaniDolgachev}
M. Artebani and I. Dolgachev,
\emph{The Hesse pencil of plane cubic curves},
L'Enseignement Math\'ematique (2) \textbf{55} (2009), no.~3--4, 235--273,
Lemma 1 (every smooth plane cubic admits a Hesse canonical form).
\url{https://doi.org/10.4171/LEM/55-3-3};
\url{https://arxiv.org/abs/math/0611590}.

\bibitem{MSNE12} I. Mori and S. P. Smith,
\emph{The classification of 3-Calabi--Yau algebras with 3 generators and 3 quadratic relations},
Math. Z. \textbf{287} (2017), 215--241.
Lemma 4.1, Theorem 4.2, Corollary 4.3; arXiv v2, pp. 15--17.
\url{https://arxiv.org/abs/1502.07403}; doi:10.1007/s00209-016-1824-5.
\bibitem{LBSVNE12} L. Le Bruyn, S. P. Smith and M. Van den Bergh,
\emph{Central extensions of three dimensional Artin--Schelter regular algebras},
Math. Z. \textbf{222} (1996), 171--212.
Theorem 2.6, Corollary 2.7, Remark 2.8, Theorem 3.1.3.
\url{https://sites.math.washington.edu/~smith/Research/Central-extns.pdf}.
\bibitem{BTNE12} R. Berger and R. Taillefer,
\emph{Poincar\'e--Birkhoff--Witt deformations of Calabi--Yau algebras},
J. Noncommut. Geom. \textbf{1} (2007), 241--270.
Theorems 3.1 and 3.6, especially the potential bimodule complex in the proof of Theorem 3.6.
\url{https://arxiv.org/abs/math/0610112}; doi:10.4171/JNCG/6.
\bibitem{BergNE12} G. M. Bergman,
\emph{The diamond lemma for ring theory}, Adv. Math. \textbf{29} (1978), 178--218, Theorem 1.2.
\url{https://doi.org/10.1016/0001-8708(78)90010-5}.
\bibitem{LP-AZ} M. Artin and J. J. Zhang,
\emph{Noncommutative projective schemes}, Adv. Math. \textbf{109} (1994), 228--287.
\url{https://doi.org/10.1006/aima.1994.1087}.
\bibitem{LP-BS} J. Burke and G. Stevenson,
\emph{The derived category of a graded Gorenstein ring},
Commutative Algebra and Noncommutative Algebraic Geometry, Vol.~2: Research Articles,
Cambridge University Press (2015), 93--124.
\url{https://doi.org/10.1017/9781316570449.006}.
本稿で確認した版はarXiv:1507.00830v1 (2015年7月3日)。
Definition 5.1、Section 4、Lemmas 6.1--6.2、Theorem 6.4。
\url{https://arxiv.org/abs/1507.00830}.
\bibitem{LP-Keller} B. Keller, \emph{On differential graded categories},
International Congress of Mathematicians, Vol. II (2006), 151--190, Theorem 3.8.
\url{https://arxiv.org/abs/math/0601185}.
\bibitem{LP-Orlov} D. Orlov,
\emph{Smooth and proper noncommutative schemes and gluing of DG categories},
Adv. Math. \textbf{302} (2016), 59--105, Section 3.4, Proposition 3.31.
\url{https://arxiv.org/abs/1402.7364}.
\bibitem{ERLowen} Wendy Lowen,
\emph{Hochschild cohomology, the characteristic morphism and derived deformations},
Compositio Mathematica \textbf{144} (2008), 1557--1580.
\url{https://arxiv.org/abs/0707.2602}. 参照：v1, Theorem 4.8, \S4.4.
\bibitem{BFKER} Matthew Ballard, David Favero and Ludmil Katzarkov,
\emph{A category of kernels for equivariant factorizations and its implications for Hodge theory},
Publications Mathématiques de l'IHÉS \textbf{120} (2014), 1--111.
\url{https://doi.org/10.1007/s10240-013-0059-9}.
参照：Theorem 5.39, pp.~82--83.
\bibitem{HRER} Daniel Huybrechts and Jørgen Vold Rennemo,
\emph{Hochschild cohomology versus the Jacobian ring and the Torelli theorem for cubic fourfolds},
Algebraic Geometry \textbf{6} (2019), 76--99.
\url{https://doi.org/10.14231/AG-2019-005}.
参照：Lemma 2.19, \S2.5, \S3.3（cubicでは$n\ge4$）.
\bibitem{JLLZER} Augustinas Jacovskis, Xun Lin, Zhiyu Liu and Shizhuo Zhang,
\emph{Infinitesimal categorical Torelli theorems for Fano threefolds},
arXiv:2203.08187v3 (2023).
\url{https://arxiv.org/abs/2203.08187}. 参照：Theorem 4.4(2).
\bibitem{HuyER} Daniel Huybrechts, \emph{The K3 category of a cubic fourfold},
Compositio Mathematica \textbf{153} (2017), 586--620.
\url{https://doi.org/10.1112/S0010437X16008137}. 参照：Proposition 3.12, \S3.5.
\bibitem{StacksER} The Stacks Project Authors, \emph{The Stacks Project}.
\url{https://stacks.math.columbia.edu/tag/02HX};
\url{https://stacks.math.columbia.edu/tag/055U}.
参照日：2026年11月10日。smoothnessの無限小判定、smooth射のétale局所section。
\bibitem{ERLinZhang} Xun Lin and Shizhuo Zhang,
\emph{Categorical Torelli and Fourier--Mukai partners of hypersurfaces},
arXiv:2611.05915v1 (2026年11月4日).
\url{https://arxiv.org/abs/2611.05915}.
参照：\S2.1, Theorem 2.2, Proposition 2.4.
幾何的な両側projectorを仮定するため、非幾何的targetには直接適用しない。

\bibitem{ExplicitUedaYoshinaga}
K. Ueda and M. Yoshinaga,
\emph{Logarithmic vector fields along smooth divisors in projective spaces},
Hokkaido Math. J. \textbf{38} (2009), 409--415.
\href{https://doi.org/10.14492/hokmj/1258553970}{doi:10.14492/hokmj/1258553970};
\href{https://arxiv.org/abs/0802.2133}{arXiv:0802.2133}.
本稿では勾配空間からの復元を扱うLemma~3を参照する。

\bibitem{ExplicitHwang}
J.-M. Hwang,
\emph{Symmetrizer group of a projective hypersurface},
J. Math. Soc. Japan \textbf{78} (2026), 55--62.
\href{https://doi.org/10.2969/jmsj/93469346}{doi:10.2969/jmsj/93469346};
\href{https://arxiv.org/abs/2501.08633}{arXiv:2501.08633}.
Theorems~1.5--1.6は勾配空間のfiberとsymmetrizer群を記述する。

% Insert inside the existing thebibliography environment.
% JLLZER is defined in papers/reconstruction/etale_references.tex.
\bibitem{PFCRVdB}
Damien Calaque, Carlo A. Rossi and Michel Van den Bergh,
\emph{C\u{a}ld\u{a}raru's conjecture and Tsygan's formality},
Annals of Mathematics (2) \textbf{176} (2012), no.~2, 865--923.
\url{https://doi.org/10.4007/annals.2012.176.2.4}.
\url{https://annals.math.princeton.edu/wp-content/uploads/annals-v176-n2-p04-p.pdf}.
参照：Theorem 1.4, p.~872（Todd補正付きHKRとHochschild加群構造）.

\end{thebibliography}

\end{document}
